\documentclass[12pt,a4paper,oneside]{book}
\usepackage{iftex}
\ifPDFTeX
  \usepackage[utf8]{inputenc}
  \usepackage[T5]{fontenc}
  \usepackage{lmodern}
\else
  \usepackage{fontspec}
  \setmainfont{lmroman10-regular.otf}[
    BoldFont=lmroman10-bold.otf,
    ItalicFont=lmroman10-italic.otf,
    BoldItalicFont=lmroman10-bolditalic.otf]
\fi
\usepackage[vietnamese]{babel}
\usepackage[a4paper,top=2.5cm,bottom=2.5cm,left=3cm,right=2cm,headheight=16pt]{geometry}
\usepackage{amsmath,amssymb,amsthm,mathtools,bm}
\usepackage{xcolor}
\usepackage{tikz}
\usepackage{booktabs,array,enumitem}
\usepackage{needspace}
\usepackage{algorithm}
\usepackage[noend]{algpseudocode}
\usepackage[theorems,breakable]{tcolorbox}
\definecolor{mainblue}{HTML}{1A5276}
\definecolor{subblue}{HTML}{EBF5FB}
\definecolor{darkgreen}{HTML}{196F3D}
\definecolor{lightgreen}{HTML}{EAFAF1}
\definecolor{darkorange}{HTML}{B9770E}
\definecolor{graybox}{HTML}{F8F9F9}
\tcbset{breakable,boxrule=0.7pt,arc=1.5mm,left=3mm,right=3mm,top=2mm,bottom=2mm,fonttitle=\bfseries}
\newtcbtheorem[number within=chapter]{dinhnghia}{Định nghĩa}{colback=lightgreen,colframe=darkgreen,coltitle=white}{dn}
\newtcbtheorem[number within=chapter]{dinhly}{Định lý}{colback=subblue,colframe=mainblue,coltitle=white}{dl}
\newtcbtheorem[number within=chapter]{menhde}{Mệnh đề}{colback=subblue,colframe=mainblue,coltitle=white}{md}
\newtcbtheorem[number within=chapter]{bode}{Bổ đề}{colback=graybox,colframe=mainblue,coltitle=white}{bd}
\newtcbtheorem[number within=chapter]{hequa}{Hệ quả}{colback=lightgreen,colframe=darkgreen,coltitle=white}{hq}
\newtcolorbox{nhanxet}{colback=graybox,colframe=darkorange,title={Nhận xét / Chú ý}}
\floatname{algorithm}{Thuật toán}
\algrenewcommand{\algorithmicrequire}{\textbf{Dữ liệu:}}
\algrenewcommand{\algorithmicensure}{\textbf{Kết quả:}}
\algrenewcommand{\algorithmicfor}{\textbf{Với}}
\algrenewcommand{\algorithmicdo}{\unskip:}
\algrenewcommand{\algorithmicwhile}{\textbf{Trong khi}}
\algrenewcommand{\algorithmicif}{\textbf{Nếu}}
\algrenewcommand{\algorithmicthen}{\unskip:}
\algrenewcommand{\algorithmicelse}{\textbf{Ngược lại:}}
\algtext*{EndFor}
\algtext*{EndIf}
\algtext*{EndWhile}
\numberwithin{algorithm}{chapter}
\setcounter{tocdepth}{3}
\setcounter{secnumdepth}{3}
\theoremstyle{definition}
\newtheorem{vidu}{Ví dụ}[chapter]
\newenvironment{chungminh}{\par\medskip\noindent\textbf{\textit{\color{mainblue}Chứng minh.}}\ }{\hfill$\blacksquare$\par\medskip}
\newcommand{\R}{\mathbb{R}}
\newcommand{\Pn}{\mathcal{P}_n}
\newcommand{\norminf}[1]{\left\lVert #1\right\rVert_{\infty}}
\newcommand{\fl}{\operatorname{fl}}
\newcommand{\lc}{\operatorname{lc}}
\usepackage{fancyhdr}
\pagestyle{fancy}
\fancyhf{}
\fancyhead[L]{\small\nouppercase{\leftmark}}
\fancyhead[R]{\small MI3042}
\fancyfoot[C]{\thepage}
\renewcommand{\headrulewidth}{0.4pt}
\usepackage{hyperref}
\usepackage{bookmark}
\hypersetup{colorlinks=true,linkcolor=mainblue,citecolor=darkgreen,urlcolor=mainblue,bookmarksnumbered=true,pdfauthor={Trần Văn Tèo},pdftitle={Phương pháp số MI3042 -- Lý thuyết}}
\setlist{itemsep=3pt,topsep=5pt}
\setlength{\emergencystretch}{2em}

\begin{document}
\hypersetup{pageanchor=false}
\begin{titlepage}
\centering
\vspace*{1.5cm}
{\scshape\Large ĐH Bách Khoa Hà Nội\par}
{\large KHOA TOÁN-TIN\par}
\vspace{2.5cm}
{\Huge\bfseries\color{mainblue} PHƯƠNG PHÁP SỐ\par}
\vspace{0.6cm}
{\LARGE\bfseries Bản note sinh viên -- Lý thuyết\par}
\vspace{1cm}
{\large Học phần MI3042\par}
\vspace{1cm}
\rule{0.6\textwidth}{1.5pt}\par
\vspace{1.5cm}
{\Large\textbf{Tác giả:} Trần Văn Tèo\par}
{\large\textbf{Giảng viên:} TS. Hà Thị Ngọc Yến\par}
\vfill
{\large Hà Nội, Năm 2026\par}
\end{titlepage}

\frontmatter
\hypersetup{pageanchor=true}
\chapter*{Ký hiệu và phạm vi}
\addcontentsline{toc}{chapter}{Ký hiệu và phạm vi}
Tài liệu được xây dựng theo toàn bộ phần lý thuyết của đề cương MI3042 \cite{outline}: xấp xỉ hàm số; tính gần đúng đạo hàm và tích phân; giải gần đúng phương trình vi phân thường, bài toán biên và giá trị riêng. Các thuật toán được đặt trong phụ lục. Lời giải bài tập được đặt trong tài liệu riêng.

Các kiến thức tiên quyết: đại số tuyến tính hữu hạn chiều; các quy tắc đạo hàm và tích phân cơ bản, công thức Newton--Leibniz, định lý giá trị trung gian, định lý giá trị lớn nhất trên tập compact, định lý Rolle và các đồng nhất thức lượng giác. Khi xét phương trình vi phân, sử dụng thêm tính đầy đủ của không gian các hàm liên tục với chuẩn đều. Các công thức và kết quả riêng của phương pháp số đều được dựng và chứng minh trong tài liệu. Chuẩn đều trên đoạn $[a,b]$ được ký hiệu
\[
 \lVert g\rVert_{\infty,[a,b]}=\max_{x\in[a,b]}|g(x)|,
 \qquad g\in C([a,b]).
\]
Khi đoạn xét đã xác định, viết $\norminf{g}$. Quy ước $\deg0=-\infty$, $x^0=1$, tích rỗng bằng $1$, tổng rỗng bằng $0$, $0!=1$ và $\operatorname{sgn}(0)=0$.

Với $p(x)=\sum_{j=0}^{n}a_jx^j$, vector hệ số là $\boldsymbol a=(a_0,\ldots,a_n)^T\in\R^{n+1}$. Các công thức truy hồi được xét trong số học chính xác; các giả thiết về làm tròn được phát biểu riêng trong phần phân tích sai số. Thuật toán được trình bày ở phụ lục.
\tableofcontents
\mainmatter

\chapter{Xấp xỉ hàm số}
\Needspace{9\baselineskip}
\section{Bài toán nội suy}
\Needspace{7\baselineskip}
\subsection{Đa thức nội suy}
Cho $n\geq0$, các mốc $x_0<\cdots<x_n$ và dữ liệu $y_0,\ldots,y_n\in\R$. Đặt
\[
 \mathcal P_n=\left\{\sum_{j=0}^{n}a_jx^j:a_j\in\R\right\}.
\]
Bài toán nội suy đa thức là tìm $p\in\mathcal P_n$ thỏa
\[
 p(x_i)=y_i,\qquad i=0,\ldots,n.
\]
Khi $y_i=f(x_i)$, ký hiệu $p=I_Xf$, với $X=(x_0,\ldots,x_n)$. Số mốc là $m=n+1$; $n$ là chặn bậc, không bắt buộc là bậc thực của $p$.

\begin{bode}{}{roots}
Với mọi đa thức $p$ và $c\in\R$, tồn tại đa thức $q$ sao cho
\[
 p(x)-p(c)=(x-c)q(x).
\]
Nếu $p\ne0$ và $\deg p=d$, thì $p$ có không quá $d$ nghiệm thực phân biệt.
\end{bode}
\begin{chungminh}
Chọn $d\geq0$ sao cho $p(x)=\sum_{j=0}^{d}a_jx^j$, và đặt
\[
 q(x)=\sum_{j=1}^{d}a_j\sum_{\ell=0}^{j-1}x^{j-1-\ell}c^\ell.
\]
Với $j\geq1$,
\begin{align*}
 (x-c)\sum_{\ell=0}^{j-1}x^{j-1-\ell}c^\ell
 &=\sum_{\ell=0}^{j-1}x^{j-\ell}c^\ell
   -\sum_{\ell=0}^{j-1}x^{j-1-\ell}c^{\ell+1}\\
 &=\sum_{\ell=0}^{j-1}x^{j-\ell}c^\ell
   -\sum_{\ell=1}^{j}x^{j-\ell}c^\ell
 =x^j-c^j.
\end{align*}
Do đó
\[
 (x-c)q(x)=\sum_{j=1}^{d}a_j(x^j-c^j)=p(x)-p(c).
\]
Nếu $p(c)=0$, $d\geq1$ và $a_d\ne0$, thì
\[
 p(x)=(x-c)q(x),\qquad [x^{d-1}]q=a_d\ne0,
 \qquad\deg q=d-1.
\]
Ở đây $[x^k]q$ là hệ số của $x^k$ trong $q$.

Đối với chặn số nghiệm, lấy $p\ne0$, $d=\deg p$ và $Z(p)=\{t\in\R:p(t)=0\}$. Quy nạp theo $d$:
\[
 d=0,\ p\ne0\ \Longrightarrow\ Z(p)=\varnothing.
\]
Giả sử chặn số nghiệm đúng với bậc $d-1$. Nếu $Z(p)=\varnothing$, chặn đúng. Nếu $c\in Z(p)$, thì
\[
 t\ne c:\quad p(t)=0
 \ \Longleftrightarrow\ (t-c)q(t)=0
 \ \Longleftrightarrow\ q(t)=0,
\]
\[
 Z(p)=\{c\}\cup\bigl(Z(q)\setminus\{c\}\bigr),
 \qquad\#Z(p)\leq1+\#Z(q)\leq1+(d-1)=d.
\]
\end{chungminh}

\Needspace{7\baselineskip}
\subsection{Sự tồn tại và duy nhất}
Với $p(x)=\sum_{j=0}^{n}a_jx^j$, đặt $A=(a_0,\ldots,a_n)^T$ và $Y=(y_0,\ldots,y_n)^T$. Khi đó
\begin{equation}\label{eq:vandermonde-system}
 \left[p(x_i)=y_i\ (0\leq i\leq n)\right]
 \ \Longleftrightarrow\ VA=Y,
 \qquad
 V=\begin{bmatrix}
 1&x_0&\cdots&x_0^n\\
 1&x_1&\cdots&x_1^n\\
 \vdots&\vdots&\ddots&\vdots\\
 1&x_n&\cdots&x_n^n
 \end{bmatrix}.
\end{equation}

\begin{dinhly}{}{vandermonde}
Với mọi $x_0,\ldots,x_n\in\R$, định thức ma trận $V$ trong \eqref{eq:vandermonde-system} là
\begin{equation}\label{eq:vandermonde-det}
 \det V=\prod_{0\leq i<j\leq n}(x_j-x_i).
\end{equation}
\end{dinhly}
\begin{chungminh}
Với $n=0$, $\det[1]=1$. Giả sử công thức đúng với $n$ hàng. Cố định $x_0,\ldots,x_{n-1}$ và đặt
\[
 D_n(t)=\det\begin{bmatrix}
 1&x_0&\cdots&x_0^n\\
 \vdots&\vdots&&\vdots\\
 1&x_{n-1}&\cdots&x_{n-1}^n\\
 1&t&\cdots&t^n
 \end{bmatrix}.
\]
Gọi $M_{n,j}$ là định thức con bỏ hàng cuối và cột $j$, với chỉ số cột bắt đầu từ $0$. Khai triển theo hàng cuối:
\[
 D_n(t)=\sum_{j=0}^{n}(-1)^{n+j}M_{n,j}t^j,
 \qquad M_{n,n}=\det V_{n-1}.
\]
Nếu $x_0,\ldots,x_{n-1}$ phân biệt, đặt
\[
 R(t)=D_n(t)-\det V_{n-1}\prod_{i=0}^{n-1}(t-x_i).
\]
Ta có
\[
 [t^n]R=M_{n,n}-\det V_{n-1}=0,\qquad\deg R\leq n-1,
\]
\[
 D_n(x_i)=0\ \text{(hai hàng bằng nhau)},\qquad
 R(x_i)=0\quad(0\leq i<n).
\]
Theo bổ đề \ref{bd:roots}, $R=0$. Bởi vậy
\begin{align*}
 \det V=D_n(x_n)
 &=\det V_{n-1}\prod_{i=0}^{n-1}(x_n-x_i)\\
 &=\left(\prod_{0\leq i<j<n}(x_j-x_i)\right)
   \left(\prod_{i=0}^{n-1}(x_n-x_i)\right)\\
 &=\prod_{0\leq i<j\leq n}(x_j-x_i).
\end{align*}
Nếu $x_i=x_j$ với $0\leq i<j<n$, gọi $R_i(V),R_j(V)$ là hai hàng tương ứng. Khi đó
\[
 R_i(V)=R_j(V)\ \Longrightarrow\ \det V=0,
 \qquad
 \prod_{0\leq k<\ell\leq n}(x_\ell-x_k)
 =(x_j-x_i)\prod_{\substack{k<\ell\\(k,\ell)\ne(i,j)}}(x_\ell-x_k)=0.
\]
\end{chungminh}

\begin{dinhly}{}{interpolation}
Với $n+1$ mốc đôi một phân biệt và mọi $y_0,\ldots,y_n\in\R$, tồn tại duy nhất $p\in\mathcal P_n$ thỏa $p(x_i)=y_i$.
\end{dinhly}
\begin{chungminh}
Theo định lý \ref{dl:vandermonde},
\[
 \det V=\prod_{i<j}(x_j-x_i)\ne0.
\]
Đặt $A=V^{-1}Y$ và $p(x)=\sum_{j=0}^{n}a_jx^j$. Khi đó
\[
 VA=VV^{-1}Y=Y
 \ \Longrightarrow\ p(x_i)=y_i\quad(0\leq i\leq n).
\]
Nếu $\widetilde p(x)=\sum_{j=0}^{n}\widetilde a_jx^j$ cũng nội suy dữ liệu, thì
\[
 V\widetilde A=Y=VA
 \ \Longrightarrow\ V(\widetilde A-A)=0
 \ \Longrightarrow\ \widetilde A-A=V^{-1}0=0
 \ \Longrightarrow\ \widetilde p=p.
\]
\end{chungminh}

\Needspace{7\baselineskip}
\subsection{Sai số nội suy}
Cho $a<b$, $f\in C^{n+1}([a,b])$, các mốc $x_i\in[a,b]$ đôi một phân biệt và $p=I_Xf$. Đặt
\[
 \omega_{n+1}(t)=\prod_{i=0}^{n}(t-x_i).
\]

\begin{bode}{}{rolle}
Cho $k\geq1$ và $g\in C^k([\alpha,\beta])$. Nếu $g$ có $k+1$ nghiệm phân biệt trong $[\alpha,\beta]$, thì tồn tại $\xi\in(\alpha,\beta)$ sao cho $g^{(k)}(\xi)=0$.
\end{bode}
\begin{chungminh}
Chọn các nghiệm $z_0^{(0)}<\cdots<z_k^{(0)}$. Giả sử ở cấp $r<k$ đã có
\[
 \alpha\leq z_0^{(r)}<\cdots<z_{k-r}^{(r)}\leq\beta,
 \qquad g^{(r)}(z_j^{(r)})=0.
\]
Do $g^{(r)}\in C^{k-r}([\alpha,\beta])$, định lý Rolle trên từng đoạn $[z_j^{(r)},z_{j+1}^{(r)}]$ cho
\[
 z_j^{(r+1)}\in(z_j^{(r)},z_{j+1}^{(r)}),\qquad
 g^{(r+1)}(z_j^{(r+1)})=0\quad(0\leq j<k-r).
\]
Các nghiệm mới thỏa
\[
 z_j^{(r+1)}<z_{j+1}^{(r)}<z_{j+1}^{(r+1)}
 \quad(0\leq j<k-r-1).
\]
Quy nạp đến $r=k$ thu được
\[
 \xi=z_0^{(k)}\in(z_0^{(0)},z_k^{(0)})\subset(\alpha,\beta),
 \qquad g^{(k)}(\xi)=0.
\]
\end{chungminh}

\begin{dinhly}{}{error}
Với mọi $x\in[a,b]$, tồn tại $\xi_x\in(a,b)$ sao cho
\begin{equation}\label{eq:error}
 f(x)-p(x)=\frac{f^{(n+1)}(\xi_x)}{(n+1)!}\omega_{n+1}(x).
\end{equation}
\end{dinhly}
\begin{chungminh}
Nếu $x=x_i$, thì
\[
 f(x_i)-p(x_i)=0,\qquad
 \frac{f^{(n+1)}(\xi)}{(n+1)!}\omega_{n+1}(x_i)=0
 \quad\text{với mọi }\xi\in(a,b).
\]
Nếu $x\notin\{x_0,\ldots,x_n\}$, đặt
\[
 \alpha=\min\{x,x_0,\ldots,x_n\},\qquad
 \beta=\max\{x,x_0,\ldots,x_n\},
\]
\[
 K=\frac{f(x)-p(x)}{\omega_{n+1}(x)},\qquad
 g(t)=f(t)-p(t)-K\omega_{n+1}(t).
\]
Ta có $a\leq\alpha<\beta\leq b$, $g\in C^{n+1}([\alpha,\beta])$ và
\[
 g(x_i)=f(x_i)-p(x_i)-K\cdot0=0,
 \qquad g(x)=f(x)-p(x)-K\omega_{n+1}(x)=0.
\]
Viết $\omega_{n+1}(t)=t^{n+1}+\sum_{j=0}^{n}\beta_jt^j$. Khi đó
\[
 p^{(n+1)}=0,\qquad
 \omega_{n+1}^{(n+1)}=(n+1)!.
\]
Theo bổ đề \ref{bd:rolle}, tồn tại $\xi_x\in(\alpha,\beta)$ sao cho
\begin{align*}
 0=g^{(n+1)}(\xi_x)
 &=f^{(n+1)}(\xi_x)-p^{(n+1)}(\xi_x)
   -K\omega_{n+1}^{(n+1)}(\xi_x)\\
 &=f^{(n+1)}(\xi_x)-K(n+1)!.
\end{align*}
Suy ra
\[
 f(x)-p(x)=K\omega_{n+1}(x)
 =\frac{f^{(n+1)}(\xi_x)}{(n+1)!}\omega_{n+1}(x),
 \qquad\xi_x\in(\alpha,\beta)\subset(a,b).
\]
\end{chungminh}

\begin{hequa}{}{uniform-error}
Đặt $M_{n+1}=\norminf{f^{(n+1)}}$. Khi đó
\begin{equation}\label{eq:uniform-error}
 \norminf{f-p}\leq\frac{M_{n+1}}{(n+1)!}\norminf{\omega_{n+1}}.
\end{equation}
\end{hequa}
\begin{chungminh}
Với mọi $x\in[a,b]$, công thức \eqref{eq:error} cho
\[
 |f(x)-p(x)|
 =\frac{|f^{(n+1)}(\xi_x)|}{(n+1)!}|\omega_{n+1}(x)|
 \leq\frac{M_{n+1}}{(n+1)!}|\omega_{n+1}(x)|
 \leq\frac{M_{n+1}}{(n+1)!}\norminf{\omega_{n+1}}.
\]
Lấy $\max_{x\in[a,b]}$ ở vế trái được \eqref{eq:uniform-error}.
\end{chungminh}

\Needspace{7\baselineskip}
\subsection{Đa thức Chebyshev và mốc tối ưu}
\begin{dinhnghia}{}{chebyshev}
Các đa thức Chebyshev loại một được xác định bởi
\begin{equation}\label{eq:chebyshev-rec}
 T_0(t)=1,\qquad T_1(t)=t,\qquad
 T_{k+1}(t)=2tT_k(t)-T_{k-1}(t)\quad(k\geq1).
\end{equation}
\end{dinhnghia}

\begin{menhde}{}{chebyshev-properties}
Với $k\geq1$, $\deg T_k=k$ và $\lc(T_k)=2^{k-1}$. Với $0\leq\theta\leq\pi$,
\[
 T_k(\cos\theta)=\cos(k\theta),\qquad
 \lVert T_k\rVert_{\infty,[-1,1]}=1.
\]
Các nghiệm của $T_k$ là
\[
 \cos\frac{(2j+1)\pi}{2k},\qquad j=0,\ldots,k-1.
\]
Tại $t_j=\cos(j\pi/k)$, $j=0,\ldots,k$, ta có $T_k(t_j)=(-1)^j$.
\end{menhde}
\begin{chungminh}
Với $k=0,1$,
\[
 T_0(\cos\theta)=1=\cos0,\qquad T_1(\cos\theta)=\cos\theta.
\]
Nếu đẳng thức lượng giác đúng ở cấp $k$ và $k-1$, thì
\begin{align*}
 T_{k+1}(\cos\theta)
 &=2\cos\theta\cos(k\theta)-\cos((k-1)\theta)\\
 &=\cos((k+1)\theta)+\cos((k-1)\theta)-\cos((k-1)\theta)\\
 &=\cos((k+1)\theta).
\end{align*}
Vì $\cos([0,\pi])=[-1,1]$,
\[
 \max_{t\in[-1,1]}|T_k(t)|
 =\max_{0\leq\theta\leq\pi}|\cos(k\theta)|=1.
\]
Về hệ số đầu,
\[
 T_1(t)=t,\qquad T_2(t)=2t^2-1.
\]
Giả sử $k\geq2$, $\deg T_k=k$, $\lc(T_k)=2^{k-1}$ và $\deg T_{k-1}=k-1$. Từ \eqref{eq:chebyshev-rec},
\[
 [t^{k+1}]T_{k+1}=2[t^k]T_k-[t^{k+1}]T_{k-1}
 =2\cdot2^{k-1}-0=2^k,
\]
\[
 \deg T_{k+1}=k+1,\qquad\lc(T_{k+1})=2^k.
\]
Đặt $\theta_j=(2j+1)\pi/(2k)$. Khi đó
\[
 0<\theta_0<\cdots<\theta_{k-1}<\pi,
 \qquad T_k(\cos\theta_j)=\cos((2j+1)\pi/2)=0.
\]
Với $0\leq j<k-1$,
\[
 \cos\theta_j-\cos\theta_{j+1}
 =2\sin\frac{\theta_j+\theta_{j+1}}2
   \sin\frac{\theta_{j+1}-\theta_j}2>0.
\]
Đặt $Z(T_k)=\{t\in\R:T_k(t)=0\}$. Các $\cos\theta_j$ phân biệt, nên bổ đề \ref{bd:roots} cho
\[
 \{\cos\theta_0,\ldots,\cos\theta_{k-1}\}\subset Z(T_k),
 \qquad k\leq\#Z(T_k)\leq\deg T_k=k.
\]
Do đó $Z(T_k)=\{\cos\theta_0,\ldots,\cos\theta_{k-1}\}$. Cuối cùng,
\[
 T_k(t_j)=\cos(k\cdot j\pi/k)=\cos(j\pi)=(-1)^j.
\]
\end{chungminh}

\begin{dinhly}{}{monic-minimax}
Với $m\geq1$, trong các đa thức thực bậc $m$ có hệ số đầu bằng $1$, đa thức duy nhất đạt chuẩn đều nhỏ nhất trên $[-1,1]$ là
\[
 q_*(t)=2^{1-m}T_m(t),\qquad
 \lVert q_*\rVert_{\infty,[-1,1]}=2^{1-m}.
\]
\end{dinhly}
\begin{chungminh}
Đặt $C=2^{1-m}$, $t_j=\cos(j\pi/m)$ và $q_*=CT_m$. Theo mệnh đề \ref{md:chebyshev-properties},
\[
 \lc(q_*)=C2^{m-1}=1,\qquad\norminf{q_*}=C,
 \qquad t_0>\cdots>t_m,\qquad q_*(t_j)=C(-1)^j.
\]
Xét $q$ có bậc $m$, hệ số đầu $1$ và $\norminf q\leq C$. Đặt $r=q-q_*$. Ta có
\[
 [t^m]r=1-1=0,\qquad\deg r\leq m-1,
\]
\[
 (-1)^j r(t_j)=(-1)^j q(t_j)-C\leq|q(t_j)|-C\leq0.
\]
Viết $r(t)=\sum_{\ell=0}^{m-1}r_\ell t^\ell$ và xét
\[
 A=\begin{bmatrix}
 1&t_0&\cdots&t_0^{m-1}&r(t_0)\\
 1&t_1&\cdots&t_1^{m-1}&r(t_1)\\
 \vdots&\vdots&&\vdots&\vdots\\
 1&t_m&\cdots&t_m^{m-1}&r(t_m)
 \end{bmatrix}.
\]
Gọi $C_\ell$ là cột $\ell$ của $A$, với chỉ số từ $0$. Khi đó
\[
 C_m=\sum_{\ell=0}^{m-1}r_\ell C_\ell
 \ \Longrightarrow\ \det A=0.
\]
Đặt
\[
 W_j=\prod_{\substack{0\leq k<\ell\leq m\\k,\ell\ne j}}(t_k-t_\ell)>0.
\]
Theo định lý \ref{dl:vandermonde}, định thức con bỏ hàng $j$ và cột cuối bằng
\[
 \prod_{\substack{k<\ell\\k,\ell\ne j}}(t_\ell-t_k)
 =(-1)^{m(m-1)/2}W_j.
\]
Khai triển $\det A$ theo cột cuối:
\[
 0=\det A=(-1)^{m+m(m-1)/2}
 \sum_{j=0}^{m}(-1)^jW_jr(t_j).
\]
Vì $(-1)^jW_jr(t_j)\leq0$, ta được
\[
 \sum_{j=0}^{m}(-1)^jW_jr(t_j)=0
 \ \Longrightarrow\ (-1)^jW_jr(t_j)=0\quad\forall j
 \ \Longrightarrow\ r(t_j)=0\quad\forall j.
\]
Bổ đề \ref{bd:roots} và $\deg r\leq m-1$ cho $r=0$, tức $q=q_*$. Bởi vậy
\[
 q\ne q_*,\ \deg q=m,\ \lc(q)=1
 \ \Longrightarrow\ \norminf q>C=\norminf{q_*}.
\]
\end{chungminh}

\begin{hequa}{}{optimal-nodes}
Cho $a<b$, $m\geq1$, $\mu=(a+b)/2$ và $\rho=(b-a)/2$. Với $m$ mốc phân biệt thuộc $[a,b]$, đặt $\omega_m(x)=\prod_{i=0}^{m-1}(x-x_i)$. Khi đó
\begin{equation}\label{eq:minimal-node-product}
 \min_X\lVert \omega_m\rVert_{\infty,[a,b]}
 =\rho^m2^{1-m}=\frac{(b-a)^m}{2^{2m-1}}.
\end{equation}
Bộ mốc duy nhất theo thứ tự tăng dần là
\begin{equation}\label{eq:optimal-nodes}
 x_i=\mu-\rho\cos\frac{(2i+1)\pi}{2m},\qquad i=0,\ldots,m-1.
\end{equation}
\end{hequa}
\begin{chungminh}
Với mỗi $X$, đặt $q_X(t)=\rho^{-m}\omega_m(\mu+\rho t)$. Vì $\rho>0$,
\[
 \deg q_X=m,\quad\lc(q_X)=\rho^{-m}\rho^m=1,\quad
 \lVert q_X\rVert_{\infty,[-1,1]}=\rho^{-m}\lVert \omega_m\rVert_{\infty,[a,b]}.
\]
Định lý \ref{dl:monic-minimax} cho
\[
 \lVert \omega_m\rVert_{\infty,[a,b]}\geq\rho^m2^{1-m},
\]
\[
 \lVert \omega_m\rVert_{\infty,[a,b]}=\rho^m2^{1-m}
 \ \Longleftrightarrow\
 \omega_m(x)=\rho^m2^{1-m}T_m\left(\frac{x-\mu}{\rho}\right).
\]
Đặt $\theta_i=(2i+1)\pi/(2m)$ và $x_i^*=\mu-\rho\cos\theta_i$. Ta có
\[
 \pi-\theta_i=\frac{(2(m-1-i)+1)\pi}{2m},
 \qquad\frac{x_i^*-\mu}{\rho}=-\cos\theta_i=\cos(\pi-\theta_i),
\]
\[
 T_m\left(\frac{x_i^*-\mu}{\rho}\right)
 =\cos\frac{(2(m-1-i)+1)\pi}{2}=0.
\]
Hơn nữa,
\[
 0<\theta_0<\cdots<\theta_{m-1}<\pi,\qquad
 a<x_i^*<b.
\]
Với $0\leq i<m-1$,
\[
 x_{i+1}^*-x_i^*
 =\rho(\cos\theta_i-\cos\theta_{i+1})
 =2\rho\sin\frac{\theta_i+\theta_{i+1}}2
         \sin\frac{\theta_{i+1}-\theta_i}2>0.
\]
Do đó $a<x_0^*<\cdots<x_{m-1}^*<b$. Đặt
\[
 Q(x)=\rho^m2^{1-m}T_m\left(\frac{x-\mu}{\rho}\right),\qquad
 H(x)=\prod_{i=0}^{m-1}(x-x_i^*)-Q(x).
\]
Ta có
\[
 [x^m]H=1-1=0,\qquad\deg H\leq m-1,
 \qquad H(x_i^*)=0\quad(0\leq i<m).
\]
Theo bổ đề \ref{bd:roots}, $H=0$. Bởi vậy
\[
 \omega_{X^*}=Q,\qquad
 \norminf{\omega_{X^*}}=\rho^m2^{1-m}.
\]
Nếu bộ $X=(x_0,\ldots,x_{m-1})$ cũng đạt chặn, thì
\[
 \omega_X=Q=\omega_{X^*},\qquad
 0=\omega_{X^*}(x_i)=\prod_{j=0}^{m-1}(x_i-x_j^*)
 \ \Longrightarrow\ x_i\in\{x_0^*,\ldots,x_{m-1}^*\}.
\]
Hai tập mốc đều có $m$ phần tử, nên
\[
 \{x_0,\ldots,x_{m-1}\}=\{x_0^*,\ldots,x_{m-1}^*\}.
\]
Với hai bộ đã sắp tăng dần,
\[
 x_i=\text{phần tử nhỏ thứ }i+1\text{ của tập mốc}=x_i^*
 \quad(0\leq i<m).
\]
\end{chungminh}

\begin{dinhly}{}{worst-case}
Cho $m\geq1$, $M\geq0$ và
\[
 \mathcal F_{m,M}=\{f\in C^m([a,b]):\norminf{f^{(m)}}\leq M\}.
\]
Với mỗi bộ $m$ mốc phân biệt $X\subset[a,b]$,
\begin{equation}\label{eq:worst-case}
 \sup_{f\in\mathcal F_{m,M}}\norminf{f-I_Xf}
 =\frac{M}{m!}\norminf{\omega_m}.
\end{equation}
Nếu $M>0$, bộ mốc tối ưu cho vế trái là bộ \eqref{eq:optimal-nodes}, và là duy nhất khi sắp tăng dần.
\end{dinhly}
\begin{chungminh}
Hệ quả \ref{hq:uniform-error} với $n=m-1$ cho
\[
 f\in\mathcal F_{m,M}\ \Longrightarrow\
 \norminf{f-I_Xf}\leq\frac{M}{m!}\norminf{\omega_m},
\]
\[
 \sup_{f\in\mathcal F_{m,M}}\norminf{f-I_Xf}
 \leq\frac{M}{m!}\norminf{\omega_m}.
\]
Đặt $f_*(x)=Mx^m/m!$. Khi đó $f_*^{(m)}=M$ và $f_*\in\mathcal F_{m,M}$. Đa thức
\[
 H(x)=f_*(x)-I_Xf_*(x)-\frac{M}{m!}\omega_m(x)
\]
thỏa
\[
 [x^m]H=\frac{M}{m!}-0-\frac{M}{m!}=0,\quad
 \deg H\leq m-1,\quad H(x_i)=0\quad(0\leq i<m).
\]
Bổ đề \ref{bd:roots} cho $H=0$. Vì vậy
\[
 \sup_{f\in\mathcal F_{m,M}}\norminf{f-I_Xf}
 \geq\norminf{f_*-I_Xf_*}
 =\frac{M}{m!}\norminf{\omega_m}.
\]
Kết hợp hai chặn được \eqref{eq:worst-case}, kể cả $M=0$. Với $M>0$, đặt
\[
 E(X)=\sup_{f\in\mathcal F_{m,M}}\norminf{f-I_Xf},\qquad
 \alpha=M/m!>0.
\]
Hệ quả \ref{hq:optimal-nodes} cho
\[
 E(X)-E(X^*)=\alpha\bigl(\norminf{\omega_X}-\norminf{\omega_{X^*}}\bigr)\geq0,
\]
\[
 E(X)=E(X^*)\ \Longleftrightarrow\
 \norminf{\omega_X}=\norminf{\omega_{X^*}}
 \ \Longleftrightarrow\ X=X^*.
\]
\end{chungminh}

\Needspace{9\baselineskip}
\begin{vidu}
Với $f(x)=x^2$ trên $[-1,1]$ và một mốc $c\in[-1,1]$, $I_{\{c\}}f=c^2$. Đặt $v=c^2\in[0,1]$. Khi đó
\[
 \norminf{f-I_{\{c\}}f}
 =\max_{0\leq u\leq1}|u-v|
 =\max\{v,1-v\}=\frac12+\left|v-\frac12\right|.
\]
Sai số nhỏ nhất bằng $1/2$ tại $c=\pm1/\sqrt2$; mốc Chebyshev $c=0$ cho sai số $1$. Bài toán tối ưu trên lớp $\mathcal F_{m,M}$ và bài toán tối ưu cho một hàm cố định có thể có nghiệm khác nhau.

\end{vidu}

Với $\theta_i=(2i+1)\pi/(2m)$, các đẳng thức đối xứng là
\[
 \theta_{m-1-i}=\pi-\theta_i,\qquad
 x_{m-1-i}=\mu+\rho\cos\theta_i.
\]
Nếu $m=2h+1$, thì $\theta_h=\pi/2$ và $x_h=\mu$.

\Needspace{9\baselineskip}
\section{Sơ đồ Horner và ứng dụng}
\Needspace{7\baselineskip}
\subsection{Biểu diễn đa thức}
Xét
\[
 p(x)=\sum_{j=0}^{n}a_jx^j,\qquad \boldsymbol a=(a_0,\ldots,a_n)^T\in\R^{n+1},\qquad n\geq0.
\]
Chỉ số $n$ là chặn bậc: $a_n$ có thể bằng $0$. Hệ số $a_j$ luôn ứng với đơn thức $x^j$.

\Needspace{7\baselineskip}
\subsection{Tính giá trị đa thức}
\begin{dinhly}{}{horner}
Với $c\in\R$, đặt
\begin{equation}\label{eq:horner}
 b_n=a_n,\qquad b_j=a_j+cb_{j+1}\quad(j=n-1,\ldots,0).
\end{equation}
Khi đó
\begin{equation}\label{eq:horner-invariant}
 b_j=\sum_{\ell=j}^{n}a_\ell c^{\ell-j}\quad(0\leq j\leq n),
 \qquad b_0=p(c).
\end{equation}
\end{dinhly}
\begin{chungminh}
Với $j=n$,
\[
 b_n=a_n=\sum_{\ell=n}^{n}a_\ell c^{\ell-n}.
\]
Giả sử \eqref{eq:horner-invariant} đúng tại $j+1$. Khi đó
\begin{align*}
 b_j&=a_j+cb_{j+1}
 =a_j+c\sum_{\ell=j+1}^{n}a_\ell c^{\ell-j-1}\\
 &=a_j+\sum_{\ell=j+1}^{n}a_\ell c^{\ell-j}
 =\sum_{\ell=j}^{n}a_\ell c^{\ell-j}.
\end{align*}
Quy nạp lùi đến $j=0$ cho $b_0=\sum_{\ell=0}^{n}a_\ell c^\ell=p(c)$.
\end{chungminh}

\Needspace{9\baselineskip}
\begin{vidu}
Sơ đồ ngang của \eqref{eq:horner} với $n\geq3$:
\begin{center}
\begin{tabular}{c|ccccc}
\toprule
$j$ & $n$ & $n-1$ & $n-2$ & $\cdots$ & $0$\\
\midrule
$a_j$ & $a_n$ & $a_{n-1}$ & $a_{n-2}$ & $\cdots$ & $a_0$\\
$cb_{j+1}$ & --- & $cb_n$ & $cb_{n-1}$ & $\cdots$ & $cb_1$\\
\midrule
$b_j$ & $b_n=a_n$ & $b_{n-1}$ & $b_{n-2}$ & $\cdots$ & $\boxed{p(c)}$\\
\bottomrule
\end{tabular}
\end{center}

Với $p(x)=x^3-2x^2+3x-4$:
\begin{center}
\begin{tabular}{c|rrrr}
\toprule
$c=2,\ j$ & $3$ & $2$ & $1$ & $0$\\
\midrule
$a_j$ & $1$ & $-2$ & $3$ & $-4$\\
$2b_{j+1}$ & --- & $2$ & $0$ & $6$\\
\midrule
$b_j$ & $1$ & $0$ & $3$ & $\boxed{2}$\\
\bottomrule
\end{tabular}
\qquad
\begin{tabular}{c|rrrr}
\toprule
$c=-1,\ j$ & $3$ & $2$ & $1$ & $0$\\
\midrule
$a_j$ & $1$ & $-2$ & $3$ & $-4$\\
$-b_{j+1}$ & --- & $-1$ & $3$ & $-6$\\
\midrule
$b_j$ & $1$ & $-3$ & $6$ & $\boxed{-10}$\\
\bottomrule
\end{tabular}
\end{center}
\[
 p(2)=2,\qquad p(-1)=-10.
\]
\end{vidu}

\Needspace{7\baselineskip}
\subsection{Chia cho \texorpdfstring{$x-c$}{x-c}}
\begin{dinhly}{}{synthetic-division}
Với hàng $b$ trong \eqref{eq:horner}, đặt
\[
 q(x)=\sum_{j=1}^{n}b_jx^{j-1},\qquad r=b_0.
\]
Khi $n=0$, đặt $q=0$. Khi đó
\begin{equation}\label{eq:synthetic-division}
 p(x)=(x-c)q(x)+r,\qquad r=p(c).
\end{equation}
Cặp $(q,r)$ là duy nhất trong các cặp đa thức $q$ và hằng số $r$.
\end{dinhly}
\begin{chungminh}
Nếu $n=0$, $(x-c)\cdot0+a_0=p(x)$. Với $n\geq1$, từ \eqref{eq:horner},
\[
 b_n=a_n,\qquad b_j-cb_{j+1}=a_j\quad(0\leq j<n).
\]
Do đó
\begin{align*}
 (x-c)q(x)+r
 &=\sum_{j=1}^{n}b_jx^j
   -\sum_{j=1}^{n}cb_jx^{j-1}+b_0\\
 &=b_nx^n+\sum_{j=1}^{n-1}(b_j-cb_{j+1})x^j+(b_0-cb_1)\\
 &=a_nx^n+\sum_{j=1}^{n-1}a_jx^j+a_0=p(x).
\end{align*}
Thay $x=c$ cho $r=p(c)$. Nếu $p=(x-c)\widetilde q+\widetilde r$, thì
\[
 \widetilde r=p(c)=r,\qquad(x-c)(q-\widetilde q)=0.
\]
Nếu $q-\widetilde q\ne0$, thì
\[
 \lc\bigl((x-c)(q-\widetilde q)\bigr)=\lc(q-\widetilde q)\ne0,
\]
mâu thuẫn với tích bằng đa thức không. Do đó $q=\widetilde q$.
\end{chungminh}

\Needspace{9\baselineskip}
\begin{vidu}
Với $p(x)=x^3-2x^2+3x-4$ và $c=2$:
\begin{center}
\begin{tabular}{c|rrr|r}
\toprule
$j$ & $3$ & $2$ & $1$ & $0$\\
\midrule
$a_j$ & $1$ & $-2$ & $3$ & $-4$\\
$2b_{j+1}$ & --- & $2$ & $0$ & $6$\\
\midrule
$b_j$ & $1$ & $0$ & $3$ & $\boxed{2}$\\
\midrule
$(q,r)$ & $[x^2]q=1$ & $[x]q=0$ & $[1]q=3$ & $r=2$\\
\bottomrule
\end{tabular}
\end{center}
\[
 q(x)=x^2+3,\qquad r=2,\qquad
 p(x)=(x-2)(x^2+3)+2.
\]
\end{vidu}

\Needspace{8\baselineskip}
\subsection{Đạo hàm bằng chia Horner lặp}
\begin{dinhly}{}{repeated-horner}
Đặt $p_0=p$. Với $s=0,\ldots,n$, xác định bằng \eqref{eq:synthetic-division}
\[
 p_s(x)=(x-c)p_{s+1}(x)+r_s,
 \qquad\deg p_s\leq n-s.
\]
Khi đó
\begin{equation}\label{eq:shifted-expansion}
 p(x)=\sum_{s=0}^{n}r_s(x-c)^s,
 \qquad p^{(k)}(c)=k!r_k\quad(0\leq k\leq n).
\end{equation}
Với $k>n$, $p^{(k)}=0$.
\end{dinhly}
\begin{chungminh}
Ta có $p_0\in\mathcal P_n$. Với $0\leq s<n$, giả sử $p_s\in\mathcal P_{n-s}$ và viết
\[
 p_s(x)=\sum_{j=0}^{n-s}a_j^{(s)}x^j.
\]
Theo công thức thương trong \eqref{eq:synthetic-division},
\[
 p_{s+1}(x)=\sum_{j=1}^{n-s}b_j^{(s)}x^{j-1}\in\mathcal P_{n-s-1}.
\]
Quy nạp theo $s$ cho
\[
 p_s\in\mathcal P_{n-s}\quad(0\leq s\leq n),\qquad p_n\in\mathcal P_0,
\]
\[
 p_n(x)=(x-c)\cdot0+p_n(c),\qquad p_{n+1}=0.
\]
Với $s=0$, $p=r_0+(x-c)p_1$. Nếu
\[
 p(x)=\sum_{j=0}^{s}r_j(x-c)^j+(x-c)^{s+1}p_{s+1}(x),
\]
thì thay $p_{s+1}=r_{s+1}+(x-c)p_{s+2}$ cho
\begin{align*}
 p(x)&=\sum_{j=0}^{s}r_j(x-c)^j
 +(x-c)^{s+1}\bigl(r_{s+1}+(x-c)p_{s+2}(x)\bigr)\\
 &=\sum_{j=0}^{s+1}r_j(x-c)^j+(x-c)^{s+2}p_{s+2}(x).
\end{align*}
Quy nạp đến $s=n$, dùng $p_{n+1}=0$, cho đẳng thức khai triển.

Với $s\geq0$, đặt $D^k=d^k/dx^k$. Công thức
\[
 D^k(x-c)^s=
 \begin{cases}
 \dfrac{s!}{(s-k)!}(x-c)^{s-k},&0\leq k\leq s,\\
 0,&k>s
 \end{cases}
\]
đúng tại $k=0$. Nếu $k<s$, lấy đạo hàm biểu thức cấp $k$ được
\[
 D^{k+1}(x-c)^s
 =\frac{s!}{(s-k)!}(s-k)(x-c)^{s-k-1}
 =\frac{s!}{(s-k-1)!}(x-c)^{s-k-1}.
\]
Ở $k=s$, biểu thức là hằng số $s!$, nên đạo hàm tiếp bằng $0$. Do đó
\[
 \left.D^k(x-c)^s\right|_{x=c}
 =\begin{cases}k!,&s=k,\\0,&s\ne k,\end{cases}
 \qquad
 p^{(k)}(c)=\sum_{s=0}^{n}r_s\left.D^k(x-c)^s\right|_{x=c}=k!r_k.
\]
Với $k>n$, mọi số hạng đạo hàm đều bằng $0$.
\end{chungminh}

\Needspace{9\baselineskip}
\begin{vidu}
Với $p(x)=x^3-2x^2+3x-4$ và $c=2$, viết $p_s(x)=\sum_{j=0}^{3-s}a_j^{(s)}x^j$. Khi đó
\begin{center}
\begin{tabular}{c|c|rrrr|c}
\toprule
$s$ & $j$ & $3$ & $2$ & $1$ & $0$ & $r_s$\\
\midrule
$0$ & $a_j^{(0)}$ & $1$ & $-2$ & $3$ & $-4$ &\\
    & $2b_{j+1}^{(0)}$   & --- & $2$ & $0$ & $6$ &\\
    & $b_j^{(0)}$        & $1$ & $0$ & $3$ & $\boxed{2}$ & $2$\\
\midrule
$1$ & $a_j^{(1)}$ & & $1$ & $0$ & $3$ &\\
    & $2b_{j+1}^{(1)}$   & & --- & $2$ & $4$ &\\
    & $b_j^{(1)}$        & & $1$ & $2$ & $\boxed{7}$ & $7$\\
\midrule
$2$ & $a_j^{(2)}$ & & & $1$ & $2$ &\\
    & $2b_{j+1}^{(2)}$   & & & --- & $2$ &\\
    & $b_j^{(2)}$        & & & $1$ & $\boxed{4}$ & $4$\\
\midrule
$3$ & $a_j^{(3)}$ & & & & $\boxed{1}$ & $1$\\
\bottomrule
\end{tabular}
\end{center}
\begin{center}
\begin{tabular}{c|rrrr}
\toprule
$k$ & $0$ & $1$ & $2$ & $3$\\
\midrule
$r_k$ & $2$ & $7$ & $4$ & $1$\\
$k!$ & $1$ & $1$ & $2$ & $6$\\
\midrule
$p^{(k)}(2)=k!r_k$ & $\boxed{2}$ & $\boxed{7}$ & $\boxed{8}$ & $\boxed{6}$\\
\bottomrule
\end{tabular}
\end{center}
\[
 p(x)=2+7(x-2)+4(x-2)^2+(x-2)^3.
\]
\end{vidu}

\Needspace{7\baselineskip}
\subsection{Horner mở rộng}
Đặt
\[
 P_i(x)=\sum_{\ell=i}^{n}a_\ell x^{\ell-i},\qquad
 P_n=a_n,\qquad P_i=xP_{i+1}+a_i.
\]

\begin{dinhly}{}{extended-horner}
Với $0\leq K\leq n$, đặt $d_0^{(n)}=a_n$, $d_j^{(n)}=0$ ($1\leq j\leq K$) và
\begin{equation}\label{eq:extended-horner}
 \begin{aligned}
 d_0^{(i)}&=cd_0^{(i+1)}+a_i,\\
 d_j^{(i)}&=cd_j^{(i+1)}+d_{j-1}^{(i+1)}\quad(1\leq j\leq K),
 \end{aligned}
 \qquad i=n-1,\ldots,0.
\end{equation}
Khi đó
\[
 d_j^{(i)}=\frac{P_i^{(j)}(c)}{j!},\qquad
 d_j^{(0)}=\frac{p^{(j)}(c)}{j!}=r_j.
\]
\end{dinhly}
\begin{chungminh}
Với một đa thức $R$, quy tắc đạo hàm tích cho
\[
 (xR)'=xR'+R.
\]
Nếu $(xR)^{(j)}=xR^{(j)}+jR^{(j-1)}$, thì
\[
 (xR)^{(j+1)}
 =xR^{(j+1)}+R^{(j)}+jR^{(j)}
 =xR^{(j+1)}+(j+1)R^{(j)}.
\]
Bởi quy nạp, công thức đúng với mọi $j\geq1$. Từ $P_i=xP_{i+1}+a_i$,
\[
 \frac{P_i^{(j)}(c)}{j!}
 =c\frac{P_{i+1}^{(j)}(c)}{j!}
  +\frac{P_{i+1}^{(j-1)}(c)}{(j-1)!}\quad(j\geq1),
 \qquad P_i(c)=cP_{i+1}(c)+a_i.
\]
Với $i=n$,
\[
 d_0^{(n)}=a_n=P_n(c),\qquad
 d_j^{(n)}=0=P_n^{(j)}(c)/j!\quad(1\leq j\leq K).
\]
Nếu $d_j^{(i+1)}=P_{i+1}^{(j)}(c)/j!$ với $0\leq j\leq K$, thì \eqref{eq:extended-horner} và các đẳng thức đạo hàm trên cho $d_j^{(i)}=P_i^{(j)}(c)/j!$. Quy nạp lùi đến $i=0$, dùng $P_0=p$ và \eqref{eq:shifted-expansion}, được $d_j^{(0)}=p^{(j)}(c)/j!=r_j$.
\end{chungminh}

\Needspace{9\baselineskip}
\begin{vidu}
Với $p(x)=x^3-2x^2+3x-4$, $c=2$, $K=3$:
\begin{center}
\begin{tabular}{c|rrrr}
\toprule
$i$ & $3$ & $2$ & $1$ & $0$\\
$a_i$ & $1$ & $-2$ & $3$ & $-4$\\
\midrule
$d_0^{(i)}$ & $1$ & $2\cdot1-2=0$ & $2\cdot0+3=3$ & $2\cdot3-4=2$\\
$d_1^{(i)}$ & $0$ & $2\cdot0+1=1$ & $2\cdot1+0=2$ & $2\cdot2+3=7$\\
$d_2^{(i)}$ & $0$ & $2\cdot0+0=0$ & $2\cdot0+1=1$ & $2\cdot1+2=4$\\
$d_3^{(i)}$ & $0$ & $2\cdot0+0=0$ & $2\cdot0+0=0$ & $2\cdot0+1=1$\\
\bottomrule
\end{tabular}
\end{center}
\[
 (d_0^{(0)},d_1^{(0)},d_2^{(0)},d_3^{(0)})=(2,7,4,1),\qquad
 (p(2),p'(2),p''(2),p'''(2))=(2,7,8,6).
\]
\end{vidu}

\Needspace{7\baselineskip}
\subsection{Nhân với \texorpdfstring{$x-c$}{x-c} bằng Horner ngược}
\begin{dinhly}{}{inverse-horner}
Cho $p(x)=\sum_{j=0}^{n}a_jx^j$ và $c\in\R$. Đặt
\begin{equation}\label{eq:inverse-horner}
 d_{n+1}=a_n,\qquad
 d_j=a_{j-1}-ca_j\quad(1\leq j\leq n),\qquad d_0=-ca_0.
\end{equation}
Khi đó
\[
 D(x)=\sum_{j=0}^{n+1}d_jx^j=(x-c)p(x).
\]
\end{dinhly}
\begin{chungminh}
Chia $D$ bằng Horner tại $c$, gọi hàng kết quả là $b_{n+1},\ldots,b_0$. Khi đó
\[
 b_{n+1}=d_{n+1}=a_n.
\]
Với $j=n,\ldots,1$, nếu $b_{j+1}=a_j$, thì
\[
 b_j=d_j+cb_{j+1}=a_{j-1}-ca_j+ca_j=a_{j-1}.
\]
Quy nạp lùi cho $b_j=a_{j-1}$ ($1\leq j\leq n+1$), và
\[
 b_0=d_0+cb_1=-ca_0+ca_0=0.
\]
Theo định lý \ref{dl:synthetic-division},
\[
 D(x)=(x-c)\sum_{j=1}^{n+1}b_jx^{j-1}+b_0
 =(x-c)\sum_{j=1}^{n+1}a_{j-1}x^{j-1}=(x-c)p(x).
\]
\end{chungminh}

\Needspace{9\baselineskip}
\begin{vidu}
Với $p(x)=x^3-2x^2+3x-4$ và $c=2$, quy ước $a_{-1}=a_4=0$:
\begin{center}
\begin{tabular}{c|rrrrr}
\toprule
$j$ & $4$ & $3$ & $2$ & $1$ & $0$\\
\midrule
$a_{j-1}$ & $1$ & $-2$ & $3$ & $-4$ & $0$\\
$-2a_j$ & $0$ & $-2$ & $4$ & $-6$ & $8$\\
\midrule
$d_j$ & $\boxed{1}$ & $\boxed{-4}$ & $\boxed{7}$ & $\boxed{-10}$ & $\boxed{8}$\\
\bottomrule
\end{tabular}
\end{center}
\[
 (x-2)p(x)=x^4-4x^3+7x^2-10x+8.
\]
\end{vidu}

\Needspace{7\baselineskip}
\subsection{Tạo đa thức từ các nhân tử tuyến tính}
\begin{menhde}{}{node-product}
Cho $z_0,\ldots,z_{m-1}\in\R$, $m\geq0$. Đặt $W_0=1$ và xây dựng bằng Horner ngược
\[
 W_{s+1}(x)=(x-z_s)W_s(x)\quad(0\leq s<m).
\]
Khi đó
\begin{equation}\label{eq:node-product}
 W_s(x)=\prod_{i=0}^{s-1}(x-z_i),\qquad\deg W_s=s,\qquad\lc(W_s)=1.
\end{equation}
\end{menhde}
\begin{chungminh}
Với $s=0$, $W_0=1=\prod_{i=0}^{-1}(x-z_i)$, $\deg W_0=0$ và $\lc(W_0)=1$. Nếu \eqref{eq:node-product} đúng ở $s$, định lý \ref{dl:inverse-horner} cho
\[
 W_{s+1}(x)=(x-z_s)\prod_{i=0}^{s-1}(x-z_i)
 =\prod_{i=0}^{s}(x-z_i),
\]
\[
 \deg W_{s+1}=1+s,\qquad\lc(W_{s+1})=1\cdot\lc(W_s)=1.
\]
Quy nạp đến $s=m$ được kết quả.
\end{chungminh}

\Needspace{9\baselineskip}
\begin{vidu}
Với $(z_0,z_1,z_2)=(-1,0,2)$:
\begin{center}
\begin{tabular}{c|c|rrrr}
\toprule
$s$ & $z_{s-1}$ & $[x^3]W_s$ & $[x^2]W_s$ & $[x]W_s$ & $[1]W_s$\\
\midrule
$0$ & --- & $0$ & $0$ & $0$ & $1$\\
$1$ & $-1$ & $0$ & $0$ & $1$ & $1$\\
$2$ & $0$ & $0$ & $1$ & $1$ & $0$\\
$3$ & $2$ & $1$ & $-1$ & $-2$ & $0$\\
\bottomrule
\end{tabular}
\end{center}
\[
 W_0=1,\quad W_1=x+1,\quad W_2=x^2+x,\quad
 W_3=x^3-x^2-2x.
\]
\end{vidu}

\Needspace{7\baselineskip}
\subsection{Nhân hai đa thức bằng truy hồi Horner}
\begin{dinhly}{}{polynomial-product}
Cho $p(x)=\sum_{j=0}^{n}a_jx^j$ và $q(x)=\sum_{k=0}^{m}\beta_kx^k$. Đặt
\begin{equation}\label{eq:polynomial-horner-product}
 R_m=\beta_mp,\qquad R_k=xR_{k+1}+\beta_kp\quad(k=m-1,\ldots,0).
\end{equation}
Khi đó
\[
 R_k(x)=p(x)\sum_{\ell=k}^{m}\beta_\ell x^{\ell-k},\qquad R_0=pq.
\]
Phép $xR_{k+1}$ là Horner ngược với $c=0$.
\end{dinhly}
\begin{chungminh}
Ở $k=m$, $R_m=\beta_mp=p\sum_{\ell=m}^{m}\beta_\ell x^{\ell-m}$. Nếu đẳng thức đúng ở $k+1$, thì
\begin{align*}
 R_k&=xR_{k+1}+\beta_kp
 =xp\sum_{\ell=k+1}^{m}\beta_\ell x^{\ell-k-1}+\beta_kp\\
 &=p\left(\sum_{\ell=k+1}^{m}\beta_\ell x^{\ell-k}+\beta_k\right)
 =p\sum_{\ell=k}^{m}\beta_\ell x^{\ell-k}.
\end{align*}
Quy nạp lùi đến $k=0$ cho $R_0=pq$.
\end{chungminh}

\Needspace{9\baselineskip}
\begin{vidu}
Với $p(x)=x^3-2x^2+3x-4$, $q(x)=x^2+x-1$:
\begin{center}
\begin{tabular}{c|rrrrrr}
\toprule
$R$ & $[x^5]R$ & $[x^4]R$ & $[x^3]R$ & $[x^2]R$ & $[x]R$ & $[1]R$\\
\midrule
$R_2=p$ & $0$ & $0$ & $1$ & $-2$ & $3$ & $-4$\\
$xR_2$ & $0$ & $1$ & $-2$ & $3$ & $-4$ & $0$\\
$\beta_1p=p$ & $0$ & $0$ & $1$ & $-2$ & $3$ & $-4$\\
\midrule
$R_1=xR_2+p$ & $0$ & $1$ & $-1$ & $1$ & $-1$ & $-4$\\
$xR_1$ & $1$ & $-1$ & $1$ & $-1$ & $-4$ & $0$\\
$\beta_0p=-p$ & $0$ & $0$ & $-1$ & $2$ & $-3$ & $4$\\
\midrule
$R_0=xR_1-p$ & $1$ & $-1$ & $0$ & $1$ & $-7$ & $4$\\
\bottomrule
\end{tabular}
\end{center}
\[
 p(x)q(x)=x^5-x^4+x^2-7x+4.
\]
\end{vidu}

\Needspace{7\baselineskip}
\subsection{Chia đa thức bằng Horner tổng quát}
\begin{dinhly}{}{generalized-horner}
Cho $p(x)=\sum_{j=0}^{n}a_jx^j$ và $g(x)=\sum_{j=0}^{d}g_jx^j$, với $n\geq d\geq1$ và $g_d\ne0$. Quy ước $g_j=0$ khi $j\notin\{0,\ldots,d\}$. Xác định theo thứ tự giảm
\begin{equation}\label{eq:generalized-horner}
 \begin{aligned}
 b_i&=\frac{1}{g_d}\left(a_i-
    \sum_{\ell=i+1}^{n}g_{i+d-\ell}b_\ell\right),&&i=n,\ldots,d,\\
 b_j&=a_j-\sum_{\ell=d}^{n}g_{j+d-\ell}b_\ell,&&j=d-1,\ldots,0.
 \end{aligned}
\end{equation}
Khi đó
\begin{equation}\label{eq:generalized-division}
 p(x)=g(x)q(x)+r(x),\qquad
 q(x)=\sum_{\ell=d}^{n}b_\ell x^{\ell-d},\quad
 r(x)=\sum_{j=0}^{d-1}b_jx^j.
\end{equation}
Cặp $(q,r)$ là duy nhất với $\deg r<d$.
\end{dinhly}
\begin{chungminh}
Với $d\leq j\leq n$, hệ số của $x^j$ trong $gq$ là
\begin{align*}
 [x^j](gq)
 &=\sum_{\ell=d}^{n}g_{j+d-\ell}b_\ell\\
 &=g_db_j+\sum_{\ell=j+1}^{n}g_{j+d-\ell}b_\ell=a_j.
\end{align*}
Các số hạng $\ell<j$ bằng $0$ do $j+d-\ell>d$. Với $0\leq j<d$,
\[
 [x^j](gq+r)=\sum_{\ell=d}^{n}g_{j+d-\ell}b_\ell+b_j=a_j.
\]
Với $j>n$, hệ số của cả $p$ và $gq+r$ bằng $0$. Do đó $p=gq+r$.

Nếu $p=g\widetilde q+\widetilde r$ với $\deg\widetilde r<d$, thì
\[
 g(q-\widetilde q)=\widetilde r-r.
\]
Nếu $h=q-\widetilde q\ne0$, đặt $e=\deg h$ và $h_e=[x^e]h\ne0$. Khi đó
\[
 [x^{d+e}](gh)=g_dh_e\ne0,\qquad
 \deg(gh)=d+e\geq d,\qquad\deg(\widetilde r-r)<d,
\]
mâu thuẫn với $gh=\widetilde r-r$. Vậy $q=\widetilde q$ và $r=\widetilde r$.

\end{chungminh}

\Needspace{9\baselineskip}
\begin{vidu}
Với $p(x)=x^3-2x^2+3x-4$ và $g(x)=x^2-x+1$, ta có $g_2=1$, $-g_1=1$, $-g_0=-1$:
\begin{center}
\begin{tabular}{c|rr|rr}
\toprule
$j$ & $3$ & $2$ & $1$ & $0$\\
\midrule
$a_j$ & $1$ & $-2$ & $3$ & $-4$\\
$b_3=1$ & --- & $-g_1b_3=1$ & $-g_0b_3=-1$ & ---\\
$b_2=-1$ & --- & --- & $-g_1b_2=-1$ & $-g_0b_2=1$\\
\midrule
$b_j$ & $1$ & $-1$ & $1$ & $-3$\\
\midrule
$(q,r)$ & $[x]q=1$ & $[1]q=-1$ & $[x]r=1$ & $[1]r=-3$\\
\bottomrule
\end{tabular}
\end{center}
\[
 q(x)=x-1,\quad r(x)=x-3,\quad
 p(x)=(x^2-x+1)(x-1)+(x-3).
\]
\end{vidu}

\Needspace{7\baselineskip}
\subsection{Sai số làm tròn khi tính giá trị}
Giả sử $a_j,c$ là dữ liệu được biểu diễn chính xác, $0<u<1$, và mỗi phép toán được dùng thỏa mô hình
\begin{equation}\label{eq:rounding-model}
 \fl(v\mathbin\circ w)=(v\mathbin\circ w)(1+\delta),
 \qquad |\delta|\leq u,\qquad\circ\in\{+,\times\}.
\end{equation}
Mô hình \eqref{eq:rounding-model} được giả thiết có hiệu lực ở mọi bước; không xét tràn và các trường hợp mất chặn sai số tương đối. Phân tích dưới đây dùng mô hình trong \cite{bindel}.

\begin{bode}{}{gamma}
Với số nguyên $s\geq0$, $su<1$ và $|\delta_i|\leq u$, ta có
\[
 \prod_{i=1}^{s}(1+\delta_i)=1+\theta_s,
 \qquad |\theta_s|\leq\gamma_s:=\frac{su}{1-su}.
\]
\end{bode}
\begin{chungminh}
Với $s=0$, tích bằng $1$, $\theta_0=0=\gamma_0$. Với $s\geq1$, do $1-u>0$,
\[
 (1-u)^s\leq\prod_{i=1}^{s}(1+\delta_i)\leq(1+u)^s.
\]
Bất đẳng thức $(1-u)^r\geq1-ru$ đúng ở $r=0$. Nếu đúng ở $r$, thì
\[
 (1-u)^{r+1}\geq(1-ru)(1-u)
 =1-(r+1)u+ru^2\geq1-(r+1)u.
\]
Do đó $(1-u)^s\geq1-su$. Mặt khác, với $0\leq j\leq s$,
\[
 \binom{s}{j}=\frac{s(s-1)\cdots(s-j+1)}{j!}\leq s^j,
\]
\[
 (1+u)^s=\sum_{j=0}^{s}\binom{s}{j}u^j
 \leq\sum_{j=0}^{s}(su)^j
 \leq\sum_{j=0}^{\infty}(su)^j=\frac{1}{1-su}.
\]
Đặt $\theta_s=\prod_{i=1}^{s}(1+\delta_i)-1$. Ta được
\[
 -su\leq\theta_s\leq\frac{1}{1-su}-1=\frac{su}{1-su},
 \qquad |\theta_s|\leq\frac{su}{1-su}.
\]
\end{chungminh}

\begin{dinhly}{}{horner-roundoff}
Giả sử mỗi bước Horner nhân và cộng với hai lần làm tròn riêng, thỏa \eqref{eq:rounding-model}. Với $n\geq1$, $2nu<1$, kết quả $\widehat p(c)$ thỏa
\begin{align}
 \widehat p(c)&=\sum_{j=0}^{n}(a_j+\Delta a_j)c^j,
 &|\Delta a_j|&\leq\gamma_{2n}|a_j|,\label{eq:backward-horner}\\
 |\widehat p(c)-p(c)|&\leq\gamma_{2n}\sum_{j=0}^{n}|a_j||c|^j.
 \label{eq:forward-horner}
\end{align}
Với $n=0$, $\widehat p(c)=a_0=p(c)$.
\end{dinhly}
\begin{chungminh}
Viết các bước làm tròn dưới dạng
\[
 \widehat b_n=a_n,\qquad
 \widehat b_i=(c\widehat b_{i+1}(1+\delta_i)+a_i)(1+\varepsilon_i),
 \qquad |\delta_i|,|\varepsilon_i|\leq u.
\]
Đặt $\alpha_i=(1+\delta_i)(1+\varepsilon_i)$ và $\eta_i=1+\varepsilon_i$. Khi đó
\[
 \widehat b_i=c\alpha_i\widehat b_{i+1}+a_i\eta_i.
\]
Ta chứng minh bằng quy nạp lùi công thức
\begin{equation}\label{eq:rounded-horner-expansion}
 \widehat b_k=a_nc^{n-k}\prod_{i=k}^{n-1}\alpha_i
 +\sum_{j=k}^{n-1}a_jc^{j-k}\eta_j\prod_{i=k}^{j-1}\alpha_i.
\end{equation}
Ở $k=n$, vế phải bằng $a_n$. Nếu công thức đúng ở $k+1$, thì
\begin{align*}
 \widehat b_k
 &=c\alpha_k\left(a_nc^{n-k-1}\prod_{i=k+1}^{n-1}\alpha_i
   +\sum_{j=k+1}^{n-1}a_jc^{j-k-1}\eta_j
       \prod_{i=k+1}^{j-1}\alpha_i\right)+a_k\eta_k\\
 &=a_nc^{n-k}\prod_{i=k}^{n-1}\alpha_i
   +\sum_{j=k+1}^{n-1}a_jc^{j-k}\eta_j
       \prod_{i=k}^{j-1}\alpha_i+a_k\eta_k\\
 &=a_nc^{n-k}\prod_{i=k}^{n-1}\alpha_i
   +\sum_{j=k}^{n-1}a_jc^{j-k}\eta_j
       \prod_{i=k}^{j-1}\alpha_i.
\end{align*}
Đặt
\[
 \lambda_j=(1+\varepsilon_j)
   \prod_{i=0}^{j-1}(1+\delta_i)(1+\varepsilon_i)\quad(0\leq j<n),
 \qquad
 \lambda_n=\prod_{i=0}^{n-1}(1+\delta_i)(1+\varepsilon_i).
\]
Công thức \eqref{eq:rounded-horner-expansion} với $k=0$ cho
\[
 \widehat p(c)=\widehat b_0=\sum_{j=0}^{n}a_jc^j\lambda_j.
\]
Số thừa số trong $\lambda_j$ là $s_j=2j+1$ khi $j<n$ và $s_n=2n$. Với $0\leq s\leq t$ và $tu<1$,
\[
 \gamma_t-\gamma_s
 =\frac{tu}{1-tu}-\frac{su}{1-su}
 =\frac{(t-s)u}{(1-tu)(1-su)}\geq0.
\]
Theo bổ đề \ref{bd:gamma},
\[
 |\lambda_j-1|\leq\gamma_{s_j}\leq\gamma_{2n}.
\]
Đặt $\Delta a_j=a_j(\lambda_j-1)$. Khi đó
\[
 a_j+\Delta a_j=a_j\lambda_j,\qquad
 |\Delta a_j|\leq\gamma_{2n}|a_j|,
\]
\[
 |\widehat p(c)-p(c)|
 =\left|\sum_{j=0}^{n}\Delta a_jc^j\right|
 \leq\sum_{j=0}^{n}|\Delta a_j||c|^j
 \leq\gamma_{2n}\sum_{j=0}^{n}|a_j||c|^j.
\]
\end{chungminh}

\begin{menhde}{}{coefficient-condition}
Nếu $p(c)\ne0$ và $\varepsilon\geq0$, thì
\[
 \max_{|\Delta a_j|\leq\varepsilon|a_j|}
 \frac{\left|\sum_{j=0}^{n}\Delta a_jc^j\right|}{|p(c)|}
 =\varepsilon\kappa_A(c),\qquad
 \kappa_A(c)=\frac{\sum_{j=0}^{n}|a_j||c|^j}{|p(c)|}.
\]
\end{menhde}
\begin{chungminh}
Với mọi nhiễu thỏa giả thiết,
\[
 \left|\sum_{j=0}^{n}\Delta a_jc^j\right|
 \leq\sum_{j=0}^{n}|\Delta a_j||c|^j
 \leq\varepsilon\sum_{j=0}^{n}|a_j||c|^j.
\]
Chọn $\Delta a_j^*=\varepsilon|a_j|\operatorname{sgn}(p(c))\operatorname{sgn}(c^j)$. Ta có
\[
 |\Delta a_j^*|\leq\varepsilon|a_j|,\qquad
 \sum_{j=0}^{n}\Delta a_j^*c^j
 =\varepsilon\operatorname{sgn}(p(c))\sum_{j=0}^{n}|a_j||c|^j.
\]
Lấy trị tuyệt đối và chia cho $|p(c)|$ cho chặn dưới bằng chặn trên.
\end{chungminh}

Với $p(c)\ne0$, \eqref{eq:forward-horner} cho
\[
 \frac{|\widehat p(c)-p(c)|}{|p(c)|}\leq\gamma_{2n}\kappa_A(c).
\]
Với $p(c)=0$, sử dụng chặn tuyệt đối \eqref{eq:forward-horner}.

\begin{hequa}{}{horner-fma}
Nếu mỗi bước dùng một phép nhân--cộng gộp với một lần làm tròn,
\[
 \widehat b_i=(c\widehat b_{i+1}+a_i)(1+\varepsilon_i),\qquad
 |\varepsilon_i|\leq u,
\]
thì với $nu<1$, các công thức \eqref{eq:backward-horner}--\eqref{eq:forward-horner} đúng khi thay $\gamma_{2n}$ bằng $\gamma_n$.
\end{hequa}
\begin{chungminh}
Đặt $\delta_i=0$ trong \eqref{eq:rounded-horner-expansion}. Khi đó
\[
 \widehat p(c)=\sum_{j=0}^{n}a_jc^j\lambda_j',\qquad
 \lambda_j'=\prod_{i=0}^{j}(1+\varepsilon_i)\ (j<n),\qquad
 \lambda_n'=\prod_{i=0}^{n-1}(1+\varepsilon_i).
\]
Theo bổ đề \ref{bd:gamma},
\[
 |\lambda_j'-1|\leq\gamma_{j+1}\leq\gamma_n\quad(j<n),
 \qquad|\lambda_n'-1|\leq\gamma_n.
\]
Đặt $\Delta a_j=a_j(\lambda_j'-1)$. Suy ra
\[
 |\Delta a_j|\leq\gamma_n|a_j|,\qquad
 |\widehat p(c)-p(c)|
 \leq\sum_{j=0}^{n}|\Delta a_j||c|^j
 \leq\gamma_n\sum_{j=0}^{n}|a_j||c|^j.
\]
\end{chungminh}

\Needspace{9\baselineskip}
\section{Đa thức nội suy Lagrange}
\Needspace{7\baselineskip}
\subsection{Từ hệ Vandermonde đến cơ sở nội suy}
Trong cơ sở $1,x,\ldots,x^n$, đa thức $p(x)=\sum_{j=0}^n a_jx^j$ được xác định bởi
\[
 \mathbf a=(a_0,\ldots,a_n)^T,\qquad
 \mathbf y=(y_0,\ldots,y_n)^T,\qquad V\mathbf a=\mathbf y.
\]
Giải hệ này cho các tọa độ của $p$ trong cơ sở đơn thức. Để đánh giá độ nhạy của các tọa độ, với $\mathbf z\in\R^s$, $s\geq1$, và $B=(b_{ij})\in\R^{s\times s}$, đặt
\[
 \|\mathbf z\|_\infty=\max_{0\leq j<s}|z_j|,\qquad
 \|B\|_\infty=\max_{0\leq i<s}\sum_{j=0}^{s-1}|b_{ij}|,
 \qquad \kappa_\infty(B)=\|B\|_\infty\|B^{-1}\|_\infty
 \quad(\det B\ne0).
\]

\begin{menhde}{}{vandermonde-condition}
Giữ cố định các mốc đôi một phân biệt. Cho $\mathbf y\ne\mathbf0$, $\mathbf a=V^{-1}\mathbf y$ và $\widetilde{\mathbf a}=V^{-1}(\mathbf y+\delta\mathbf y)$. Khi đó
\[
 \frac{\|\widetilde{\mathbf a}-\mathbf a\|_\infty}{\|\mathbf a\|_\infty}
 \leq\kappa_\infty(V)
       \frac{\|\delta\mathbf y\|_\infty}{\|\mathbf y\|_\infty}.
\]
Với hai mốc $0,h$, $0<h\leq1$, ta có
\[
 V_h=\begin{pmatrix}1&0\\1&h\end{pmatrix},\qquad
 V_h^{-1}=\begin{pmatrix}1&0\\-1/h&1/h\end{pmatrix},\qquad
 \kappa_\infty(V_h)=\frac{2(1+h)}h.
\]
\end{menhde}
\begin{chungminh}
Với mọi $B$ và $\mathbf z$,
\[
 \|B\mathbf z\|_\infty
 =\max_i\left|\sum_jb_{ij}z_j\right|
 \leq\max_i\sum_j|b_{ij}||z_j|
 \leq\left(\max_i\sum_j|b_{ij}|\right)\|\mathbf z\|_\infty
 =\|B\|_\infty\|\mathbf z\|_\infty.
\]
Vì $\mathbf y\ne\mathbf0$ và $V\mathbf a=\mathbf y$, $\mathbf a\ne\mathbf0$. Ta có
\[
 \widetilde{\mathbf a}-\mathbf a
 =V^{-1}(\mathbf y+\delta\mathbf y)-V^{-1}\mathbf y
 =V^{-1}\delta\mathbf y,
 \qquad
 \|\mathbf y\|_\infty=\|V\mathbf a\|_\infty
 \leq\|V\|_\infty\|\mathbf a\|_\infty.
\]
Do đó
\begin{align*}
 \frac{\|\widetilde{\mathbf a}-\mathbf a\|_\infty}{\|\mathbf a\|_\infty}
 &\leq\|V^{-1}\|_\infty\|\delta\mathbf y\|_\infty
        \frac{\|V\|_\infty}{\|\mathbf y\|_\infty}\\
 &=\kappa_\infty(V)\frac{\|\delta\mathbf y\|_\infty}{\|\mathbf y\|_\infty}.
\end{align*}
Đối với $V_h$,
\[
 \begin{pmatrix}1&0\\1&h\end{pmatrix}
 \begin{pmatrix}1&0\\-1/h&1/h\end{pmatrix}
 =\begin{pmatrix}1&0\\0&1\end{pmatrix},
 \qquad
 \begin{pmatrix}1&0\\-1/h&1/h\end{pmatrix}
 \begin{pmatrix}1&0\\1&h\end{pmatrix}
 =\begin{pmatrix}1&0\\0&1\end{pmatrix},
\]
\[
 \|V_h\|_\infty=\max\{1,1+h\}=1+h,
 \qquad
 \|V_h^{-1}\|_\infty=\max\{1,2/h\}=2/h,
\]
\[
 \kappa_\infty(V_h)=(1+h)\frac2h=\frac{2(1+h)}h.
\]
\end{chungminh}

\Needspace{9\baselineskip}
\begin{vidu}
Một thiết bị ghi nhiệt độ theo bộ đếm thời gian $x$ tính bằng giây. Xét ba số đo giả định sau; các nhiệt độ được coi là số hữu tỉ chính xác trong ví dụ:
\begin{center}
\begin{tabular}{c|ccc}
\toprule
$i$ & $0$ & $1$ & $2$\\
\midrule
$x_i$ & $M$ & $M+1$ & $M+2$\\
$y_i$ (${ }^{\circ}\mathrm C$) & $20.0$ & $20.1$ & $20.4$\\
\bottomrule
\end{tabular}
\end{center}
Ở đây $M=10^9$. Cần tìm đa thức nội suy bậc không quá $2$ và ước lượng nhiệt độ tại $c=M+3/2$.

Trong cơ sở $1,x,x^2$, hệ phương trình là
\[
 \begin{pmatrix}
 1&M&M^2\\
 1&M+1&M^2+2M+1\\
 1&M+2&M^2+4M+4
 \end{pmatrix}
 \begin{pmatrix}a_0\\a_1\\a_2\end{pmatrix}
 =\begin{pmatrix}20\\201/10\\102/5\end{pmatrix}.
\]
Theo \eqref{eq:vandermonde-det}, $\det V=(1)(2)(1)=2\ne0$.

Trừ phương trình đầu khỏi hai phương trình sau, rồi trừ hai lần phương trình thứ hai mới khỏi phương trình thứ ba mới:
\[
 \left[
 \begin{array}{ccc|c}
 1&M&M^2&20\\
 0&1&2M+1&1/10\\
 0&2&4M+4&2/5
 \end{array}\right]
 \ \xrightarrow{\ R_2\gets R_2-2R_1\ }\
 \left[
 \begin{array}{ccc|c}
 1&M&M^2&20\\
 0&1&2M+1&1/10\\
 0&0&2&1/5
 \end{array}\right].
\]
Chỉ số hàng ở đây bắt đầu từ $0$. Thế ngược cho
\[
 a_2=\frac{1/5}{2}=\frac1{10},\qquad
 a_1=\frac1{10}-(2M+1)\frac1{10}=-\frac M5,
\]
\[
 a_0=20-Ma_1-M^2a_2
 =20+\frac{M^2}5-\frac{M^2}{10}
 =\frac{M^2}{10}+20.
\]

Xét mô hình lưu số với $53$ chữ số nhị phân có nghĩa và làm tròn tới số gần nhất, như định dạng double trong \cite{oracle-754}. Trên khoảng chuẩn hóa $[2^e,2^{e+1})$ đang xét, tập số lưu được là
\[
 \mathbb F_e=\{k2^{e-52}:k\in\mathbb Z,\ 2^{52}\leq k<2^{53}\}.
\]
Hai số liên tiếp thỏa $(k+1)2^{e-52}-k2^{e-52}=2^{e-52}$. Trong ví dụ này,
\[
 2^{59}<M^2<M^2+4M+4<2^{60},\qquad 2^{59-52}=128,
\]
\[
 M^2=128\cdot7812500000000000,
 \qquad 2M=128\cdot15625000,
 \qquad 4M=128\cdot31250000,
\]
\[
 0<1<64,\qquad0<4<64.
\]
Do đó, khi chỉ làm tròn các phần tử $x_i^2$ của ma trận,
\[
 \widehat{x_0^2}=M^2,\qquad
 \widehat{x_1^2}=M^2+2M,\qquad
 \widehat{x_2^2}=M^2+4M.
\]
Các $x_i$ lưu được chính xác: chúng thuộc $[2^{29},2^{30})$ và là bội của $2^{29-52}=2^{-23}$. Ma trận đã lưu là
\[
 \widehat V=
 \begin{pmatrix}
 1&M&M^2\\
 1&M+1&M^2+2M\\
 1&M+2&M^2+4M
 \end{pmatrix}.
\]
Nếu $C_j$ là cột $j$ của $\widehat V$, thì
\[
 C_2=2MC_1-M^2C_0.
\]
Ngoài ra,
\[
 \alpha C_0+\beta C_1=0
 \ \Longrightarrow\ \beta((M+1)-M)=0
 \ \Longrightarrow\ \beta=0
 \ \Longrightarrow\ \alpha=0,
 \qquad\operatorname{rank}\widehat V=2.
\]
Ngay cả khi giải chính xác hệ với ma trận đã lưu và giữ nguyên vector nhiệt độ hữu tỉ, các bước khử Gauss cho
\[
 \left[
 \begin{array}{ccc|c}
 1&M&M^2&20\\
 0&1&2M&1/10\\
 0&2&4M&2/5
 \end{array}\right]
 \ \xrightarrow{\ R_2\gets R_2-2R_1\ }\
 \left[
 \begin{array}{ccc|c}
 1&M&M^2&20\\
 0&1&2M&1/10\\
 0&0&0&1/5
 \end{array}\right].
\]
Phương trình cuối là $0=1/5$; hệ đã lưu vô nghiệm.

Đổi tọa độ thời gian $t=x-M$ cho ba mốc $0,1,2$. Tìm $q(t)=b_0+b_1t+b_2t^2$ từ
\[
 \begin{pmatrix}1&0&0\\1&1&1\\1&2&4\end{pmatrix}
 \begin{pmatrix}b_0\\b_1\\b_2\end{pmatrix}
 =\begin{pmatrix}20\\201/10\\102/5\end{pmatrix}.
\]
Khử Gauss cho
\[
 b_0=20,\qquad b_1+b_2=\frac1{10},\qquad
 2b_1+4b_2=\frac25,
\]
\[
 2b_2=\frac25-2\frac1{10}=\frac15,
 \qquad b_2=\frac1{10},\qquad b_1=0.
\]
Các phần tử của ma trận mới là $0,1,2,4$ và đều lưu được chính xác trong mô hình đã chọn. Đặt
\[
 B=\begin{pmatrix}1&0&0\\1&1&1\\1&2&4\end{pmatrix},\qquad
 B^{-1}=\begin{pmatrix}1&0&0\\-3/2&2&-1/2\\1/2&-1&1/2\end{pmatrix}.
\]
Ta có
\[
 BB^{-1}
 =\begin{pmatrix}
 1&0&0\\
 1-3/2+1/2&2-1&-1/2+1/2\\
 1-3+2&4-4&-1+2
\end{pmatrix}=I.
\]
\[
 \|B\|_\infty=\max\{1,3,7\}=7,
 \qquad\|B^{-1}\|_\infty=\max\{1,4,2\}=4,
\]
\[
 \kappa_\infty(B)=\|B\|_\infty\|B^{-1}\|_\infty=7\cdot4=28.
\]
Vì $q(x_i-M)=y_i$ và $q(x-M)\in\mathcal P_2$, tính duy nhất của đa thức nội suy cho $p(x)=q(x-M)$. Tại $c=M+3/2$, dùng Horner với $t=3/2$:
\begin{center}
\begin{tabular}{c|rrr}
\toprule
$j$ & $2$ & $1$ & $0$\\
\midrule
$b_j$ & $1/10$ & $0$ & $20$\\
$(3/2)d_{j+1}$ & --- & $3/20$ & $9/40$\\
\midrule
$d_j$ & $1/10$ & $3/20$ & $\boxed{809/40}$\\
\bottomrule
\end{tabular}
\end{center}
\[
 p(M+3/2)=q(3/2)=\frac{809}{40}=20.225\ {}^{\circ}\mathrm C.
\]
\end{vidu}

Việc làm tròn các lũy thừa trong tọa độ thời gian gốc đã biến một ma trận khả nghịch thành ma trận hạng $2$; đổi gốc thời gian tránh mất thông tin này trong ví dụ. Nếu chỉ cần giá trị của đa thức, một biểu diễn theo các mốc và dữ liệu còn cho phép tính trực tiếp mà không cần tạo trước các tọa độ trong cơ sở đơn thức. Phần tiếp theo xây dựng biểu diễn đó.

Ta xét các đa thức nhận giá trị $1$ tại một mốc và $0$ tại các mốc còn lại.

\Needspace{7\baselineskip}
\subsection{Dạng tích}
Cho các mốc $x_0,\ldots,x_n$ đôi một phân biệt và $X=\{x_0,\ldots,x_n\}$. Đặt
\[
 N_i(x)=\prod_{\substack{0\leq j\leq n\\j\ne i}}(x-x_j),
 \qquad
 D_i=N_i(x_i)=\prod_{\substack{0\leq j\leq n\\j\ne i}}(x_i-x_j)\ne0.
\]
Các đa thức cơ sở Lagrange là
\begin{equation}\label{eq:lagrange-basis}
 L_i(x)=\frac{N_i(x)}{D_i}
 =\prod_{\substack{0\leq j\leq n\\j\ne i}}
   \frac{x-x_j}{x_i-x_j},\qquad i=0,\ldots,n.
\end{equation}

\begin{dinhly}{}{lagrange}
Các đa thức $L_0,\ldots,L_n$ tạo thành một cơ sở của $\mathcal P_n$. Đa thức nội suy dữ liệu $(x_i,y_i)$ là
\begin{equation}\label{eq:lagrange}
 p(x)=\sum_{i=0}^ny_iL_i(x).
\end{equation}
Với mọi $v\in\mathcal P_n$,
\begin{equation}\label{eq:cardinal}
 v(x)=\sum_{i=0}^nv(x_i)L_i(x).
\end{equation}
\end{dinhly}
\begin{chungminh}
Với $0\leq i,k\leq n$,
\[
 L_i(x_k)=
 \begin{cases}
 \displaystyle\prod_{j\ne i}\frac{x_i-x_j}{x_i-x_j}=1,&k=i,\\[2mm]
 \displaystyle\frac{x_k-x_k}{x_i-x_k}
  \prod_{\substack{j\ne i\\j\ne k}}\frac{x_k-x_j}{x_i-x_j}=0,&k\ne i.
 \end{cases}
\]
Đặt
\[
 C_{ji}=[x^j]L_i,\qquad
 C=(C_{ji})_{0\leq j,i\leq n}\in\R^{(n+1)\times(n+1)},\qquad
 L_i(x)=\sum_{j=0}^n C_{ji}x^j.
\]
Với $V=(x_k^j)_{0\leq k,j\leq n}$ trong \eqref{eq:vandermonde-system},
\[
 (VC)_{ki}=\sum_{j=0}^n V_{kj}C_{ji}
 =\sum_{j=0}^n x_k^jC_{ji}
 =L_i(x_k)=\delta_{ki}\quad(0\leq k,i\leq n).
\]
Do đó
\[
 VC=I_{n+1}.
\]
Định lý \ref{dl:vandermonde} cho $\det V\ne0$, nên
\[
 C=V^{-1}(VC)=V^{-1}I_{n+1}=V^{-1},\qquad CV=I_{n+1}.
\]

Gọi $p$ là đa thức nội suy đã tồn tại theo định lý \ref{dl:interpolation}, với vector hệ số $\mathbf a_p$ thỏa $V\mathbf a_p=\mathbf y$. Đặt
\[
 \mathbf a=C\mathbf y,\qquad
 a_j=\sum_{i=0}^n C_{ji}y_i,\qquad
 Q(x)=\sum_{i=0}^n y_iL_i(x).
\]
Khi đó
\begin{align*}
 Q(x)&=\sum_{i=0}^n y_i\sum_{j=0}^n C_{ji}x^j
 =\sum_{j=0}^n\left(\sum_{i=0}^n C_{ji}y_i\right)x^j
 =\sum_{j=0}^n a_jx^j,\\
 V\mathbf a&=VC\mathbf y=I_{n+1}\mathbf y=\mathbf y.
\end{align*}
Vì vậy
\[
 \sum_{j=0}^n a_jx_k^j=(V\mathbf a)_k=y_k\quad(0\leq k\leq n),
\]
\[
 V(\mathbf a-\mathbf a_p)=\mathbf y-\mathbf y=\mathbf0,
 \qquad
 \mathbf a-\mathbf a_p=V^{-1}\mathbf0=\mathbf0,
 \qquad Q=p.
\]
Đây là \eqref{eq:lagrange}.

Nếu $\sum_{i=0}^n\gamma_iL_i=0$, với $\boldsymbol\gamma=(\gamma_0,\ldots,\gamma_n)^T$, thì
\[
 \sum_{i=0}^n C_{ji}\gamma_i
 =[x^j]\left(\sum_{i=0}^n\gamma_iL_i\right)=0\quad(0\leq j\leq n),
 \qquad C\boldsymbol\gamma=\mathbf0,
\]
\[
 \boldsymbol\gamma=VC\boldsymbol\gamma=V\mathbf0=\mathbf0.
\]
Với $v(x)=\sum_{j=0}^n\beta_jx^j\in\mathcal P_n$, đặt
\[
 \boldsymbol\beta=(\beta_0,\ldots,\beta_n)^T,\qquad
 \mathbf z=(v(x_0),\ldots,v(x_n))^T.
\]
Ta có
\[
 z_k=v(x_k)=\sum_{j=0}^n x_k^j\beta_j=(V\boldsymbol\beta)_k,
 \qquad \mathbf z=V\boldsymbol\beta,
 \qquad \boldsymbol\beta=CV\boldsymbol\beta=C\mathbf z.
\]
Do đó
\begin{align*}
 v(x)&=\sum_{j=0}^n\beta_jx^j
 =\sum_{j=0}^n\left(\sum_{i=0}^n C_{ji}z_i\right)x^j\\
 &=\sum_{i=0}^n z_i\sum_{j=0}^n C_{ji}x^j
 =\sum_{i=0}^n v(x_i)L_i(x).
\end{align*}
Suy ra \eqref{eq:cardinal} và $\mathcal P_n=\operatorname{span}\{L_0,\ldots,L_n\}$. Kết hợp với tính độc lập tuyến tính đã chứng minh, $L_0,\ldots,L_n$ tạo thành một cơ sở của $\mathcal P_n$.
\end{chungminh}

\Needspace{9\baselineskip}
\begin{vidu}
Với $\mathbf x=(-1,0,2)^T$ và $\mathbf y=(6,3,9)^T$,
\begin{center}
\begin{tabular}{c|c|r|c}
\toprule
$x_i$ & $N_i(x)$ & $D_i$ & $L_i(x)$\\
\midrule
$-1$ & $x(x-2)$ & $3$ & $x(x-2)/3$\\
$0$ & $(x+1)(x-2)$ & $-2$ & $-(x+1)(x-2)/2$\\
$2$ & $(x+1)x$ & $6$ & $(x+1)x/6$\\
\bottomrule
\end{tabular}
\end{center}
Do đó
\[
 p(x)=2x(x-2)-\frac32(x+1)(x-2)+\frac32(x+1)x.
\]
\[
 p(-1)=6,\qquad p(0)=3,\qquad p(2)=9.
\]
\end{vidu}

Dạng tích xác định riêng từng $L_i$ và cho phép dựng hệ số của $p$ từ các tích đa thức. Với các mốc cố định, $L_i$ không phụ thuộc $\mathbf y$; dữ liệu khác được thay trực tiếp vào \eqref{eq:lagrange}.

\Needspace{7\baselineskip}
\subsection{Dạng dùng đa thức tích các mốc}
Đặt
\[
 \omega_{n+1}(x)=\prod_{j=0}^n(x-x_j),\qquad w_i=\frac1{D_i}.
\]

\begin{menhde}{}{lagrange-node-product}
Với mọi $i=0,\ldots,n$,
\[
 D_i=\omega_{n+1}'(x_i).
\]
Với $x\notin\{x_0,\ldots,x_n\}$,
\begin{equation}\label{eq:lagrange-omega}
 p(x)=\omega_{n+1}(x)\sum_{i=0}^n\frac{y_i}{D_i(x-x_i)}
     =\omega_{n+1}(x)\sum_{i=0}^n\frac{w_iy_i}{x-x_i}.
\end{equation}
Vế phải của \eqref{eq:lagrange-omega} có giới hạn $y_k$ khi $x\to x_k$.
\end{menhde}
\begin{chungminh}
Ta có
\[
 \omega_{n+1}(x)=(x-x_i)N_i(x),\qquad
 \omega_{n+1}'(x)=N_i(x)+(x-x_i)N_i'(x),
\]
\[
 \omega_{n+1}'(x_i)=N_i(x_i)+(x_i-x_i)N_i'(x_i)=D_i.
\]
Với $x\notin X$,
\[
 \frac{\omega_{n+1}(x)}{D_i(x-x_i)}
 =\frac{(x-x_i)N_i(x)}{D_i(x-x_i)}
 =\frac{N_i(x)}{D_i}=L_i(x),
\]
\[
 \omega_{n+1}(x)\sum_{i=0}^n\frac{y_i}{D_i(x-x_i)}
 =\sum_{i=0}^ny_i\frac{N_i(x)}{D_i}
 =\sum_{i=0}^ny_iL_i(x)=p(x).
\]
Với mỗi $k$,
\begin{align*}
 \lim_{\substack{x\to x_k\\x\notin X}}
 \left(\omega_{n+1}(x)\sum_{i=0}^n\frac{y_i}{D_i(x-x_i)}\right)
 &=\sum_{i=0}^ny_i\frac{N_i(x_k)}{D_i}\\
 &=\sum_{i=0}^ny_i\delta_{ik}=y_k.
\end{align*}
\end{chungminh}

\Needspace{9\baselineskip}
\begin{vidu}
Với $\mathbf x=(-1,0,2)^T$ và $\mathbf y=(6,3,9)^T$,
\[
 \omega_3(x)=(x+1)x(x-2)=x^3-x^2-2x,
 \qquad
 (w_0,w_1,w_2)=\left(\frac13,-\frac12,\frac16\right),
\]
\[
 p(x)=\omega_3(x)\left(\frac2{x+1}-\frac{3}{2x}+\frac{3}{2(x-2)}\right),
 \qquad x\notin\{-1,0,2\}.
\]
Tại $x=1$,
\begin{center}
\begin{tabular}{c|rrr|c}
\toprule
$i$ & $0$ & $1$ & $2$ & Tổng\\
\midrule
$x_i$ & $-1$ & $0$ & $2$ & ---\\
$w_iy_i/(1-x_i)$ & $1$ & $-3/2$ & $-3/2$ & $-2$\\
\bottomrule
\end{tabular}
\end{center}
\[
 \omega_3(1)=(1+1)\cdot1\cdot(1-2)=-2,
 \qquad p(1)=(-2)(-2)=4.
\]
\end{vidu}

Dạng \eqref{eq:lagrange-omega} dùng chung một đa thức $\omega_{n+1}$ cho tất cả các số hạng. Với $X$ cố định, $\omega_{n+1}$ và $w_i$ có thể dùng lại khi thay dữ liệu $\mathbf y$. Biểu diễn \eqref{eq:lagrange-omega} được gọi là dạng barycentric thứ nhất \cite{berrut}.


\begin{menhde}{}{lagrange-evaluation-cost}
Cho $n\geq1$, $x\notin X$. Giả sử các số $w_i$, $\eta_i=w_iy_i$ đã được xác định. Đặt $t_i=x-x_i$. Xét việc tính hai biểu thức
\[
 \sum_{i=0}^n\eta_i\prod_{j\ne i}t_j,
 \qquad
 \left(\prod_{j=0}^nt_j\right)\sum_{i=0}^n\frac{\eta_i}{t_i}.
\]
Nếu mỗi tích được tính riêng bằng các phép nhân liên tiếp, mỗi tổng bằng các phép cộng liên tiếp, không giản lược theo giá trị đặc biệt của dữ liệu, thì số phép toán sau khi đã có $t_i$ là
\begin{center}
\begin{tabular}{c|ccc}
\toprule
Biểu thức & Nhân & Chia & Cộng\\
\midrule
Dạng tích & $n(n+1)$ & $0$ & $n$\\
Dạng $\omega_{n+1}$ & $n+1$ & $n+1$ & $n$\\
\bottomrule
\end{tabular}
\end{center}
Việc xác định $t_0,\ldots,t_n$ cần $n+1$ phép trừ ở cả hai trường hợp.
\end{menhde}
\begin{chungminh}
Tích $N_i(x)=\prod_{j\ne i}t_j$ có $n$ thừa số, nên cần $n-1$ phép nhân; nhân với $\eta_i$ cần thêm một phép nhân. Do đó
\[
 C_{\times,1}=\sum_{i=0}^n((n-1)+1)=n(n+1),
 \qquad C_{\div,1}=0,\qquad C_{+,1}=(n+1)-1=n.
\]
Tích $\prod_{j=0}^nt_j$ có $n+1$ thừa số. Các thương $\eta_i/t_i$ và phép nhân cuối cho
\[
 C_{\times,2}=((n+1)-1)+1=n+1,
 \qquad C_{\div,2}=\sum_{i=0}^n1=n+1,
 \qquad C_{+,2}=(n+1)-1=n.
\]
Mỗi $t_i=x-x_i$ cần một phép trừ, nên $C_- =\sum_{i=0}^n1=n+1$.
\end{chungminh}

\Needspace{7\baselineskip}
\subsection{Dạng thương barycentric}

\begin{menhde}{}{barycentric-ratio}
Với $x\notin X$, đặt
\[
 S_0(x)=\sum_{i=0}^n\frac{w_i}{x-x_i},\qquad
 S_1(x)=\sum_{i=0}^n\frac{w_iy_i}{x-x_i}.
\]
Khi đó
\begin{equation}\label{eq:barycentric-denominator}
 S_0(x)=\frac1{\omega_{n+1}(x)}\ne0,
\end{equation}
\begin{equation}\label{eq:barycentric-ratio}
 p(x)=\frac{S_1(x)}{S_0(x)}
 =\frac{\displaystyle\sum_{i=0}^nw_iy_i/(x-x_i)}
        {\displaystyle\sum_{i=0}^nw_i/(x-x_i)}.
\end{equation}
Thay tất cả $w_i$ bằng $\alpha w_i$, với $\alpha\ne0$, không thay đổi thương trong \eqref{eq:barycentric-ratio}. Thương này có giới hạn $y_k$ khi $x\to x_k$.
\end{menhde}
\begin{chungminh}
Áp dụng \eqref{eq:cardinal} cho $v(x)=1$:
\[
 1=\sum_{i=0}^nL_i(x).
\]
Với $x\notin X$,
\[
 1=\sum_{i=0}^n\frac{\omega_{n+1}(x)w_i}{x-x_i}
   =\omega_{n+1}(x)S_0(x),
 \qquad S_0(x)=\frac1{\omega_{n+1}(x)}\ne0,
\]
\[
 \frac{S_1(x)}{S_0(x)}
 =\omega_{n+1}(x)S_1(x)=p(x).
\]
Với $\alpha\ne0$,
\[
 \frac{\displaystyle\sum_{i=0}^n\alpha w_iy_i/(x-x_i)}
      {\displaystyle\sum_{i=0}^n\alpha w_i/(x-x_i)}
 =\frac{\alpha S_1(x)}{\alpha S_0(x)}
 =\frac{S_1(x)}{S_0(x)}.
\]
Với $x\notin X$, nhân tử và mẫu với $x-x_k$ cho
\[
 \frac{S_1(x)}{S_0(x)}
 =\frac{\displaystyle w_ky_k+
       \sum_{\substack{0\leq i\leq n\\i\ne k}}w_iy_i\frac{x-x_k}{x-x_i}}
      {\displaystyle w_k+
       \sum_{\substack{0\leq i\leq n\\i\ne k}}w_i\frac{x-x_k}{x-x_i}}.
\]
Do $x_k-x_i\ne0$ khi $i\ne k$ và $w_k\ne0$,
\[
 \lim_{\substack{x\to x_k\\x\notin X}}\frac{S_1(x)}{S_0(x)}
 =\frac{w_ky_k+\sum_{i\ne k}w_iy_i\cdot0}
       {w_k+\sum_{i\ne k}w_i\cdot0}=y_k.
\]
\end{chungminh}

\Needspace{9\baselineskip}
\begin{vidu}
Với $\mathbf x=(-1,0,2)^T$, $\mathbf y=(6,3,9)^T$ và $x=1$,
\begin{center}
\begin{tabular}{c|rrr|r}
\toprule
$i$ & $0$ & $1$ & $2$ & Tổng\\
\midrule
$x_i$ & $-1$ & $0$ & $2$ & ---\\
$L_i(1)$ & $-1/3$ & $1$ & $1/3$ & $1$\\
$y_iL_i(1)$ & $-2$ & $3$ & $3$ & $4$\\
$w_i/(1-x_i)$ & $1/6$ & $-1/2$ & $-1/6$ & $-1/2$\\
$w_iy_i/(1-x_i)$ & $1$ & $-3/2$ & $-3/2$ & $-2$\\
\bottomrule
\end{tabular}
\end{center}
\[
 p(1)=\frac{-2}{-1/2}=4.
\]
Với $\alpha=6$, các trọng số mới là $(2,-3,1)$. Tại $x=1$,
\[
 \widetilde S_1(1)=6S_1(1)=-12,
 \qquad\widetilde S_0(1)=6S_0(1)=-3,
 \qquad\frac{\widetilde S_1(1)}{\widetilde S_0(1)}=4.
\]
\end{vidu}

Dạng thương \eqref{eq:barycentric-ratio} là dạng barycentric thứ hai \cite{berrut}. Nó không cần tính thừa số $\omega_{n+1}(x)$, và cho phép nhân chung các trọng số với một hằng khác $0$. Đối với dạng \eqref{eq:lagrange-omega}, nếu giữ nguyên $\omega_{n+1}$ thì
\[
 \omega_{n+1}(x)\sum_{i=0}^n\frac{\alpha w_iy_i}{x-x_i}
 =\alpha p(x).
\]

\begin{menhde}{}{barycentric-quotient-error}
Giữ cố định $x\notin X$, đặt
\[
 \ell(x)=\sum_{i=0}^n|L_i(x)|,\qquad
 A(x)=\sum_{i=0}^n|y_iL_i(x)|.
\]
Cho $\varepsilon\geq0$, $\varepsilon\ell(x)<1$ và hai giá trị $\widehat S_0,\widehat S_1$ thỏa
\[
 |\widehat S_0-S_0(x)|\leq\varepsilon\sum_{i=0}^n\left|\frac{w_i}{x-x_i}\right|,
 \qquad
 |\widehat S_1-S_1(x)|\leq\varepsilon\sum_{i=0}^n\left|\frac{w_iy_i}{x-x_i}\right|.
\]
Khi đó $\widehat S_0\ne0$, và thương chính xác $\widehat q=\widehat S_1/\widehat S_0$ thỏa
\[
 |\widehat q-p(x)|
 \leq\frac{\varepsilon\bigl(A(x)+|p(x)|\ell(x)\bigr)}{1-\varepsilon\ell(x)}.
\]
\end{menhde}
\begin{chungminh}
Đặt $\delta_0=\widehat S_0-S_0(x)$ và $\delta_1=\widehat S_1-S_1(x)$. Từ \eqref{eq:barycentric-denominator} và $L_i(x)=\omega_{n+1}(x)w_i/(x-x_i)$,
\[
 \frac{\sum_i|w_i/(x-x_i)|}{|S_0(x)|}=\sum_i|L_i(x)|=\ell(x),
 \qquad
 \frac{\sum_i|w_iy_i/(x-x_i)|}{|S_0(x)|}=\sum_i|y_iL_i(x)|=A(x).
\]
Do đó
\[
 |\delta_0|\leq\varepsilon|S_0(x)|\ell(x),\qquad
 |\delta_1|\leq\varepsilon|S_0(x)|A(x),
\]
\[
 |\widehat S_0|=|S_0(x)+\delta_0|
 \geq|S_0(x)|-|\delta_0|
 \geq|S_0(x)|\bigl(1-\varepsilon\ell(x)\bigr)>0.
\]
Vì $S_1(x)=p(x)S_0(x)$,
\begin{align*}
 |\widehat q-p(x)|
 &=\left|\frac{S_1(x)+\delta_1-p(x)(S_0(x)+\delta_0)}{S_0(x)+\delta_0}\right|\\
 &=\frac{|\delta_1-p(x)\delta_0|}{|S_0(x)+\delta_0|}\\
 &\leq\frac{|\delta_1|+|p(x)||\delta_0|}{|S_0(x)|\bigl(1-\varepsilon\ell(x)\bigr)}\\
 &\leq\frac{\varepsilon\bigl(A(x)+|p(x)|\ell(x)\bigr)}{1-\varepsilon\ell(x)}.
\end{align*}
\end{chungminh}

\Needspace{7\baselineskip}
\subsection{Dựng hệ số bằng các hệ thức Horner}

\begin{menhde}{}{lagrange-horner}
Viết $\omega_{n+1}(x)=\sum_{j=0}^{n+1}c_jx^j$, với $c_{n+1}=1$. Với mỗi $i$, đặt
\[
 b_{i,n}=c_{n+1},\qquad
 b_{i,j}=c_{j+1}+x_ib_{i,j+1}\quad(j=n-1,\ldots,0),
\]
\[
 Q_i(x)=\sum_{j=0}^nb_{i,j}x^j,\qquad
 r_i=c_0+x_ib_{i,0},\qquad\lambda_i=Q_i(x_i).
\]
Khi đó
\[
 r_i=0,\qquad Q_i(x)=N_i(x),\qquad
 \lambda_i=D_i=\omega_{n+1}'(x_i)\ne0.
\]
Đặt $\mathbf q_i=(b_{i,0},\ldots,b_{i,n})^T$. Đa thức $p(x)=\sum_{j=0}^na_jx^j$ có
\begin{equation}\label{eq:lagrange-coefficients}
 \begin{aligned}
 p(x)&=\sum_{i=0}^n\frac{y_i}{\lambda_i}Q_i(x),\\
 a_j&=\sum_{i=0}^n\frac{y_i}{\lambda_i}b_{i,j},\qquad
 \mathbf a=\sum_{i=0}^n\frac{y_i}{\lambda_i}\mathbf q_i.
 \end{aligned}
\end{equation}
\end{menhde}
\begin{chungminh}
Các hệ thức xác định $b_{i,j}$ và $r_i$ cho
\begin{align*}
 (x-x_i)Q_i(x)+r_i
 &=b_{i,n}x^{n+1}
   +\sum_{j=1}^n(b_{i,j-1}-x_ib_{i,j})x^j+r_i-x_ib_{i,0}\\
 &=c_{n+1}x^{n+1}+\sum_{j=1}^nc_jx^j+c_0=\omega_{n+1}(x).
\end{align*}
Thay $x=x_i$:
\[
 r_i=\omega_{n+1}(x_i)=0.
\]
Mặt khác, $\omega_{n+1}(x)=(x-x_i)N_i(x)+0$. Theo tính duy nhất trong định lý \ref{dl:synthetic-division},
\[
 Q_i=N_i,\qquad
 \lambda_i=Q_i(x_i)=N_i(x_i)=D_i=\omega_{n+1}'(x_i)\ne0.
\]
Vì $Q_i/\lambda_i=N_i/D_i=L_i$, công thức \eqref{eq:lagrange} cho
\begin{align*}
 p(x)&=\sum_{i=0}^n\frac{y_i}{\lambda_i}Q_i(x)
 =\sum_{i=0}^n\frac{y_i}{\lambda_i}\sum_{j=0}^nb_{i,j}x^j\\
 &=\sum_{j=0}^n\left(\sum_{i=0}^n\frac{y_i}{\lambda_i}b_{i,j}\right)x^j.
\end{align*}
Với $j=0,\ldots,n$,
\[
 a_j=[x^j]p=\sum_{i=0}^n\frac{y_i}{\lambda_i}b_{i,j}.
\]
Do đó
\[
 \mathbf a=\begin{pmatrix}a_0\\\vdots\\a_n\end{pmatrix}
 =\sum_{i=0}^n\frac{y_i}{\lambda_i}
   \begin{pmatrix}b_{i,0}\\\vdots\\b_{i,n}\end{pmatrix}
 =\sum_{i=0}^n\frac{y_i}{\lambda_i}\mathbf q_i.
\]
\end{chungminh}

\Needspace{9\baselineskip}
\begin{vidu}
Với $\mathbf x=(-1,0,2)^T$, $\mathbf y=(6,3,9)^T$ và $\omega_3(x)=x^3-x^2-2x$, các hệ thức chia Horner cho
\begin{center}
\begin{tabular}{c|rrr|r|r|r}
\toprule
$x_i$ & $[x^2]Q_i$ & $[x]Q_i$ & $[x^0]Q_i$ & $r_i$ & $\lambda_i$ & $y_i/\lambda_i$\\
\midrule
$-1$ & $1$ & $-2$ & $0$ & $0$ & $3$ & $2$\\
$0$ & $1$ & $-1$ & $-2$ & $0$ & $-2$ & $-3/2$\\
$2$ & $1$ & $1$ & $0$ & $0$ & $6$ & $3/2$\\
\bottomrule
\end{tabular}
\end{center}
Phép tổ hợp tuyến tính các $Q_i$ được biểu diễn bởi
\begin{center}
\begin{tabular}{c|rrr}
\toprule
Đa thức & $[x^2]$ & $[x]$ & $[x^0]$\\
\midrule
$(y_0/\lambda_0)Q_0$ & $2$ & $-4$ & $0$\\
$(y_1/\lambda_1)Q_1$ & $-3/2$ & $3/2$ & $3$\\
$(y_2/\lambda_2)Q_2$ & $3/2$ & $3/2$ & $0$\\
\midrule
$p$ & $\boxed{2}$ & $\boxed{-1}$ & $\boxed{3}$\\
\bottomrule
\end{tabular}
\end{center}
\[
 \mathbf a
 =2\begin{pmatrix}0\\-2\\1\end{pmatrix}
  -\frac32\begin{pmatrix}-2\\-1\\1\end{pmatrix}
  +\frac32\begin{pmatrix}0\\1\\1\end{pmatrix}
 =\begin{pmatrix}3\\-1\\2\end{pmatrix},\qquad
 p(x)=2x^2-x+3.
\]
\end{vidu}

\Needspace{7\baselineskip}
\subsection{Độ nhạy đối với nhiễu dữ liệu}

\begin{menhde}{}{data-sensitivity}
Cho $a<b$ và các mốc cố định $x_i\in[a,b]$ đôi một phân biệt. Gọi $p,\widetilde p$ là các đa thức nội suy dữ liệu $y_i$ và $y_i+\delta y_i$. Khi đó
\[
 \widetilde p(x)-p(x)=\sum_{i=0}^n\delta y_iL_i(x),\qquad
 \norminf{\widetilde p-p}\leq\Lambda_n\max_i|\delta y_i|,
\]
trong đó
\[
 \Lambda_n=\max_{x\in[a,b]}\sum_{i=0}^n|L_i(x)|.
\]
Hằng số $\Lambda_n$ là nhỏ nhất để bất đẳng thức đúng với mọi nhiễu dữ liệu.
\end{menhde}
\begin{chungminh}
Từ \eqref{eq:lagrange},
\[
 \widetilde p-p
 =\sum_{i=0}^n(y_i+\delta y_i)L_i-\sum_{i=0}^ny_iL_i
 =\sum_{i=0}^n\delta y_iL_i.
\]
Do đó
\begin{align*}
 \norminf{\widetilde p-p}
 &=\max_{x\in[a,b]}\left|\sum_{i=0}^n\delta y_iL_i(x)\right|\\
 &\leq\max_{x\in[a,b]}\sum_{i=0}^n|\delta y_i||L_i(x)|\\
 &\leq\left(\max_i|\delta y_i|\right)
       \max_{x\in[a,b]}\sum_{i=0}^n|L_i(x)|
 =\Lambda_n\max_i|\delta y_i|.
\end{align*}
Hàm $\sum_i|L_i|$ liên tục trên $[a,b]$, nên chọn được $x_*\in[a,b]$ thỏa
\[
 \sum_{i=0}^n|L_i(x_*)|=\Lambda_n.
\]
Áp dụng \eqref{eq:cardinal} cho $v=1$,
\[
 \sum_{i=0}^nL_i(x_*)=1
 \ \Longrightarrow\ \exists i:\ L_i(x_*)\ne0.
\]
Với $\varepsilon>0$, chọn $\delta y_i^*=\varepsilon\operatorname{sgn}(L_i(x_*))$. Khi đó
\[
 \max_i|\delta y_i^*|=\varepsilon,\qquad
 (\widetilde p-p)(x_*)
 =\varepsilon\sum_{i=0}^n\operatorname{sgn}(L_i(x_*))L_i(x_*)
 =\varepsilon\sum_{i=0}^n|L_i(x_*)|=\varepsilon\Lambda_n,
\]
\[
 \varepsilon\Lambda_n
 =|(\widetilde p-p)(x_*)|
 \leq\norminf{\widetilde p-p}
 \leq\Lambda_n\max_i|\delta y_i^*|=\varepsilon\Lambda_n.
\]
Nếu $C$ thỏa chặn với mọi nhiễu dữ liệu, áp dụng cho $\delta\mathbf y^*$ cho
\[
 \varepsilon\Lambda_n=\norminf{\widetilde p-p}
 \leq C\max_i|\delta y_i^*|=C\varepsilon
 \ \Longrightarrow\ C\geq\Lambda_n.
\]
\end{chungminh}

\Needspace{9\baselineskip}
\begin{vidu}
Với hai mốc $0,h$, $0<h\leq1$, và đoạn xét $[0,h]$,
\[
 L_0(x)=\frac{x-h}{-h}=1-\frac{x}{h},\qquad
 L_1(x)=\frac{x}{h}.
\]
Với mọi $x\in[0,h]$,
\[
 L_0(x)\geq0,\qquad L_1(x)\geq0,
 \qquad|L_0(x)|+|L_1(x)|=1-\frac{x}{h}+\frac{x}{h}=1,
\]
\[
 \Lambda_1=1,\qquad
 \|\widetilde p-p\|_{\infty,[0,h]}
 \leq\max\{|\delta y_0|,|\delta y_1|\}.
\]
Với $\delta y_0=\delta y_1=\varepsilon>0$,
\[
 \widetilde p(x)-p(x)=\varepsilon(L_0(x)+L_1(x))=\varepsilon,
 \qquad
 \|\widetilde p-p\|_{\infty,[0,h]}=\varepsilon.
\]
Trái lại, các tọa độ trong cơ sở $1,x$ thỏa
\[
 a_0=y_0,\qquad a_1=\frac{y_1-y_0}{h},\qquad
 \kappa_\infty(V_h)=\frac{2(1+h)}h\longrightarrow+\infty\quad(h\downarrow0).
\]
Đổi biến $t=x/h$ cho
\[
 p(ht)=y_0+(y_1-y_0)t,\qquad
 \widetilde V=\begin{pmatrix}1&0\\1&1\end{pmatrix},\qquad
 \kappa_\infty(\widetilde V)=4.
\]
Nếu giữ cùng hai mốc nhưng xét trên đoạn $[0,1]$, thì
\[
 |L_0(x)|+|L_1(x)|=
 \begin{cases}
 1,&0\leq x\leq h,\\
 2x/h-1,&h\leq x\leq1,
 \end{cases}
 \qquad
 \Lambda_1=\frac2h-1.
\]
\end{vidu}

Mệnh đề \ref{md:vandermonde-condition} đo nhiễu của hệ số trong cơ sở đơn thức; mệnh đề \ref{md:data-sensitivity} đo nhiễu của giá trị đa thức trên đoạn đã chọn. Ví dụ trên cho $\kappa_\infty(V_h)\to\infty$ nhưng $\Lambda_1=1$ trên $[0,h]$.

% BEGIN EXTENSION NEWTON
% REVISED SLIDE EXAMPLES
\Needspace{9\baselineskip}
\section{Đa thức nội suy Newton}
Các biểu diễn và cách tổ chức bảng theo các bài giảng \cite{newton-slides,newton-slides-2022}; các kết quả được chứng minh cho những mốc đôi một phân biệt.
\Needspace{7\baselineskip}
\subsection{Trường hợp mốc bất kỳ}
Cho các mốc $x_0,\ldots,x_n$ đôi một phân biệt và $y_i=f(x_i)$. Ký hiệu $f$ ở đây chỉ cần xác định trên tập các mốc; chưa giả thiết tính liên tục.
\begin{dinhnghia}{}{newton-dd}
Tỷ sai phân được xác định bởi
\begin{equation}\label{eq:newton-dd}
 f[x_i]=y_i,\qquad
 f[x_i,\ldots,x_{i+k}]
 =\frac{f[x_{i+1},\ldots,x_{i+k}]-f[x_i,\ldots,x_{i+k-1}]}{x_{i+k}-x_i}
 \quad(k\geq1).
\end{equation}
\end{dinhnghia}
\begin{menhde}{}{newton-dd-symmetry}
Với $k\geq0$,
\begin{equation}\label{eq:newton-dd-explicit}
 f[x_0,\ldots,x_k]
 =\sum_{i=0}^{k}\frac{y_i}{\prod_{\substack{0\leq j\leq k\\j\ne i}}(x_i-x_j)}.
\end{equation}
Do đó tỷ sai phân không phụ thuộc thứ tự các cặp $(x_i,y_i)$.
\end{menhde}
\begin{chungminh}
Quy nạp theo $k$. Tại $k=0$,
\[
 f[x_0]=y_0=\frac{y_0}{\prod_{j\in\varnothing}(x_0-x_j)}.
\]
Giả sử công thức đúng ở cấp $k-1$. Đặt
\[
 A_i=\prod_{\substack{0\leq j\leq k\\j\ne i}}(x_i-x_j)\ne0.
\]
Với $1\leq i\leq k-1$, hai hệ số của $y_i$ trong hai tỷ sai phân cấp $k-1$ là
\[
 \left(\prod_{\substack{1\leq j\leq k\\j\ne i}}(x_i-x_j)\right)^{-1}
 =\frac{x_i-x_0}{A_i},\qquad
 \left(\prod_{\substack{0\leq j\leq k-1\\j\ne i}}(x_i-x_j)\right)^{-1}
 =\frac{x_i-x_k}{A_i}.
\]
Do đó \eqref{eq:newton-dd} cho
\begin{align*}
 f[x_0,\ldots,x_k]
 ={}&-\frac{y_0}{(x_k-x_0)\prod_{j=1}^{k-1}(x_0-x_j)}
 +\frac{y_k}{(x_k-x_0)\prod_{j=1}^{k-1}(x_k-x_j)}\\
 &+\sum_{i=1}^{k-1}\frac{y_i}{x_k-x_0}
       \left(\frac{x_i-x_0}{A_i}-\frac{x_i-x_k}{A_i}\right)\\
 ={}&\frac{y_0}{A_0}+\frac{y_k}{A_k}
 +\sum_{i=1}^{k-1}\frac{y_i((x_i-x_0)-(x_i-x_k))}{(x_k-x_0)A_i}
 =\sum_{i=0}^{k}\frac{y_i}{A_i}.
\end{align*}
Với mọi hoán vị $\pi$ của $\{0,\ldots,k\}$,
\[
 f[x_{\pi(0)},\ldots,x_{\pi(k)}]
 =\sum_{i=0}^{k}\frac{y_{\pi(i)}}{\prod_{j\ne i}(x_{\pi(i)}-x_{\pi(j)})}
 =\sum_{\ell=0}^{k}\frac{y_\ell}{\prod_{j\ne\ell}(x_\ell-x_j)}
 =f[x_0,\ldots,x_k].
\]
\end{chungminh}

\begin{menhde}{}{newton-dd-derivative}
Cho $k\geq1$, $f\in C^k([a,b])$ và $k+1$ mốc phân biệt trong $[a,b]$. Tồn tại $\xi$ nằm giữa mốc nhỏ nhất và mốc lớn nhất sao cho
\begin{equation}\label{eq:newton-dd-derivative}
 f[x_0,\ldots,x_k]=\frac{f^{(k)}(\xi)}{k!}.
\end{equation}
Nếu $f\in\mathcal P_r$ và $k>r$, thì $f[x_0,\ldots,x_k]=0$.
\end{menhde}
\begin{chungminh}
Đặt $A_i=\prod_{j\ne i}(x_i-x_j)$, $0\leq i\leq k$. Theo \eqref{eq:lagrange} và \eqref{eq:newton-dd-explicit},
\[
 [x^k]p_k=\sum_{i=0}^{k}y_i
 \left[x^k\right]\frac{\prod_{j\ne i}(x-x_j)}{A_i}
 =\sum_{i=0}^{k}\frac{y_i}{A_i}=f[x_0,\ldots,x_k]=d.
\]
Vì $p_k\in\mathcal P_k$,
\[
 p_k^{(k)}(x)=k!d,
 \qquad (f-p_k)(x_i)=y_i-y_i=0\quad(0\leq i\leq k).
\]
Bổ đề \ref{bd:rolle} cho
\[
 \exists\xi\in\bigl(\min_i x_i,\max_i x_i\bigr):\quad
 0=(f-p_k)^{(k)}(\xi)=f^{(k)}(\xi)-k!d,
 \qquad d=\frac{f^{(k)}(\xi)}{k!}.
\]
Nếu $f\in\mathcal P_r$, $r<k$, thì
\[
 f^{(k)}\equiv0\quad\Longrightarrow\quad
 f[x_0,\ldots,x_k]=0/k!=0.
\]
\end{chungminh}

\Needspace{7\baselineskip}
Trong cơ sở đơn thức, thêm mốc thường đòi hỏi giải lại hệ. Ta tìm một cơ sở mà phần hiệu chỉnh tại bước $k$ triệt tiêu ở mọi mốc trước đó. Tích $W_k$ có đúng tính chất này. Đặt
\[
 W_0(x)=1,\qquad W_k(x)=\prod_{j=0}^{k-1}(x-x_j),\qquad
 d_k=f[x_0,\ldots,x_k].
\]
\begin{dinhly}{}{newton-form}
Đa thức nội suy duy nhất trên các mốc $x_0,\ldots,x_n$ có biểu diễn
\begin{equation}\label{eq:newton-form}
 p_n(x)=\sum_{k=0}^{n}d_kW_k(x).
\end{equation}
Với $0\leq k\leq n$, tổng đến chỉ số $k$ là đa thức nội suy trên $x_0,\ldots,x_k$.
\end{dinhly}
\begin{chungminh}
Theo định lý \ref{dl:interpolation}, với mỗi $k$ tồn tại duy nhất $p_k\in\mathcal P_k$ thỏa $p_k(x_j)=y_j$, $0\leq j\leq k$. Đặt $H_k=p_k-p_{k-1}$, $k\geq1$. Khi đó
\[
 \deg H_k\leq k,\qquad H_k(x_j)=0\quad(0\leq j<k).
\]
Đặt $Q_0=H_k$. Giả sử $H_k=W_rQ_r$, $0\leq r<k$. Các mốc phân biệt cho
\[
 W_r(x_r)=\prod_{j=0}^{r-1}(x_r-x_j)\ne0,
 \qquad Q_r(x_r)=H_k(x_r)/W_r(x_r)=0.
\]
Bổ đề \ref{bd:roots} cho $Q_r=(x-x_r)Q_{r+1}$, nên
\[
 H_k=W_r(x)(x-x_r)Q_{r+1}(x)=W_{r+1}(x)Q_{r+1}(x).
\]
Sau $k$ bước, $H_k=W_kQ_k$, $\deg Q_k\leq0$; viết $Q_k=c_k$. Theo \eqref{eq:lagrange} và \eqref{eq:newton-dd-explicit},
\[
 c_k=[x^k]H_k=[x^k]p_k-[x^k]p_{k-1}
 =\sum_{i=0}^{k}\frac{y_i}{\prod_{j\ne i}(x_i-x_j)}-0=d_k.
\]
Vì $p_0=d_0=y_0$,
\[
 p_k=p_0+\sum_{j=1}^{k}(p_j-p_{j-1})
 =d_0+\sum_{j=1}^{k}d_jW_j
 =\sum_{j=0}^{k}d_jW_j\qquad(0\leq k\leq n).
\]
\end{chungminh}

\begin{menhde}{}{newton-append}
Thêm một mốc $x_{n+1}\notin\{x_0,\ldots,x_n\}$ và giá trị $y_{n+1}$. Khi đó
\begin{equation}\label{eq:newton-append}
 p_{n+1}(x)=p_n(x)+d_{n+1}W_{n+1}(x),\qquad
 d_{n+1}=\frac{y_{n+1}-p_n(x_{n+1})}{W_{n+1}(x_{n+1})}.
\end{equation}
Các hệ số $d_0,\ldots,d_n$ được giữ nguyên.
\end{menhde}
\begin{chungminh}
Ký hiệu hệ số của bảng mới là $\widetilde d_k$. Với $k\leq n$, \eqref{eq:newton-dd-explicit} cho
\[
 \widetilde d_k
 =\sum_{i=0}^{k}\frac{y_i}{\prod_{\substack{0\leq j\leq k\\j\ne i}}(x_i-x_j)}
 =d_k.
\]
Định lý \ref{dl:newton-form} cho
\[
 p_{n+1}(x)-p_n(x)
 =\sum_{k=0}^{n}(\widetilde d_k-d_k)W_k(x)
   +\widetilde d_{n+1}W_{n+1}(x)
 =\widetilde d_{n+1}W_{n+1}(x).
\]
Tại $x_{n+1}$,
\[
 y_{n+1}-p_n(x_{n+1})
 =\widetilde d_{n+1}\prod_{j=0}^{n}(x_{n+1}-x_j),
 \qquad\prod_{j=0}^{n}(x_{n+1}-x_j)\ne0,
\]
\[
 d_{n+1}=\widetilde d_{n+1}
 =\frac{y_{n+1}-p_n(x_{n+1})}{W_{n+1}(x_{n+1})}.
\]
\end{chungminh}

\begin{menhde}{}{newton-remainder}
Nếu $f$ còn xác định tại $x\notin\{x_0,\ldots,x_n\}$, thì
\begin{equation}\label{eq:newton-remainder}
 f(x)-p_n(x)=f[x_0,\ldots,x_n,x]W_{n+1}(x).
\end{equation}
Nếu $f\in C^{n+1}([a,b])$, các mốc và $x$ thuộc $[a,b]$, thì công thức sai số là \eqref{eq:error}.
\end{menhde}
\begin{chungminh}
Với $x\notin\{x_0,\ldots,x_n\}$, gọi $P$ là đa thức nội suy trên các mốc $x_0,\ldots,x_n,x$, với giá trị cuối $f(x)$. Mệnh đề \ref{md:newton-append} cho
\[
 P(z)-p_n(z)=f[x_0,\ldots,x_n,x]W_{n+1}(z),
 \qquad P(x)=f(x).
\]
Thay $z=x$:
\[
 f(x)-p_n(x)=f[x_0,\ldots,x_n,x]W_{n+1}(x).
\]
Nếu $f\in C^{n+1}([a,b])$, mệnh đề \ref{md:newton-dd-derivative} cho
\[
 \exists\xi_x\in(a,b):\quad
 f[x_0,\ldots,x_n,x]=\frac{f^{(n+1)}(\xi_x)}{(n+1)!},
\]
\[
 f(x)-p_n(x)=\frac{f^{(n+1)}(\xi_x)}{(n+1)!}W_{n+1}(x).
\]
Với $x=x_i$,
\[
 f(x_i)-p_n(x_i)=y_i-y_i=0,
 \qquad W_{n+1}(x_i)=\prod_{j=0}^{n}(x_i-x_j)=0;
\]
vì vậy \eqref{eq:error} đúng với mọi $\xi_x\in(a,b)$.
\end{chungminh}

\Needspace{7\baselineskip}
Đặt $D_i^{(0)}=y_i$ và
\begin{equation}\label{eq:newton-table}
 D_i^{(k)}=\frac{D_{i+1}^{(k-1)}-D_i^{(k-1)}}{x_{i+k}-x_i},
 \qquad0\leq i\leq n-k.
\end{equation}
Quy nạp trực tiếp trong \eqref{eq:newton-dd} cho
\[
 D_i^{(k)}=f[x_i,\ldots,x_{i+k}],\qquad d_k=D_0^{(k)}.
\]
\Needspace{7\baselineskip}
Với mốc cách đều, tên tiến và lùi đi kèm hai loại sai phân. Với mốc bất kỳ, hai tên này chỉ thứ tự kết nạp mốc từ đầu bảng hoặc từ cuối bảng. Đặt $D_i^{(k)}=f[x_i,\ldots,x_{i+k}]$ như trong \eqref{eq:newton-table}.
\begin{menhde}{}{newton-two-directions}
Cho $n+1$ mốc đôi một phân biệt. Đa thức nội suy có hai biểu diễn
\begin{align}
 p_n(x)&=\sum_{k=0}^{n}d_k^+\prod_{j=0}^{k-1}(x-x_j),
 &d_k^+&=D_0^{(k)},\label{eq:newton-arbitrary-forward}\\
 p_n(x)&=\sum_{k=0}^{n}d_k^-\prod_{j=0}^{k-1}(x-x_{n-j}),
 &d_k^-&=D_{n-k}^{(k)}.\label{eq:newton-arbitrary-backward}
\end{align}
\end{menhde}
\begin{chungminh}
Đặt $z_j^+=x_j$, $z_j^-=x_{n-j}$, $0\leq j\leq n$. Tính đối xứng cho
\[
 f[z_0^+,\ldots,z_k^+]=f[x_0,\ldots,x_k]=D_0^{(k)},
\]
\[
 f[z_0^-,\ldots,z_k^-]
 =f[x_n,\ldots,x_{n-k}]
 =f[x_{n-k},\ldots,x_n]=D_{n-k}^{(k)}.
\]
Áp dụng \eqref{eq:newton-form} với từng $\sigma\in\{+,-\}$:
\[
 p_n(x)=\sum_{k=0}^{n}f[z_0^\sigma,\ldots,z_k^\sigma]
                   \prod_{j=0}^{k-1}(x-z_j^\sigma).
\]
Thay $z_j^+=x_j$ hoặc $z_j^-=x_{n-j}$ được \eqref{eq:newton-arbitrary-forward} hoặc \eqref{eq:newton-arbitrary-backward}.
\end{chungminh}

\Needspace{7\baselineskip}
Giả sử thiếu giá trị $y_j$ tại một mốc $x_j$. Điều kiện nội suy tại các mốc còn lại chưa xác định được giá trị này: với mỗi $z\in\R$, định lý \ref{dl:interpolation} cho một đa thức $p_z\in\mathcal P_n$ qua toàn bộ bảng khi đặt $y_j=z$.
\begin{menhde}{}{newton-missing-data}
Cho $n\geq1$, $x_0,\ldots,x_n$ đôi một phân biệt, các giá trị $y_i$ đã biết với $i\ne j$, và giả thiết toàn bộ dữ liệu thuộc một đa thức bậc không quá $n-1$. Đặt
\[
 A_i=\prod_{\substack{0\leq\ell\leq n\\\ell\ne i}}(x_i-x_\ell)\ne0.
\]
Giá trị duy nhất tương thích với giả thiết là
\begin{equation}\label{eq:newton-missing-value}
 y_j=-A_j\sum_{i\ne j}\frac{y_i}{A_i}.
\end{equation}
Nó cũng bằng $P(x_j)$, trong đó $P\in\mathcal P_{n-1}$ nội suy $n$ cặp đã biết.
\end{menhde}
\begin{chungminh}
Với giá trị còn thiếu $y_j=z$, gọi $p_z\in\mathcal P_n$ là đa thức nội suy toàn bộ bảng. Theo \eqref{eq:newton-dd-explicit} và \eqref{eq:newton-form},
\[
 [x^n]p_z=\frac z{A_j}+\sum_{i\ne j}\frac{y_i}{A_i}.
\]
Vì $A_j\ne0$ và $\deg p_z\leq n$,
\begin{align*}
 \deg p_z\leq n-1
 &\ \Longleftrightarrow\ [x^n]p_z=0\\
 &\ \Longleftrightarrow\ \frac z{A_j}+\sum_{i\ne j}\frac{y_i}{A_i}=0\\
 &\ \Longleftrightarrow\ z=-A_j\sum_{i\ne j}\frac{y_i}{A_i}.
\end{align*}
Mỗi vế cuối xác định đúng một số $z$. Theo định lý \ref{dl:interpolation}, tồn tại duy nhất $P\in\mathcal P_{n-1}$ thỏa $P(x_i)=y_i$ với $i\ne j$. Với $z_0=P(x_j)$,
\[
 \deg P\leq n-1,\qquad
 P(x_i)=\begin{cases}y_i,&i\ne j,\\z_0,&i=j,\end{cases}
 \quad\Longrightarrow\quad p_{z_0}=P,
\]
\[
 [x^n]p_{z_0}=0
 \quad\Longrightarrow\quad
 P(x_j)=z_0=-A_j\sum_{i\ne j}y_i/A_i.
\]
\end{chungminh}
Khi giả thiết bậc là $r<n-1$, cần kiểm tra các dữ liệu đã biết cùng thuộc một đa thức của $\mathcal P_r$. Điều kiện một tỷ sai phân cấp $n$ bằng $0$ chỉ bảo đảm bậc không quá $n-1$.

\Needspace{7\baselineskip}

\begin{menhde}{}{newton-nested}
Đặt $R_n(x)=d_n$ và $R_i(x)=d_i+(x-x_i)R_{i+1}(x)$, $i=n-1,\ldots,0$. Khi đó
\begin{equation}\label{eq:newton-nested}
 R_0(x)=p_n(x).
\end{equation}
Tại $c\in\R$, đặt $v_n=d_n$, $u_n=0$ và
\begin{equation}\label{eq:newton-derivative}
 u_i=v_{i+1}+(c-x_i)u_{i+1},\qquad
 v_i=d_i+(c-x_i)v_{i+1}.
\end{equation}
Ta có $v_0=p_n(c)$, $u_0=p_n'(c)$.
\end{menhde}
\begin{chungminh}
Tại $i=n$,
\[
 R_n(x)=d_n=\sum_{k=n}^{n}d_k\prod_{j=n}^{k-1}(x-x_j).
\]
Giả sử công thức tổng đúng tại $i+1$. Khi đó
\begin{align*}
 R_i(x)
 &=d_i+(x-x_i)\sum_{k=i+1}^{n}d_k\prod_{j=i+1}^{k-1}(x-x_j)\\
 &=d_i+\sum_{k=i+1}^{n}d_k\prod_{j=i}^{k-1}(x-x_j)
 =\sum_{k=i}^{n}d_k\prod_{j=i}^{k-1}(x-x_j).
\end{align*}
Quy nạp giảm cho $R_0=\sum_{k=0}^nd_kW_k=p_n$. Đạo hàm của truy hồi là
\[
 R_i'(c)=R_{i+1}(c)+(c-x_i)R_{i+1}'(c).
\]
Đặt $e_i=v_i-R_i(c)$, $\varepsilon_i=u_i-R_i'(c)$. Từ dữ liệu đầu và hai truy hồi,
\[
 e_n=0,\qquad\varepsilon_n=0,\qquad
 e_i=(c-x_i)e_{i+1},\qquad
 \varepsilon_i=e_{i+1}+(c-x_i)\varepsilon_{i+1}.
\]
Quy nạp giảm cho $e_i=\varepsilon_i=0$, nên
\[
 v_0=R_0(c)=p_n(c),\qquad u_0=R_0'(c)=p_n'(c).
\]
\end{chungminh}

\Needspace{7\baselineskip}
Để tính các đạo hàm cấp ba, cấp bốn hoặc cao hơn, giữ nguyên dạng lồng $R_i(x)=d_i+(x-x_i)R_{i+1}(x)$ và dùng các giá trị đạo hàm đã chia cho giai thừa.
\begin{menhde}{}{newton-derivative-jet}
Cho số nguyên $s\geq0$, $c\in\R$. Đặt
\[
 J_n^{(0)}=d_n,\qquad J_n^{(r)}=0\quad(1\leq r\leq s),
\]
và với $i=n-1,\ldots,0$,
\begin{equation}\label{eq:newton-derivative-jet}
 J_i^{(0)}=d_i+(c-x_i)J_{i+1}^{(0)},\qquad
 J_i^{(r)}=(c-x_i)J_{i+1}^{(r)}+J_{i+1}^{(r-1)}\quad(1\leq r\leq s).
\end{equation}
Khi đó
\[
 p_n^{(r)}(c)=r!J_0^{(r)}\quad(0\leq r\leq s).
\]
Nếu các phép tính được thực hiện cho $Q(t)=p_n(x_c+ht)$ tại $t_c=(c-x_c)/h$, thì
\begin{equation}\label{eq:newton-jet-scaling}
 p_n^{(r)}(c)=h^{-r}Q^{(r)}(t_c)=h^{-r}r!J_0^{(r)}.
\end{equation}
\end{menhde}
\begin{chungminh}
Với $r\geq1$, công thức Leibniz cho
\begin{align*}
 R_i^{(r)}(x)
 &=\sum_{k=0}^{r}\binom rk (x-x_i)^{(k)}R_{i+1}^{(r-k)}(x)\\
 &=(x-x_i)R_{i+1}^{(r)}(x)+rR_{i+1}^{(r-1)}(x),
\end{align*}
vì $(x-x_i)'=1$ và $(x-x_i)^{(k)}=0$ với $k\geq2$. Do đó
\[
 \frac{R_i^{(r)}(c)}{r!}
 =(c-x_i)\frac{R_{i+1}^{(r)}(c)}{r!}
 +\frac{R_{i+1}^{(r-1)}(c)}{(r-1)!}.
\]
Ở cấp $r=0$, $R_i(c)=d_i+(c-x_i)R_{i+1}(c)$. Vì
\[
 R_n^{(0)}(c)=d_n,\qquad R_n^{(r)}(c)=0\ (r\geq1),
\]
quy nạp giảm theo $i$ trong các truy hồi cho
\[
 J_i^{(r)}=R_i^{(r)}(c)/r!,\qquad
 r!J_0^{(r)}=R_0^{(r)}(c)=p_n^{(r)}(c).
\]
Với $Q(t)=p_n(x_c+ht)$, $h\ne0$,
\[
 Q'(t)=h p_n'(x_c+ht),
\]
\[
 Q^{(r)}(t)=h^rp_n^{(r)}(x_c+ht)
 \quad\Longrightarrow\quad
 Q^{(r+1)}(t)=h^{r+1}p_n^{(r+1)}(x_c+ht).
\]
Suy ra $p_n^{(r)}(c)=h^{-r}Q^{(r)}((c-x_c)/h)$.
\end{chungminh}


\Needspace{9\baselineskip}
\begin{vidu}
Cho dữ liệu $(x_i,y_i)=(-1,6),(0,3),(2,9)$. Bảng tỷ sai phân là
\begin{center}
\begin{tabular}{c|rrr}\toprule
$x_i$&$D_i^{(0)}$&$D_i^{(1)}$&$D_i^{(2)}$\\\midrule
$-1$&$6$&$-3$&$2$\\
$0$&$3$&$3$&\\
$2$&$9$&&\\\bottomrule
\end{tabular}
\end{center}
Do đó $p_2(x)=6+(x+1)(-3+2x)$. Tại $c=1$, truy hồi \eqref{eq:newton-derivative} cho bảng ngang
\begin{center}
\begin{tabular}{c|rrr}\toprule
$i$&$2$&$1$&$0$\\\midrule
$d_i$&$2$&$-3$&$6$\\
$c-x_i$& & $1$&$2$\\
$v_i$&$2$&$-1$&$4$\\
$u_i$&$0$&$2$&$3$\\\bottomrule
\end{tabular}
\end{center}
Vậy $p_2(1)=4$, $p_2'(1)=3$. Chuyển sang hệ số đơn thức bằng Horner ngược: $R_2=2$, $R_1=-3+xR_2$, $R_0=6+(x+1)R_1$.
\begin{center}
\begin{tabular}{c|rrr}\toprule
Đa thức&$a_0$&$a_1$&$a_2$\\\midrule
$R_2$&$2$&&\\
$R_1$&$-3$&$2$&\\
$R_0$&$3$&$-1$&$2$\\\bottomrule
\end{tabular}
\end{center}
\end{vidu}

\Needspace{7\baselineskip}
Sau khi có các tỷ sai phân $d_k$, phần còn lại trong \eqref{eq:newton-form} là tổ hợp tuyến tính của các tích $W_k$. Nếu cần hệ số đơn thức cho nhiều bộ giá trị trên cùng tập mốc, có thể dựng các tích này một lần thành bảng nhân.
Đặt $B_{j,k}$ là hệ số của $x^j$ trong $W_k(x)$, với quy ước $B_{j,k}=0$ khi $j<0$ hoặc $j>k$.
\begin{menhde}{}{newton-product-table}
Bảng nhân $B=(B_{j,k})_{0\leq j,k\leq n}$ được dựng bởi
\begin{equation}\label{eq:newton-product-table}
 B_{0,0}=1,\qquad
 B_{j,k+1}=B_{j-1,k}-x_kB_{j,k}\quad(0\leq j\leq k+1).
\end{equation}
Nếu $\boldsymbol d=(d_0,\ldots,d_n)^T$ và $\boldsymbol a$ là vector hệ số đơn thức của $p_n$, thì
\begin{equation}\label{eq:newton-matrix-product}
 \boldsymbol a=B\boldsymbol d,\qquad
 a_j=\sum_{k=j}^{n}B_{j,k}d_k.
\end{equation}
Ma trận $B$ tam giác trên, có đường chéo bằng $1$ và khả nghịch. Với ma trận Vandermonde $V$ trong \eqref{eq:vandermonde-system}, đặt $T=VB$. Khi đó
\[
 T_{i,k}=W_k(x_i),\quad T_{i,k}=0\ (k>i),\quad
 T_{i,i}=\prod_{\nu=0}^{i-1}(x_i-x_\nu)\ne0,\qquad
 T\boldsymbol d=\boldsymbol y.
\]
\end{menhde}
\begin{chungminh}
Từ $W_0=1$ suy ra $B_{0,0}=1$. Với quy ước $B_{j,k}=0$ khi $j<0$ hoặc $j>k$,
\begin{align*}
 W_{k+1}(x)
 &=(x-x_k)\sum_{j=0}^{k}B_{j,k}x^j\\
 &=\sum_{j=1}^{k+1}B_{j-1,k}x^j-x_k\sum_{j=0}^{k}B_{j,k}x^j\\
 &=\sum_{j=0}^{k+1}(B_{j-1,k}-x_kB_{j,k})x^j,
\end{align*}
\[
 B_{j,k+1}=[x^j]W_{k+1}=B_{j-1,k}-x_kB_{j,k}.
\]
Mặt khác,
\begin{align*}
 p_n(x)
 &=\sum_{k=0}^{n}d_k\sum_{j=0}^{k}B_{j,k}x^j\\
 &=\sum_{j=0}^{n}\left(\sum_{k=j}^{n}B_{j,k}d_k\right)x^j,
\end{align*}
\[
 a_j=\sum_{k=j}^{n}B_{j,k}d_k,
 \qquad\boldsymbol a=B\boldsymbol d.
\]
Truy hồi đường chéo cho $B_{k+1,k+1}=B_{k,k}$; bởi vậy
\[
 B_{j,k}=0\ (j>k),\qquad B_{k,k}=B_{0,0}=1,
 \qquad\det B=\prod_{k=0}^{n}B_{k,k}=1\ne0.
\]
Với $T=VB$,
\[
 T_{i,k}=\sum_{j=0}^{n}x_i^jB_{j,k}=W_k(x_i),
\]
\[
 k>i\ \Longrightarrow\ W_k(x_i)=\prod_{j=0}^{k-1}(x_i-x_j)=0,
 \qquad T_{i,i}=\prod_{j=0}^{i-1}(x_i-x_j)\ne0,
\]
\[
 T\boldsymbol d=VB\boldsymbol d=V\boldsymbol a=\boldsymbol y.
\]
\end{chungminh}
Hệ tam giác này cũng biểu thị điều kiện nội suy tại từng mốc:
\[
 d_0=y_0,\qquad
 d_i=\frac{y_i-\sum_{k=0}^{i-1}d_kW_k(x_i)}{W_i(x_i)}.
\]


\Needspace{9\baselineskip}
\begin{vidu}\label{vidu:newton-slide-six}
Xét bảng sáu mốc trong bài giảng Newton \cite{newton-slides}. Các số thập phân ở hai cột đầu là dữ liệu chính xác. Để đọc cả hai chiều, đặt $\widehat D_i^{(k)}=D_{i-k}^{(k)}$ với $k\leq i$. Bảng tỷ sai phân được trình bày theo dạng tam giác dưới; các tỷ sai phân có biểu diễn thập phân vô hạn được làm tròn ở đây.
\begin{center}
\small\setlength{\tabcolsep}{4pt}
\begin{tabular}{c|r|rrrrr}\toprule
$x_i$&$y_i$&$\widehat D_i^{(1)}$&$\widehat D_i^{(2)}$&$\widehat D_i^{(3)}$&$\widehat D_i^{(4)}$&$\widehat D_i^{(5)}$\\\midrule
$1.2$&$3.172$&&&&&\\
$1.3$&$3.695$&$5.23$&&&&\\
$1.5$&$3.974$&$1.395$&$-12.78333333$&&&\\
$1.7$&$4.152$&$0.89$&$-1.2625$&$23.04166667$&&\\
$1.8$&$3.769$&$-3.83$&$-15.73333333$&$-28.94166667$&$-86.63888889$&\\
$2$&$3.112$&$-3.285$&$1.81666667$&$35.1$&$91.48809524$&$222.65873016$\\\bottomrule
\end{tabular}
\end{center}
Các hệ số Newton theo thứ tự $x_0,x_1,\ldots,x_5$ nằm trên đường chéo $\widehat D_k^{(k)}$; theo thứ tự $x_5,x_4,\ldots,x_0$ chúng nằm trên hàng cuối:
\begin{center}
\small\setlength{\tabcolsep}{4pt}
\begin{tabular}{c|rrrrrr}\toprule
$k$&$0$&$1$&$2$&$3$&$4$&$5$\\\midrule
$d_k^+$&$3.172$&$5.23$&$-12.78333333$&$23.04166667$&$-86.63888889$&$222.65873016$\\
$d_k^-$&$3.112$&$-3.285$&$1.81666667$&$35.1$&$91.48809524$&$222.65873016$\\\bottomrule
\end{tabular}
\end{center}
Dựng $W_{k+1}=(x-x_k)W_k$ bằng Horner ngược. Khác với bảng tỷ sai phân đã làm tròn, mọi hệ số trong bảng nhân sau là chính xác:
\begin{center}
\small\setlength{\tabcolsep}{5pt}
\begin{tabular}{c|rrrrrr}\toprule
Tích cơ sở&$x^5$&$x^4$&$x^3$&$x^2$&$x$&$1$\\\midrule
$W_0$&$0$&$0$&$0$&$0$&$0$&$1$\\
$W_1$&$0$&$0$&$0$&$0$&$1$&$-1.2$\\
$W_2$&$0$&$0$&$0$&$1$&$-2.5$&$1.56$\\
$W_3$&$0$&$0$&$1$&$-4$&$5.31$&$-2.34$\\
$W_4$&$0$&$1$&$-5.7$&$12.11$&$-11.367$&$3.978$\\
$W_5$&$1$&$-7.5$&$22.37$&$-33.165$&$24.4386$&$-7.1604$\\\bottomrule
\end{tabular}
\end{center}
Nhân từng hàng với tỷ sai phân chính xác tương ứng rồi cộng theo cột, tức $\boldsymbol a=B\boldsymbol d^+$, được
\begin{align*}
p_5(x)={}&\frac{28055}{126}x^5-\frac{221329}{126}x^4
 +\frac{13854353}{2520}x^3\\
 &-\frac{107586659}{12600}x^2+\frac{138302461}{21000}x
 -\frac{1411157}{700}.
\end{align*}
Tại $c=1.6$, dùng \eqref{eq:newton-derivative-jet} cho $J_{i,s}=R_i^{(s)}(c)/s!$. Bảng ngang ghi từng bước lồng; các kết quả được tính bằng số hữu tỷ trước khi làm tròn để trình bày.
\begin{center}
\scriptsize\setlength{\tabcolsep}{4pt}
\begin{tabular}{c|rrrrrr}\toprule
$i$&$5$&$4$&$3$&$2$&$1$&$0$\\\midrule
$c-x_i$&&$-0.2$&$-0.1$&$0.1$&$0.3$&$0.4$\\
$J_{i,0}$&$222.6587302$&$-131.1706349$&$36.1587302$&$-9.1674603$&$2.4797619$&$4.1639048$\\
$J_{i,1}$&$0$&$222.6587302$&$-153.4365079$&$20.8150794$&$-2.9229365$&$1.3105873$\\
$J_{i,2}$&$0$&$0$&$222.6587302$&$-131.1706349$&$-18.5361111$&$-10.3373810$\\
$J_{i,3}$&$0$&$0$&$0$&$222.6587302$&$-64.3730159$&$-44.2853175$\\
$J_{i,4}$&$0$&$0$&$0$&$0$&$222.6587302$&$24.6904762$\\
$J_{i,5}$&$0$&$0$&$0$&$0$&$0$&$222.6587302$\\\bottomrule
\end{tabular}
\end{center}
Do $p_5^{(s)}(c)=s!J_{0,s}$, kết quả là
\begin{center}
\small\setlength{\tabcolsep}{4pt}
\begin{tabular}{c|rrrrrr}\toprule
$s$&$0$&$1$&$2$&$3$&$4$&$5$\\\midrule
$p_5^{(s)}(1.6)$&$4.1639048$&$1.3105873$&$-20.6747619$&$-265.7119048$&$592.5714286$&$26719.0476190$\\\bottomrule
\end{tabular}
\end{center}
Đây là giá trị và các đạo hàm của đa thức nội suy; sai số so với một hàm chưa biết cần thêm giả thiết về hàm ấy.
\end{vidu}


\Needspace{7\baselineskip}
\subsection{Trường hợp mốc cách đều}
\begin{dinhnghia}{}{generalized-binomial}
Với $t\in\R$ và số nguyên $k\geq0$, hệ số nhị thức mở rộng được định nghĩa bởi
\begin{equation}\label{eq:generalized-binomial}
 \binom t0=1,\qquad
 \binom tk=\frac{t(t-1)\cdots(t-k+1)}{k!}\quad(k\geq1).
\end{equation}
Ký hiệu $C_t^k$ cũng được dùng cho $\binom tk$. Quy ước $\binom tk=0$ với $k$ nguyên âm. Khi $t=n\geq0$ là số nguyên, định nghĩa cho
\[
 \binom nk=\frac{n!}{k!(n-k)!}\quad(0\leq k\leq n),\qquad
 \binom nk=0\quad(k>n).
\]
\end{dinhnghia}

\begin{menhde}{}{generalized-binomial-identities}
Với $t\in\R$, $k\geq1$ nguyên,
\begin{align}
 \binom{t+1}k-\binom tk&=\binom t{k-1},\label{eq:generalized-binomial-pascal}\\
 k\binom tk&=(t-k+1)\binom t{k-1}
 =t\binom{t-1}{k-1},\label{eq:generalized-binomial-recurrence}\\
 \binom{-t}k&=(-1)^k\binom{t+k-1}k.\label{eq:generalized-binomial-negative}
\end{align}
Đặt $\delta_tu(t)=u(t+1)-u(t)$. Khi $r\geq0$ nguyên,
\begin{equation}\label{eq:generalized-binomial-difference}
 \delta_t^r\binom tk=
 \begin{cases}\binom t{k-r},&0\leq r\leq k,\\0,&r>k.\end{cases}
\end{equation}
\end{menhde}
\begin{chungminh}
Với $k\geq1$, tách tích chung cho
\begin{align*}
 k!\left(\binom{t+1}k-\binom tk\right)
 &=\bigl[(t+1)-(t-k+1)\bigr]\prod_{j=0}^{k-2}(t-j)\\
 &=k\prod_{j=0}^{k-2}(t-j)=k!\binom t{k-1}.
\end{align*}
Tách lần lượt nhân tử cuối và nhân tử đầu của tích cho
\[
 k\binom tk
 =\frac{\prod_{j=0}^{k-1}(t-j)}{(k-1)!}
 =(t-k+1)\binom t{k-1}
 =t\binom{t-1}{k-1}.
\]
Đổi chỉ số $\ell=k-1-j$ cho
\[
 \binom{-t}k
 =\frac{\prod_{j=0}^{k-1}(-t-j)}{k!}
 =\frac{(-1)^k\prod_{j=0}^{k-1}(t+j)}{k!}
 =(-1)^k\binom{t+k-1}k.
\]
Công thức sai phân đúng với $r=0$. Với $0\leq r<k$, bước quy nạp là
\[
 \delta_t^{r+1}\binom tk
 =\binom{t+1}{k-r}-\binom t{k-r}
 =\binom t{k-r-1}.
\]
Tại $r=k$, kết quả bằng $\binom t0=1$; do đó $\delta_t^{k+1}\binom tk=1-1=0$ và mọi sai phân cấp cao hơn đều bằng $0$.
\end{chungminh}

Cho $h\ne0$ và $x_i=x_0+ih$. Đặt
\[
 \Delta^0y_i=y_i,\quad\Delta^{k+1}y_i=\Delta^ky_{i+1}-\Delta^ky_i,
 \qquad
 \nabla^0y_i=y_i,\quad\nabla^{k+1}y_i=\nabla^ky_i-\nabla^ky_{i-1}.
\]
Chỉ dùng các sai phân mà toàn bộ chỉ số dữ liệu thuộc tập mốc đã cho.
\begin{menhde}{}{newton-equal-dd}
Với các chỉ số thích hợp,
\begin{equation}\label{eq:newton-equal-dd}
 \Delta^ky_i=\sum_{j=0}^{k}(-1)^{k-j}\binom kj y_{i+j},\qquad
 \nabla^ky_{i+k}=\Delta^ky_i,
 \qquad
 f[x_i,\ldots,x_{i+k}]=\frac{\Delta^ky_i}{k!h^k}.
\end{equation}
\end{menhde}
\begin{chungminh}
Quy nạp theo $k$. Tại $k=0$,
\[
 \Delta^0y_i=y_i=\sum_{j=0}^{0}(-1)^{0-j}\binom0j y_{i+j}.
\]
Quy ước $\binom kj=0$ nếu $j\notin\{0,\ldots,k\}$. Từ công thức ở cấp $k$,
\begin{align*}
 \Delta^{k+1}y_i
 &=\sum_{j=0}^{k}(-1)^{k-j}\binom kj y_{i+1+j}
   -\sum_{j=0}^{k}(-1)^{k-j}\binom kj y_{i+j}\\
 &=\sum_{j=0}^{k+1}(-1)^{k+1-j}
   \left(\binom k{j-1}+\binom kj\right)y_{i+j}\\
 &=\sum_{j=0}^{k+1}(-1)^{k+1-j}\binom{k+1}j y_{i+j}.
\end{align*}
Đẳng thức lùi đúng tại $k=0$. Nếu đúng ở cấp $k$, thì
\begin{align*}
 \nabla^{k+1}y_{i+k+1}
 &=\nabla^ky_{i+k+1}-\nabla^ky_{i+k}\\
 &=\Delta^ky_{i+1}-\Delta^ky_i=\Delta^{k+1}y_i.
\end{align*}
Tại $k=0$, $f[x_i]=y_i=\Delta^0y_i/(0!h^0)$. Giả sử công thức tỷ sai phân đúng ở cấp $k-1$, $k\geq1$. Vì $x_{i+k}-x_i=kh\ne0$,
\begin{align*}
 f[x_i,\ldots,x_{i+k}]
 &=\frac{f[x_{i+1},\ldots,x_{i+k}]-f[x_i,\ldots,x_{i+k-1}]}{x_{i+k}-x_i}\\
 &=\frac{\Delta^{k-1}y_{i+1}-\Delta^{k-1}y_i}{(k-1)!h^{k-1}kh}
 =\frac{\Delta^ky_i}{k!h^k}.
\end{align*}
\end{chungminh}

\Needspace{7\baselineskip}
Bảng tỷ sai phân cũ chứa $D_i^{(k)}$ với $0\leq k\leq n$, $0\leq i\leq n-k$. Sau khi thêm cặp $(x_{n+1},y_{n+1})$, chỉ có một phần tử mới ở mỗi cột.
\begin{menhde}{}{newton-table-update}
Giả sử $x_{n+1}$ khác mọi mốc cũ. Đặt $E_0=y_{n+1}$ và
\begin{equation}\label{eq:newton-dd-update}
 E_k=\frac{E_{k-1}-D_{n+1-k}^{(k-1)}}{x_{n+1}-x_{n+1-k}},
 \qquad1\leq k\leq n+1.
\end{equation}
Khi đó $E_k=f[x_{n+1-k},\ldots,x_{n+1}]$; bảng mới bổ sung đúng $D_{n+1-k}^{(k)}=E_k$ và giữ nguyên các phần tử cũ. Hệ số Newton tiến mới là $d_{n+1}=E_{n+1}$.

Nếu $x_i=x_0+ih$, $h\ne0$, đặt $F_0=y_{n+1}$ và
\begin{equation}\label{eq:newton-difference-update}
 F_k=F_{k-1}-\Delta^{k-1}y_{n+1-k},\qquad1\leq k\leq n+1.
\end{equation}
Khi đó $F_k=\Delta^ky_{n+1-k}=\nabla^ky_{n+1}$.
\end{menhde}
\begin{chungminh}
Tập chỉ số của bảng cũ và bảng mới là
\[
 \mathcal T_n=\{(i,k):0\leq k\leq n,\ 0\leq i\leq n-k\},
\]
\[
 \mathcal T_{n+1}\setminus\mathcal T_n
 =\{(n+1-k,k):0\leq k\leq n+1\}.
\]
Với $(i,k)\in\mathcal T_n$, công thức \eqref{eq:newton-dd-explicit} chỉ chứa các cặp có chỉ số $i,\ldots,i+k\leq n$, nên
\[
 (D_i^{(k)})_{\text{mới}}
 =\sum_{j=i}^{i+k}\frac{y_j}{\prod_{\substack{i\leq\ell\leq i+k\\\ell\ne j}}(x_j-x_\ell)}
 =(D_i^{(k)})_{\text{cũ}}.
\]
Tại $k=0$, $E_0=y_{n+1}=f[x_{n+1}]$. Giả sử
$E_{k-1}=f[x_{n+2-k},\ldots,x_{n+1}]$. Vì $x_{n+1}\ne x_{n+1-k}$,
\begin{align*}
 E_k
 &=\frac{f[x_{n+2-k},\ldots,x_{n+1}]-f[x_{n+1-k},\ldots,x_n]}
 {x_{n+1}-x_{n+1-k}}\\
 &=f[x_{n+1-k},\ldots,x_{n+1}]
 =(D_{n+1-k}^{(k)})_{\text{mới}}.
\end{align*}
Do đó
\[
 E_{n+1}=f[x_0,\ldots,x_{n+1}]=d_{n+1}.
\]
Với sai phân, $F_0=y_{n+1}=\Delta^0y_{n+1}$. Nếu
$F_{k-1}=\Delta^{k-1}y_{n+2-k}$, thì
\begin{align*}
 F_k
 &=\Delta^{k-1}y_{n+2-k}-\Delta^{k-1}y_{n+1-k}
 =\Delta^ky_{n+1-k}
 =\nabla^ky_{n+1},
\end{align*}
trong đó đẳng thức cuối là \eqref{eq:newton-equal-dd} với chỉ số đầu $n+1-k$.
\end{chungminh}


\begin{dinhly}{}{newton-equal-form}
Đặt $t=(x-x_0)/h$ và $s=(x-x_n)/h$. Đa thức nội suy thỏa
\begin{align}
 p_n(x_0+ht)&=\sum_{k=0}^{n}\binom tk\Delta^ky_0,\label{eq:newton-forward}\\
 p_n(x_n+hs)&=\sum_{k=0}^{n}\binom{s+k-1}k\nabla^ky_n
 =\sum_{k=0}^{n}(-1)^k\binom{-s}k\nabla^ky_n.\label{eq:newton-backward}
\end{align}
\end{dinhly}
\begin{chungminh}
Trong thứ tự mốc $x_0,x_1,\ldots,x_n$, \eqref{eq:newton-equal-dd} và định nghĩa hệ số nhị thức mở rộng cho
\[
 d_kW_k(x_0+ht)
 =\frac{\Delta^ky_0}{k!h^k}\prod_{j=0}^{k-1}h(t-j)
 =\binom tk\Delta^ky_0.
\]
Thế vào \eqref{eq:newton-form} được \eqref{eq:newton-forward}. Trong thứ tự $x_n,x_{n-1},\ldots,x_0$, tính đối xứng cho
\[
 f[x_n,\ldots,x_{n-k}]=f[x_{n-k},\ldots,x_n]
 =\frac{\Delta^ky_{n-k}}{k!h^k}=\frac{\nabla^ky_n}{k!h^k}.
\]
Vì
\[
 \frac1{k!h^k}\prod_{j=0}^{k-1}(x_n+hs-x_{n-j})
 =\frac1{k!}\prod_{j=0}^{k-1}(s+j)
 =\binom{s+k-1}k=(-1)^k\binom{-s}k,
\]
\eqref{eq:newton-form} theo thứ tự này cho \eqref{eq:newton-backward}, kể cả $k=0$ theo quy ước tích rỗng.
\end{chungminh}

\Needspace{9\baselineskip}
\begin{vidu}
Một bảng vị trí giả lập, đơn vị mét, được lấy mỗi giây:
\begin{center}
\begin{tabular}{c|rrrr}\toprule
$i$&$0$&$1$&$2$&$3$\\\midrule
$x_i$ (s)&$0$&$1$&$2$&$3$\\
$y_i$ (m)&$0$&$1$&$4$&$9$\\
$\Delta y_i$&$1$&$3$&$5$&\\
$\Delta^2y_i$&$2$&$2$&&\\
$\Delta^3y_i$&$0$&&&\\\bottomrule
\end{tabular}
\end{center}
Với $c=3/2$, $h=1$, hệ số Newton là $(0,1,1,0)^T$. Tính theo dạng lồng:
\begin{center}
\begin{tabular}{c|rrrr}\toprule
$i$&$3$&$2$&$1$&$0$\\\midrule
$d_i$&$0$&$1$&$1$&$0$\\
$c-x_i$&&$-1/2$&$1/2$&$3/2$\\
$v_i$&$0$&$1$&$3/2$&$9/4$\\\bottomrule
\end{tabular}
\end{center}
Đa thức nội suy cho vị trí $9/4$ m tại thời điểm $3/2$ s. Nếu dữ liệu chỉ là phép đo, đẳng thức này nói về đa thức nội suy; sai lệch với vị trí thực cần giả thiết và chặn đạo hàm trong \eqref{eq:uniform-error}.
\end{vidu}

\Needspace{7\baselineskip}
Khi cần tính gần $x_k$, không bắt buộc lấy toàn bộ bảng. Chọn một số nguyên $r$ và dùng các mốc liền nhau về phía phải hoặc phía trái của $x_k$. Đặt $t=(x-x_k)/h$.
\begin{menhde}{}{newton-from-k}
Với $0\leq r\leq n-k$, đa thức nội suy trên $x_k,\ldots,x_{k+r}$ là
\begin{equation}\label{eq:newton-forward-k}
 P_{k,r}^{+}(x_k+ht)=\sum_{j=0}^{r}\binom tj\Delta^jy_k.
\end{equation}
Với $0\leq r\leq k$, đa thức nội suy trên $x_{k-r},\ldots,x_k$ là
\begin{equation}\label{eq:newton-backward-k}
 P_{k,r}^{-}(x_k+ht)=\sum_{j=0}^{r}\binom{t+j-1}j\nabla^jy_k
 =\sum_{j=0}^{r}(-1)^j\binom{-t}j\nabla^jy_k.
\end{equation}
Đặt $q_{k,r}^{\pm}(t)=P_{k,r}^{\pm}(x_k+ht)$. Với $s\geq0$ nguyên,
\[
 (P_{k,r}^{\pm})^{(s)}(x_k+ht)=h^{-s}(q_{k,r}^{\pm})^{(s)}(t).
\]
\end{menhde}
\begin{chungminh}
Cho thứ tự mốc $x_k,x_{k+1},\ldots,x_{k+r}$. Hệ số Newton cấp $j$ và tích cơ sở tương ứng thỏa
\begin{align*}
 f[x_k,\ldots,x_{k+j}]&=\frac{\Delta^jy_k}{j!h^j},\\
 \prod_{\nu=0}^{j-1}(x_k+ht-x_{k+\nu})
 &=h^j\prod_{\nu=0}^{j-1}(t-\nu)=h^jj!\binom tj.
\end{align*}
Thế vào \eqref{eq:newton-form} cho \eqref{eq:newton-forward-k}. Với thứ tự $x_k,x_{k-1},\ldots,x_{k-r}$,
\begin{align*}
 f[x_k,\ldots,x_{k-j}]&=\frac{\nabla^jy_k}{j!h^j},\\
 \prod_{\nu=0}^{j-1}(x_k+ht-x_{k-\nu})
 &=h^j\prod_{\nu=0}^{j-1}(t+\nu)=h^jj!\binom{t+j-1}j.
\end{align*}
Thế vào \eqref{eq:newton-form} và dùng \eqref{eq:generalized-binomial-negative} cho \eqref{eq:newton-backward-k}. Với đạo hàm, $q_{k,r}^{\pm}=P_{k,r}^{\pm}\circ(t\mapsto x_k+ht)$, nên quy nạp theo $s$ cho
\[
 (q_{k,r}^{\pm})^{(s+1)}(t)
 =\frac d{dt}\left[h^s(P_{k,r}^{\pm})^{(s)}(x_k+ht)\right]
 =h^{s+1}(P_{k,r}^{\pm})^{(s+1)}(x_k+ht).
\]
Chia cho $h^s\ne0$ được công thức đã nêu.
\end{chungminh}
Với mốc cách đều, có thể dựng bảng nhân trong biến $t$: Newton tiến dùng các nhân tử $t,t-1,\ldots$, Newton lùi dùng $t,t+1,\ldots$. Vector nhân với bảng lần lượt là $(\Delta^jy_k/j!)_{j=0}^{r}$ và $(\nabla^jy_k/j!)_{j=0}^{r}$. Hệ số thu được là hệ số theo $t$; khi tính đạo hàm theo $x$, phải nhân thêm $h^{-s}$.

Đa thức $P_{k,r}^{\pm}$ nội suy đúng tập mốc đã chọn. Nó không phải biểu diễn của đa thức trên toàn bộ $n+1$ mốc khi tập mốc được chọn nhỏ hơn, trừ trường hợp các dữ liệu còn lại cũng nằm trên cùng đa thức này.

\Needspace{7\baselineskip}
Cho bảng $x_i=a+ih$, $0\leq i<M$, $h>0$. Muốn lấy đúng $k$ mốc, $1\leq k\leq M$, xét độ xa hình học $|x_i-c|$. Khi bằng nhau, ưu tiên chỉ số nhỏ hơn.
\begin{menhde}{}{newton-nearest-window}
Đặt $q=(c-a)/h$ và
\begin{equation}\label{eq:newton-nearest-window}
 s=\min\bigl\{M-k,\max\{0,\lceil q-k/2\rceil\}\bigr\}.
\end{equation}
Các chỉ số $s,s+1,\ldots,s+k-1$ là $k$ chỉ số gần $c$ nhất theo quy tắc đã nêu. Với mọi tập $J\subset\{0,\ldots,M-1\}$ có $k$ phần tử,
\[
 \max_{s\leq i<s+k}|c-x_i|\leq\max_{i\in J}|c-x_i|,
 \qquad
 \prod_{i=s}^{s+k-1}|c-x_i|\leq\prod_{i\in J}|c-x_i|.
\]
\end{menhde}
\begin{chungminh}
Đặt $I_s=\{s,\ldots,s+k-1\}$. Nếu $k=M$, $s=0$ và mọi tập $J$ có $k$ phần tử đều bằng $I_s$; hai bất đẳng thức trở thành đẳng thức. Xét $k<M$. Từ định nghĩa $s$,
\begin{align*}
 s>0&\ \Longrightarrow\ \left\lceil q-\frac k2\right\rceil\geq s
 \ \Longrightarrow\ q>s+\frac{k-2}{2},\\
 s<M-k&\ \Longrightarrow\ \left\lceil q-\frac k2\right\rceil\leq s
 \ \Longrightarrow\ q\leq s+\frac k2.
\end{align*}
Đặt
\[
 R=\max\{|s-q|,|s+k-1-q|\}
 =\max\{q-s,s+k-1-q\}.
\]
Với $i\in I_s$, $s\leq i\leq s+k-1$, nên
\[
 |i-q|=\max\{i-q,q-i\}
 \leq\max\{s+k-1-q,q-s\}=R.
\]
Nếu $s>0$, ta có
\begin{align*}
 |s-1-q|^2-|s-q|^2&=2(q-s)+1>k-1\geq0,\\
 |s-1-q|^2-|s+k-1-q|^2&=k(2q-2s-k+2)>0.
\end{align*}
Do $q>s+(k-2)/2\geq s-1/2>s-1$, với $j<s$,
\[
 |j-q|=q-j\geq q-s+1=|s-1-q|>R.
\]
Nếu $s<M-k$, ta có
\begin{align*}
 |s+k-q|^2-|s-q|^2&=k(2s+k-2q)\geq0,\\
 |s+k-q|^2-|s+k-1-q|^2&=2(s+k-q)-1\geq k-1\geq0.
\end{align*}
Do $q\leq s+k/2<s+k$, với $j\geq s+k$,
\[
 |j-q|=j-q\geq s+k-q=|s+k-q|\geq R.
\]
Suy ra
\[
 i\in I_s,\ j\notin I_s\quad\Longrightarrow\quad |i-q|\leq|j-q|.
\]
Ở phía trái bất đẳng thức nghiêm ngặt; ở phía phải $i<j$. Như vậy cửa sổ thỏa cả quy tắc chọn khoảng cách và quy tắc ưu tiên chỉ số khi bằng nhau.

Với một tập $J$ có $k$ phần tử, đặt
\[
 I_s\setminus J=\{a_1<\cdots<a_r\},\qquad
 J\setminus I_s=\{b_1<\cdots<b_r\};
 \quad r=k-|I_s\cap J|.
\]
Ánh xạ
\[
 T(i)=\begin{cases}i,&i\in I_s\cap J,\\b_\nu,&i=a_\nu,\quad1\leq\nu\leq r\end{cases}
\]
thỏa $T(I_s)=J$ và là song ánh. Các so sánh trên cho $|i-q|\leq|T(i)-q|$ với mọi $i\in I_s$. Vì $|c-x_i|=h|q-i|$ và $h>0$,
\begin{align*}
 \max_{i\in I_s}|c-x_i|
 &=h\max_{i\in I_s}|i-q|
 \leq h\max_{i\in I_s}|T(i)-q|
 =\max_{j\in J}|c-x_j|,\\
 \prod_{i\in I_s}|c-x_i|
 &=h^k\prod_{i\in I_s}|i-q|
 \leq h^k\prod_{i\in I_s}|T(i)-q|
 =\prod_{j\in J}|c-x_j|.
\end{align*}
\end{chungminh}
Đây là tối ưu của các khoảng cách và tích nhân tử hình học tại $c$. Muốn suy ra chặn sai số hàm, cần một chặn đạo hàm hợp lệ; muốn suy ra chặn sai số đạo hàm, cần một lập luận sai số đạo hàm riêng. Số mốc được chọn không được suy ra chỉ từ độ gần.

\Needspace{9\baselineskip}

\begin{vidu}\label{vidu:newton-100}
Để tách sai số xấp xỉ khỏi sai số phép đo, xét một bộ dữ liệu mô phỏng đường đáp ứng chuẩn hóa
\[
 f(x)=\frac1{1+x^2},\qquad x_i=-5+\frac{10i}{99},\qquad
 y_i=\frac{9801}{9801+25(2i-99)^2},\quad0\leq i\leq99.
\]
Các $y_i$ được định nghĩa bằng số hữu tỷ chính xác. Bảng dưới chỉ làm tròn để hiển thị; không dùng các số đã làm tròn này để dựng đa thức.
Bộ 100 giá trị là:
\begin{center}
\small\setlength{\tabcolsep}{4pt}
\begin{tabular}{rr|rr|rr|rr|rr}\toprule
$i$&$y_i$&$i$&$y_i$&$i$&$y_i$&$i$&$y_i$&$i$&$y_i$\\\midrule
0&0.038462&20&0.101223&40&0.520610&60&0.470614&80&0.095316\\
1&0.040000&21&0.107673&41&0.575649&61&0.425649&81&0.089896\\
2&0.041631&22&0.114731&42&0.635356&62&0.385472&82&0.084912\\
3&0.043362&23&0.122473&43&0.698774&63&0.349711&83&0.080319\\
4&0.045202&24&0.130984&44&0.764151&64&0.317946&84&0.076079\\
5&0.047160&25&0.140363&45&0.828767&65&0.289748&85&0.072158\\
6&0.049245&26&0.150724&46&0.888899&66&0.264706&86&0.068526\\
7&0.051469&27&0.162198&47&0.940054&67&0.242443&87&0.065155\\
8&0.053844&28&0.174937&48&0.977558&68&0.222618&88&0.062021\\
9&0.056384&29&0.189114&49&0.997456&69&0.204930&89&0.059104\\
10&0.059104&30&0.204930&50&0.997456&70&0.189114&90&0.056384\\
11&0.062021&31&0.222618&51&0.977558&71&0.174937&91&0.053844\\
12&0.065155&32&0.242443&52&0.940054&72&0.162198&92&0.051469\\
13&0.068526&33&0.264706&53&0.888899&73&0.150724&93&0.049245\\
14&0.072158&34&0.289748&54&0.828767&74&0.140363&94&0.047160\\
15&0.076079&35&0.317946&55&0.764151&75&0.130984&95&0.045202\\
16&0.080319&36&0.349711&56&0.698774&76&0.122473&96&0.043362\\
17&0.084912&37&0.385472&57&0.635356&77&0.114731&97&0.041631\\
18&0.089896&38&0.425649&58&0.575649&78&0.107673&98&0.040000\\
19&0.095316&39&0.470614&59&0.520610&79&0.101223&99&0.038462\\
\bottomrule\end{tabular}
\end{center}
Tính giá trị và đạo hàm tại $c=9/2$, một điểm nằm trong $[-5,5]$. Dựng bốn đa thức bằng cùng công thức Newton: $p_{99}$ dùng đủ 100 mốc; $P_{93,2}^+$ dùng các chỉ số $93,94,95$; $P_{92,4}^+$ dùng $92,\ldots,96$; $P_{91,6}^+$ dùng $91,\ldots,97$. Trong từng trường hợp,
\[
 d_j=f[x_\ell,\ldots,x_{\ell+j}]
 =\frac{\Delta^jy_\ell}{j!h^j},\qquad h=\frac{10}{99},
\]
Bảng nhân $B$ của từng tập mốc được dựng bằng \eqref{eq:newton-product-table}; tích $B\boldsymbol d$ cho hệ số đơn thức của đa thức tương ứng. Để tính giá trị và đạo hàm, có thể giữ trực tiếp cơ sở Newton và dùng
\[
 v_r=d_r,\quad u_r=0,\qquad
 u_j=v_{j+1}+(c-x_{\ell+j})u_{j+1},\quad
 v_j=d_j+(c-x_{\ell+j})v_{j+1}.
\]
Kết quả tính bằng số hữu tỷ rồi làm tròn để trình bày là
\begin{center}
\small\setlength{\tabcolsep}{4pt}
\begin{tabular}{c|rr}\toprule
Tập chỉ số mốc&$P(c)$&$P'(c)$\\\midrule
$0,\ldots,99$&$2.62333898674\cdot10^6$&$5.88103586782\cdot10^8$\\
$93,94,95$&$0.0470587357968$&$-0.0199480578110$\\
$92,\ldots,96$&$0.0470588237423$&$-0.0199307539996$\\
$91,\ldots,97$&$0.0470588235285$&$-0.0199307960326$\\\midrule
Hàm đã cho&$f(c)=4/85$&$f'(c)=-144/7225$\\\bottomrule
\end{tabular}
\end{center}
Các sai số tuyệt đối tương ứng:
\begin{center}
\small
\begin{tabular}{c|rr}\toprule
Số mốc&$|f(c)-P(c)|$&$|f'(c)-P'(c)|$\\\midrule
$100$&$2.62333893968\cdot10^6$&$5.88103586802\cdot10^8$\\
$3$&$8.77326192947\cdot10^{-8}$&$1.72619632275\cdot10^{-5}$\\
$5$&$2.12929017330\cdot10^{-10}$&$4.18481648091\cdot10^{-8}$\\
$7$&$9.40799612127\cdot10^{-13}$&$1.84817056407\cdot10^{-10}$\\\bottomrule
\end{tabular}
\end{center}
Tại tâm $c=0$, cùng $p_{99}$ cho $|f(0)-p_{99}(0)|\approx4.36017431239\cdot10^{-13}$ và $p_{99}'(0)=0$. Cả $0$ và $9/2$ đều nằm trong đoạn nội suy, nhưng các sai số rất khác nhau. Đây là sai số của đa thức trong số học chính xác, không phải hậu quả của làm tròn khi tính.


Đặt
\[
 r_j=\frac{5(2j+1)}{99},\quad0\leq j\leq49,\qquad
 \Omega(x)=\prod_{j=0}^{49}(x^2-r_j^2),\qquad
 D=\prod_{j=0}^{49}(1+r_j^2)>0.
\]
Các mốc thỏa $x_{99-i}=-x_i$ và $f(-x_i)=f(x_i)$. Đặt
$\widetilde p(x)=p_{99}(-x)$. Khi đó
\[
 \deg\widetilde p\leq99,\qquad
 \widetilde p(x_i)=p_{99}(x_{99-i})=f(x_{99-i})=f(x_i).
\]
Tính duy nhất cho $\widetilde p=p_{99}$. Viết
$p_{99}(x)=\sum_{\ell=0}^{99}a_\ell x^\ell$. So sánh hệ số cho
\[
 (-1)^\ell a_\ell=a_\ell
 \ \Longrightarrow\ a_{2j+1}=0\quad(0\leq j\leq49),
\]
\[
 p_{99}(x)=Q(x^2),\qquad Q(z)=\sum_{j=0}^{49}a_{2j}z^j.
\]
Đặt $R(x)=1-(1+x^2)p_{99}(x)$, $K=[x^{100}]R$. Vì $\Omega$ là đa thức đơn nhất bậc $100$,
\[
 \deg(R-K\Omega)\leq99,\qquad
 R(x_i)=1-(1+x_i^2)f(x_i)=0,\qquad\Omega(x_i)=0.
\]
Bổ đề \ref{bd:roots} suy ra $R-K\Omega=0$. Đặt
\[
 U(z)=\prod_{j=0}^{49}(z-r_j^2),\qquad
 H(z)=1-(1+z)Q(z)-KU(z).
\]
Từ $H(x^2)=R(x)-K\Omega(x)=0$,
\[
 [z^j]H=[x^{2j}]H(x^2)=0\quad(j\geq0),\qquad H=0.
\]
Thay $z=-1$ được
\[
 1=KU(-1)=K(-1)^{50}\prod_{j=0}^{49}(1+r_j^2)=KD,
 \qquad K=\frac1D.
\]
Vì $D>0$, $1+x^2>0$,
\begin{equation}\label{eq:newton-100-exact-error}
 f(x)-p_{99}(x)
 =\frac{1-(1+x^2)p_{99}(x)}{1+x^2}
 =\frac{\Omega(x)}{D(1+x^2)}.
\end{equation}
Lấy đạo hàm,
\begin{align}
 f'(x)-p_{99}'(x)
 &=\frac1D\left(\frac{\Omega'(x)}{1+x^2}
       -\frac{2x\Omega(x)}{(1+x^2)^2}\right)\notag\\
 &=\frac{(1+x^2)\Omega'(x)-2x\Omega(x)}{D(1+x^2)^2}.
\label{eq:newton-100-deriv-error}
\end{align}
Tại $c=9/2$, mỗi nhân tử thỏa
\[
 \frac{c^2-r_j^2}{1+r_j^2}
 =\frac{793881-100(2j+1)^2}{4[9801+25(2j+1)^2]}.
\]
Với $0\leq b\leq4$, đặt
\[
 P_b=\prod_{j=10b}^{10b+9}[793881-100(2j+1)^2],\qquad
 Q_b=\prod_{j=10b}^{10b+9}4[9801+25(2j+1)^2]>0,
\]
\[
 C_b=P_b/Q_b,\qquad
 \Delta_b^-=S_bP_b-L_bQ_b,\qquad
 \Delta_b^+=U_bQ_b-S_bP_b.
\]
Các số nguyên sau cho $S_b>0$ và $\Delta_b^-,\Delta_b^+>0$:
\begin{center}
\small
\begin{tabular}{c|rrr}\toprule
$b$&$S_b$&$L_b$&$U_b$\\\midrule
$0$&$1$&$671500000000$&$671510000000$\\
$1$&$1$&$23892000$&$23893000$\\
$2$&$100$&$53411$&$53412$\\
$3$&$10^7$&$24920$&$24921$\\
$4$&$10^{16}$&$-26104$&$-26103$\\\bottomrule
\end{tabular}
\end{center}
\begin{center}
\scriptsize
\begin{tabular}{cc|r}\toprule
$b$&Dấu&$\Delta_b^{\text{dấu}}$\\\midrule
$0$&$-$&$564460969348834046711847698088542484475980255307743801$\\
$0$&$+$&$684724826470029336161660449100001655679319830452256199$\\
$1$&$-$&$979724127227905110223191947116561091663793560633151801$\\
$1$&$+$&$197806234275437721862288427185902333616177604575424199$\\
$2$&$-$&$1930419243763990891963147862293447081401390661444727364$\\
$2$&$+$&$1835720507968248292494252144689366161274950628990881212$\\
$3$&$-$&$712457450876316338714507918277174321045663279055092286080$\\
$3$&$+$&$920949731233869418149642945013771182134689377734248922496$\\
$4$&$-$&$9123098148032443477488874790691938794181524716887951867904$\\
$4$&$+$&$179139523686347218419752530116283002436411363097803930140672$\\\bottomrule
\end{tabular}
\end{center}
Do đó
\[
 C_b-\frac{L_b}{S_b}=\frac{\Delta_b^-}{S_bQ_b}>0,
 \qquad
 \frac{U_b}{S_b}-C_b=\frac{\Delta_b^+}{S_bQ_b}>0.
\]
Đặc biệt $C_0,C_1,C_2,C_3>0$, $C_4<0$ và
\begin{align*}
 |f(c)-p_{99}(c)|
 &=\frac4{85}\prod_{b=0}^{4}|C_b|\\
 &>\frac4{85}(671500000000)(23892000)
    \frac{53411}{100}\frac{24920}{10^7}\frac{26103}{10^{16}}\\
 &=\frac{409853629424422953}{156250000000}\\
 &=2600000+\frac{3603629424422953}{156250000000}
 >2.6\cdot10^6.
\end{align*}
\end{vidu}

\Needspace{9\baselineskip}
\begin{vidu}
Giữ nguyên 100 mốc, chỉ làm nhiễu một giá trị: $\widetilde y_{50}=y_{50}+\varepsilon$, $\widetilde y_i=y_i$ với $i\ne50$. Theo tính tuyến tính của nội suy và \eqref{eq:cardinal},
\[
 \widetilde p_{99}(c)-p_{99}(c)=\varepsilon L_{50}(c),\qquad
 L_{50}(c)=\prod_{j\ne50}\frac{c-x_j}{x_{50}-x_j}.
\]
Tại $c=9/2$, các tích hữu tỷ cho
\[
 L_{50}(c)\approx9.19268262604\cdot10^{17},\qquad
 \sum_{i=0}^{99}|L_i(c)|\approx1.15685939400\cdot10^{19}.
\]
Với $\varepsilon=10^{-6}$, độ thay đổi giá trị nội suy xấp xỉ $9.19\cdot10^{11}$. Trong khi đó, đa thức cục bộ dùng các chỉ số $92,\ldots,96$ không đổi vì không sử dụng mốc $50$. Chặn sai số thực tế cần kiểm soát cả sai số xấp xỉ lẫn độ nhạy dữ liệu; điều kiện $c\in[a,b]$ không kiểm soát riêng được hai đại lượng này.
\end{vidu}

\Needspace{9\baselineskip}
\section{Các công thức nội suy trung tâm}
Các công thức dưới đây là các biểu diễn hữu hạn của cùng đa thức nội suy, thu được bằng cách sắp thứ tự mốc trong công thức Newton. Chỉ số âm là nhãn của mốc $x_j=x_0+jh$, không biểu thị một loại dữ liệu khác.
\Needspace{7\baselineskip}
\subsection{Công thức Gauss I và Gauss II}
Cho $m\geq1$, $h\ne0$ và các mốc $x_{-m},\ldots,x_m$. Đặt
\[
 t=\frac{x-x_0}{h},\qquad A_0(t)=t,\qquad
 A_r(t)=t\prod_{j=1}^{r}(t^2-j^2)\quad(r\geq1).
\]
\begin{dinhly}{}{central-gauss}
Đa thức nội suy $p_{2m}$ trên $2m+1$ mốc có hai biểu diễn
\begin{align}
 p_{2m}(x_0+ht)
 &=y_0+\sum_{r=1}^{m}\binom{t+r-1}{2r}\Delta^{2r}y_{-r}
 +\sum_{r=0}^{m-1}\binom{t+r}{2r+1}\Delta^{2r+1}y_{-r},
 \label{eq:central-gauss1}\\
 p_{2m}(x_0+ht)
 &=y_0+\sum_{r=1}^{m}\binom{t+r}{2r}\Delta^{2r}y_{-r}
 +\sum_{r=0}^{m-1}\binom{t+r}{2r+1}\Delta^{2r+1}y_{-r-1}.
 \label{eq:central-gauss2}
\end{align}
\end{dinhly}
\begin{chungminh}
Hai thứ tự chỉ số mốc được xác định bởi
\[
 z_0^{\rm I}=z_0^{\rm II}=0,\qquad
 z_k^{\rm I}=(-1)^{k+1}\lceil k/2\rceil,
 \qquad z_k^{\rm II}=-z_k^{\rm I}\quad(1\leq k\leq2m).
\]
Với $Q_k^\sigma=\{z_0^\sigma,\ldots,z_{k-1}^\sigma\}$, các tập chỉ số trong tích cơ sở và trong tỷ sai phân là
\begin{center}
\small
\begin{tabular}{c|c|c|c}\toprule
Thứ tự&Cấp $k$&$Q_k^\sigma$&$Q_{k+1}^\sigma$\\\midrule
I&$2r$, $1\leq r\leq m$&$\{1-r,\ldots,r\}$&$\{-r,\ldots,r\}$\\
II&$2r$, $1\leq r\leq m$&$\{-r,\ldots,r-1\}$&$\{-r,\ldots,r\}$\\
I&$2r+1$, $0\leq r<m$&$\{-r,\ldots,r\}$&$\{-r,\ldots,r+1\}$\\
II&$2r+1$, $0\leq r<m$&$\{-r,\ldots,r\}$&$\{-r-1,\ldots,r\}$\\\bottomrule
\end{tabular}
\end{center}
Tính đối xứng của tỷ sai phân và \eqref{eq:newton-equal-dd} cho
\[
 d_{2r}^{\rm I}=d_{2r}^{\rm II}=\frac{\Delta^{2r}y_{-r}}{(2r)!h^{2r}},\qquad
 d_{2r+1}^{\rm I}=\frac{\Delta^{2r+1}y_{-r}}{(2r+1)!h^{2r+1}},\qquad
 d_{2r+1}^{\rm II}=\frac{\Delta^{2r+1}y_{-r-1}}{(2r+1)!h^{2r+1}}.
\]
Trong biến $t$, các tích của bảng bằng
\begin{align*}
 h^{-2r}W_{2r}^{\rm I}(x_0+ht)
 &=\prod_{j=1-r}^{r}(t-j)
 =(2r)!\binom{t+r-1}{2r}=(t-r)A_{r-1}(t),\\
 h^{-2r}W_{2r}^{\rm II}(x_0+ht)
 &=\prod_{j=-r}^{r-1}(t-j)
 =(2r)!\binom{t+r}{2r}=(t+r)A_{r-1}(t),\\
 h^{-2r-1}W_{2r+1}^{\rm I}(x_0+ht)
 &=h^{-2r-1}W_{2r+1}^{\rm II}(x_0+ht)
 =\prod_{j=-r}^{r}(t-j)\\
 &=(2r+1)!\binom{t+r}{2r+1}=A_r(t).
\end{align*}
Thế các $d_k^\sigma W_k^\sigma$ vào \eqref{eq:newton-form} được hai công thức.
\end{chungminh}

\Needspace{9\baselineskip}
\begin{vidu}\label{vidu:central-five-table}
Cho $x_j=2+j/10$, $-2\leq j\leq2$, và bảng dữ liệu sau. Dạng bảng sai phân và bảng nhân theo cấu trúc bài giảng nội suy trung tâm \cite{central-slides}.
\begin{center}
\begin{tabular}{c|rrrrr}\toprule
$j$&$-2$&$-1$&$0$&$1$&$2$\\\midrule
$x_j$&$1.8$&$1.9$&$2$&$2.1$&$2.2$\\
$y_j$&$3$&$1$&$1$&$3$&$31$\\
$\Delta y_j$&$-2$&$0$&$2$&$28$&\\
$\Delta^2y_j$&$2$&$2$&$26$&&\\
$\Delta^3y_j$&$0$&$24$&&&\\
$\Delta^4y_j$&$24$&&&&\\\bottomrule
\end{tabular}
\end{center}
Đặt $q(t)=p_4(2+t/10)$. Trong biến $t$, hệ số ở cấp $k$ là sai phân được chọn chia cho $k!$. Hai đường trích hệ số là
\begin{center}
\small\setlength{\tabcolsep}{5pt}
\begin{tabular}{c|c|c|r|c|c|r}\toprule
&\multicolumn{3}{c|}{Gauss I}&\multicolumn{3}{c}{Gauss II}\\
$k$&Mốc thêm&Đại lượng chọn&$d_k^{\rm I}$&Mốc thêm&Đại lượng chọn&$d_k^{\rm II}$\\\midrule
$0$&$0$&$y_0$&$1$&$0$&$y_0$&$1$\\
$1$&$1$&$\Delta y_0$&$2$&$-1$&$\Delta y_{-1}$&$0$\\
$2$&$-1$&$\Delta^2y_{-1}$&$1$&$1$&$\Delta^2y_{-1}$&$1$\\
$3$&$2$&$\Delta^3y_{-1}$&$4$&$-2$&$\Delta^3y_{-2}$&$0$\\
$4$&$-2$&$\Delta^4y_{-2}$&$1$&$2$&$\Delta^4y_{-2}$&$1$\\\bottomrule
\end{tabular}
\end{center}
Cột ``Mốc thêm'' ghi chỉ số $j$, tức tọa độ của mốc trong biến $t$. Với từng thứ tự, bảng nhân được dựng bằng Horner ngược từ $W_0=1$:
\begin{center}
\small\setlength{\tabcolsep}{6pt}
\begin{tabular}{c|c|rrrrr|r}\toprule
&Tích cơ sở&$t^4$&$t^3$&$t^2$&$t$&$1$&Hệ số\\\midrule
Gauss I&$1$&$0$&$0$&$0$&$0$&$1$&$1$\\
&$t$&$0$&$0$&$0$&$1$&$0$&$2$\\
&$t(t-1)$&$0$&$0$&$1$&$-1$&$0$&$1$\\
&$t(t-1)(t+1)$&$0$&$1$&$0$&$-1$&$0$&$4$\\
&$t(t-1)(t+1)(t-2)$&$1$&$-2$&$-1$&$2$&$0$&$1$\\\cmidrule{2-8}
&Tổng có trọng số&$1$&$2$&$0$&$-1$&$1$&\\\midrule
Gauss II&$1$&$0$&$0$&$0$&$0$&$1$&$1$\\
&$t$&$0$&$0$&$0$&$1$&$0$&$0$\\
&$t(t+1)$&$0$&$0$&$1$&$1$&$0$&$1$\\
&$t(t+1)(t-1)$&$0$&$1$&$0$&$-1$&$0$&$0$\\
&$t(t+1)(t-1)(t+2)$&$1$&$2$&$-1$&$-2$&$0$&$1$\\\cmidrule{2-8}
&Tổng có trọng số&$1$&$2$&$0$&$-1$&$1$&\\\bottomrule
\end{tabular}
\end{center}
Hai phép nhân bảng cho cùng vector hệ số, đúng với tính duy nhất đã chứng minh:
\[
 q(t)=t^4+2t^3-t+1,\qquad
 p_4(x)=[10(x-2)]^4+2[10(x-2)]^3-10(x-2)+1.
\]
Tại $c=81/40$, $t=1/4$, các số hạng theo cấp là
\begin{center}
\begin{tabular}{c|rrrrr|r}\toprule
&$0$&$1$&$2$&$3$&$4$&Tổng\\\midrule
Gauss I&$1$&$1/2$&$-3/16$&$-15/16$&$105/256$&$201/256$\\
Gauss II&$1$&$0$&$5/16$&$0$&$-135/256$&$201/256$\\\bottomrule
\end{tabular}
\end{center}
\end{vidu}



\Needspace{7\baselineskip}
\subsection{Công thức Stirling}
\begin{dinhly}{}{central-stirling}
Với các mốc $x_{-m},\ldots,x_m$,
\begin{equation}\label{eq:central-stirling}
\begin{split}
 p_{2m}(x_0+ht)=y_0
 &+\sum_{r=1}^{m}\frac{t}{2r}\binom{t+r-1}{2r-1}\Delta^{2r}y_{-r}\\
 &+\sum_{r=0}^{m-1}\binom{t+r}{2r+1}
 \frac{\Delta^{2r+1}y_{-r-1}+\Delta^{2r+1}y_{-r}}2.
\end{split}
\end{equation}
\end{dinhly}
\begin{chungminh}
Theo \eqref{eq:generalized-binomial-recurrence}, với $r\geq1$,
\[
 \binom{t+r-1}{2r}=\frac{t-r}{2r}\binom{t+r-1}{2r-1},\qquad
 \binom{t+r}{2r}=\frac{t+r}{2r}\binom{t+r-1}{2r-1}.
\]
Do đó
\[
 \frac12\left[\binom{t+r-1}{2r}+\binom{t+r}{2r}\right]
 =\frac{t}{2r}\binom{t+r-1}{2r-1}
 =\frac{tA_{r-1}(t)}{(2r)!}.
\]
Lấy trung bình từng số hạng của \eqref{eq:central-gauss1} và \eqref{eq:central-gauss2}:
\begin{align*}
 \frac{p_{2m}(x_0+ht)+p_{2m}(x_0+ht)}2
 =y_0&+\sum_{r=1}^{m}\frac{t}{2r}\binom{t+r-1}{2r-1}\Delta^{2r}y_{-r}\\
 &+\sum_{r=0}^{m-1}\binom{t+r}{2r+1}
 \frac{\Delta^{2r+1}y_{-r-1}+\Delta^{2r+1}y_{-r}}2.
\end{align*}
Vế trái bằng $p_{2m}(x_0+ht)$, cho \eqref{eq:central-stirling}.
\end{chungminh}

\Needspace{9\baselineskip}
\begin{vidu}
Áp dụng Stirling cho dữ liệu của ví dụ \ref{vidu:central-five-table}. Bảng nhân của các tích chẵn và lẻ thu được từ \eqref{eq:central-stirling}:
\begin{center}
\small
\begin{tabular}{c|rrrrr|c}\toprule
Tích cơ sở&$t^4$&$t^3$&$t^2$&$t$&$1$&Hệ số\\\midrule
$1$&$0$&$0$&$0$&$0$&$1$&$y_0=1$\\
$t$&$0$&$0$&$0$&$1$&$0$&$(\Delta y_{-1}+\Delta y_0)/2=1$\\
$t^2$&$0$&$0$&$1$&$0$&$0$&$\Delta^2y_{-1}/2!=1$\\
$t(t^2-1)$&$0$&$1$&$0$&$-1$&$0$&$(\Delta^3y_{-2}+\Delta^3y_{-1})/(2\cdot3!)=2$\\
$t^2(t^2-1)$&$1$&$0$&$-1$&$0$&$0$&$\Delta^4y_{-2}/4!=1$\\\midrule
Tổng có trọng số&$1$&$2$&$0$&$-1$&$1$&\\\bottomrule
\end{tabular}
\end{center}
Tại $t=1/4$, năm số hạng là $1,1/4,1/16,-15/32,-15/256$; tổng bằng $201/256$. Để tính các đạo hàm, áp dụng Horner mở rộng cho vector hệ số của $q$. Ký hiệu $H_{j,s}$ là đạo hàm chuẩn hóa cấp $s$ của đa thức đuôi sau bước $j$:
\begin{center}
\begin{tabular}{c|rrrrr}\toprule
Bậc hệ số đang kết nạp $j$&$4$&$3$&$2$&$1$&$0$\\\midrule
$a_j$&$1$&$2$&$0$&$-1$&$1$\\
$H_{j,0}$&$1$&$9/4$&$9/16$&$-55/64$&$201/256$\\
$H_{j,1}$&$0$&$1$&$5/2$&$19/16$&$-9/16$\\
$H_{j,2}$&$0$&$0$&$1$&$11/4$&$15/8$\\
$H_{j,3}$&$0$&$0$&$0$&$1$&$3$\\
$H_{j,4}$&$0$&$0$&$0$&$0$&$1$\\\bottomrule
\end{tabular}
\end{center}
Hệ thức $p_4^{(s)}(c)=h^{-s}q^{(s)}(1/4)=10^s s!H_{0,s}$ cho
\begin{center}
\begin{tabular}{c|rrrrr}\toprule
$s$&$0$&$1$&$2$&$3$&$4$\\\midrule
$q^{(s)}(1/4)$&$201/256$&$-9/16$&$15/4$&$18$&$24$\\
$p_4^{(s)}(81/40)$&$201/256$&$-45/8$&$375$&$18000$&$240000$\\\bottomrule
\end{tabular}
\end{center}
\end{vidu}



\Needspace{7\baselineskip}
\subsection{Công thức Bessel}
Cho $m\geq0$, các mốc $x_{-m},\ldots,x_{m+1}$ và đặt
\[
 t=\frac{x-x_0}{h},\qquad
 B_0(t)=1,\qquad B_r(t)=\prod_{j=0}^{r-1}(t+j)(t-j-1).
\]
\begin{dinhly}{}{central-bessel}
Đa thức nội suy trên $2m+2$ mốc thỏa
\begin{equation}\label{eq:central-bessel}
\begin{split}
 p_{2m+1}(x_0+ht)
 =&\sum_{r=0}^{m}\binom{t+r-1}{2r}
 \frac{\Delta^{2r}y_{-r}+\Delta^{2r}y_{1-r}}2\\
 &+\sum_{r=0}^{m}\frac{t-1/2}{2r+1}\binom{t+r-1}{2r}\Delta^{2r+1}y_{-r}.
\end{split}
\end{equation}
\end{dinhly}
\begin{chungminh}
Chọn $z_0^{\rm I}=0$,
\[
 z_k^{\rm I}=(-1)^{k+1}\lceil k/2\rceil\quad(1\leq k\leq2m+1),\qquad
 z_k^{\rm II}=1-z_k^{\rm I}\quad(0\leq k\leq2m+1).
\]
Với $Q_k^\sigma=\{z_0^\sigma,\ldots,z_{k-1}^\sigma\}$,
\[
 Q_{2r}^{\rm I}=Q_{2r}^{\rm II}=\{1-r,\ldots,r\},\qquad
 Q_{2r+1}^{\rm I}=\{-r,\ldots,r\},\qquad
 Q_{2r+1}^{\rm II}=\{1-r,\ldots,r+1\}.
\]
Tại $r=0$, $Q_0^\sigma=\varnothing$. Do đó
\begin{align*}
 d_{2r}^{\rm I}&=\frac{\Delta^{2r}y_{-r}}{(2r)!h^{2r}},&
 d_{2r}^{\rm II}&=\frac{\Delta^{2r}y_{1-r}}{(2r)!h^{2r}},\\
 d_{2r+1}^{\rm I}&=\frac{\Delta^{2r+1}y_{-r}}{(2r+1)!h^{2r+1}},&
 d_{2r+1}^{\rm II}&=\frac{\Delta^{2r+1}y_{-r}}{(2r+1)!h^{2r+1}},
\end{align*}
và
\begin{align*}
 h^{-2r}W_{2r}^{\rm I}(x_0+ht)
 &=h^{-2r}W_{2r}^{\rm II}(x_0+ht)
 =\prod_{j=1-r}^{r}(t-j)=B_r(t),\\
 h^{-2r-1}W_{2r+1}^{\rm I}(x_0+ht)&=(t+r)B_r(t),\\
 h^{-2r-1}W_{2r+1}^{\rm II}(x_0+ht)&=(t-r-1)B_r(t).
\end{align*}
Mỗi thứ tự chứa toàn bộ chỉ số $-m,\ldots,m+1$, nên định lý \ref{dl:newton-form} cho hai tổng đều bằng $p_{2m+1}$. Lấy trung bình, dùng
\[
 \frac{(t+r)+(t-r-1)}2=t-\frac12,\qquad
 \frac{B_r(t)}{(2r)!}=\binom{t+r-1}{2r},
\]
thu được \eqref{eq:central-bessel}.
\end{chungminh}

\Needspace{9\baselineskip}
\begin{vidu}\label{vidu:central-six-table}
Cho sáu mốc $x_j=2+j/10$, $-2\leq j\leq3$. Dựng bảng sai phân trước khi chọn hai mốc giữa $x_0,x_1$:
\begin{center}
\begin{tabular}{c|rrrrrr}\toprule
$j$&$-2$&$-1$&$0$&$1$&$2$&$3$\\\midrule
$x_j$&$1.8$&$1.9$&$2$&$2.1$&$2.2$&$2.3$\\
$y_j$&$-45$&$-1$&$1$&$3$&$47$&$295$\\
$\Delta y_j$&$44$&$2$&$2$&$44$&$248$&\\
$\Delta^2y_j$&$-42$&$0$&$42$&$204$&&\\
$\Delta^3y_j$&$42$&$42$&$162$&&&\\
$\Delta^4y_j$&$0$&$120$&&&&\\
$\Delta^5y_j$&$120$&&&&&\\\bottomrule
\end{tabular}
\end{center}
Với $q(t)=p_5(2+t/10)$, công thức Bessel chọn các hệ số
\begin{center}
\small
\begin{tabular}{c|c|r}\toprule
$k$&Đại lượng lấy từ bảng sai phân&Hệ số\\\midrule
$0$&$(y_0+y_1)/2$&$2$\\
$1$&$\Delta y_0$&$2$\\
$2$&$(\Delta^2y_{-1}+\Delta^2y_0)/(2\cdot2!)$&$21/2$\\
$3$&$\Delta^3y_{-1}/3!$&$7$\\
$4$&$(\Delta^4y_{-2}+\Delta^4y_{-1})/(2\cdot4!)$&$5/2$\\
$5$&$\Delta^5y_{-2}/5!$&$1$\\\bottomrule
\end{tabular}
\end{center}
Các hàng chẵn là hệ số của $B_r(t)$; các hàng lẻ là hệ số của $(t-1/2)B_r(t)$. Dựng chúng bằng các phép nhân Horner ngược:
\begin{center}
\small\setlength{\tabcolsep}{5pt}
\begin{tabular}{c|rrrrrr|r}\toprule
Tích cơ sở&$t^5$&$t^4$&$t^3$&$t^2$&$t$&$1$&Hệ số\\\midrule
$1$&$0$&$0$&$0$&$0$&$0$&$1$&$2$\\
$t-1/2$&$0$&$0$&$0$&$0$&$1$&$-1/2$&$2$\\
$t(t-1)$&$0$&$0$&$0$&$1$&$-1$&$0$&$21/2$\\
$(t-1/2)t(t-1)$&$0$&$0$&$1$&$-3/2$&$1/2$&$0$&$7$\\
$(t+1)t(t-1)(t-2)$&$0$&$1$&$-2$&$-1$&$2$&$0$&$5/2$\\
$(t-1/2)(t+1)t(t-1)(t-2)$&$1$&$-5/2$&$0$&$5/2$&$-1$&$0$&$1$\\\midrule
Tổng có trọng số&$1$&$0$&$2$&$0$&$-1$&$1$&\\\bottomrule
\end{tabular}
\end{center}
Do đó $q(t)=t^5+2t^3-t+1$. Với $c=103/50$, $t=3/5$, bảng tính từng số hạng cho
\begin{center}
\begin{tabular}{c|rrrrrr}\toprule
Cấp&$0$&$1$&$2$&$3$&$4$&$5$\\\midrule
Số hạng&$2$&$1/5$&$-63/25$&$-21/125$&$168/125$&$168/3125$\\\bottomrule
\end{tabular}
\end{center}
Tổng bằng $2843/3125$. Horner mở rộng trên vector $(1,-1,0,2,0,1)^T$ cho các đạo hàm, rồi đổi biến theo \eqref{eq:newton-jet-scaling}:
\begin{center}
\small
\begin{tabular}{c|rrrrrr}\toprule
$s$&$0$&$1$&$2$&$3$&$4$&$5$\\\midrule
$q^{(s)}(3/5)$&$2843/3125$&$226/125$&$288/25$&$168/5$&$72$&$120$\\
$p_5^{(s)}(103/50)$&$2843/3125$&$452/25$&$1152$&$33600$&$720000$&$12000000$\\\bottomrule
\end{tabular}
\end{center}
\end{vidu}


\Needspace{7\baselineskip}
Các công thức trung tâm sử dụng những thứ tự mốc khác nhau. Khi cần toàn bộ đa thức và nhiều đạo hàm, có thể dựng bảng nhân trong biến $t=(x-x_0)/h$, rồi nhân bảng này với các sai phân đã chia cho giai thừa.

\begin{menhde}{}{central-product-coefficients}
Cho $N\geq0$, $h\ne0$ và $N+1$ mốc đều $x_j=x_0+jh$ có chỉ số trong một tập các số nguyên liên tiếp. Chọn một hoán vị $z_0,\ldots,z_N$ của tập chỉ số ấy sao cho mỗi tập đầu $\{z_0,\ldots,z_k\}$ cũng gồm các số nguyên liên tiếp. Đặt
\[
 \ell_k=\min\{z_0,\ldots,z_k\},\qquad
 d_k=\frac{\Delta^ky_{\ell_k}}{k!},\qquad
 W_k(t)=\prod_{i=0}^{k-1}(t-z_i).
\]
Nếu $B$ là bảng nhân của các $W_k$ theo \eqref{eq:newton-product-table} với $x_i$ thay bằng $z_i$, thì vector hệ số của $q(t)=p_N(x_0+ht)$ là
\[
 \boldsymbol a=B\boldsymbol d.
\]
Với hai thứ tự dùng để dựng Stirling hoặc Bessel, ta cũng có
\[
 \boldsymbol a=\frac{B^{(1)}\boldsymbol d^{(1)}+
 B^{(2)}\boldsymbol d^{(2)}}2.
\]
\end{menhde}
\begin{chungminh}
Với $S_k=\{z_0,\ldots,z_k\}$, giả thiết cho
\[
 |S_k|=k+1,\quad\min S_k=\ell_k,
 \quad S_k=\{\ell_k,\ell_k+1,\ldots,\ell_k+k\}.
\]
Từ tính đối xứng của tỷ sai phân và \eqref{eq:newton-equal-dd},
\[
 f[x_{z_0},\ldots,x_{z_k}]
 =f[x_{\ell_k},\ldots,x_{\ell_k+k}]
 =\frac{\Delta^ky_{\ell_k}}{k!h^k}.
\]
Với mọi $k\geq0$,
\[
 \prod_{i=0}^{k-1}(x_0+ht-x_{z_i})
 =\prod_{i=0}^{k-1}h(t-z_i)=h^kW_k(t).
\]
Công thức Newton cho
\begin{align*}
 q(t)=p_N(x_0+ht)
 &=\sum_{k=0}^{N}\frac{\Delta^ky_{\ell_k}}{k!h^k}h^kW_k(t)\\
 &=\sum_{k=0}^{N}d_k\sum_{j=0}^{k}B_{j,k}t^j
 =\sum_{j=0}^{N}\left(\sum_{k=j}^{N}B_{j,k}d_k\right)t^j.
\end{align*}
So sánh hệ số cho
\[
 a_j=\sum_{k=j}^{N}B_{j,k}d_k,\qquad\boldsymbol a=B\boldsymbol d.
\]
Đối với hai thứ tự của cùng tập mốc cuối, gọi $q_1,q_2$ là các đa thức vừa dựng. Khi đó
\[
 \deg(q_1-q_2)\leq N,\qquad
 (q_1-q_2)(z_i)=y_{z_i}-y_{z_i}=0\quad(0\leq i\leq N).
\]
Bổ đề \ref{bd:roots} cho $q_1=q_2=q$. So sánh từng hệ số,
\[
 (B^{(1)}\boldsymbol d^{(1)})_j=a_j
 =(B^{(2)}\boldsymbol d^{(2)})_j,
\]
\[
 \frac{(B^{(1)}\boldsymbol d^{(1)})_j
       +(B^{(2)}\boldsymbol d^{(2)})_j}{2}=a_j
 \quad(0\leq j\leq N).
\]
\end{chungminh}

\begin{menhde}{}{central-error}
Nếu $f\in C^{2m+1}([a,b])$ và các mốc $x_{-m},\ldots,x_m$ cùng $x$ thuộc $[a,b]$, thì sai số của Gauss I, Gauss II và Stirling thỏa
\[
 f(x)-p_{2m}(x)=\frac{f^{(2m+1)}(\xi_x)}{(2m+1)!}
 h^{2m+1}t\prod_{j=1}^{m}(t^2-j^2).
\]
Với Bessel, giả sử $f\in C^{2m+2}([a,b])$ và các mốc $x_{-m},\ldots,x_{m+1}$ cùng $x$ thuộc $[a,b]$. Khi đó
\[
 f(x)-p_{2m+1}(x)=\frac{f^{(2m+2)}(\eta_x)}{(2m+2)!}
 h^{2m+2}B_{m+1}(t).
\]
Các điểm $\xi_x,\eta_x$ thuộc $(a,b)$.
\end{menhde}
\begin{chungminh}
Vì $x-x_j=h(t-j)$,
\begin{align*}
 \prod_{j=-m}^{m}(x-x_j)
 &=h^{2m+1}\prod_{j=-m}^{m}(t-j)\\
 &=h^{2m+1}t\prod_{j=1}^{m}(t-j)(t+j)
 =h^{2m+1}t\prod_{j=1}^{m}(t^2-j^2),
\end{align*}
\begin{align*}
 \prod_{j=-m}^{m+1}(x-x_j)
 &=h^{2m+2}\prod_{j=-m}^{m+1}(t-j)\\
 &=h^{2m+2}\prod_{j=0}^{m}(t+j)(t-j-1)
 =h^{2m+2}B_{m+1}(t).
\end{align*}
Nếu $x\notin\{x_{-m},\ldots,x_m\}$, định lý \ref{dl:error} cho
\[
 f(x)-p_{2m}(x)
 =\frac{f^{(2m+1)}(\xi_x)}{(2m+1)!}
   \prod_{j=-m}^{m}(x-x_j)
 =\frac{f^{(2m+1)}(\xi_x)}{(2m+1)!}
   h^{2m+1}t\prod_{j=1}^{m}(t^2-j^2),\quad\xi_x\in(a,b).
\]
Nếu $x=x_i$, $-m\leq i\leq m$, thì
\[
 f(x_i)-p_{2m}(x_i)=0,\qquad
 \prod_{j=-m}^{m}(x_i-x_j)=0,
\]
nên đẳng thức vẫn đúng với mọi $\xi_x\in(a,b)$.
Tương tự, nếu $x\notin\{x_{-m},\ldots,x_{m+1}\}$,
\[
 f(x)-p_{2m+1}(x)
 =\frac{f^{(2m+2)}(\eta_x)}{(2m+2)!}
   \prod_{j=-m}^{m+1}(x-x_j)
 =\frac{f^{(2m+2)}(\eta_x)}{(2m+2)!}
   h^{2m+2}B_{m+1}(t),\quad\eta_x\in(a,b).
\]
Tại $x=x_i$, $-m\leq i\leq m+1$, hai vế đều bằng $0$ với mọi $\eta_x\in(a,b)$.
\end{chungminh}
Với $|t|\leq1/2$, công thức Stirling biểu diễn quanh mốc $x_0$; với $|t-1/2|\leq1/2$, Bessel biểu diễn quanh trung điểm $(x_0+x_1)/2$. Muốn giảm số mốc, phải dựng lại đa thức trên tập mốc được chọn và dùng chặn sai số của chính tập đó. Việc bỏ số hạng của một công thức trung bình không tự động cho đa thức nội suy trên tập mốc nhỏ hơn.

\Needspace{9\baselineskip}
\section{Nội suy ngược}
Các phương pháp và cách khoanh nhánh theo bài giảng \cite{inverse-slides}.
\Needspace{7\baselineskip}
\subsection{Bài toán nội suy ngược}
Cho dữ liệu $(x_i,y_i)$ và giá trị đích $y_\ast$. Bài toán là xấp xỉ một nghiệm của $f(x)=y_\ast$. Để xác định nhánh nghiệm, chọn $I=[a,b]$, $a<b$, trên đó $f$ đơn điệu nghiêm ngặt và $y_\ast\in f(I)$. Nếu chỉ có bảng dữ liệu, tính đơn điệu của các $y_i$ là điều kiện cho bảng, chưa chứng minh tính đơn điệu của hàm giữa các mốc.
\Needspace{7\baselineskip}
Giả sử $n\geq1$, $x_0<\cdots<x_n$. Với giá trị đích $y_\ast$, đặt
\[
\begin{aligned}
 Z&=\{i\in\{0,\ldots,n\}:y_i=y_\ast\},\\
 C&=\{i\in\{0,\ldots,n-1\}:(y_i-y_\ast)(y_{i+1}-y_\ast)<0\},\\
 F&=\{i\in\{0,\ldots,n-1\}:y_i=y_{i+1}=y_\ast\}.
\end{aligned}
\]
Tập $Z$ ghi các nghiệm tại mốc; tập $C$ ghi các khoảng đổi dấu liền nhau; tập $F$ ghi các đoạn hằng của nội suy tuyến tính có giá trị đích.
Định nghĩa $P_1$ trên từng đoạn bởi
\[
 P_1(x)=(1-\theta)y_i+\theta y_{i+1},\qquad
 \theta=\frac{x-x_i}{x_{i+1}-x_i},\quad x\in[x_i,x_{i+1}].
\]
\begin{menhde}{}{inv-sampled-branches}
Nếu $f$ liên tục và $f(x_i)=y_i$, mỗi $i\in C$ cho ít nhất một nghiệm trong $(x_i,x_{i+1})$; mỗi $i\in Z$ là một nghiệm tại mốc. Nếu còn biết $f$ đơn điệu nghiêm ngặt trên $[x_i,x_{i+1}]$, nghiệm trong khoảng đổi dấu đó là duy nhất.

Đối với hàm nội suy tuyến tính từng đoạn $P_1$, toàn bộ tập nghiệm được mô tả bởi các mốc $x_i$, $i\in Z$; các đoạn $[x_i,x_{i+1}]$, $i\in F$; và với $i\in C$, điểm duy nhất
\[
 \widehat x_i=x_i-\frac{(y_i-y_\ast)(x_{i+1}-x_i)}{y_{i+1}-y_i}.
\]
\end{menhde}
\begin{chungminh}
Với mỗi $0\leq i<n$, đặt
\[
 a=y_i-y_\ast,\qquad b=y_{i+1}-y_\ast,\qquad
 h_i=x_{i+1}-x_i>0,\qquad H(x)=f(x)-y_\ast.
\]
Nếu $i\in C$, thì
\[
 ab<0,\qquad\min(a,b)<0<\max(a,b),\qquad
 H(x_i)=a\ne0,\quad H(x_{i+1})=b\ne0.
\]
Định lý giá trị trung gian cho $\zeta\in[x_i,x_{i+1}]$ với
$H(\zeta)=0$. Vì $H(x_i),H(x_{i+1})\ne0$,
\[
 \zeta\in(x_i,x_{i+1}),\qquad f(\zeta)=y_\ast.
\]
Nếu $f$ đơn điệu nghiêm ngặt trên đoạn ấy, chọn
$\sigma\in\{-1,1\}$ để $\sigma f$ tăng nghiêm ngặt. Với hai nghiệm $u<v$,
\[
 0=\sigma\bigl(f(v)-f(u)\bigr)>0,
\]
mâu thuẫn. Với $i\in Z$,
\[
 f(x_i)-y_\ast=y_i-y_\ast=0.
\]

Ánh xạ $\theta=(x-x_i)/h_i$ đưa $[x_i,x_{i+1}]$ song ánh lên $[0,1]$, với ánh xạ ngược $x=x_i+h_i\theta$. Khi đó
\[
 P_1(x)-y_\ast=(1-\theta)a+\theta b=a+(b-a)\theta,
\]
\[
 \Theta_i=\{\theta\in[0,1]:a+(b-a)\theta=0\}.
\]
Năm trường hợp sau phủ mọi $(a,b)\in\R^2$:
\[
 a=b=0\ \Longrightarrow\ a+(b-a)\theta=0\quad(\theta\in[0,1]),
 \qquad\Theta_i=[0,1].
\]
\[
 a=0\ne b\ \Longrightarrow\
 a+(b-a)\theta=b\theta=0\ \Longleftrightarrow\ \theta=0,
 \qquad\Theta_i=\{0\}.
\]
\[
 b=0\ne a\ \Longrightarrow\
 a+(b-a)\theta=a(1-\theta)=0\ \Longleftrightarrow\ \theta=1,
 \qquad\Theta_i=\{1\}.
\]
Nếu $ab>0$, đặt $\eta=\operatorname{sgn}(a)=\operatorname{sgn}(b)$. Với $0\leq\theta\leq1$,
\[
 \eta\bigl(a+(b-a)\theta\bigr)
 =(1-\theta)|a|+\theta|b|
 \geq\min(|a|,|b|)>0,
 \qquad\Theta_i=\varnothing.
\]
Nếu $ab<0$, thì $b-a\ne0$ và
\[
 \theta_\ast=-\frac a{b-a}
 =\frac{|a|}{|a|+|b|}\in(0,1),\qquad
 a+(b-a)\theta=(b-a)(\theta-\theta_\ast),
\]
\[
 \Theta_i=\{\theta_\ast\},\qquad
 x_i+h_i\theta_\ast
 =x_i-\frac{(y_i-y_\ast)(x_{i+1}-x_i)}{y_{i+1}-y_i}
 =\widehat x_i.
\]
Do $[x_0,x_n]=\bigcup_{i=0}^{n-1}[x_i,x_{i+1}]$,
\begin{align*}
 \{x\in[x_0,x_n]:P_1(x)=y_\ast\}
 &=\bigcup_{i=0}^{n-1}\{x_i+h_i\theta:\theta\in\Theta_i\}\\
 &=\{x_i:i\in Z\}
   \cup\bigcup_{i\in F}[x_i,x_{i+1}]
   \cup\{\widehat x_i:i\in C\}.
\end{align*}
\end{chungminh}
Các khoảng đổi dấu của bảng không bảo đảm phát hiện mọi nghiệm của một hàm liên tục chưa biết. Chẳng hạn $f(x)=\cos(2\pi x)$ trên $[0,1]$ có $f(0)=f(1)=1$, nhưng $f(1/4)=f(3/4)=0$.

Đặt $\delta_i=\operatorname{sgn}(y_{i+1}-y_i)$, $0\leq i<n$. Một khối đơn điệu tăng của bảng là một đoạn chỉ số mốc $[a,b]$ có $\delta_a=\cdots=\delta_{b-1}=1$; khối giảm dùng dấu $-1$; khối hằng dùng dấu $0$. Khối cực đại không mở rộng được sang cạnh liền kề có cùng dấu. Các khối liên tiếp có thể chung mốc ở ranh giới.
\begin{menhde}{}{inv-sampled-monotonicity}
Trong mỗi khối tăng hoặc giảm, các giá trị $y_i$ đôi một phân biệt; vì vậy có thể lập tỷ sai phân cho các cặp đảo $(y_i,x_i)$. Tính đơn điệu này là của bảng và của nội suy tuyến tính trên khối; nó chưa suy ra tính đơn điệu của một hàm trơn phù hợp với bảng.
\end{menhde}
\begin{chungminh}
Trong một khối tăng hoặc giảm $[a,b]$, chọn
$\sigma\in\{-1,1\}$ sao cho
\[
 \sigma(y_{\ell+1}-y_\ell)=|y_{\ell+1}-y_\ell|>0
 \quad(a\leq\ell<b).
\]
Với $a\leq i<j\leq b$,
\[
 \sigma(y_j-y_i)
 =\sum_{\ell=i}^{j-1}\sigma(y_{\ell+1}-y_\ell)
 =\sum_{\ell=i}^{j-1}|y_{\ell+1}-y_\ell|>0.
\]
Do đó $y_j-y_i\ne0$; mọi mẫu số trong tỷ sai phân của các mốc đảo đều khác $0$.
Đặt
\[
 m_\ell=\frac{y_{\ell+1}-y_\ell}{x_{\ell+1}-x_\ell},
 \qquad\sigma m_\ell=|m_\ell|>0\quad(a\leq\ell<b).
\]
Với $x_a\leq u<v\leq x_b$, lập dãy
$u=z_0<z_1<\cdots<z_q=v$ bằng cách thêm mọi mốc $x_i\in(u,v)$.
Mỗi $[z_j,z_{j+1}]$ nằm trong một đoạn $[x_{\ell_j},x_{\ell_j+1}]$, nên
\[
 P_1(z_{j+1})-P_1(z_j)=m_{\ell_j}(z_{j+1}-z_j).
\]
Cộng các hiệu cho
\begin{align*}
 \sigma\bigl(P_1(v)-P_1(u)\bigr)
 &=\sum_{j=0}^{q-1}|m_{\ell_j}|(z_{j+1}-z_j)\\
 &\geq\min_{a\leq\ell<b}|m_\ell|\sum_{j=0}^{q-1}(z_{j+1}-z_j)\\
 &=\min_{a\leq\ell<b}|m_\ell|(v-u)>0.
\end{align*}

Với $f(x)=x+\sin(2\pi x)$,
\[
 f(0)=0<1=f(1),\qquad
 \lim_{\varepsilon\downarrow0}
 \frac{f(1/2+\varepsilon)-f(1/2)}{\varepsilon}
 =1-2\pi<0.
\]
Theo định nghĩa giới hạn, tồn tại $\delta\in(0,1/2)$ để
\[
 0<\varepsilon<\delta\ \Longrightarrow\
 \frac{f(1/2+\varepsilon)-f(1/2)}{\varepsilon}
 <\frac{1-2\pi}{2}<0,
\]
\[
 f(1/2+\varepsilon)<f(1/2).
\]
Hàm này không tăng trên $[0,1]$; đồng thời $f(0)<f(1)$ loại khả năng giảm. Bảng hai mốc tăng vì thế không đủ suy ra tính đơn điệu của hàm giữa các mốc.
\end{chungminh}

\Needspace{9\baselineskip}
\begin{vidu}
Với bảng chín mốc trong bài giảng nội suy ngược \cite{inverse-slides} và giá trị đích $y_\ast=3.5$:
\begin{center}
\small\setlength{\tabcolsep}{4pt}
\begin{tabular}{c|rrrrrrrrr}\toprule
$i$&$0$&$1$&$2$&$3$&$4$&$5$&$6$&$7$&$8$\\\midrule
$x_i$&$1.2$&$1.4$&$1.6$&$1.8$&$2.0$&$2.2$&$2.4$&$2.6$&$2.8$\\
$y_i$&$4.02$&$3.98$&$4.23$&$3.67$&$2.99$&$1.24$&$0.87$&$2.08$&$4.54$\\
$\operatorname{sgn}(y_i-3.5)$&$+$&$+$&$+$&$+$&$-$&$-$&$-$&$-$&$+$\\
$\delta_i$&$-$&$+$&$-$&$-$&$-$&$-$&$+$&$+$&\\\bottomrule
\end{tabular}
\end{center}
Ta có $Z=F=\varnothing$, $C=\{3,7\}$. Hai khoảng đổi dấu là $(1.8,2.0)$ và $(2.6,2.8)$. Các khối cực đại của bảng là $[0,1]$ giảm, $[1,2]$ tăng, $[2,6]$ giảm và $[6,8]$ tăng. Hai khoảng đổi dấu thuộc hai khối khác nhau, nên phép nội suy hàm ngược cần chọn và xử lý riêng từng nhánh. Khi không có thêm giả thiết về hàm giữa các mốc, bảng này chưa chứng minh mỗi khoảng có đúng một nghiệm.
\end{vidu}

\Needspace{7\baselineskip}
\subsection{Phương pháp hàm ngược}

\begin{menhde}{}{inv-function}
Giả sử $I=[a,b]$, $a<b$, $n\geq0$ nguyên, $f\in C^{n+1}(I)$ và $|f'(x)|\geq\mu>0$ trên $I$. Khi đó $f$ có hàm ngược $g:f(I)\to I$ thuộc $C^{n+1}$, và
\begin{equation}\label{eq:inv-derivatives}
\begin{aligned}
 g'(y)&=\frac1{f'(g(y))},\qquad
 g''(y)=-\frac{f''(g(y))}{(f'(g(y)))^3},\\
 g'''(y)&=\frac{3(f''(g(y)))^2-f'(g(y))f'''(g(y))}{(f'(g(y)))^5},
\end{aligned}
\end{equation}
trong đó chỉ dùng các công thức có đạo hàm được giả thiết tồn tại.
Ngoài ra, $|g(y)-g(z)|\leq|y-z|/\mu$.
\end{menhde}
\begin{chungminh}
Vì $|f'|\geq\mu>0$, $f'$ không triệt tiêu. Nếu có $u,v\in I$ với
$f'(u)f'(v)<0$, tính liên tục của $f'$ và định lý giá trị trung gian cho
\[
 \exists\theta\in(\min(u,v),\max(u,v)):\quad f'(\theta)=0,
\]
mâu thuẫn với $|f'(\theta)|\geq\mu$. Do đó tồn tại
$\sigma\in\{-1,1\}$ sao cho $\sigma f'(x)=|f'(x)|\geq\mu$ trên $I$.
Với $u<v$, định lý giá trị trung bình cho $\theta\in(u,v)$ và
\[
 \sigma\bigl(f(v)-f(u)\bigr)
 =\sigma f'(\theta)(v-u)=|f'(\theta)|(v-u)
 \geq\mu(v-u)>0.
\]
Đặc biệt $f$ đơn ánh. Đặt
\[
 A=\min(f(a),f(b)),\qquad B=\max(f(a),f(b)).
\]
Tính đơn điệu vừa chứng minh cho $f(I)\subset[A,B]$; định lý giá trị trung gian cho $[A,B]\subset f(I)$. Vì vậy
\[
 f(I)=[A,B],\qquad
 \forall y\in[A,B]\ \exists!x\in I:\ f(x)=y.
\]
Định nghĩa $g(y)$ là $x$ duy nhất ấy. Với $y,z\in[A,B]$,
\[
 |y-z|=|f(g(y))-f(g(z))|
 \geq\mu|g(y)-g(z)|,
\]
\[
 |g(y)-g(z)|\leq\frac{|y-z|}{\mu}.
\]
Với $z\ne y$, tính đơn ánh cho $g(z)\ne g(y)$, nên tồn tại
$\eta_{y,z}$ nằm giữa $g(y),g(z)$ sao cho
\[
 \frac{g(z)-g(y)}{z-y}
 =\frac1{[f(g(z))-f(g(y))]/[g(z)-g(y)]}
 =\frac1{f'(\eta_{y,z})}.
\]
Mặt khác,
\[
 |\eta_{y,z}-g(y)|\leq|g(z)-g(y)|\leq\frac{|z-y|}{\mu}
 \longrightarrow0\quad(z\to y).
\]
Tính liên tục của $f'$ và $|f'|\geq\mu$ suy ra
\[
 g'(y)=\lim_{z\to y}\frac{g(z)-g(y)}{z-y}
 =\frac1{f'(g(y))}.
\]
Ở hai đầu $A,B$, lấy giới hạn một phía. Hàm $1/(f'\circ g)$ liên tục, nên $g\in C^1([A,B])$.
Nếu $1\leq r\leq n$ và $g\in C^r([A,B])$, thì
\[
 f'\in C^n(I),\quad g\in C^r([A,B]),\quad |f'\circ g|\geq\mu
 \ \Longrightarrow\ \frac1{f'\circ g}\in C^r([A,B]).
\]
\[
 g'=\frac1{f'\circ g}\in C^r([A,B])
 \ \Longrightarrow\ g\in C^{r+1}([A,B]).
\]
Quy nạp cho $g\in C^{n+1}([A,B])$.
Với $x=g(y)$, quy tắc dây chuyền cho
\begin{align*}
 g''(y)
 &=-\frac{f''(g(y))g'(y)}{(f'(g(y)))^2}
 =-\frac{f''(x)}{(f'(x))^3},\\
 g'''(y)
 &=-\frac{f'''(x)g'(y)}{(f'(x))^3}
   +\frac{3(f''(x))^2g'(y)}{(f'(x))^4}\\
 &=-\frac{f'''(x)}{(f'(x))^4}
   +\frac{3(f''(x))^2}{(f'(x))^5}\\
 &=\frac{3(f''(x))^2-f'(x)f'''(x)}{(f'(x))^5}.
\end{align*}
Các phép lấy đạo hàm cuối chỉ áp dụng khi $n+1\geq2$ hoặc $n+1\geq3$ tương ứng.
\end{chungminh}

\Needspace{7\baselineskip}
Giả sử các $y_i$ đôi một phân biệt. Dựng đa thức $q_n$ nội suy các cặp $(y_i,x_i)$, rồi đặt $\widehat x=q_n(y_\ast)$. Khi $f$ thỏa điều kiện ở mệnh đề \ref{md:inv-function}, đây là nội suy hàm ngược $g$.
\begin{dinhly}{}{inv-interpolation}
Giả sử $f$ thỏa các giả thiết của mệnh đề \ref{md:inv-function}, các mốc $x_0,\ldots,x_n\in I$ đôi một phân biệt, $y_i=f(x_i)$ và $y_\ast\in f(I)$. Khi đó
\begin{equation}\label{eq:inv-interpolation-error}
 g(y_\ast)-q_n(y_\ast)
 =\frac{g^{(n+1)}(\zeta)}{(n+1)!}\prod_{i=0}^{n}(y_\ast-y_i),
 \qquad\zeta\in\operatorname{int}(f(I)).
\end{equation}
\end{dinhly}
\begin{chungminh}
Mệnh đề \ref{md:inv-function} cho $J=f(I)=[A,B]$, $g\in C^{n+1}(J)$ và
\[
 |y_i-y_j|=|f(x_i)-f(x_j)|\geq\mu|x_i-x_j|>0
 \quad(i\ne j),\qquad g(y_i)=x_i.
\]
Chọn hoán vị $\pi$ sao cho $y_{\pi(0)}<\cdots<y_{\pi(n)}$.
Đa thức $q_n$ thỏa
\[
 \deg q_n\leq n,\qquad q_n(y_{\pi(i)})=x_{\pi(i)}=g(y_{\pi(i)}).
\]
Nếu $y_\ast\notin\{y_0,\ldots,y_n\}$, định lý \ref{dl:error} áp dụng cho $g$ trên $J$ cho
\begin{align*}
 g(y_\ast)-q_n(y_\ast)
 &=\frac{g^{(n+1)}(\zeta)}{(n+1)!}
   \prod_{i=0}^{n}(y_\ast-y_{\pi(i)})\\
 &=\frac{g^{(n+1)}(\zeta)}{(n+1)!}
   \prod_{i=0}^{n}(y_\ast-y_i),\qquad\zeta\in(A,B).
\end{align*}
Nếu $y_\ast=y_j$, thì
\[
 g(y_j)-q_n(y_j)=x_j-x_j=0,
 \qquad\prod_{i=0}^{n}(y_j-y_i)=0,
\]
nên đẳng thức đúng với mọi $\zeta\in(A,B)$.
\end{chungminh}

\Needspace{9\baselineskip}
\begin{vidu}
Cho $f(x)=x^2$ trên $[1,3]$ với bảng $(1,1),(2,4),(3,9)$. Tỷ sai phân của bảng đảo là
\[
 d_0=1,\quad d_1=\frac{2-1}{4-1}=\frac13,\quad
 d_2=\frac{(3-2)/(9-4)-(2-1)/(4-1)}{9-1}=-\frac1{60}.
\]
Tại $y_\ast=2$, Horner dạng Newton cho
\begin{center}
\begin{tabular}{c|rrr}\toprule
$i$&$2$&$1$&$0$\\\midrule
$d_i$&$-1/60$&$1/3$&$1$\\
$y_\ast-y_i$&&$-2$&$1$\\
$v_i$&$-1/60$&$11/30$&$41/30$\\\bottomrule
\end{tabular}
\end{center}
Vậy $\widehat x=41/30$. Nội suy thuận cho đúng $p_2(x)=x^2$ theo tính duy nhất, nhưng
\[
 p_2(\widehat x)-2=\frac{1681}{900}-2=-\frac{119}{900}\ne0.
\]
Nội suy hàm ngược và giải phương trình của đa thức nội suy thuận là hai phép xấp xỉ khác nhau.
\end{vidu}

\Needspace{7\baselineskip}
\subsection{Phương pháp lặp}
Cho $I=[\alpha,\beta]$, $\alpha\leq\beta$, và $\varphi:I\to I$. Điểm bất động của $\varphi$ được xấp xỉ bằng $t_{k+1}=\varphi(t_k)$.
\begin{dinhly}{}{inv-contraction}
Giả sử $|\varphi(u)-\varphi(v)|\leq q|u-v|$ với $0\leq q<1$ và mọi $u,v\in I$. Khi đó tồn tại duy nhất $t_\ast\in I$ thỏa $\varphi(t_\ast)=t_\ast$. Với $t_0\in I$,
\begin{equation}\label{eq:inv-iteration-bound}
 |t_k-t_\ast|\leq\frac{q^k}{1-q}|t_1-t_0|\quad(k\geq0),\qquad
 |t_k-t_\ast|\leq\frac{q}{1-q}|t_k-t_{k-1}|\quad(k\geq1).
\end{equation}
\end{dinhly}
\begin{chungminh}
Từ $t_0\in I$ và $\varphi(I)\subset I$, quy nạp cho
\[
 t_j\in I\ \Longrightarrow\ t_{j+1}=\varphi(t_j)\in I,
 \qquad t_j\in I\quad(j\geq0).
\]
Đặt $\delta_j=|t_{j+1}-t_j|$. Với $j\geq1$,
\[
 \delta_j=|\varphi(t_j)-\varphi(t_{j-1})|
 \leq q\delta_{j-1},\qquad
 \delta_j\leq q^j\delta_0.
\]
Với $\ell>k$,
\begin{align*}
 |t_\ell-t_k|
 &\leq\sum_{j=k}^{\ell-1}|t_{j+1}-t_j|
 \leq\delta_0\sum_{j=k}^{\ell-1}q^j\\
 &=\frac{q^k(1-q^{\ell-k})}{1-q}\delta_0
 \leq\frac{q^k}{1-q}\delta_0.
\end{align*}
Vì $q^k\to0$, với mọi $\varepsilon>0$ tồn tại $K$ sao cho
\[
 \ell>k\geq K\ \Longrightarrow\ |t_\ell-t_k|<\varepsilon.
\]
Tính đầy đủ của $\R$ cho $t_k\to t_\ast\in\R$. Từ
$\alpha\leq t_k\leq\beta$ và phép lấy giới hạn,
$\alpha\leq t_\ast\leq\beta$, tức $t_\ast\in I$. Khi đó
\[
 |\varphi(t_\ast)-t_\ast|
 \leq|\varphi(t_\ast)-\varphi(t_k)|+|t_{k+1}-t_\ast|
 \leq q|t_\ast-t_k|+|t_{k+1}-t_\ast|\longrightarrow0,
\]
\[
 \varphi(t_\ast)=t_\ast.
\]
Nếu $s_\ast\in I$ cũng thỏa $\varphi(s_\ast)=s_\ast$, thì
\[
 |s_\ast-t_\ast|
 =|\varphi(s_\ast)-\varphi(t_\ast)|\leq q|s_\ast-t_\ast|,
\]
\[
 0\leq(1-q)|s_\ast-t_\ast|\leq0
 \ \Longrightarrow\ s_\ast=t_\ast.
\]
Lấy $\ell\to\infty$ trong chặn tổng hữu hạn cho
\[
 |t_k-t_\ast|\leq\frac{q^k}{1-q}|t_1-t_0|.
\]
Với $k\geq1$, từ $\delta_j\leq q^{j-k+1}\delta_{k-1}$ ($j\geq k$),
\begin{align*}
 |t_\ell-t_k|
 &\leq\delta_{k-1}\sum_{j=k}^{\ell-1}q^{j-k+1}\\
 &=\frac{q(1-q^{\ell-k})}{1-q}|t_k-t_{k-1}|.
\end{align*}
Cho $\ell\to\infty$ được
\[
 |t_k-t_\ast|\leq\frac q{1-q}|t_k-t_{k-1}|.
\]
Khi $q=0$, dùng $q^0=1$ trong các tổng trên; khi đó $\delta_j=0$ với $j\geq1$.
\end{chungminh}

\Needspace{7\baselineskip}
Cho $h>0$, mốc cách đều, $A=y_{k+1}-y_k\ne0$ và $y_\ast$ nằm giữa $y_k,y_{k+1}$. Một đa thức tiến được dựng trên các mốc $x_k,\ldots,x_{k+d_+}$, $d_+\geq1$; một đa thức lùi được dựng trên các mốc $x_{k+1-d_-},\ldots,x_{k+1}$, $d_-\geq1$. Các tập mốc phải thuộc bảng đã cho. Đặt
\[
 t=\frac{x-x_k}h\in[0,1],\quad
 s=\frac{x-x_{k+1}}h\in[-1,0],\quad
 \alpha=\frac{y_\ast-y_k}{A}\in[0,1],
\]
\[
 H_+(t)=\sum_{r=2}^{d_+}\frac{\Delta^ry_k}{r!}\prod_{j=0}^{r-1}(t-j),\qquad
 H_-(s)=\sum_{r=2}^{d_-}\frac{\nabla^ry_{k+1}}{r!}\prod_{j=0}^{r-1}(s+j).
\]
\begin{menhde}{}{inv-two-iterations}
\label{md:inv-forward-iteration}
Các phương trình nội suy thuận trên khoảng $[x_k,x_{k+1}]$ tương đương
\begin{align}
 P_{k,d_+}^{+}(x_k+ht)=y_\ast
 &\ \Longleftrightarrow\ t=\varphi(t),
 &\varphi(t)&=\alpha-H_+(t)/A,\label{eq:inv-forward-k}\\
 P_{k+1,d_-}^{-}(x_{k+1}+hs)=y_\ast
 &\ \Longleftrightarrow\ s=\psi(s),
 &\psi(s)&=\alpha-1-H_-(s)/A.\label{eq:inv-backward-k}
\end{align}
Nếu
\[
 \varphi([0,1])\subset[0,1],\quad
 \max_{[0,1]}|H_+'|\leq q_+|A|,\quad 0\leq q_+<1,
\]
thì lặp $t_0=\alpha$, $t_{j+1}=\varphi(t_j)$ hội tụ đến nghiệm duy nhất của đa thức tiến trong $[0,1]$. Nếu
\[
 \psi([-1,0])\subset[-1,0],\quad
 \max_{[-1,0]}|H_-'|\leq q_-|A|,\quad 0\leq q_-<1,
\]
thì lặp $s_0=\alpha-1$, $s_{j+1}=\psi(s_j)$ hội tụ đến nghiệm duy nhất của đa thức lùi trong $[-1,0]$. Chặn sai số theo $x$ lần lượt là
\[
 |x_j-x_*^+|\leq h\frac{q_+}{1-q_+}|t_j-t_{j-1}|,
 \quad x_j=x_k+ht_j,
\]
\[
 |x_j-x_*^-|\leq h\frac{q_-}{1-q_-}|s_j-s_{j-1}|,
 \quad x_j=x_{k+1}+hs_j.
\]
\end{menhde}
\begin{chungminh}
Từ $\alpha=(y_\ast-y_k)/A$ và $y_{k+1}-y_k=A$,
\[
 y_k-y_\ast=-A\alpha,\qquad
 y_{k+1}-y_\ast=A(1-\alpha),\qquad
 \nabla y_{k+1}=A.
\]
Công thức \eqref{eq:newton-forward-k} cho
\begin{align*}
 P_{k,d_+}^{+}(x_k+ht)-y_\ast
 &=y_k-y_\ast+At+H_+(t)\\
 &=A(t-\alpha)+H_+(t)
 =A\bigl(t-\varphi(t)\bigr).
\end{align*}
Công thức \eqref{eq:newton-backward-k} cho
\begin{align*}
 P_{k+1,d_-}^{-}(x_{k+1}+hs)-y_\ast
 &=y_{k+1}-y_\ast+As+H_-(s)\\
 &=A(s-\alpha+1)+H_-(s)
 =A\bigl(s-\psi(s)\bigr).
\end{align*}
Vì $A\ne0$,
\[
 P_{k,d_+}^{+}(x_k+ht)=y_\ast\ \Longleftrightarrow\ t=\varphi(t),
\]
\[
 P_{k+1,d_-}^{-}(x_{k+1}+hs)=y_\ast\ \Longleftrightarrow\ s=\psi(s).
\]
Với $u\ne v$ thuộc $[0,1]$, định lý giá trị trung bình cho
$\xi_+$ nằm giữa $u,v$ và
\[
 |\varphi(u)-\varphi(v)|
 =\frac{|H_+(u)-H_+(v)|}{|A|}
 =\frac{|H_+'(\xi_+)|}{|A|}|u-v|
 \leq q_+|u-v|.
\]
Với $u\ne v$ thuộc $[-1,0]$, tương tự,
\[
 |\psi(u)-\psi(v)|
 =\frac{|H_-(u)-H_-(v)|}{|A|}
 =\frac{|H_-'(\xi_-)|}{|A|}|u-v|
 \leq q_-|u-v|.
\]
Các bất đẳng thức cũng đúng khi $u=v$. Kết hợp với các giả thiết về ảnh, định lý \ref{dl:inv-contraction} cho các điểm bất động duy nhất
\[
 t_\ast\in[0,1],\quad s_\ast\in[-1,0],\quad
 t_j\to t_\ast,\quad s_j\to s_\ast,
\]
trong đó $t_0=\alpha\in[0,1]$, $s_0=\alpha-1\in[-1,0]$.
Phép đổi biến cho
\[
 x_\ast^+=x_k+ht_\ast\in[x_k,x_{k+1}],\qquad
 x_\ast^-=x_{k+1}+hs_\ast\in[x_k,x_{k+1}].
\]
Với $j\geq1$, từ \eqref{eq:inv-iteration-bound},
\[
 |x_k+ht_j-x_\ast^+|
 =h|t_j-t_\ast|
 \leq h\frac{q_+}{1-q_+}|t_j-t_{j-1}|,
\]
\[
 |x_{k+1}+hs_j-x_\ast^-|
 =h|s_j-s_\ast|
 \leq h\frac{q_-}{1-q_-}|s_j-s_{j-1}|.
\]
\end{chungminh}
Hai đa thức được dựng từ hai tập mốc khác nhau có thể cho hai nghiệm xấp xỉ khác nhau. Khi hai tập mốc trùng nhau, tính duy nhất cho hai đa thức bằng nhau; nếu mỗi phép lặp thỏa các điều kiện trên, hai nghiệm theo $x$ trùng nhau.

\begin{menhde}{}{inv-coefficient-certificate}
Viết $H_+(t)=\sum_{r=0}^{d_+}a_r^+t^r$, $H_-(s)=\sum_{r=0}^{d_-}a_r^-s^r$, các hệ số được dựng bằng Horner ngược. Đặt $\delta=\min(\alpha,1-\alpha)$. Với một trong hai dấu $\sigma\in\{+,-\}$, các điều kiện đủ kiểm tra được từ hệ số là
\begin{equation}\label{eq:inv-coefficient-certificate}
 \sum_r|a_r^\sigma|\leq|A|\delta,
 \qquad q_\sigma=\frac{\sum_{r\geq1}r|a_r^\sigma|}{|A|}<1.
\end{equation}
Khi đó phép lặp tương ứng có ảnh nằm trong đoạn và hệ số co không quá $q_\sigma$.
\end{menhde}
\begin{chungminh}
Đặt $I_+=[0,1]$, $I_-=[-1,0]$. Với $z\in I_\sigma$, $|z|\leq1$, nên
\begin{align*}
 |H_\sigma(z)|
 &=\left|\sum_{r=0}^{d_\sigma}a_r^\sigma z^r\right|
 \leq\sum_{r=0}^{d_\sigma}|a_r^\sigma||z|^r
 \leq\sum_{r=0}^{d_\sigma}|a_r^\sigma|
 \leq|A|\delta,\\
 |H_\sigma'(z)|
 &=\left|\sum_{r=1}^{d_\sigma}r a_r^\sigma z^{r-1}\right|
 \leq\sum_{r=1}^{d_\sigma}r|a_r^\sigma||z|^{r-1}
 \leq|A|q_\sigma.
\end{align*}
Vì $\delta=\min(\alpha,1-\alpha)$,
\[
 0\leq\delta\leq\alpha,\qquad\delta\leq1-\alpha.
\]
Với dạng tiến,
\[
 |\varphi(t)-\alpha|=\frac{|H_+(t)|}{|A|}\leq\delta,
\]
\[
 0\leq\alpha-\delta\leq\varphi(t)
 \leq\alpha+\delta\leq1.
\]
Với dạng lùi,
\[
 |\psi(s)-(\alpha-1)|=\frac{|H_-(s)|}{|A|}\leq\delta,
\]
\[
 -1\leq\alpha-1-\delta\leq\psi(s)
 \leq\alpha-1+\delta\leq0.
\]
Với hai điểm phân biệt trong đoạn tương ứng, định lý giá trị trung bình cho
\[
 |\varphi(u)-\varphi(v)|
 =\frac{|H_+'(\xi)|}{|A|}|u-v|\leq q_+|u-v|
 \quad(u,v\in I_+),
\]
hoặc
\[
 |\psi(u)-\psi(v)|
 =\frac{|H_-'(\xi)|}{|A|}|u-v|\leq q_-|u-v|
 \quad(u,v\in I_-).
\]
Khi hai điểm trùng nhau, hai vế bằng $0$. Mỗi dấu chỉ sử dụng điều kiện hệ số của chính dấu ấy.
\end{chungminh}
Các điều kiện \eqref{eq:inv-coefficient-certificate} là điều kiện đủ. Không thỏa chúng chưa kết luận phép lặp phân kỳ: có thể dùng chặn chính xác hơn của $H_\sigma,H_\sigma'$ trên đoạn. Chặn lặp đánh giá sai số đối với nghiệm của đa thức. Để đánh giá nghiệm của hàm, cần kết hợp sai số lặp với sai số nội suy và một chặn dưới cho $|f'|$.


\begin{menhde}{}{inv-residual}
Giả sử $f$ thỏa mệnh đề \ref{md:inv-function}, $f(x_\ast)=y_\ast$, $p$ là một đa thức và $\widehat x\in I$. Khi đó
\begin{equation}\label{eq:inv-residual}
 |\widehat x-x_\ast|
 \leq\frac{|p(\widehat x)-y_\ast|+\|f-p\|_{\infty,I}}{\mu}.
\end{equation}
\end{menhde}
\begin{chungminh}
Nếu $\widehat x=x_\ast$, vế trái của \eqref{eq:inv-residual} bằng $0$ và vế phải không âm. Nếu $\widehat x\ne x_\ast$, định lý giá trị trung bình cho $\theta$ nằm giữa $\widehat x,x_\ast$ và
\begin{align*}
 \mu|\widehat x-x_\ast|
 &\leq|f'(\theta)||\widehat x-x_\ast|\\
 &=|f(\widehat x)-f(x_\ast)|=|f(\widehat x)-y_\ast|\\
 &=|f(\widehat x)-p(\widehat x)+p(\widehat x)-y_\ast|\\
 &\leq|f(\widehat x)-p(\widehat x)|+|p(\widehat x)-y_\ast|\\
 &\leq\|f-p\|_{\infty,I}+|p(\widehat x)-y_\ast|.
\end{align*}
Chia cho $\mu>0$ được \eqref{eq:inv-residual}.
\end{chungminh}

\Needspace{9\baselineskip}
\begin{vidu}
Cho $p(x)=x+x^2/10$, bảng mốc $0,1,2$ có $\boldsymbol y=(0,11/10,12/5)^T$, và $y_\ast=1/2$. Công thức sai phân cho $\Delta y_0=11/10$, $\Delta^2y_0=1/5$, nên
\[
 \varphi(t)=\frac5{11}-\frac{t(t-1)}{11}.
\]
Trên $I=[0,1]$, $-1/4\leq t(t-1)\leq0$ và $|2t-1|\leq1$, nên
\[
 \varphi(I)\subset[5/11,21/44]\subset I,\qquad
 |\varphi'(t)|\leq1/11.
\]
Từ $t_0=5/11$, một bước cho $t_1=635/1331$; chặn hậu nghiệm là
\[
 |t_1-t_\ast|\leq\frac1{10}\left|\frac{635}{1331}-\frac5{11}\right|
 =\frac3{1331}.
\]
Giá trị nghiệm của đa thức được xác định bởi $(t+5)^2=30$, $t\in[0,1]$, nên $t_\ast=\sqrt{30}-5$. Nếu $p$ chỉ nội suy một hàm chưa biết, chặn này phải được kết hợp với \eqref{eq:inv-residual} để đánh giá sai số đối với nghiệm của hàm.
\end{vidu}

% END EXTENSION NEWTON

% BEGIN EXTENSION SPLINE LS
\Needspace{9\baselineskip}
\section{Hàm ghép trơn}
Các bậc spline và điều kiện biên trong bài giảng \cite{spline-slides} được đặt trong cấu trúc ghép đa thức tổng quát.
Mỗi khoảng giữa hai vị trí đo có thể được mô tả bằng một đa thức riêng. Ta yêu cầu các đa thức khớp giá trị ở mốc chung; nếu cần đường cong có đạo hàm liên tục, các đạo hàm ở mốc chung cũng phải bằng nhau. Bài toán bắt đầu từ cách ghép các đoạn thẳng, rồi tăng bậc đa thức và bổ sung các điều kiện khớp đạo hàm.
% BEGIN SPL-GENERAL-QUADRATIC
\Needspace{7\baselineskip}
Cho $N\geq1$, $x_0<\cdots<x_N$, $h_i=x_{i+1}-x_i$ và các giá trị $y_i$.
Với $p\geq1$, ta tìm các đa thức $S_i\in\mathcal P_p$ sao cho
\[
 S_i(x_i)=y_i,\quad S_i(x_{i+1})=y_{i+1},\qquad
 S_{i-1}^{(k)}(x_i)=S_i^{(k)}(x_i)\quad(1\leq k<p).
\]
Như vậy hàm ghép thuộc $C^{p-1}([x_0,x_N])$. Điều kiện $S_i\in\mathcal P_p$ có nghĩa là bậc không quá $p$; bậc có thể thấp hơn khi dữ liệu khiến hệ số đầu bằng $0$.
Các giá trị nội suy chưa xác định duy nhất các đa thức khi $p\geq2$. Ta bắt đầu bằng cách cho thêm các đạo hàm tại đầu trái.
\begin{dinhly}{}{spl-initial-jet}
Với mọi $\alpha_1,\ldots,\alpha_{p-1}\in\R$, tồn tại duy nhất hàm ghép trên thỏa
\[
 S^{(k)}(x_0)=\alpha_k\quad(1\leq k<p).
\]
Đặt $s_i^{(k)}=S^{(k)}(x_i)$, $s_i^{(0)}=y_i$. Các hệ số trong biến $t=x-x_i$ được xác định bởi
\begin{equation}\label{eq:spl-initial-jet-coefficients}
 S_i(x_i+t)=\sum_{k=0}^{p-1}\frac{s_i^{(k)}}{k!}t^k+c_it^p,
 \qquad c_i=\frac{y_{i+1}-\sum_{k=0}^{p-1}s_i^{(k)}h_i^k/k!}{h_i^p},
\end{equation}
\begin{equation}\label{eq:spl-initial-jet-propagation}
 s_{i+1}^{(k)}=\sum_{j=k}^{p-1}\frac{s_i^{(j)}}{(j-k)!}h_i^{j-k}
 +\frac{p!}{(p-k)!}c_ih_i^{p-k}\quad(1\leq k<p).
\end{equation}
Tập các hàm ghép nội suy khi chưa thêm điều kiện biên là một tập affine có số chiều $p-1$.
\end{dinhly}
\begin{chungminh}
Một đa thức $P(t)=\sum_{k=0}^pa_kt^k$ thỏa các điều kiện ở hai đầu khi và chỉ khi
\[
 k!a_k=s_i^{(k)}\quad(0\leq k<p),\qquad
 h_i^pa_p=y_{i+1}-\sum_{k=0}^{p-1}\frac{s_i^{(k)}}{k!}h_i^k.
\]
Vì $h_i^p>0$, các hệ số duy nhất là $a_k=s_i^{(k)}/k!$ với $k<p$ và $a_p=c_i$ trong \eqref{eq:spl-initial-jet-coefficients}. Với $1\leq k<p$,
\[
 S_i^{(k)}(x_i+t)
 =\sum_{j=k}^{p-1}\frac{s_i^{(j)}}{(j-k)!}t^{j-k}
 +\frac{p!}{(p-k)!}c_it^{p-k}.
\]
Thay $t=h_i$ được \eqref{eq:spl-initial-jet-propagation}. Bắt đầu từ $s_0^{(0)}=y_0$, $s_0^{(k)}=\alpha_k$, tại mỗi bước các công thức xác định duy nhất $c_i$, $S_i$ và $s_{i+1}^{(k)}$.

Đặt $D_k$ là hàm ghép từ các $S_i^{(k)}$ với $0\leq k<p$. Các điều kiện đã dựng cho
\[
 S_{i-1}^{(k)}(x_i)=s_i^{(k)}=S_i^{(k)}(x_i)
 \quad(0\leq k<p),
\]
nên mỗi $D_k$ liên tục. Với $k<p-1$ và $u$ đủ nhỏ,
\[
 D_k(x_i+u)-D_k(x_i)=\int_0^uD_{k+1}(x_i+t)\,dt,
\]
\[
 \lim_{u\to0}\frac{D_k(x_i+u)-D_k(x_i)}u
 =D_{k+1}(x_i).
\]
Vì vậy $D_k'=D_{k+1}$ ở mốc cũng như trong từng đoạn; lặp theo $k$ cho $S=D_0\in C^{p-1}$. Nếu $\widetilde S$ thỏa cùng dữ liệu và cùng $\boldsymbol\alpha$, các hệ số đầu đoạn $0$ bằng nhau. Quy nạp theo $i$ cho các $s_i^{(k)}$, $c_i$ và $S_i$ bằng nhau, tức $\widetilde S=S$.

Gọi $S^{\boldsymbol\alpha}$ là hàm ứng với các đạo hàm đầu $\boldsymbol\alpha$, và đặt $F_k=S^{\boldsymbol e_k}-S^{\boldsymbol0}$. Khi đó
\[
 F_k(x_i)=0,\qquad F_k^{(j)}(x_0)=\delta_{kj}
 \quad(1\leq j,k<p).
\]
Hàm
\[
 \widehat S=S^{\boldsymbol0}+\sum_{k=1}^{p-1}\alpha_kF_k
\]
là spline cùng bậc, cùng độ trơn và có
\[
 \widehat S(x_i)=y_i,\qquad
 \widehat S^{(j)}(x_0)
 =0+\sum_{k=1}^{p-1}\alpha_k\delta_{kj}=\alpha_j.
\]
Tính duy nhất vừa chứng minh cho $\widehat S=S^{\boldsymbol\alpha}$. Nếu $\sum_{k=1}^{p-1}c_kF_k=0$, thì
\[
 0=\left(\sum_{k=1}^{p-1}c_kF_k\right)^{(j)}(x_0)
 =\sum_{k=1}^{p-1}c_k\delta_{kj}=c_j
 \quad(1\leq j<p).
\]
Do đó
\[
 \{S^{\boldsymbol\alpha}:\boldsymbol\alpha\in\R^{p-1}\}
 =S^{\boldsymbol0}+\operatorname{span}\{F_1,\ldots,F_{p-1}\},
 \qquad
 \dim\operatorname{span}\{F_1,\ldots,F_{p-1}\}=p-1.
\]
Với $p=1$, tổng và họ $F_k$ rỗng, không gian hiệu là $\{0\}$ và số chiều bằng $0$.
\end{chungminh}

\Needspace{7\baselineskip}

\subsection{Nội suy tuyến tính từng đoạn}
Cho $N\geq1$, $a=x_0<x_1<\cdots<x_N=b$ và $y_0,\ldots,y_N\in\R$. Đặt
\[
 h_i=x_{i+1}-x_i>0,\qquad \delta_i=\frac{y_{i+1}-y_i}{h_i}\quad(0\leq i<N),\qquad h=\max_{0\leq i<N}h_i.
\]
\begin{dinhly}{}{spl-linear}
Hàm liên tục $L$ nhận giá trị $y_i$ tại $x_i$ và là đa thức bậc không quá $1$ trên mỗi $[x_i,x_{i+1}]$ được xác định duy nhất bởi
\begin{equation}\label{eq:spl-linear}
 L(x)=y_i+\delta_i(x-x_i)
 =\frac{x_{i+1}-x}{h_i}y_i+\frac{x-x_i}{h_i}y_{i+1},\qquad x\in[x_i,x_{i+1}].
\end{equation}
Nếu $y_i=f(x_i)$ với $f\in C^2([a,b])$, thì
\[
 |f(x)-L(x)|\leq\frac{\|f''\|_{\infty,[x_i,x_{i+1}]}}2(x-x_i)(x_{i+1}-x),\qquad
 \|f-L\|_\infty\leq\frac{h^2}8\|f''\|_\infty.
\]
\end{dinhly}
\begin{chungminh}
Viết $L_i(x)=u_i+v_i(x-x_i)$. Hai điều kiện tại mốc cho
\[
 u_i=y_i,\qquad u_i+v_ih_i=y_{i+1},\qquad v_i=\delta_i.
\]
Tại mốc trong,
\[
 L_{i-1}(x_i)=y_{i-1}+\delta_{i-1}h_{i-1}=y_i=L_i(x_i).
\]
Định lý \ref{dl:error} với hai mốc $x_i,x_{i+1}$ cho, khi $x$ không trùng mốc,
\[
 f(x)-L_i(x)=\frac{f''(\xi_x)}2(x-x_i)(x-x_{i+1}),\qquad \xi_x\in(x_i,x_{i+1}).
\]
Tại hai mốc, $f-L_i=0$. Với $t=x-x_i\in[0,h_i]$,
\[
 t(h_i-t)=\frac{h_i^2}4-\left(t-\frac{h_i}2\right)^2\leq\frac{h_i^2}4.
\]
Với mọi $x\in[x_i,x_{i+1}]$, kể cả hai đầu mút,
\[
 |f(x)-L_i(x)|
 \leq\frac{\|f''\|_{\infty,[x_i,x_{i+1}]}}2t(h_i-t)
 \leq\frac{h_i^2}8\|f''\|_{\infty,[x_i,x_{i+1}]}.
\]
Suy ra
\[
 \|f-L\|_\infty
 =\max_{0\leq i<N}\|f-L_i\|_{\infty,[x_i,x_{i+1}]}
 \leq\max_i\frac{h_i^2}8\|f''\|_{\infty,[x_i,x_{i+1}]}
 \leq\frac{h^2}8\|f''\|_\infty.
\]
\end{chungminh}
\begin{vidu}
Ba số đo chiều cao mặt đường tại các vị trí $0,1,2$ là $0,1,0$ (đơn vị được chọn thống nhất). Khi đó
\[
 L(x)=\begin{cases}x,&0\leq x\leq1,\\2-x,&1\leq x\leq2,\end{cases}\qquad
 L'(1^-)=1,\quad L'(1^+)=-1.
\]
Horner trên đoạn đầu tại $x=1/2$ cho
\begin{center}\begin{tabular}{c|rr}\toprule
$j$&$1$&$0$\\\midrule
$a_j$&$1$&$0$\\
$(1/2)b_{j+1}$&---&$1/2$\\\midrule
$b_j$&$1$&$\boxed{1/2}$\\\bottomrule
\end{tabular}\end{center}
\end{vidu}

\subsection{Spline bậc hai}
Khi $p=2$, đặt $m_i=S'(x_i)$ và $\delta_i=(y_{i+1}-y_i)/h_i$. Đạo hàm trên mỗi đoạn là một đa thức bậc không quá một; điều kiện $C^1$ được biểu thị bằng cùng một $m_i$ ở hai phía mốc.
\begin{menhde}{}{spl-quadratic}
Các spline bậc hai nội suy dữ liệu được xác định bởi
\begin{equation}\label{eq:spl-quadratic}
 \begin{aligned}
 m_{i+1}&=2\delta_i-m_i,\\
 S_i(x_i+t)&=y_i+t\left(m_i+\frac{m_{i+1}-m_i}{2h_i}t\right),
 \qquad0\leq t\leq h_i.
 \end{aligned}
\end{equation}
Nếu cho $m_0=\alpha$, thì
\begin{equation}\label{eq:spl-quadratic-explicit-slopes}
 m_i=(-1)^i\alpha+2\sum_{j=0}^{i-1}(-1)^{i-1-j}\delta_j.
\end{equation}
Điều kiện thêm $m_N=\beta$ cùng với $m_0=\alpha$ khả thi khi và chỉ khi $\beta$ bằng vế phải của \eqref{eq:spl-quadratic-explicit-slopes} tại $i=N$.
\end{menhde}
\begin{chungminh}
Từ \eqref{eq:spl-initial-jet-coefficients},
\[
 S_i(x_i+t)=y_i+m_it+\frac{\delta_i-m_i}{h_i}t^2,
 \qquad S_i'(x_{i+1})=m_i+2(\delta_i-m_i)=2\delta_i-m_i.
\]
Do đó $(\delta_i-m_i)/h_i=(m_{i+1}-m_i)/(2h_i)$, cho \eqref{eq:spl-quadratic}. Công thức tường minh đúng với $i=0$ bởi tổng rỗng bằng $0$. Nếu đúng tại $i$, thì
\begin{align*}
 m_{i+1}&=2\delta_i-(-1)^i\alpha
       -2\sum_{j=0}^{i-1}(-1)^{i-1-j}\delta_j\\
 &=(-1)^{i+1}\alpha+2\sum_{j=0}^{i}(-1)^{i-j}\delta_j.
\end{align*}
Quy nạp cho \eqref{eq:spl-quadratic-explicit-slopes}. Với nghiệm duy nhất có $m_0=\alpha$,
\[
 S'(x_N)=m_N=(-1)^N\alpha+
 2\sum_{j=0}^{N-1}(-1)^{N-1-j}\delta_j.
\]
Vì vậy
\[
 S'(x_N)=\beta
 \ \Longleftrightarrow\
 \beta=(-1)^N\alpha+2\sum_{j=0}^{N-1}(-1)^{N-1-j}\delta_j.
\]
\end{chungminh}
\Needspace{9\baselineskip}
\begin{vidu}
Với $\boldsymbol x=(0,1,2,3)^T$, $\boldsymbol y=(0,2,3,5)^T$ và $S'(0)=1$, ta có
\begin{center}
\begin{tabular}{c|rrrr}\toprule
$i$&$0$&$1$&$2$&$3$\\\midrule
$y_i$&$0$&$2$&$3$&$5$\\
$\delta_i$&$2$&$1$&$2$&\\
$m_i$&$1$&$3$&$-1$&$5$\\\bottomrule
\end{tabular}
\end{center}
Các hệ số trong biến $t=x-x_i$ là
\begin{center}
\begin{tabular}{c|rrr}\toprule
Đoạn&$[t^2]S_i$&$[t]S_i$&$[1]S_i$\\\midrule
$[0,1]$&$1$&$1$&$0$\\
$[1,2]$&$-2$&$3$&$2$\\
$[2,3]$&$3$&$-1$&$3$\\\bottomrule
\end{tabular}
\end{center}
Tại $x=3/2$, dùng đoạn $[1,2]$, $t=1/2$. Horner đồng thời cho giá trị và đạo hàm:
\begin{center}
\begin{tabular}{c|rrr}\toprule
$j$&$2$&$1$&$0$\\\midrule
$a_j$&$-2$&$3$&$2$\\
$t v_{j+1}$&&$-1$&$1$\\
$v_j$&$-2$&$2$&$\boxed{3}$\\
$u_j$&$0$&$-2$&$\boxed{1}$\\\bottomrule
\end{tabular}
\end{center}
Ở đây $v_j=a_j+tv_{j+1}$, $u_j=v_{j+1}+tu_{j+1}$; vì $dt/dx=1$, $S(3/2)=3$ và $S'(3/2)=1$.
\end{vidu}
% END SPL-GENERAL-QUADRATIC

\subsection{Spline bậc ba và các điều kiện biên}
Khi yêu cầu cả giá trị, đạo hàm cấp một và đạo hàm cấp hai liên tục qua các mốc, ta tìm các đoạn đa thức bậc ba. Một spline bậc ba trên các mốc đã cho là hàm $S\in C^2([a,b])$ sao cho $S|_{[x_i,x_{i+1}]}\in\mathcal P_3$. Điều kiện nội suy là $S(x_i)=y_i$. Đặt $M_i=S''(x_i)$. Hai lựa chọn điều kiện biên được xét là
\[
 M_0=M_N=0\quad\text{(tự nhiên)},\qquad
 S'(a)=\alpha,\quad S'(b)=\beta\quad\text{(ấn định đạo hàm ở biên)}.
\]
\begin{menhde}{}{spl-local-cubic}
Với các số $M_0,\ldots,M_N$, đa thức trên đoạn $[x_i,x_{i+1}]$ có giá trị $y_i,y_{i+1}$ và đạo hàm cấp hai $M_i,M_{i+1}$ tại hai đầu được xác định duy nhất bởi
\begin{equation}\label{eq:spl-cubic}
 \begin{aligned}
 S_i(x)&=\frac{M_i(x_{i+1}-x)^3+M_{i+1}(x-x_i)^3}{6h_i}\\
 &\quad+\left(y_i-\frac{M_ih_i^2}6\right)\frac{x_{i+1}-x}{h_i}
 +\left(y_{i+1}-\frac{M_{i+1}h_i^2}6\right)\frac{x-x_i}{h_i}.
 \end{aligned}
\end{equation}
Với $t=x-x_i$, công thức Horner tương ứng là
\begin{equation}\label{eq:spl-horner}
 \begin{aligned}
 S_i(x)&=a_i+t\bigl(b_i+t(c_i+td_i)\bigr),\\
 a_i&=y_i,\quad b_i=\delta_i-\frac{h_i}6(2M_i+M_{i+1}),\quad
 c_i=\frac{M_i}2,\quad d_i=\frac{M_{i+1}-M_i}{6h_i}.
 \end{aligned}
\end{equation}
Các đạo hàm tại hai đầu là
\[
 S_i'(x_i)=\delta_i-\frac{h_i}6(2M_i+M_{i+1}),\qquad
 S_i'(x_{i+1})=\delta_i+\frac{h_i}6(M_i+2M_{i+1}).
\]
\end{menhde}
\begin{chungminh}
Viết $S_i=a_i+b_it+c_it^2+d_it^3$. Bốn điều kiện cho
\[
 a_i=y_i,\quad 2c_i=M_i,\quad 2c_i+6d_ih_i=M_{i+1},\quad
 a_i+b_ih_i+c_ih_i^2+d_ih_i^3=y_{i+1}.
\]
Do đó
\[
 d_i=\frac{M_{i+1}-M_i}{6h_i},\qquad
 b_i=\delta_i-\frac{M_ih_i}2-\frac{(M_{i+1}-M_i)h_i}6
 =\delta_i-\frac{h_i}6(2M_i+M_{i+1}).
\]
Với $x=x_i+t$, vế phải của \eqref{eq:spl-cubic} bằng
\begin{align*}
 &\frac{M_i(h_i-t)^3+M_{i+1}t^3}{6h_i}
 +\left(y_i-\frac{M_ih_i^2}6\right)\frac{h_i-t}{h_i}
 +\left(y_{i+1}-\frac{M_{i+1}h_i^2}6\right)\frac{t}{h_i}\\
 &=\frac{M_i(h_i^3-3h_i^2t+3h_it^2-t^3)+M_{i+1}t^3}{6h_i}
 +y_i-\frac{M_ih_i^2}6
 +\left(\delta_i+\frac{(M_i-M_{i+1})h_i}6\right)t\\
 &=y_i+\left(\delta_i-\frac{h_i}6(2M_i+M_{i+1})\right)t
 +\frac{M_i}2t^2+\frac{M_{i+1}-M_i}{6h_i}t^3.
\end{align*} Ngoài ra,
\[
 S_i'=b_i+2c_it+3d_it^2,\qquad
 S_i''=M_i+\frac{M_{i+1}-M_i}{h_i}t.
\]
\[
 S_i'(x_i)=b_i=\delta_i-\frac{h_i}6(2M_i+M_{i+1}),
\]
\begin{align*}
 S_i'(x_{i+1})
 &=b_i+2c_ih_i+3d_ih_i^2\\
 &=\delta_i-\frac{h_i}6(2M_i+M_{i+1})
   +h_iM_i+\frac{h_i}2(M_{i+1}-M_i)\\
 &=\delta_i+\frac{h_i}6(M_i+2M_{i+1}).
\end{align*}
\end{chungminh}

\begin{dinhly}{}{spl-existence}
Với mọi dữ liệu $y_0,\ldots,y_N$, tồn tại duy nhất spline nội suy tự nhiên. Với mọi $\alpha,\beta\in\R$, tồn tại duy nhất spline nội suy có đạo hàm biên $\alpha,\beta$.

Trong cả hai trường hợp, các mốc trong thỏa hệ tam đường chéo
\begin{equation}\label{eq:spl-interior-system}
 h_{i-1}M_{i-1}+2(h_{i-1}+h_i)M_i+h_iM_{i+1}
 =6(\delta_i-\delta_{i-1}),\qquad 1\leq i<N.
\end{equation}
Trường hợp tự nhiên có $M_0=M_N=0$. Trường hợp ấn định đạo hàm có hai phương trình bổ sung
\begin{equation}\label{eq:spl-clamped-system}
 2h_0M_0+h_0M_1=6(\delta_0-\alpha),\qquad
 h_{N-1}M_{N-1}+2h_{N-1}M_N=6(\beta-\delta_{N-1}).
\end{equation}
\end{dinhly}
\begin{chungminh}
Các công thức cục bộ cho
\[
 S_{i-1}(x_i)=S_i(x_i)=y_i,\qquad
 S_{i-1}''(x_i)=S_i''(x_i)=M_i.
\]
Đối với đạo hàm cấp một,
\begin{align*}
 S_{i-1}'(x_i)=S_i'(x_i)
 &\ \Longleftrightarrow\
 \delta_{i-1}+\frac{h_{i-1}}6(M_{i-1}+2M_i)
 =\delta_i-\frac{h_i}6(2M_i+M_{i+1})\\
 &\ \Longleftrightarrow\
 h_{i-1}M_{i-1}+2(h_{i-1}+h_i)M_i+h_iM_{i+1}
 =6(\delta_i-\delta_{i-1}).
\end{align*}
Tại hai đầu,
\begin{align*}
 S_0'(a)=\alpha&\ \Longleftrightarrow\ \delta_0-\frac{h_0}6(2M_0+M_1)=\alpha\\
 &\ \Longleftrightarrow\ 2h_0M_0+h_0M_1=6(\delta_0-\alpha).
\end{align*}
\begin{align*}
 S_{N-1}'(b)=\beta&\ \Longleftrightarrow\ \delta_{N-1}+\frac{h_{N-1}}6(M_{N-1}+2M_N)=\beta\\
 &\ \Longleftrightarrow\ h_{N-1}M_{N-1}+2h_{N-1}M_N=6(\beta-\delta_{N-1}).
\end{align*}
Gọi $T$ là ma trận của hệ ấn định đạo hàm, có
\[
 T_{00}=2h_0,\quad T_{NN}=2h_{N-1},\quad
 T_{ii}=2(h_{i-1}+h_i)\ (1\leq i<N),\quad
 T_{i,i+1}=T_{i+1,i}=h_i.
\]
Đặt $h_{\min}=\min_i h_i>0$. Với $\mathbf z=(z_0,\ldots,z_N)^T$,
\begin{align*}
 \mathbf z^TT\mathbf z
 &=\sum_{i=0}^{N-1}2h_i(z_i^2+z_iz_{i+1}+z_{i+1}^2)\\
 &=\sum_{i=0}^{N-1}h_i\bigl((z_i+z_{i+1})^2+z_i^2+z_{i+1}^2\bigr)\\
 &\geq h_{\min}\sum_{i=0}^{N-1}(z_i^2+z_{i+1}^2)
 \geq h_{\min}\sum_{i=0}^Nz_i^2.
\end{align*}
Do đó
\[
 T\mathbf z=0\ \Longrightarrow\
 0=\mathbf z^TT\mathbf z\geq h_{\min}\|\mathbf z\|_2^2
 \ \Longrightarrow\ \mathbf z=0,
\]
\[
 \operatorname{rank}T=(N+1)-\dim\ker T=N+1.
\]
Hệ có duy nhất $\mathbf M=T^{-1}\mathbf R$.

Với điều kiện tự nhiên và $N\geq2$, đặt $\mathbf z=(0,\mathbf z_\circ^T,0)^T$ cho ma trận con chính $T_\circ$. Khi đó
\[
 \mathbf z_\circ^TT_\circ\mathbf z_\circ
 =\mathbf z^TT\mathbf z
 \geq h_{\min}\|\mathbf z_\circ\|_2^2,
\]
\[
 T_\circ\mathbf z_\circ=0\ \Longrightarrow\ \mathbf z_\circ=0,
 \qquad \operatorname{rank}T_\circ=N-1.
\]
Vì vậy $\mathbf M_\circ=T_\circ^{-1}\mathbf R_\circ$ duy nhất, với $M_0=M_N=0$. Nếu $N=1$, chỉ có hai moment biên $M_0=M_1=0$.

Các đẳng thức ghép ở đầu chứng minh cho $S,S',S''$ liên tục tại mọi mốc. Lập luận tích phân ở định lý \ref{dl:spl-initial-jet} cho $S\in C^2$; các công thức biên cho đúng điều kiện đã chọn. Nếu $\widetilde S$ là nghiệm khác thì moment $\widetilde{\mathbf M}$ thỏa cùng hệ, nên
\[
 T(\widetilde{\mathbf M}-\mathbf M)=0
 \ \Longrightarrow\ \widetilde{\mathbf M}=\mathbf M
\]
trong trường hợp ấn định đạo hàm; trong trường hợp tự nhiên dùng $T_\circ$ và các moment biên bằng nhau. Mệnh đề \ref{md:spl-local-cubic} cho $\widetilde S_i=S_i$ trên mọi đoạn, tức $\widetilde S=S$.
\end{chungminh}
\begin{vidu}
Với các số đo $(x_i,y_i)=(0,0),(1,1),(2,0)$ và điều kiện tự nhiên,
\[
 h_0=h_1=1,\quad\delta_0=1,\quad\delta_1=-1,\quad M_0=M_2=0,\quad4M_1=-12,\quad M_1=-3.
\]
Các hệ số theo \eqref{eq:spl-horner} là
\begin{center}\begin{tabular}{c|rrrr}\toprule
$i$&$a_i$&$b_i$&$c_i$&$d_i$\\\midrule
$0$&$0$&$3/2$&$0$&$-1/2$\\
$1$&$1$&$0$&$-3/2$&$1/2$\\\bottomrule
\end{tabular}\end{center}
\[
 S(x)=\begin{cases}x(3/2-x^2/2),&0\leq x\leq1,\\
 1+(x-1)^2(-3/2+(x-1)/2),&1\leq x\leq2.
 \end{cases}
\]
Đặt $a_{0,0}=a_0$, $a_{0,1}=b_0$, $a_{0,2}=c_0$, $a_{0,3}=d_0$. Horner tại $x=1/2$ trên đoạn đầu cho
\begin{center}\begin{tabular}{c|rrrr}\toprule
$j$&$3$&$2$&$1$&$0$\\\midrule
$a_{0,j}$&$-1/2$&$0$&$3/2$&$0$\\
$(1/2)v_{j+1}$&---&$-1/4$&$-1/8$&$11/16$\\\midrule
$v_j$&$-1/2$&$-1/4$&$11/8$&$\boxed{11/16}$\\\bottomrule
\end{tabular}\end{center}
Ta có $S(1/2)=S(3/2)=11/16$ và $S'(1)=0$.

Với một đoạn $[0,1]$, dữ liệu $0,1$ và đạo hàm biên $0,3$, hệ \eqref{eq:spl-clamped-system} là
\[
 2M_0+M_1=6,\qquad M_0+2M_1=12.
\]
Khử phương trình cho $M_1=6$, $M_0=0$; \eqref{eq:spl-horner} cho $S(x)=x^3$.
\end{vidu}

% BEGIN SPL-PRESCRIBED-MOMENTS
% Ghép sau định lý và ví dụ tồn tại duy nhất spline bậc ba, trước phần giải hệ.
\Needspace{7\baselineskip}
Các phương trình ghép tại mốc trong đoạn chỉ xác định $N-1$ điều kiện cho $N+1$ moment. Ngoài việc cho hai đạo hàm bậc nhất, có thể cho hai đạo hàm bậc hai tại hai đầu đoạn. Spline tự nhiên là trường hợp cả hai giá trị này bằng $0$.
\begin{dinhly}{}{spl-prescribed-moments}
Cho $A,B\in\R$. Có duy nhất spline bậc ba nội suy dữ liệu và thỏa
\[
 S''(x_0)=A,\qquad S''(x_N)=B.
\]
Với $N\geq2$, đặt $M_0=A$, $M_N=B$. Các moment còn lại là nghiệm của hệ đối xứng xác định dương
\begin{equation}\label{eq:spl-prescribed-moment-system}
 h_{i-1}M_{i-1}+2(h_{i-1}+h_i)M_i+h_iM_{i+1}
 =6(\delta_i-\delta_{i-1}),\quad1\leq i<N.
\end{equation}
Vế phải sau khi loại hai moment biên là
\[
 R_i=6(\delta_i-\delta_{i-1})
 -\delta_{i1}h_0A-\delta_{i,N-1}h_{N-1}B.
\]
Ở đây $\delta_{ij}$ là ký hiệu Kronecker; $\delta_i$ vẫn là độ dốc dây cung đã định nghĩa.
\end{dinhly}
\begin{chungminh}
Theo mệnh đề \ref{md:spl-local-cubic},
\begin{align*}
 S_{i-1}'(x_i)=S_i'(x_i)
 &\ \Longleftrightarrow\
 \delta_{i-1}+\frac{h_{i-1}}6(M_{i-1}+2M_i)
 =\delta_i-\frac{h_i}6(2M_i+M_{i+1})\\
 &\ \Longleftrightarrow\
 h_{i-1}M_{i-1}+2(h_{i-1}+h_i)M_i+h_iM_{i+1}
 =6(\delta_i-\delta_{i-1}).
\end{align*}
Thế $M_0=A$, $M_N=B$ vào các hàng trong cho
\[
 H\mathbf M_\circ=\mathbf R,\qquad
 R_i=6(\delta_i-\delta_{i-1})
 -\delta_{i1}h_0A-\delta_{i,N-1}h_{N-1}B.
\]
Với $N\geq2$, đặt $z_0=z_N=0$ và $h_{\min}=\min_i h_i>0$. Khi đó
\begin{align*}
 \mathbf z_\circ^TH\mathbf z_\circ
 &=\sum_{i=0}^{N-1}h_i\bigl((z_i+z_{i+1})^2+z_i^2+z_{i+1}^2\bigr)\\
 &\geq h_{\min}\sum_{i=1}^{N-1}z_i^2,
\end{align*}
\[
 H\mathbf z_\circ=0\ \Longrightarrow\ \mathbf z_\circ=0,
 \qquad \operatorname{rank}H=N-1,\qquad
 \mathbf M_\circ=H^{-1}\mathbf R.
\]
Từ các moment này, công thức \eqref{eq:spl-cubic} cho các đoạn có cùng giá trị, đạo hàm cấp một và cấp hai tại các mốc. Phép ghép đã chứng minh trong định lý \ref{dl:spl-initial-jet} cho $S\in C^2$, $S(x_i)=y_i$, $S''(x_0)=A$, $S''(x_N)=B$. Với hai nghiệm,
\[
 H(\widetilde{\mathbf M}_\circ-\mathbf M_\circ)=0
 \ \Longrightarrow\ \widetilde{\mathbf M}_\circ=\mathbf M_\circ,
 \qquad \widetilde M_0=M_0,\quad\widetilde M_N=M_N.
\]
Mệnh đề \ref{md:spl-local-cubic} cho $\widetilde S_i=S_i$ trên từng đoạn. Nếu $N=1$, chính mệnh đề đó xác định duy nhất đa thức từ $y_0,y_1,A,B$.
\end{chungminh}
\Needspace{9\baselineskip}
\begin{vidu}
Cho $\boldsymbol x=(0,1,2)^T$, $\boldsymbol y=(0,1,8)^T$, $A=0$, $B=12$. Phương trình duy nhất là
\[
 0+4M_1+12=6(7-1),\qquad M_1=6.
\]
Bảng hệ số theo $t=x-x_i$:
\begin{center}
\begin{tabular}{c|rrrr}\toprule
Đoạn&$[t^3]S_i$&$[t^2]S_i$&$[t]S_i$&$[1]S_i$\\\midrule
$[0,1]$&$1$&$0$&$0$&$0$\\
$[1,2]$&$1$&$3$&$3$&$1$\\\bottomrule
\end{tabular}
\end{center}
Trên đoạn thứ hai, $S_i(1+t)=1+t(3+t(3+t))$. Horner tại $t=1/2$ cho lần lượt $1,7/2,19/4,27/8$; vậy $S(3/2)=27/8$. Hai đoạn là các hạn chế của $x^3$, nên $S''(0)=0$, $S''(2)=12$.
\end{vidu}
% END SPL-PRESCRIBED-MOMENTS

\begin{dinhly}{}{spl-thomas}
Cho ma trận đối xứng xác định dương $T\in\R^{q\times q}$, $q\geq1$, có $T_{ii}=d_i$, $T_{i,i+1}=T_{i+1,i}=e_i$ và các phần tử còn lại bằng $0$. Với $T\mathbf z=\mathbf r$, đặt
\begin{equation}\label{eq:spl-thomas}
 \begin{aligned}
 D_0&=d_0,&\rho_0&=r_0,\\
 D_i&=d_i-e_{i-1}^2/D_{i-1},&
 \rho_i&=r_i-e_{i-1}\rho_{i-1}/D_{i-1}\quad(1\leq i<q).
 \end{aligned}
\end{equation}
Mọi $D_i>0$. Nghiệm duy nhất được xác định bởi
\[
 z_{q-1}=\rho_{q-1}/D_{q-1},\qquad
 z_i=(\rho_i-e_iz_{i+1})/D_i\quad(i=q-2,\ldots,0).
\]
\end{dinhly}
\begin{chungminh}
Với một ma trận đối xứng xác định dương chia khối
\[
 A=\begin{pmatrix}d&\mathbf v^T\\\mathbf v&B\end{pmatrix},
\]
ta có $d=(1,\mathbf0^T)A(1,\mathbf0^T)^T>0$. Với $\mathbf u\ne0$, đặt $\mathbf w=(-\mathbf v^T\mathbf u/d,\mathbf u^T)^T\ne0$. Khi đó
\[
 \mathbf w^TA\mathbf w
 =\frac{(\mathbf v^T\mathbf u)^2}{d}-2\frac{(\mathbf v^T\mathbf u)^2}{d}+\mathbf u^TB\mathbf u
 =\mathbf u^T\left(B-\frac{\mathbf v\mathbf v^T}{d}\right)\mathbf u>0.
\]
Vì vậy $B-\mathbf v\mathbf v^T/d$ xác định dương. Đặt $T^{(0)}=T$. Nếu $i<q-1$, ma trận còn lại ở bước $i$ có dạng
\[
 (T^{(i)})_{00}=D_i,\qquad
 (T^{(i)})_{jj}=d_{i+j}\quad(1\leq j<q-i),
\]
\[
 (T^{(i)})_{j,j+1}=(T^{(i)})_{j+1,j}=e_{i+j}
 \quad(0\leq j<q-i-1),\qquad
 (T^{(i)})_{jk}=0\quad(|j-k|>1).
\]
Viết
\[
 T^{(i)}=
 \begin{pmatrix}D_i&\mathbf v_i^T\\\mathbf v_i&B_i\end{pmatrix},
 \qquad \mathbf v_i=(e_i,0,\ldots,0)^T.
\]
Nếu $T^{(i)}$ xác định dương thì $D_i>0$, và phép tính chia khối ở trên cho ma trận xác định dương
\[
 T^{(i+1)}
 =B_i-\frac{\mathbf v_i\mathbf v_i^T}{D_i},\qquad
 (T^{(i+1)})_{00}=d_{i+1}-\frac{e_i^2}{D_i}=D_{i+1}.
\]
Các phần tử khác của $B_i$ giữ nguyên vì
\[
 (\mathbf v_i\mathbf v_i^T)_{jk}
 =e_i^2\delta_{j0}\delta_{k0}.
\]
Quy nạp theo $i$ cho $T^{(i)}$ xác định dương, và
\[
 D_i=(1,\mathbf0^T)T^{(i)}(1,\mathbf0^T)^T>0
 \quad(0\leq i<q).
\]

Quy ước $e_{q-1}=0$, $z_q=0$ trong các công thức ở hàng cuối. Khử phương trình $i-1$ khỏi phương trình $i$ bằng hệ số $e_{i-1}/D_{i-1}$ cho
\[
 d_i z_i+e_{i-1}z_{i-1}+e_i z_{i+1}
 -\frac{e_{i-1}}{D_{i-1}}(D_{i-1}z_{i-1}+e_{i-1}z_i)
 =r_i-\frac{e_{i-1}}{D_{i-1}}\rho_{i-1},
\]
\[
 D_i z_i+e_i z_{i+1}=\rho_i.
\]
Hệ đã khử là
\[
 D_i z_i+e_i z_{i+1}=\rho_i\quad(0\leq i<q),\qquad e_{q-1}=0.
\]
Vì $D_i>0$, hàng cuối và các hàng đi ngược cho
\[
 z_{q-1}=\rho_{q-1}/D_{q-1},\qquad
 z_i=(\rho_i-e_i z_{i+1})/D_i.
\]
Mỗi phép khử thay phương trình thứ $i$ bằng hiệu của nó với $e_{i-1}/D_{i-1}$ lần phương trình thứ $i-1$; phép ngược là cộng lại cùng bội phương trình. Do đó hệ trước và sau khử có cùng tập nghiệm; các công thức thế ngược cho đúng một nghiệm.
\end{chungminh}

\begin{dinhly}{}{spl-energy}
Với spline tự nhiên $S$, mọi $g\in C^2([a,b])$ thỏa $g(x_i)=y_i$ đều có
\begin{equation}\label{eq:spl-energy}
 \int_a^b(g'')^2\,dx
 =\int_a^b(S'')^2\,dx+\int_a^b(g''-S'')^2\,dx.
\end{equation}
Do đó $S$ là hàm duy nhất có năng lượng $\int_a^b(g'')^2dx$ nhỏ nhất trong lớp trên. Với spline ấn định đạo hàm, cùng kết quả đúng trong lớp bổ sung $g'(a)=\alpha$, $g'(b)=\beta$.
\end{dinhly}
\begin{chungminh}
Đặt $E=g-S$, nên $E(x_i)=0$. Trên mỗi đoạn, tích phân từng phần hai lần và $S_i^{(4)}=0$ cho
\[
 \int_{x_i}^{x_{i+1}}S_i''E''\,dx
 =[S_i''E'-S_i'''E]_{x_i}^{x_{i+1}}
 =S_i''(x_{i+1})E'(x_{i+1})-S_i''(x_i)E'(x_i).
\]
Cộng theo $i$:
\begin{align*}
 \int_a^bS''E''\,dx
 &=S''(b)E'(b)-S''(a)E'(a)\\
 &\quad+\sum_{i=1}^{N-1}
 \bigl(S_{i-1}''(x_i)-S_i''(x_i)\bigr)E'(x_i)\\
 &=S''(b)E'(b)-S''(a)E'(a)=0.
\end{align*}
Đẳng thức cuối đúng do $S''(a)=S''(b)=0$ trong trường hợp tự nhiên, hoặc do $E'(a)=E'(b)=0$ trong trường hợp ấn định đạo hàm. Vì $g''=S''+E''$,
\[
 \int_a^b(g'')^2dx=\int_a^b(S'')^2dx+2\int_a^bS''E''dx+\int_a^b(E'')^2dx,
\]
cho \eqref{eq:spl-energy}, và
\[
 \int_a^b(g'')^2dx-\int_a^b(S'')^2dx
 =\int_a^b(E'')^2dx\geq0.
\]
Nếu hai năng lượng bằng nhau mà $E''(c)\ne0$, tính liên tục cho một đoạn $I\subseteq[a,b]$, $|I|>0$, trên đó $|E''|\geq|E''(c)|/2$. Khi đó
\[
 0=\int_a^b(E'')^2dx
 \geq\int_I(E'')^2dx
 \geq\frac{|I||E''(c)|^2}4>0,
\]
mâu thuẫn. Vì vậy $E''=0$ và
\[
 E'(x)=E'(a)+\int_a^xE''(t)dt=E'(a),\qquad
 E(x)=E(a)+(x-a)E'(a).
\]
Từ $E(a)=E(b)=0$,
\[
 0=(b-a)E'(a),\qquad b-a>0
 \ \Longrightarrow\ E'(a)=0
 \ \Longrightarrow\ E(x)=0.
\]
Do đó $g=S$ là trường hợp duy nhất có năng lượng bằng năng lượng của $S$.
\end{chungminh}
\begin{vidu}
Với spline tự nhiên của các số đo $0,1,0$ tại $0,1,2$,
\[
 S''(x)=\begin{cases}-3x,&0\leq x\leq1,\\3(x-2),&1\leq x\leq2,\end{cases}
\]
\[
 \int_0^2(S'')^2dx=9\left[\frac{x^3}3\right]_0^1+9\left[\frac{(x-2)^3}3\right]_1^2=6.
\]
Mọi $g\in C^2([0,2])$ có cùng ba giá trị có năng lượng $6+\int_0^2(g''-S'')^2dx$.
\end{vidu}

\begin{bode}{}{spl-dirichlet-bound}
Nếu $u\in C^2([l,r])$, $l<r$, $u(l)=u(r)=0$, thì
\[
 |u(x)|\leq\frac{(x-l)(r-x)}2\|u''\|_{\infty,[l,r]},\qquad
 \|u\|_{\infty,[l,r]}\leq\frac{(r-l)^2}8\|u''\|_{\infty,[l,r]}.
\]
\end{bode}
\begin{chungminh}
Tích phân từng phần cho
\[
 \int_l^x(x-t)u''(t)dt=u(x)-u(l)-(x-l)u'(l),
\]
\[
 u'(l)=-\frac1{r-l}\int_l^r(r-t)u''(t)dt.
\]
Thế vào công thức đầu:
\[
 u(x)=-\frac{r-x}{r-l}\int_l^x(t-l)u''(t)dt
      -\frac{x-l}{r-l}\int_x^r(r-t)u''(t)dt.
\]
Hai nhân tử trước tích phân không âm, nên
\[
 |u(x)|\leq\|u''\|_\infty\left(\frac{(r-x)(x-l)^2}{2(r-l)}+
 \frac{(x-l)(r-x)^2}{2(r-l)}\right)
 =\frac{(x-l)(r-x)}2\|u''\|_\infty.
\]
\[
 (x-l)(r-x)
 =\frac{(r-l)^2}4-\left(x-\frac{l+r}2\right)^2
 \leq\frac{(r-l)^2}4,
\]
\[
 \|u\|_{\infty,[l,r]}
 \leq\max_{x\in[l,r]}\frac{(x-l)(r-x)}2\|u''\|_{\infty,[l,r]}
 \leq\frac{(r-l)^2}8\|u''\|_{\infty,[l,r]}.
\]
\end{chungminh}

\begin{dinhly}{}{spl-error}
Cho $f\in C^4([a,b])$, $y_i=f(x_i)$, $F=\|f^{(4)}\|_\infty$. Nếu $S$ là spline ấn định đạo hàm với $\alpha=f'(a)$, $\beta=f'(b)$, thì
\begin{equation}\label{eq:spl-error-clamped}
 \|f-S\|_\infty\leq\frac7{64}h^4F.
\end{equation}
Nếu $S$ là spline tự nhiên và $B=\max\{|f''(a)|,|f''(b)|\}$, thì
\begin{equation}\label{eq:spl-error-natural}
 \|f-S\|_\infty\leq\frac{h^2}8B+\frac7{64}h^4F.
\end{equation}
Đặc biệt, spline tự nhiên thỏa chặn \eqref{eq:spl-error-clamped} khi $f''(a)=f''(b)=0$.
\end{dinhly}
\begin{chungminh}
Đặt $z_i=f''(x_i)$ và $E_i=M_i-z_i$. Tích phân từng phần lặp cho các công thức Taylor có dư tích phân
\[
 f(x+t)=\sum_{j=0}^3\frac{f^{(j)}(x)}{j!}t^j+R_x(t),\qquad
 |R_x(t)|\leq F|t|^4/24,
\]
\[
 f''(x+t)=f''(x)+tf'''(x)+U_x(t),\qquad |U_x(t)|\leq F|t|^2/2.
\]
Ở đây, với $t$ sao cho $x,x+t\in[a,b]$,
\[
 R_x(t)=\frac16\int_0^t(t-s)^3f^{(4)}(x+s)ds,\qquad
 U_x(t)=\int_0^t(t-s)f^{(4)}(x+s)ds.
\]
Các lần tích phân từng phần cho
\begin{align*}
 R_x(t)
 &=-\frac{t^3}6f^{(3)}(x)
   +\frac12\int_0^t(t-s)^2f^{(3)}(x+s)ds\\
 &=-\frac{t^3}6f^{(3)}(x)-\frac{t^2}2f''(x)
   +\int_0^t(t-s)f''(x+s)ds\\
 &=-\frac{t^3}6f^{(3)}(x)-\frac{t^2}2f''(x)-tf'(x)
   +\int_0^tf'(x+s)ds\\
 &=f(x+t)-f(x)-tf'(x)-\frac{t^2}2f''(x)-\frac{t^3}6f^{(3)}(x),\\
 U_x(t)&=-tf^{(3)}(x)+\int_0^tf^{(3)}(x+s)ds
        =f''(x+t)-f''(x)-tf^{(3)}(x).
\end{align*}
Đổi biến $s=tu$ cho
\[
 R_x(t)=\frac{t^4}6\int_0^1(1-u)^3f^{(4)}(x+tu)du,\qquad
 U_x(t)=t^2\int_0^1(1-u)f^{(4)}(x+tu)du,
\]
\[
 |R_x(t)|\leq\frac{F|t|^4}6\int_0^1(1-u)^3du=\frac{F|t|^4}{24},
 \qquad |U_x(t)|\leq F|t|^2\int_0^1(1-u)du=\frac{F|t|^2}2.
\] Với mốc trong, Taylor cho
\[
 \delta_i=f'(x_i)+\frac{h_i}2z_i+\frac{h_i^2}6f^{(3)}(x_i)
              +\frac{R_{x_i}(h_i)}{h_i},
\]
\[
 \delta_{i-1}=f'(x_i)-\frac{h_{i-1}}2z_i+\frac{h_{i-1}^2}6f^{(3)}(x_i)
              -\frac{R_{x_i}(-h_{i-1})}{h_{i-1}}.
\]
Đặt
\[
 v_i=6\left(\frac{R_{x_i}(h_i)}{h_i}
             +\frac{R_{x_i}(-h_{i-1})}{h_{i-1}}\right),\qquad
 u_i=h_{i-1}U_{x_i}(-h_{i-1})+h_iU_{x_i}(h_i).
\]
Khi đó
\[
 6(\delta_i-\delta_{i-1})
 =3(h_{i-1}+h_i)z_i+(h_i^2-h_{i-1}^2)f'''(x_i)+v_i,
 \qquad |v_i|\leq\frac F4(h_{i-1}^3+h_i^3),
\]
\begin{align*}
 h_{i-1}z_{i-1}+2(h_{i-1}+h_i)z_i+h_iz_{i+1}
 &=3(h_{i-1}+h_i)z_i\\
 &\quad+(h_i^2-h_{i-1}^2)f'''(x_i)+u_i,\\
 |u_i|&\leq\frac F2(h_{i-1}^3+h_i^3).
\end{align*}
Trừ hai hệ thức và dùng \eqref{eq:spl-interior-system}:
\[
 h_{i-1}E_{i-1}+2(h_{i-1}+h_i)E_i+h_iE_{i+1}=v_i-u_i,
 \quad |v_i-u_i|\leq\frac{3F}4h^2(h_{i-1}+h_i).
\]
Trong trường hợp ấn định đạo hàm chính xác, đặt
\begin{align*}
 v_0&=6R_a(h_0)/h_0,& u_0&=h_0U_a(h_0),\\
 v_N&=6R_b(-h_{N-1})/h_{N-1},& u_N&=h_{N-1}U_b(-h_{N-1}).
\end{align*}
Taylor tại $a$ cho
\[
 6(\delta_0-f'(a))=3h_0z_0+h_0^2f'''(a)+v_0,\quad |v_0|\leq Fh_0^3/4,
\]
\[
 2h_0z_0+h_0z_1=3h_0z_0+h_0^2f'''(a)+u_0,\quad |u_0|\leq Fh_0^3/2.
\]
Taylor tại $b$ với $t=-h_{N-1}$ cho
\[
 6(f'(b)-\delta_{N-1})=3h_{N-1}z_N-h_{N-1}^2f'''(b)+v_N,
 \quad |v_N|\leq Fh_{N-1}^3/4,
\]
\[
 h_{N-1}z_{N-1}+2h_{N-1}z_N=3h_{N-1}z_N-h_{N-1}^2f'''(b)+u_N,
 \quad |u_N|\leq Fh_{N-1}^3/2.
\]
Do đó hai hàng biên cho
\[
 |2h_0E_0+h_0E_1|\leq3Fh_0h^2/4,\qquad
 |h_{N-1}E_{N-1}+2h_{N-1}E_N|\leq3Fh_{N-1}h^2/4.
\]
Đặt $E=\max_i|E_i|$. Nếu cực đại tại mốc trong $i$, thì
\[
 (h_{i-1}+h_i)E
 \leq2(h_{i-1}+h_i)|E_i|-h_{i-1}|E_{i-1}|-h_i|E_{i+1}|
 \leq|v_i-u_i|
 \leq\frac{3F}4h^2(h_{i-1}+h_i).
\]
Nếu $|E_0|=E$ trong trường hợp ấn định đạo hàm,
\[
 h_0E
 \leq2h_0|E_0|-h_0|E_1|
 \leq|2h_0E_0+h_0E_1|
 \leq3Fh_0h^2/4.
\]
Nếu $|E_N|=E$ thì
\[
 h_{N-1}E
 \leq2h_{N-1}|E_N|-h_{N-1}|E_{N-1}|
 \leq|2h_{N-1}E_N+h_{N-1}E_{N-1}|
 \leq3Fh_{N-1}h^2/4.
\]
Trong trường hợp tự nhiên,
\[
 |E_0|=|z_0|=|f''(a)|\leq B,\qquad
 |E_N|=|z_N|=|f''(b)|\leq B.
\]
Do $h_i>0$, chia các bất đẳng thức tại một chỉ số đạt cực đại cho
\[
 E\leq3Fh^2/4\quad\text{(ấn định đạo hàm)},\qquad
 E\leq\max\{B,3Fh^2/4\}\quad\text{(tự nhiên)}.
\]

Gọi $T_i$ là đa thức \eqref{eq:spl-cubic} với $M_i,M_{i+1}$ được thay bằng $z_i,z_{i+1}$. Khi đó
\[
 (f-T_i)(x_i)=(f-T_i)(x_{i+1})=0,\qquad
 T_i''(x)=\frac{x_{i+1}-x}{h_i}f''(x_i)+\frac{x-x_i}{h_i}f''(x_{i+1}).
\]
Định lý \ref{dl:spl-linear} áp dụng cho $f''$ cho $\|(f-T_i)''\|_\infty\leq Fh_i^2/8$. Bổ đề \ref{bd:spl-dirichlet-bound} suy ra
\[
 \|f-T_i\|_{\infty,[x_i,x_{i+1}]}\leq Fh_i^4/64.
\]
Mặt khác,
\[
 (S_i-T_i)''(x)=\frac{x_{i+1}-x}{h_i}E_i+\frac{x-x_i}{h_i}E_{i+1},\qquad
 \|(S_i-T_i)''\|_\infty\leq E,
\]
\[
 (S_i-T_i)(x_i)=(S_i-T_i)(x_{i+1})=0,\qquad
 \|S_i-T_i\|_\infty\leq h_i^2E/8.
\]
Trên mỗi đoạn,
\[
 \|f-S_i\|_{\infty,[x_i,x_{i+1}]}
 \leq\|f-T_i\|_{\infty,[x_i,x_{i+1}]}
   +\|T_i-S_i\|_{\infty,[x_i,x_{i+1}]}
 \leq\frac{Fh_i^4}{64}+\frac{h_i^2E}8.
\]
Lấy cực đại theo $i$ cho
\[
 \|f-S\|_\infty\leq\frac{Fh^4}{64}+\frac{h^2E}8.
\]
Với điều kiện ấn định đạo hàm,
\[
 \|f-S\|_\infty
 \leq\frac{Fh^4}{64}+\frac{h^2}8\frac{3Fh^2}4
 =\frac7{64}Fh^4.
\]
Với điều kiện tự nhiên,
\[
 \|f-S\|_\infty
 \leq\frac{Fh^4}{64}+\frac{h^2}8\max\{B,3Fh^2/4\}
 \leq\frac{h^2B}8+\frac7{64}Fh^4.
\]
Nếu $f''(a)=f''(b)=0$, thì $B=0$ và chặn cuối bằng chặn trong trường hợp ấn định đạo hàm.
\end{chungminh}
\begin{vidu}
Với $f(x)=x^2$ trên $[0,1]$, lấy $x_i=i/N$, $h=1/N$, và spline tự nhiên $S_N$. Đặt $E_i=M_i-2$. Hệ moment cho
\[
 E_{i-1}+4E_i+E_{i+1}=0\quad(1\leq i<N),\qquad E_0=E_N=-2.
\]
Đặt $\lambda=-2+\sqrt3$, nên $-1<\lambda<0$ và $\lambda^2+4\lambda+1=0$. Các số
\[
 E_i=-2\frac{\lambda^i+\lambda^{N-i}}{1+\lambda^N}
\]
đúng ở hai biên và thỏa truy hồi bởi phương trình của $\lambda$; tính duy nhất của hệ cho đây là các sai số moment. Trên đoạn đầu, $(f-S_N)''$ tuyến tính với hai giá trị $-E_0,-E_1$. Công thức tích phân trong bổ đề \ref{bd:spl-dirichlet-bound} tại $h/2$ cho
\[
 (f-S_N)(h/2)=\frac{h^2}{16}(E_0+E_1),\qquad
 \lim_{N\to\infty}\frac{|(f-S_N)(h/2)|}{h^2}=\frac{1+\lambda}8>0.
\]
Vì thế không tồn tại hằng số $C$ độc lập $N$ sao cho $\|f-S_N\|_\infty\leq Ch^4$ với dữ liệu này.
\end{vidu}

\Needspace{9\baselineskip}
% BEGIN SPL-QUARTIC
% Ghép sau spline bậc ba, trước phần giải hệ hoặc sau phần sai số spline bậc ba.
\Needspace{7\baselineskip}
\subsection{Spline bậc bốn}
Với bậc bốn, ta phải ghép cả đạo hàm cấp $1,2,3$ tại mọi mốc trong đoạn. Đặt
\[
 m_i=S'(x_i),\qquad v_i=S''(x_i),\qquad w_i=S'''(x_i).
\]
Dữ liệu nội suy còn để tự do ba tham số. Một lựa chọn đủ xác định chúng là cho $m_0,v_0,w_0$.
\begin{menhde}{}{spl-quartic}
Với $S'(x_0)=\alpha$, $S''(x_0)=\beta$, $S'''(x_0)=\gamma$, có duy nhất spline bậc bốn thuộc $C^3$ nội suy dữ liệu. Các hệ thức dựng từng đoạn là
\begin{equation}\label{eq:spl-quartic}
 \begin{aligned}
 c_i&=\frac{y_{i+1}-y_i-h_im_i-h_i^2v_i/2-h_i^3w_i/6}{h_i^4},\\
 S_i(x_i+t)&=y_i+t\left(m_i+t\left(\frac{v_i}2+t\left(\frac{w_i}6+c_it\right)\right)\right),\\
 m_{i+1}&=m_i+h_iv_i+\frac{h_i^2}2w_i+4h_i^3c_i,\\
 v_{i+1}&=v_i+h_iw_i+12h_i^2c_i,\\
 w_{i+1}&=w_i+24h_ic_i.
 \end{aligned}
\end{equation}
\end{menhde}
\begin{chungminh}
Trong \eqref{eq:spl-initial-jet-coefficients}, thay $p=4$ cho
\[
 S_i(x_i+t)=y_i+m_it+\frac{v_i}2t^2+\frac{w_i}6t^3+c_it^4.
\]
Điều kiện $S_i(x_{i+1})=y_{i+1}$ cho công thức $c_i$. Ba đạo hàm là
\[
 S_i'=m_i+v_it+\frac{w_i}2t^2+4c_it^3,\qquad
 S_i''=v_i+w_it+12c_it^2,\qquad S_i'''=w_i+24c_it.
\]
\[
 S_i'(x_{i+1})=m_i+h_iv_i+\frac{h_i^2}2w_i+4h_i^3c_i=m_{i+1},
\]
\[
 S_i''(x_{i+1})=v_i+h_iw_i+12h_i^2c_i=v_{i+1},\qquad
 S_i'''(x_{i+1})=w_i+24h_ic_i=w_{i+1}.
\]
Các công thức này chính là \eqref{eq:spl-initial-jet-propagation} với $p=4$, $s_i^{(1)}=m_i$, $s_i^{(2)}=v_i$, $s_i^{(3)}=w_i$. Định lý \ref{dl:spl-initial-jet} cho tồn tại và duy nhất trong $C^3$ với dữ liệu đầu $(\alpha,\beta,\gamma)$.
\end{chungminh}
\Needspace{9\baselineskip}
\begin{vidu}
Cho $\boldsymbol x=(0,1,2)^T$, $\boldsymbol y=(0,1,16)^T$ và $m_0=v_0=w_0=0$. Truy hồi cho
\begin{center}
\begin{tabular}{c|rrr}\toprule
$i$&$0$&$1$&$2$\\\midrule
$y_i$&$0$&$1$&$16$\\
$m_i$&$0$&$4$&$32$\\
$v_i$&$0$&$12$&$48$\\
$w_i$&$0$&$24$&$48$\\
$c_i$&$1$&$1$&\\\bottomrule
\end{tabular}
\end{center}
Trên $[1,2]$, hệ số giảm bậc trong biến $t=x-1$ là $(1,4,6,4,1)$. Tại $t=1/2$:
\begin{center}
\begin{tabular}{c|rrrrr}\toprule
$j$&$4$&$3$&$2$&$1$&$0$\\\midrule
$a_j$&$1$&$4$&$6$&$4$&$1$\\
$t q_{j+1}$&&$1/2$&$9/4$&$33/8$&$65/16$\\
$q_j$&$1$&$9/2$&$33/4$&$65/8$&$\boxed{81/16}$\\
$u_j$&$0$&$1$&$5$&$43/4$&$\boxed{27/2}$\\\bottomrule
\end{tabular}
\end{center}
Vì vậy $S(3/2)=81/16$, $S'(3/2)=27/2$. Trên mỗi đoạn, các hệ số cũng cho đúng hạn chế của $x^4$.
\end{vidu}

\begin{menhde}{}{spl-jet-local-error}
Cho $f\in C^{p+1}([x_i,x_{i+1}])$, $y_j=f(x_j)$ ở hai đầu và $s_i^{(k)}$ là các đạo hàm đầu trái được dùng trong \eqref{eq:spl-initial-jet-coefficients}. Đặt
\[
 e_{i,k}=s_i^{(k)}-f^{(k)}(x_i),\quad
 F_i=\max_{[x_i,x_{i+1}]}|f^{(p+1)}|,\qquad1\leq k<p.
\]
Khi $0\leq t\leq h_i$,
\begin{equation}\label{eq:spl-jet-local-error}
 \left|f(x_i+t)-S_i(x_i+t)\right|
 \leq\frac{F_i}{(p+1)!}t^p(h_i-t)
 +\sum_{k=1}^{p-1}\frac{|e_{i,k}|}{k!}
       \left(t^k-h_i^{k-p}t^p\right).
\end{equation}
Đặc biệt, nếu các đạo hàm đầu trái đều đúng, thì sai số trên đoạn không quá
\[
 \frac{p^p}{(p+1)^{p+1}(p+1)!}F_i h_i^{p+1}.
\]
\end{menhde}
\begin{chungminh}
Gọi $T_i$ là đa thức cục bộ dùng các đạo hàm đúng của $f$ tại $x_i$ và giá trị $f(x_{i+1})$. Với $0<t<h_i$, đặt
\[
 K=\frac{f(x_i+t)-T_i(x_i+t)}{t^p(t-h_i)},\qquad
 R(z)=f(x_i+z)-T_i(x_i+z)-Kz^p(z-h_i).
\]
Ta có
\[
 R^{(k)}(0)=0\quad(0\leq k<p),\qquad R(t)=R(h_i)=0.
\]
Đặt $a_0=t$, $b_0=h_i$. Ta chứng minh bằng quy nạp rằng với mỗi $0\leq j<p$ có
\[
 0<a_j<b_j\leq h_i,\qquad
 R^{(j)}(0)=R^{(j)}(a_j)=R^{(j)}(b_j)=0.
\]
Khẳng định đúng với $j=0$. Nếu $j<p-1$, Rolle trên hai đoạn $[0,a_j]$ và $[a_j,b_j]$ cho
\[
 a_{j+1}\in(0,a_j),\quad b_{j+1}\in(a_j,b_j),\qquad
 R^{(j+1)}(a_{j+1})=R^{(j+1)}(b_{j+1})=0.
\]
Đồng thời $R^{(j+1)}(0)=0$ vì $j+1<p$, hoàn tất bước quy nạp. Tại $j=p-1$, áp dụng Rolle thêm hai lần:
\[
 A\in(0,a_{p-1}),\quad B\in(a_{p-1},b_{p-1}),\qquad
 R^{(p)}(A)=R^{(p)}(B)=0,
\]
\[
 \xi\in(A,B),\qquad R^{(p+1)}(\xi)=0.
\]
Lập luận này cũng áp dụng với $p=1$, khi $a_{p-1}=t$, $b_{p-1}=h_i$. Vì $\deg T_i\leq p$,
\[
 0=R^{(p+1)}(\xi)=f^{(p+1)}(x_i+\xi)-(p+1)!K,
\]
cho $|f-T_i|\leq F_it^p(h_i-t)/(p+1)!$. Công thức vẫn đúng ở $t=0,h_i$ vì cả hai vế bằng $0$.

Trừ công thức của $S_i$ và $T_i$, rồi dùng điều kiện bằng nhau ở $t=h_i$, được
\[
 S_i(x_i+t)-T_i(x_i+t)
 =\sum_{k=1}^{p-1}\frac{e_{i,k}}{k!}
   \left(t^k-h_i^{k-p}t^p\right).
\]
Với $0\leq t\leq h_i$ và $1\leq k<p$,
\[
 t^k-h_i^{k-p}t^p
 =t^k\left(1-\left(\frac{t}{h_i}\right)^{p-k}\right)\geq0.
\]
Do đó
\begin{align*}
 |f(x_i+t)-S_i(x_i+t)|
 &\leq|f(x_i+t)-T_i(x_i+t)|+|T_i(x_i+t)-S_i(x_i+t)|\\
 &\leq\frac{F_i}{(p+1)!}t^p(h_i-t)
 +\sum_{k=1}^{p-1}\frac{|e_{i,k}|}{k!}
   \left(t^k-h_i^{k-p}t^p\right).
\end{align*}
Đặt $\psi(t)=t^p(h_i-t)$ và $t_*=ph_i/(p+1)$. Khi $0<t<h_i$,
\[
 \psi'(t)=t^{p-1}(ph_i-(p+1)t)
 \begin{cases}>0,&t<t_*,\\=0,&t=t_*,\\<0,&t>t_*.\end{cases}
\]
Vì $\psi(0)=\psi(h_i)=0$,
\[
 \max_{0\leq t\leq h_i}\psi(t)
 =\psi(t_*)
 =\left(\frac{ph_i}{p+1}\right)^p\frac{h_i}{p+1}
 =\frac{p^ph_i^{p+1}}{(p+1)^{p+1}}.
\]
Nếu mọi $e_{i,k}=0$, thế cực đại này vào chặn đã chứng minh cho kết luận cuối.
\end{chungminh}
Chặn này chứa sai số của các đạo hàm được truyền từ đoạn trước. Chỉ cho đạo hàm đúng tại $x_0$ không làm các $e_{i,k}$ tự động bằng $0$ tại những mốc còn lại.
% END SPL-QUARTIC

\section{Phương pháp bình phương tối thiểu}
Các mô hình trong bài giảng \cite{ls-slides} được xây dựng từ cùng điều kiện trực giao của phần dư.
\subsection{Dữ liệu rời rạc và phương trình chuẩn}
Với nhiều số đo, việc buộc một mô hình ít hệ số đi qua mọi điểm có thể dẫn tới một hệ không có nghiệm. Thay điều kiện bằng nhau bởi bài toán cực tiểu, xét $m\geq1$ số đo $(x_i,y_i)$, $0\leq i<m$, và mô hình
\[
 p_{\mathbf a}(x)=\sum_{j=0}^{r}a_j\phi_j(x),\qquad
 \mathbf a=(a_0,\ldots,a_r)^T\in\R^{r+1}.
\]
Cho trọng số $w_i>0$. Độ sai lệch cần cực tiểu là
\begin{equation}\label{eq:ls-discrete-objective}
 J(\mathbf a)=\sum_{i=0}^{m-1}w_i\left(p_{\mathbf a}(x_i)-y_i\right)^2.
\end{equation}
Đặt
\[
 A_{ij}=\phi_j(x_i),\quad A\in\R^{m\times(r+1)},\quad
 W=\operatorname{diag}(w_0,\ldots,w_{m-1}),\quad \mathbf y=(y_0,\ldots,y_{m-1})^T,
\]
\[
 \|\mathbf u\|_W^2=\mathbf u^TW\mathbf u=\sum_iw_iu_i^2,\quad
 G=A^TWA,\quad\mathbf b=A^TW\mathbf y.
\]
\begin{dinhly}{}{ls-normal-equations}
Bài toán \eqref{eq:ls-discrete-objective} luôn có nghiệm. Một vector $\mathbf a_*$ là nghiệm khi và chỉ khi
\begin{equation}\label{eq:ls-normal-equations}
 G\mathbf a_*=\mathbf b,\qquad
 G_{jk}=\sum_{i=0}^{m-1}w_i\phi_j(x_i)\phi_k(x_i),\quad
 b_j=\sum_{i=0}^{m-1}w_iy_i\phi_j(x_i).
\end{equation}
Nghiệm duy nhất khi và chỉ khi $\operatorname{rank}A=r+1$; khi đó $G$ xác định dương và $\mathbf a_*=G^{-1}\mathbf b$. Nếu hạng không đầy đủ, mọi nghiệm là $\mathbf a_*+\ker A$, nhưng vector giá trị $A\mathbf a_*$ vẫn duy nhất. Với mọi $\mathbf a$,
\begin{equation}\label{eq:ls-discrete-pythagoras}
 J(\mathbf a)=J(\mathbf a_*)+\|A(\mathbf a-\mathbf a_*)\|_W^2.
\end{equation}
\end{dinhly}
\begin{chungminh}
Khai triển các bình phương cho
\[
 J(\mathbf a)=\mathbf a^TG\mathbf a-2\mathbf a^T\mathbf b+\mathbf y^TW\mathbf y,
 \qquad\mathbf z^TG\mathbf z=\sum_iw_i(A\mathbf z)_i^2\geq0.
\]
Với mọi $t\in\R$,
\[
 J(\mathbf a+t\mathbf z)-J(\mathbf a)
 =2t\mathbf z^T(G\mathbf a-\mathbf b)+t^2\mathbf z^TG\mathbf z.
\]
Nếu $c=\mathbf z^T(G\mathbf a-\mathbf b)\ne0$, đặt $g=\mathbf z^TG\mathbf z\geq0$, $t=-c/(g+1)$. Khi đó
\[
 J(\mathbf a+t\mathbf z)-J(\mathbf a)
 =-\frac{c^2(g+2)}{(g+1)^2}<0.
\]
Bởi vậy, tại một cực tiểu,
\[
 \mathbf z^T(G\mathbf a-\mathbf b)=0\quad\text{với mọi }\mathbf z\in\R^{r+1}.
\]
Chọn $\mathbf z=G\mathbf a-\mathbf b$ cho
\[
 \|G\mathbf a-\mathbf b\|_2^2=0
 \ \Longrightarrow\ G\mathbf a=\mathbf b.
\]

Vì mọi $w_i>0$,
\[
 \sum_iw_i(A\mathbf z)_i^2=0
 \ \Longleftrightarrow\
 (A\mathbf z)_i=0\text{ với mọi }i
 \ \Longleftrightarrow\ A\mathbf z=0.
\]
Do đó
\[
 G\mathbf z=0\ \Longrightarrow\
 \mathbf z^TG\mathbf z=\sum_iw_i(A\mathbf z)_i^2=0
 \ \Longrightarrow\ A\mathbf z=0
 \ \Longrightarrow\ G\mathbf z=A^TWA\mathbf z=0,
\]
\[
 \ker G=\ker A,
 \qquad \mathbf z\in\ker G\ \Longrightarrow\ \mathbf z^T\mathbf b=(A\mathbf z)^TW\mathbf y=0.
\]
Với $\mathbf v=G\mathbf u\in\operatorname{im}G$ và $\mathbf z\in\ker G$,
\[
 \mathbf z^T\mathbf v=\mathbf z^TG\mathbf u=(G\mathbf z)^T\mathbf u=0.
\]
Vì vậy $\operatorname{im}G\subseteq(\ker G)^\perp$. Công thức hạng--không gian nghiệm cho
\[
 \dim\operatorname{im}G=(r+1)-\dim\ker G=\dim(\ker G)^\perp,
\]
 nên $\operatorname{im}G=(\ker G)^\perp$. Vì $\mathbf b\perp\ker G$, tồn tại $\mathbf a_*$ với $G\mathbf a_*=\mathbf b$. Thế $\mathbf a=\mathbf a_*+\mathbf z$ vào khai triển $J$:
\[
 J(\mathbf a_*+\mathbf z)-J(\mathbf a_*)
 =2\mathbf z^T(G\mathbf a_*-\mathbf b)+\mathbf z^TG\mathbf z
 =\|A\mathbf z\|_W^2.
\]
Đây là \eqref{eq:ls-discrete-pythagoras}. Đặc biệt,
\[
 J(\mathbf a)\geq J(\mathbf a_*),\qquad
 J(\mathbf a)=J(\mathbf a_*)
 \ \Longleftrightarrow\
 A(\mathbf a-\mathbf a_*)=0.
\]
Tập cực tiểu là $\mathbf a_*+\ker A$. Nếu $\widetilde{\mathbf a}_*$ là cực tiểu khác,
\[
 A(\widetilde{\mathbf a}_*-\mathbf a_*)=0
 \ \Longrightarrow\ A\widetilde{\mathbf a}_*=A\mathbf a_*,
\]
nên vector giá trị duy nhất. Mặt khác,
\[
 \dim\ker A=(r+1)-\operatorname{rank}A,
\]
\[
 \mathbf a_*\text{ duy nhất}
 \ \Longleftrightarrow\ \ker A=\{0\}
 \ \Longleftrightarrow\ \operatorname{rank}A=r+1.
\]
Trong trường hợp hạng đầy đủ, với $\mathbf z\ne0$,
\[
 A\mathbf z\ne0
 \ \Longrightarrow\
 \mathbf z^TG\mathbf z=\sum_iw_i(A\mathbf z)_i^2>0.
\]
Do đó $\ker G=\{0\}$, $\operatorname{rank}G=r+1$ và $\mathbf a_*=G^{-1}\mathbf b$.
\end{chungminh}

\begin{hequa}{}{ls-polynomial-rank}
Nếu $x_0,\ldots,x_{m-1}$ đôi một phân biệt và $\phi_j(x)=x^j$, thì bài toán đa thức bậc không quá $r$ có nghiệm duy nhất khi và chỉ khi $r<m$. Khi $r=m-1$, nghiệm là đa thức nội suy. Khi $r\geq m$, sai lệch nhỏ nhất bằng $0$ và mọi nghiệm là
\[
 p(x)=I_X\mathbf y(x)+\omega_m(x)q(x),\qquad
 \omega_m(x)=\prod_{i=0}^{m-1}(x-x_i),\qquad q\in\mathcal P_{r-m},
\]
trong đó $I_X\mathbf y$ là đa thức nội suy bậc không quá $m-1$ của dữ liệu.
\end{hequa}
\begin{chungminh}
Nếu $r<m$ và $A\mathbf a=0$, đa thức $p_{\mathbf a}$ có $m>r$ nghiệm phân biệt. Bổ đề \ref{bd:roots} cho $p_{\mathbf a}=0$, tức $\mathbf a=0$. Do đó $A$ có hạng $r+1$. Nếu $r\geq m$, $\operatorname{rank}A\leq m<r+1$, nên nghiệm không duy nhất.

Đa thức nội suy có sai lệch $0$, nên trong trường hợp $r\geq m-1$, mọi cực tiểu thỏa $p(x_i)=y_i$. Đặt $v=p-I_X\mathbf y$. Chia Horner $v$ cho $x-x_0$ cho phần dư $v(x_0)=0$, nên $v=(x-x_0)v_1$. Với $i\geq1$,
\[
 0=v(x_i)=(x_i-x_0)v_1(x_i),\qquad x_i-x_0\ne0\quad\Longrightarrow\quad v_1(x_i)=0.
\]
Đặt $v_0=v$. Quy nạp cho, sau $k$ phép chia,
\[
 v(x)=\left(\prod_{j=0}^{k-1}(x-x_j)\right)v_k(x),\qquad
 v_k(x_i)=0\quad(k\leq i<m).
\]
Nếu $k<m$, phần dư khi chia $v_k$ cho $x-x_k$ là $v_k(x_k)=0$, nên
\[
 v_k=(x-x_k)v_{k+1},\qquad
 0=v_k(x_i)=(x_i-x_k)v_{k+1}(x_i)
 \ \Longrightarrow\ v_{k+1}(x_i)=0\quad(i>k).
\]
Tại $k=m$, đặt $q=v_m$ được $v=\omega_mq$. Với $r\geq m$, nếu $v\ne0$ thì
\[
 \deg q=\deg v-m\leq r-m;
\]
nếu $v=0$, chọn $q=0\in\mathcal P_{r-m}$. Ngược lại, với $q\in\mathcal P_{r-m}$,
\[
 \deg(I_X\mathbf y+\omega_mq)\leq r,\qquad
 (I_X\mathbf y+\omega_mq)(x_i)=y_i+0=y_i.
\]
Khi $r=m-1$, $\deg v\leq m-1$ và $v(x_i)=0$ tại $m$ mốc phân biệt; bổ đề \ref{bd:roots} cho $v=0$.
\end{chungminh}

% BEGIN LS-ERROR-MEASURES
% Ghép sau định lý phương trình chuẩn.
\Needspace{7\baselineskip}
Sau khi tìm hệ số, các sai lệch trên bảng được tính từ chính hàm thực nghiệm đã tìm được. Với $m$ quan sát và trọng số $w_i>0$, đặt
\begin{equation}\label{eq:ls-mean-square-error}
 r_i=y_i-g(x_i),\qquad
 Q=\sum_{i=0}^{m-1}w_ir_i^2,\qquad
 E_2=\frac{Q}{\sum_{i=0}^{m-1}w_i},\qquad E=\sqrt{E_2}.
\end{equation}
Nếu $w_i=1$, $E_2=m^{-1}\sum_i r_i^2$ là sai số trung bình phương và $E$ là căn sai số trung bình phương. Không thay mẫu số $m$ bằng số tham số khi tính các đại lượng này.
\begin{menhde}{}{ls-error-measures}
Ba tiêu chuẩn $Q,E_2,E$ có cùng tập điểm cực tiểu. Với thiết kế $A$ và nghiệm $\boldsymbol a_*$ của phương trình chuẩn, giá trị tối thiểu là
\begin{equation}\label{eq:ls-minimum-residual}
 Q_{\min}=\boldsymbol y^TW\boldsymbol y
 -\boldsymbol a_*^TA^TW\boldsymbol y.
\end{equation}
Nếu $A$ có hạng cột đầy đủ và $G=A^TWA$, thì
\[
 Q_{\min}=\boldsymbol y^TW\boldsymbol y
 -(A^TW\boldsymbol y)^TG^{-1}(A^TW\boldsymbol y).
\]
\end{menhde}
\begin{chungminh}
Đặt $s=\sum_iw_i>0$. Với $Q_1,Q_2\geq0$,
\[
 Q_1\leq Q_2\ \Longleftrightarrow\ Q_1/s\leq Q_2/s
 \ \Longleftrightarrow\ \sqrt{Q_1/s}\leq\sqrt{Q_2/s}.
\]
Do đó các cực tiểu trùng nhau. Khai triển bình phương tại $\boldsymbol a_*$ cho
\begin{align*}
 Q_{\min}
 &=\boldsymbol y^TW\boldsymbol y
 -2\boldsymbol a_*^TA^TW\boldsymbol y
 +\boldsymbol a_*^TG\boldsymbol a_*\\
 &=\boldsymbol y^TW\boldsymbol y-\boldsymbol a_*^TA^TW\boldsymbol y,
\end{align*}
vì $G\boldsymbol a_*=A^TW\boldsymbol y$. Nếu $G$ khả nghịch, đặt $\boldsymbol b=A^TW\boldsymbol y$. Vì $G^T=G$,
\[
 \boldsymbol a_*^T\boldsymbol b
 =(G^{-1}\boldsymbol b)^T\boldsymbol b
 =\boldsymbol b^T(G^{-1})^T\boldsymbol b
 =\boldsymbol b^TG^{-1}\boldsymbol b.
\]
Thế vào biểu thức của $Q_{\min}$ cho công thức còn lại.
\end{chungminh}
% END LS-ERROR-MEASURES

\begin{menhde}{}{ls-projection}
Giả sử $\operatorname{rank}A=r+1$. Đặt $P=AG^{-1}A^TW$. Khi đó
\[
 P^2=P,\qquad P^TW=WP,\qquad A^TW(\mathbf y-P\mathbf y)=0,
\]
\[
 \|\mathbf y\|_W^2=\|P\mathbf y\|_W^2+\|\mathbf y-P\mathbf y\|_W^2,
 \qquad\|P\mathbf e\|_W\leq\|\mathbf e\|_W.
\]
Nếu dữ liệu là $\mathbf y=\mathbf f+\boldsymbol\eta$, thì
\[
 \|P\mathbf y-P\mathbf f\|_W\leq\|\boldsymbol\eta\|_W,
\]
\begin{equation}\label{eq:ls-approx-noise}
 \|\mathbf f-P\mathbf y\|_W^2
 =\min_{\mathbf a}\|\mathbf f-A\mathbf a\|_W^2+\|P\boldsymbol\eta\|_W^2
 \leq\min_{\mathbf a}\|\mathbf f-A\mathbf a\|_W^2+\|\boldsymbol\eta\|_W^2.
\end{equation}
\end{menhde}
\begin{chungminh}
Vì $G^T=G$ và $(G^{-1})^T=G^{-1}$,
\[
 P^2=AG^{-1}(A^TWA)G^{-1}A^TW=P,\qquad
 P^TW=WA G^{-1}A^TW=WP,
\]
\[
 A^TW(I-P)=A^TW-GG^{-1}A^TW=0.
\]
Với mọi $\mathbf z$,
\[
 (P\mathbf z)^TW(\mathbf z-P\mathbf z)
 =\mathbf z^TWP(I-P)\mathbf z=0.
\]
\begin{align*}
 \|\mathbf z\|_W^2
 &=\|P\mathbf z\|_W^2
 +2(P\mathbf z)^TW(I-P)\mathbf z+\|(I-P)\mathbf z\|_W^2\\
 &=\|P\mathbf z\|_W^2+\|(I-P)\mathbf z\|_W^2.
\end{align*}
Đặt $\mathbf z=\mathbf e$:
\[
 0\leq\|P\mathbf e\|_W^2
 \leq\|P\mathbf e\|_W^2+\|(I-P)\mathbf e\|_W^2
 =\|\mathbf e\|_W^2
 \ \Longrightarrow\ \|P\mathbf e\|_W\leq\|\mathbf e\|_W.
\] Từ tính tuyến tính,
\[
 P\mathbf y-P\mathbf f=P\boldsymbol\eta,
 \qquad\mathbf f-P\mathbf y=(I-P)\mathbf f-P\boldsymbol\eta.
\]
Hai số hạng cuối trực giao theo $W$ vì
\[
 ((I-P)\mathbf f)^TWP\boldsymbol\eta
 =\mathbf f^T(I-P)^TWP\boldsymbol\eta
 =\mathbf f^TW(I-P)P\boldsymbol\eta=0.
\]
Với $\mathbf a_{\mathbf f}=G^{-1}A^TW\mathbf f$, ta có $A\mathbf a_{\mathbf f}=P\mathbf f$. Định lý \ref{dl:ls-normal-equations} cho
\[
 \|(I-P)\mathbf f\|_W^2
 =\|\mathbf f-A\mathbf a_{\mathbf f}\|_W^2
 =\min_{\mathbf a}\|\mathbf f-A\mathbf a\|_W^2.
\]
Vì hai số hạng trực giao,
\begin{align*}
 \|\mathbf f-P\mathbf y\|_W^2
 &=\|(I-P)\mathbf f\|_W^2+\|P\boldsymbol\eta\|_W^2\\
 &=\min_{\mathbf a}\|\mathbf f-A\mathbf a\|_W^2+\|P\boldsymbol\eta\|_W^2\\
 &\leq\min_{\mathbf a}\|\mathbf f-A\mathbf a\|_W^2+\|\boldsymbol\eta\|_W^2.
\end{align*}
Ngoài ra,
\[
 \|P\mathbf y-P\mathbf f\|_W
 =\|P\boldsymbol\eta\|_W\leq\|\boldsymbol\eta\|_W.
\]
\end{chungminh}

\begin{dinhly}{}{ls-cholesky}
Với ma trận đối xứng xác định dương $G\in\R^{q\times q}$, tồn tại duy nhất ma trận tam giác dưới $L$ có đường chéo dương sao cho $G=LL^T$. Các phần tử thỏa
\begin{equation}\label{eq:ls-cholesky}
 L_{kk}=\sqrt{G_{kk}-\sum_{j=0}^{k-1}L_{kj}^2},\qquad
 L_{ik}=\frac{G_{ik}-\sum_{j=0}^{k-1}L_{ij}L_{kj}}{L_{kk}}\quad(i>k).
\end{equation}
Hệ $G\mathbf a=\mathbf b$ được giải bằng
\[
\begin{aligned}
 z_i&=\frac{b_i-\sum_{j=0}^{i-1}L_{ij}z_j}{L_{ii}}\quad(i=0,\ldots,q-1),\\
 a_i&=\frac{z_i-\sum_{j=i+1}^{q-1}L_{ji}a_j}{L_{ii}}\quad(i=q-1,\ldots,0).
\end{aligned}
\]
\end{dinhly}
\begin{chungminh}
Với $q=1$, $G_{00}>0$ và $L_{00}=\sqrt{G_{00}}$. Với $q>1$, viết
\[
 G=\begin{pmatrix}d&\mathbf v^T\\\mathbf v&B\end{pmatrix},\qquad d>0.
\]
Phép chia khối trong chứng minh định lý \ref{dl:spl-thomas} cho $H=B-\mathbf v\mathbf v^T/d$ xác định dương. Theo quy nạp, $H=KK^T$ với $K$ tam giác dưới có đường chéo dương. Khi đó
\[
 L=\begin{pmatrix}\sqrt d&0\\\mathbf v/\sqrt d&K\end{pmatrix},\qquad
 LL^T=\begin{pmatrix}d&\mathbf v^T\\\mathbf v&\mathbf v\mathbf v^T/d+KK^T\end{pmatrix}=G.
\]
Nếu $\widetilde L$ là một phân tích khác có đường chéo dương, viết
\[
 \widetilde L=
 \begin{pmatrix}\ell&0\\\mathbf u&\widetilde K\end{pmatrix},
 \qquad \ell>0.
\]
So sánh từng khối của $\widetilde L\widetilde L^T=G$ cho
\[
 \ell^2=d,\qquad \ell\mathbf u=\mathbf v,\qquad
 \mathbf u\mathbf u^T+\widetilde K\widetilde K^T=B.
\]
Do đó
\[
 \ell=\sqrt d,\qquad\mathbf u=\mathbf v/\sqrt d,\qquad
 \widetilde K\widetilde K^T=B-\mathbf v\mathbf v^T/d=H.
\]
Giả thiết quy nạp cho $\widetilde K=K$, nên $\widetilde L=L$. Từ tích ma trận,
\[
 G_{kk}=\sum_{j=0}^{k-1}L_{kj}^2+L_{kk}^2,\qquad
 G_{ik}=\sum_{j=0}^{k-1}L_{ij}L_{kj}+L_{ik}L_{kk}\quad(i>k),
\]
cho \eqref{eq:ls-cholesky}; số dưới căn bằng $L_{kk}^2>0$. Các công thức thế cho
\[
 (L\mathbf z)_i=\sum_{j<i}L_{ij}z_j+L_{ii}z_i=b_i,
 \qquad (L^T\mathbf a)_i=L_{ii}a_i+\sum_{j>i}L_{ji}a_j=z_i.
\]
Do đó $G\mathbf a=LL^T\mathbf a=L\mathbf z=\mathbf b$.
\end{chungminh}

\subsection{Các mô hình xấp xỉ}
\begin{menhde}{}{ls-affine}
Đặt
\[
 s_k=\sum_iw_ix_i^k\ (k=0,1,2),\qquad t_k=\sum_iw_iy_ix_i^k\ (k=0,1),\qquad D=s_0s_2-s_1^2.
\]
Nếu có ít nhất hai vị trí $x_i$ khác nhau, thì
\[
 D=\sum_{i<j}w_iw_j(x_i-x_j)^2>0.
\]
Nghiệm $p(x)=a+bx$ là
\begin{equation}\label{eq:ls-affine}
 a=\frac{s_2t_0-s_1t_1}D,\qquad b=\frac{s_0t_1-s_1t_0}D.
\end{equation}
Với $\bar x=s_1/s_0$, $\bar y=t_0/s_0$, cùng nghiệm có biểu diễn
\[
 b=\frac{\sum_iw_i(x_i-\bar x)(y_i-\bar y)}{\sum_iw_i(x_i-\bar x)^2},\qquad
 a=\bar y-b\bar x.
\]
\end{menhde}
\begin{chungminh}
Tính trực tiếp cho
\begin{align*}
 D&=\sum_{i,j}w_iw_j(x_j^2-x_ix_j)
 =\frac12\sum_{i,j}w_iw_j(x_i^2+x_j^2-2x_ix_j)\\
 &=\sum_{i<j}w_iw_j(x_i-x_j)^2.
\end{align*}
Chọn một cặp $i<j$ có $x_i\ne x_j$. Khi đó
\[
 D=\sum_{k<l}w_kw_l(x_k-x_l)^2
 \geq w_iw_j(x_i-x_j)^2>0.
\]
Hệ phương trình chuẩn là
\[
 s_0a+s_1b=t_0,\qquad s_1a+s_2b=t_1.
\]
\[
 s_0(s_1a+s_2b)-s_1(s_0a+s_1b)
 =Db=s_0t_1-s_1t_0,
\]
\[
 s_2(s_0a+s_1b)-s_1(s_1a+s_2b)
 =Da=s_2t_0-s_1t_1.
\]
Chia cho $D>0$ được \eqref{eq:ls-affine}. Ngoài ra,
\begin{align*}
 \sum_iw_i(x_i-\bar x)^2
 &=s_2-2\bar x s_1+\bar x^2s_0
 =s_2-\frac{s_1^2}{s_0}=\frac D{s_0}>0,\\
 \sum_iw_i(x_i-\bar x)(y_i-\bar y)
 &=t_1-\bar x t_0-\bar y s_1+\bar x\bar y s_0
 =t_1-\frac{s_1t_0}{s_0}.
\end{align*}
Do đó
\[
 \frac{\sum_iw_i(x_i-\bar x)(y_i-\bar y)}
      {\sum_iw_i(x_i-\bar x)^2}
 =\frac{s_0t_1-s_1t_0}D=b,\qquad
 a=\frac{t_0-s_1b}{s_0}=\bar y-b\bar x.
\]
\end{chungminh}
\begin{vidu}
Xét các số đo minh họa để hiệu chuẩn cảm biến: nhiệt độ $x$ là $0,10,20,30$ và tín hiệu $y$ là $1,2,3,5$, với trọng số bằng $1$. Khi đó
\[
 s_0=4,\ s_1=60,\ s_2=1400,\ t_0=11,\ t_1=230,\ D=2000,
\]
\[
 b=\frac{4\cdot230-60\cdot11}{2000}=\frac{13}{100},\qquad
 a=\frac{1400\cdot11-60\cdot230}{2000}=\frac45.
\]
\begin{center}\begin{tabular}{c|rrrr}\toprule
$x_i$&$0$&$10$&$20$&$30$\\\midrule
$y_i$&$1$&$2$&$3$&$5$\\
$p(x_i)$&$4/5$&$21/10$&$17/5$&$47/10$\\
$y_i-p(x_i)$&$1/5$&$-1/10$&$-2/5$&$3/10$\\\bottomrule
\end{tabular}\end{center}
\[
 \sum_i(y_i-p(x_i))=0,\quad\sum_ix_i(y_i-p(x_i))=-1-8+9=0,
\]
\[
 J=\frac1{25}+\frac1{100}+\frac4{25}+\frac9{100}=\frac3{10}.
\]
Horner tại $x=15$ cho $v=(13/100)\cdot15+4/5=11/4$.
\end{vidu}

% BEGIN LS-SPECIFIC-BASES
% Ghép sau mô hình affine, trước phần xấp xỉ tích phân.
\Needspace{7\baselineskip}
Bậc hai vẫn tuyến tính theo ba tham số, dù phụ thuộc phi tuyến vào biến $x$. Đặt
\[
 g(x)=a_0+a_1x+a_2x^2,\qquad
 s_k=\sum_iw_ix_i^k\ (0\leq k\leq4),\quad
 t_k=\sum_iw_ix_i^ky_i\ (0\leq k\leq2).
\]
\begin{menhde}{}{ls-quadratic-moments}
Các hệ số cực tiểu thỏa
\begin{equation}\label{eq:ls-quadratic-moments}
 \begin{pmatrix}s_0&s_1&s_2\\s_1&s_2&s_3\\s_2&s_3&s_4\end{pmatrix}
 \begin{pmatrix}a_0\\a_1\\a_2\end{pmatrix}
 =\begin{pmatrix}t_0\\t_1\\t_2\end{pmatrix}.
\end{equation}
Hệ có nghiệm duy nhất nếu bảng có ít nhất ba hoành độ phân biệt.
\end{menhde}
\begin{chungminh}
Với $Q=\sum_iw_i(a_0+a_1x_i+a_2x_i^2-y_i)^2$ và $j=0,1,2$,
\[
 \frac12\frac{\partial Q}{\partial a_j}
 =\sum_iw_ix_i^j(a_0+a_1x_i+a_2x_i^2-y_i)
 =\sum_{k=0}^2s_{j+k}a_k-t_j.
\]
Với $A_{ij}=x_i^j$, định lý \ref{dl:ls-normal-equations} cho
\[
 \boldsymbol a\text{ cực tiểu}
 \ \Longleftrightarrow\ A^TWA\boldsymbol a=A^TW\boldsymbol y
 \ \Longleftrightarrow\ \sum_{k=0}^2s_{j+k}a_k=t_j\quad(j=0,1,2).
\]
Đây là \eqref{eq:ls-quadratic-moments}. Nếu bảng có ít nhất ba hoành độ phân biệt và $A\boldsymbol z=0$, đặt $q(x)=z_0+z_1x+z_2x^2$. Khi đó
\[
 q(x_i)=(A\boldsymbol z)_i=0\quad\text{với mọi }i,\qquad
 \deg q\leq2.
\]
Bổ đề \ref{bd:roots} cho $q=0$, nên $z_j=[x^j]q=0$ với $j=0,1,2$. Vì vậy $\ker A=\{0\}$, $\operatorname{rank}A=3$. Với $\boldsymbol z\ne0$,
\[
 \boldsymbol z^TA^TWA\boldsymbol z
 =\sum_iw_i(A\boldsymbol z)_i^2>0.
\]
Định lý \ref{dl:ls-normal-equations} cho nghiệm duy nhất.
\end{chungminh}

\begin{menhde}{}{ls-nonpolynomial-bases}
Với $x_i\ne0$, mô hình $g(x)=ax+b/x$ có phương trình chuẩn
\begin{equation}\label{eq:ls-reciprocal-model}
 \begin{pmatrix}\sum_iw_ix_i^2&\sum_iw_i\\
                 \sum_iw_i&\sum_iw_i/x_i^2\end{pmatrix}
 \begin{pmatrix}a\\b\end{pmatrix}
 =\begin{pmatrix}\sum_iw_ix_iy_i\\\sum_iw_iy_i/x_i\end{pmatrix}.
\end{equation}
Hệ có nghiệm duy nhất khi và chỉ khi có hai hoành độ với bình phương khác nhau.
Mô hình $g(x)=ax^2+b\cos x+c$ có phương trình chuẩn
\begin{equation}\label{eq:ls-cosine-model}
 \begin{pmatrix}
 \sum_iw_ix_i^4&\sum_iw_ix_i^2\cos x_i&\sum_iw_ix_i^2\\
 \sum_iw_ix_i^2\cos x_i&\sum_iw_i\cos^2x_i&\sum_iw_i\cos x_i\\
 \sum_iw_ix_i^2&\sum_iw_i\cos x_i&\sum_iw_i
 \end{pmatrix}
 \begin{pmatrix}a\\b\\c\end{pmatrix}
 =\begin{pmatrix}\sum_iw_ix_i^2y_i\\\sum_iw_iy_i\cos x_i\\\sum_iw_iy_i\end{pmatrix}.
\end{equation}
Nghiệm là duy nhất khi và chỉ khi ba vector $(x_i^2)_i$, $(\cos x_i)_i$, $(1)_i$ độc lập tuyến tính.
\end{menhde}
\begin{chungminh}
Với các vector giá trị cơ sở $\mathbf v_j=(\phi_j(x_i))_i$, đặt $A=(\mathbf v_0,\ldots,\mathbf v_r)$. Ta có
\[
 (A^TWA)_{jk}=\sum_iw_i\phi_j(x_i)\phi_k(x_i),\qquad
 (A^TW\boldsymbol y)_j=\sum_iw_i\phi_j(x_i)y_i.
\]
Thay $(\phi_0,\phi_1)=(x,1/x)$ hoặc $(\phi_0,\phi_1,\phi_2)=(x^2,\cos x,1)$ cho từng phần tử trong \eqref{eq:ls-reciprocal-model} hoặc \eqref{eq:ls-cosine-model}. Định lý \ref{dl:ls-normal-equations} cho hai hệ phương trình chuẩn.
Với mô hình ba tham số,
\[
 \ker A=\{0\}
 \ \Longleftrightarrow\
 \left(c_0\mathbf v_0+c_1\mathbf v_1+c_2\mathbf v_2=0
       \ \Longrightarrow\ c_0=c_1=c_2=0\right)
 \ \Longleftrightarrow\ \operatorname{rank}A=3.
\]
Đây chính là điều kiện độc lập tuyến tính trong phát biểu.

Trong mô hình nghịch đảo, định thức Gram bằng
\begin{align*}
 D&=\left(\sum_iw_ix_i^2\right)\left(\sum_iw_i/x_i^2\right)
       -\left(\sum_iw_i\right)^2\\
  &=\sum_{i<j}w_iw_j\left(\frac{x_i^2}{x_j^2}+
                          \frac{x_j^2}{x_i^2}-2\right)
   =\sum_{i<j}w_iw_j\left(\frac{x_i}{x_j}-\frac{x_j}{x_i}\right)^2.
\end{align*}
Vì $x_ix_j\ne0$,
\[
 \frac{x_i}{x_j}-\frac{x_j}{x_i}
 =\frac{x_i^2-x_j^2}{x_ix_j}.
\]
Do đó
\[
 D>0
 \ \Longleftrightarrow\
 \exists i<j:\left(\frac{x_i^2-x_j^2}{x_ix_j}\right)^2>0
 \ \Longleftrightarrow\
 \exists i<j:x_i^2\ne x_j^2.
\]
Nếu $D>0$, ma trận Gram $2\times2$ khả nghịch, nên nghiệm duy nhất. Nếu mọi $x_i^2=r^2>0$, thì
\[
 (x_i)_i=r^2(1/x_i)_i,\qquad
 A(1,-r^2)^T=0,\qquad (1,-r^2)^T\ne0.
\]
Khi đó $\ker A\ne\{0\}$ và định lý \ref{dl:ls-normal-equations} cho nghiệm không duy nhất.
\end{chungminh}
% END LS-SPECIFIC-BASES

% BEGIN LS-TRANSFORMED-MODELS
% Ghép sau các mô hình cơ sở cụ thể, trước phần xấp xỉ tích phân.
\Needspace{7\baselineskip}
\subsection{Các mô hình qua phép biến đổi dữ liệu}
Một mô hình có thể phi tuyến theo tham số nhưng trở thành tuyến tính sau một phép biến đổi. Khi dùng phép biến đổi này, trước hết phải xác định miền mà nó hợp lệ và tiêu chuẩn sai lệch sẽ được cực tiểu hóa.
\begin{menhde}{}{ls-exponential-transform}
Cho $s\in\{-1,1\}$, $sy_i>0$ với mọi $i$, và xét
\[
 g(x)=a\exp\left(\sum_{j=1}^rb_j\varphi_j(x)\right),\qquad sa>0.
\]
Đặt $z_i=\ln(sy_i)$, $c=\ln(sa)$ và $\varphi_0=1$. Việc cực tiểu hóa
\begin{equation}\label{eq:ls-exponential-objective}
 Q_{\log}(c,\boldsymbol b)
 =\sum_iw_i\left(z_i-c-\sum_{j=1}^rb_j\varphi_j(x_i)\right)^2
\end{equation}
tương đương bài toán bình phương tối thiểu tuyến tính với các cơ sở $1,\varphi_1,\ldots,\varphi_r$. Nếu thiết kế $A_{ij}=\varphi_j(x_i)$ có hạng cột đầy đủ, các tham số được xác định duy nhất bởi
\[
 (A^TWA)\boldsymbol\theta=A^TW\boldsymbol z,
 \quad\boldsymbol\theta=(c,b_1,\ldots,b_r)^T,\qquad a=s e^c.
\]
\end{menhde}
\begin{chungminh}
Với $sa>0$,
\[
 \ln(sg(x))=\ln(sa)+\sum_{j=1}^rb_j\varphi_j(x).
\]
Vì $s^2=1$,
\[
 s e^{\ln(sa)}=s(sa)=a\quad(sa>0),\qquad
 \ln\bigl(s(s e^c)\bigr)=\ln(e^c)=c\quad(c\in\R).
\]
Do đó $a\mapsto c$ và $c\mapsto a=s e^c$ là hai ánh xạ ngược nhau. Với $\boldsymbol\theta=(c,b_1,\ldots,b_r)^T$,
\[
 z_i-\ln(sg(x_i))=z_i-(A\boldsymbol\theta)_i,
\]
\[
 Q_{\log}=\sum_iw_i(z_i-(A\boldsymbol\theta)_i)^2
 =\|W^{1/2}(\boldsymbol z-A\boldsymbol\theta)\|_2^2.
\]
Định lý \ref{dl:ls-normal-equations} cho
\[
 \boldsymbol\theta\text{ cực tiểu}
 \ \Longleftrightarrow\
 A^TWA\boldsymbol\theta=A^TW\boldsymbol z.
\]
Nếu $\operatorname{rank}A=r+1$, vector $\boldsymbol\theta$ duy nhất; hai ánh xạ ngược nhau ở trên cho duy nhất $(a,b_1,\ldots,b_r)$.
\end{chungminh}
\Needspace{9\baselineskip}
\begin{vidu}
Dữ liệu dương $\boldsymbol x=(0,1,2,3)^T$, $\boldsymbol y=(1,e,e^2,e^4)^T$ được ghép bằng $g(x)=a e^{bx+cx^2}$, $w_i=1$. Bảng sau phép biến đổi:
\begin{center}
\begin{tabular}{c|rrrr}\toprule
$i$&$0$&$1$&$2$&$3$\\\midrule
$x_i$&$0$&$1$&$2$&$3$\\
$z_i=\ln y_i$&$0$&$1$&$2$&$4$\\
$x_i z_i$&$0$&$1$&$4$&$12$\\
$x_i^2z_i$&$0$&$1$&$8$&$36$\\\bottomrule
\end{tabular}
\end{center}
Hệ moment và nghiệm được kiểm tra trực tiếp bởi
\[
 \begin{pmatrix}4&6&14\\6&14&36\\14&36&98\end{pmatrix}
 \begin{pmatrix}1/20\\11/20\\1/4\end{pmatrix}
 =\begin{pmatrix}7\\17\\45\end{pmatrix}.
\]
Vậy $a=e^{1/20}$, $b=11/20$, $c=1/4$. Giá trị dự đoán và phần dư trên thang logarit:
\begin{center}
\begin{tabular}{c|rrrr}\toprule
$i$&$0$&$1$&$2$&$3$\\\midrule
$\ln g(x_i)$&$1/20$&$17/20$&$43/20$&$79/20$\\
$z_i-\ln g(x_i)$&$-1/20$&$3/20$&$-3/20$&$1/20$\\\bottomrule
\end{tabular}
\end{center}
Suy ra $Q_{\log}=1/20$, $E_{2,\log}=1/80$. Tại $x=3/2$, Horner cho
\begin{center}
\begin{tabular}{c|rrr}\toprule
$j$&$2$&$1$&$0$\\\midrule
Hệ số&$1/4$&$11/20$&$1/20$\\
$(3/2)v_{j+1}$&&$3/8$&$111/80$\\
$v_j$&$1/4$&$37/40$&$\boxed{23/16}$\\\bottomrule
\end{tabular}
\end{center}
Do đó $g(3/2)=e^{23/16}$. Sai số trên thang dữ liệu gốc phải được tính riêng từ $y_i-g(x_i)$.
\end{vidu}

\begin{hequa}{}{ls-power-transform}
Cho $x_i>0$, $sy_i>0$ và mô hình $g(x)=ax^b$ với $sa>0$, $b\in\R$. Đặt $u_i=\ln x_i$, $z_i=\ln(sy_i)$. Khi có ít nhất hai hoành độ phân biệt,
\begin{equation}\label{eq:ls-power-transform}
 b=\frac{\sum_iw_i(u_i-\overline u)(z_i-\overline z)}
         {\sum_iw_i(u_i-\overline u)^2},
 \qquad c=\overline z-b\overline u,\qquad a=s e^c,
\end{equation}
trong đó $\overline u=\sum_iw_iu_i/\sum_iw_i$, $\overline z=\sum_iw_iz_i/\sum_iw_i$. Các tham số này cực tiểu hóa $\sum_iw_i(\ln(sy_i)-\ln(sg(x_i)))^2$.
\end{hequa}
\begin{chungminh}
Vì $x_i>0$, $x_i^b=e^{b\ln x_i}$ và
\[
 \ln(sg(x_i))=c+bu_i.
\]
Ta có $e^{u_i}=x_i$, nên
\[
 x_i\ne x_j\ \Longrightarrow\ u_i\ne u_j.
\]
Đặt $s_0=\sum_iw_i>0$. Mệnh đề \ref{md:ls-affine} cho
\[
 \sum_iw_i(u_i-\overline u)^2
 =\frac1{s_0}\sum_{i<j}w_iw_j(u_i-u_j)^2>0.
\]
Phương trình chuẩn của mô hình $c+bu$ là
\[
 s_0c+\left(\sum_iw_iu_i\right)b=\sum_iw_iz_i,
 \qquad
 \left(\sum_iw_iu_i\right)c+\left(\sum_iw_iu_i^2\right)b
 =\sum_iw_iu_iz_i.
\]
Từ phương trình đầu $c=\overline z-b\overline u$; thế vào phương trình sau cho
\[
 b\sum_iw_i(u_i-\overline u)^2
 =\sum_iw_i(u_i-\overline u)(z_i-\overline z).
\]
Chia cho mẫu số dương và dùng $a=s e^c$ được \eqref{eq:ls-power-transform}. Mặt khác,
\[
 \sum_iw_i(\ln(sy_i)-\ln(sg(x_i)))^2
 =\sum_iw_i(z_i-c-bu_i)^2.
\]
Mệnh đề \ref{md:ls-exponential-transform} cho các tham số vừa tính cực tiểu hóa đúng tiêu chuẩn này.
\end{chungminh}
\Needspace{9\baselineskip}
\begin{vidu}
Với $\boldsymbol x=(1,2,4)^T$, $\boldsymbol y=(3,12,48)^T$, $w_i=1$, bảng biến đổi là
\begin{center}
\begin{tabular}{c|rrr}\toprule
$i$&$0$&$1$&$2$\\\midrule
$u_i=\ln x_i$&$0$&$\ln2$&$2\ln2$\\
$z_i=\ln y_i$&$\ln3$&$\ln3+2\ln2$&$\ln3+4\ln2$\\\bottomrule
\end{tabular}
\end{center}
Ta có $\overline u=\ln2$, $\overline z=\ln3+2\ln2$, tử số $4(\ln2)^2$ và mẫu số $2(\ln2)^2$. Vì vậy $b=2$, $a=3$ và mọi phần dư bằng $0$.
\end{vidu}

\begin{menhde}{}{ls-logarithmic-transform}
Với mô hình $g(x)=\ln H_{\boldsymbol b}(x)$, $H_{\boldsymbol b}(x)=\sum_{j=1}^rb_j\varphi_j(x)$, điều kiện xác định là $H_{\boldsymbol b}(x)>0$. Đặt $z_i=e^{y_i}$. Bài toán tuyến tính
\begin{equation}\label{eq:ls-logarithmic-transform}
 \min_{\boldsymbol b}\sum_iw_i\left(z_i-\sum_{j=1}^rb_j\varphi_j(x_i)\right)^2
\end{equation}
có phương trình chuẩn $A^TWA\boldsymbol b=A^TW\boldsymbol z$. Nghiệm của bài toán này chỉ xác định được $g(x_i)$ nếu $(A\boldsymbol b)_i>0$ với mọi $i$; tính dương trên cả đoạn phải được kiểm tra riêng.
\end{menhde}
\begin{chungminh}
Với $H_{\boldsymbol b}(x_i)>0$,
\[
 y_i=g(x_i)\ \Longleftrightarrow\ e^{y_i}=H_{\boldsymbol b}(x_i).
\]
Với $A_{ij}=\varphi_j(x_i)$,
\[
 H_{\boldsymbol b}(x_i)=(A\boldsymbol b)_i,\qquad
 \sum_iw_i(z_i-H_{\boldsymbol b}(x_i))^2
 =\|\boldsymbol z-A\boldsymbol b\|_W^2.
\]
Định lý \ref{dl:ls-normal-equations} cho
\[
 \boldsymbol b\text{ cực tiểu của \eqref{eq:ls-logarithmic-transform}}
 \ \Longleftrightarrow\
 A^TWA\boldsymbol b=A^TW\boldsymbol z.
\]
Miền của logarit thực cho các tương đương
\[
 g(x_i)\in\R\text{ với mọi }i
 \ \Longleftrightarrow\
 (A\boldsymbol b)_i>0\text{ với mọi }i,
\]
\[
 g(x)\in\R\text{ với mọi }x\in[a,b]
 \ \Longleftrightarrow\
 H_{\boldsymbol b}(x)>0\text{ với mọi }x\in[a,b].
\]
Điều kiện thứ nhất không kéo theo điều kiện thứ hai: trên $[0,1]$, đa thức
\[
 H(x)=(x-\tfrac12)^2-\tfrac18
\]
thỏa $H(0)=H(1)=1/8>0$ nhưng $H(1/2)=-1/8<0$.
\end{chungminh}
\Needspace{9\baselineskip}
\begin{vidu}
Các $x_i=(0,1,2)$, $y_i=(0,0,\ln10)$ cho $z_i=(1,1,10)$. Với cơ sở $(1,x)$, công thức affine cho $H(x)=-1/2+(9/2)x$:
\[
 \sum_i z_i=12,\quad\sum_i x_i z_i=21,\quad
 \overline x=1,\quad\overline z=4,
 \quad b=\frac{21-3\cdot1\cdot4}{5-3}=\frac92,
 \quad a=4-\frac92=-\frac12.
\]
Mặc dù cả ba $z_i$ dương, $H(0)=-1/2<0$, nên $\ln H(0)$ không xác định trong $\R$. Hệ tuyến tính sau phép biến đổi đã có nghiệm, nhưng nghiệm này không cho một mô hình logarit hợp lệ trên bảng.
\end{vidu}

\Needspace{7\baselineskip}
Một phép biến đổi đơn ánh bảo toàn đẳng thức giữa các giá trị; nó không bảo toàn tổng bình phương sai lệch. Để đánh giá mô hình trên đơn vị đo ban đầu, phải tính thêm các $r_i=y_i-g(x_i)$ trong \eqref{eq:ls-mean-square-error}.
\begin{menhde}{}{ls-transform-objectives}
Giả sử $sy_i>0$, $sg(x_i)>0$ và $0<\rho\leq sy_i,sg(x_i)\leq R$. Đặt
\[
 Q_y=\sum_iw_i(y_i-g(x_i))^2,\qquad
 Q_{\log}=\sum_iw_i(\ln(sy_i)-\ln(sg(x_i)))^2.
\]
Khi đó
\begin{equation}\label{eq:ls-transform-comparison}
 \rho^2Q_{\log}\leq Q_y\leq R^2Q_{\log}.
\end{equation}
Nghiệm cực tiểu của $Q_{\log}$ không nhất thiết cực tiểu $Q_y$.
\end{menhde}
\begin{chungminh}
Nếu $y_i\ne g(x_i)$, định lý giá trị trung bình cho một $\xi_i$ nằm giữa $sy_i$ và $sg(x_i)$. Nếu $y_i=g(x_i)$, chọn $\xi_i=sy_i=sg(x_i)$. Trong cả hai trường hợp,
\[
 \ln(sy_i)-\ln(sg(x_i))=\frac{s(y_i-g(x_i))}{\xi_i},
 \qquad\rho\leq\xi_i\leq R.
\]
Do đó
\[
 \rho^2(\ln(sy_i)-\ln(sg(x_i)))^2
 \leq(y_i-g(x_i))^2
 \leq R^2(\ln(sy_i)-\ln(sg(x_i)))^2.
\]
Nhân với $w_i$ và cộng được \eqref{eq:ls-transform-comparison}. Với mô hình hằng $g(x)=a>0$, hai giá trị $y_0=1$, $y_1=4$, $w_0=w_1=1$ cho
\[
 Q_{\log}(a)=(\ln a)^2+(\ln4-\ln a)^2
 =2(\ln a-\ln2)^2+2(\ln2)^2,
\]
\[
 Q_y(a)=(1-a)^2+(4-a)^2=2(a-5/2)^2+9/2.
\]
\[
 Q_{\log}(a)-Q_{\log}(2)=2(\ln a-\ln2)^2\geq0,
 \qquad
 Q_{\log}(a)=Q_{\log}(2)\ \Longleftrightarrow\ a=2,
\]
\[
 Q_y(a)-Q_y(5/2)=2(a-5/2)^2\geq0,
 \qquad
 Q_y(a)=Q_y(5/2)\ \Longleftrightarrow\ a=5/2.
\]
Vì $2\ne5/2$, hai tiêu chuẩn có các cực tiểu khác nhau.
\end{chungminh}

Nếu chọn cực tiểu trực tiếp sai lệch dữ liệu gốc của mô hình mũ, phương trình theo tham số trở thành phi tuyến. Đặt $\boldsymbol v_i=(1,\varphi_1(x_i),\ldots,\varphi_r(x_i))^T$, $g_i=s e^{\boldsymbol v_i^T\boldsymbol\theta}$. Ta có
\[
 \frac{\partial g_i}{\partial\theta_j}=g_i v_{i,j},\qquad
 \frac{\partial^2g_i}{\partial\theta_j\partial\theta_k}=g_i v_{i,j}v_{i,k}.
\]
Với $Q_y(\boldsymbol\theta)=\sum_iw_i(g_i-y_i)^2$, đạo hàm từng số hạng cho
\begin{align*}
 \frac12\nabla Q_y(\boldsymbol\theta)
 &=\sum_iw_i(g_i-y_i)g_i\boldsymbol v_i,\\
 \frac12\nabla^2 Q_y(\boldsymbol\theta)
 &=\sum_iw_i(2g_i^2-y_ig_i)\boldsymbol v_i\boldsymbol v_i^T.
\end{align*}
Các hệ số $g_i$ phụ thuộc vào $\boldsymbol\theta$; hệ phương trình điểm dừng này khác $A^TWA\boldsymbol\theta=A^TW\boldsymbol z$. Ngay với một quan sát $y=1$, $s=1$ và mô hình hằng $g=e^c$, đạo hàm bậc hai bằng $2e^c(2e^c-1)<0$ khi $e^c<1/2$. Vì vậy không thể áp dụng kết luận cực tiểu toàn cục của bài toán bình phương tối thiểu tuyến tính cho hệ điểm dừng phi tuyến này.
% END LS-TRANSFORMED-MODELS

\subsection{Xấp xỉ theo chuẩn tích phân}
Đối với hàm $f\in C([a,b])$, thay tổng sai lệch tại hữu hạn mốc bởi
\[
 J_\rho(p)=\int_a^b\rho(x)(f(x)-p(x))^2dx,
 \qquad\rho\in C([a,b]),\quad\rho(x)>0.
\]
Với $u,v\in C([a,b])$, đặt
\[
 \langle u,v\rangle_\rho=\int_a^b\rho uv\,dx,\qquad\|u\|_\rho^2=\langle u,u\rangle_\rho.
\]
\begin{dinhly}{}{ls-continuous}
Cho các hàm $\phi_0,\ldots,\phi_r$ liên tục, độc lập tuyến tính trên $[a,b]$. Trong $E=\operatorname{span}\{\phi_0,\ldots,\phi_r\}$, tồn tại duy nhất $p_*=\sum_ja_j\phi_j$ cực tiểu $J_\rho$. Các hệ số giải hệ
\begin{equation}\label{eq:ls-continuous-normal}
 H\mathbf a=\mathbf b,\qquad
 H_{jk}=\langle\phi_j,\phi_k\rangle_\rho,\quad b_j=\langle f,\phi_j\rangle_\rho.
\end{equation}
Ma trận $H$ xác định dương. Với mọi $v\in E$,
\[
 \langle f-p_*,v\rangle_\rho=0,\qquad
 \|f-v\|_\rho^2=\|f-p_*\|_\rho^2+\|p_*-v\|_\rho^2.
\]
Phép chiếu $f\mapsto p_*$ tuyến tính và thỏa $\|p_*\|_\rho\leq\|f\|_\rho$.
\end{dinhly}
\begin{chungminh}
Với $\mathbf z\ne0$, độc lập tuyến tính cho $u=\sum_jz_j\phi_j\ne0$. Chọn $c$ sao cho $u(c)\ne0$. Tính liên tục và $\rho(c)>0$ cho một đoạn $I\subseteq[a,b]$ độ dài dương với $|u|\geq|u(c)|/2$ và $\rho\geq\rho(c)/2$. Do đó
\[
 \mathbf z^TH\mathbf z=\int_a^b\rho u^2dx\geq|I|\frac{\rho(c)|u(c)|^2}8>0.
\]
\[
 H\mathbf z=0\ \Longrightarrow\
 \mathbf z^TH\mathbf z=0\ \Longrightarrow\ \mathbf z=0,
 \qquad \operatorname{rank}H=(r+1)-\dim\ker H=r+1.
\]
Do đó $H$ khả nghịch. Đặt $\mathbf a=H^{-1}\mathbf b$. Với $v=\sum_jc_j\phi_j$,
\[
 \langle f-p_*,v\rangle_\rho
 =\sum_jc_j\left(b_j-\sum_kH_{jk}a_k\right)=0.
\]
Vì $p_*-v\in E$,
\[
 \|f-v\|_\rho^2
 =\|f-p_*\|_\rho^2+2\langle f-p_*,p_*-v\rangle_\rho+\|p_*-v\|_\rho^2
 =\|f-p_*\|_\rho^2+\|p_*-v\|_\rho^2.
\]
\[
 J_\rho(v)\geq J_\rho(p_*),\qquad
 J_\rho(v)=J_\rho(p_*)
 \ \Longleftrightarrow\ \|p_*-v\|_\rho^2=0.
\]
Nếu $p_*-v$ không đồng nhất bằng $0$, lập luận trên với $u=p_*-v$ cho $\|p_*-v\|_\rho^2>0$, mâu thuẫn. Vì vậy $v=p_*$ là cực tiểu duy nhất.

Viết $\mathbf b[f]_j=\langle f,\phi_j\rangle_\rho$, $\mathbf a[f]=H^{-1}\mathbf b[f]$. Với $\lambda,\mu\in\R$,
\[
 \mathbf b[\lambda f_1+\mu f_2]
 =\lambda\mathbf b[f_1]+\mu\mathbf b[f_2],\qquad
 \mathbf a[\lambda f_1+\mu f_2]
 =\lambda\mathbf a[f_1]+\mu\mathbf a[f_2],
\]
\[
 p_*[\lambda f_1+\mu f_2]
 =\sum_j(\lambda a_j[f_1]+\mu a_j[f_2])\phi_j
 =\lambda p_*[f_1]+\mu p_*[f_2].
\]
Đặt $v=0$ trong đồng nhất thức chuẩn đã chứng minh:
\[
 \|f\|_\rho^2=\|f-p_*\|_\rho^2+\|p_*\|_\rho^2
 \geq\|p_*\|_\rho^2
 \ \Longrightarrow\ \|p_*\|_\rho\leq\|f\|_\rho.
\]
\end{chungminh}

Xét một trong hai tích vô hướng
\[
 \langle u,v\rangle=\sum_{i=0}^{m-1}w_iu(x_i)v(x_i),\qquad
 \langle u,v\rangle=\int_a^b\rho(x)u(x)v(x)dx.
\]
Trong trường hợp rời rạc, các mốc phân biệt và $0\leq r<m$; trong trường hợp tích phân, $r\geq0$. Các điều kiện trên bảo đảm $\langle p,p\rangle>0$ với mọi $p\in\mathcal P_r\setminus\{0\}$.
\begin{dinhly}{}{ls-orthogonal-polynomials}
Đặt $\pi_{-1}=0$, $\pi_0=1$, $\nu_k=\langle\pi_k,\pi_k\rangle$. Với $0\leq k<r$, đặt
\begin{equation}\label{eq:ls-orthogonal-recurrence}
 \alpha_k=\frac{\langle x\pi_k,\pi_k\rangle}{\nu_k},\qquad
 \beta_0=0,\quad\beta_k=\frac{\nu_k}{\nu_{k-1}}\ (k\geq1),\qquad
 \pi_{k+1}=(x-\alpha_k)\pi_k-\beta_k\pi_{k-1}.
\end{equation}
Các $\pi_k$ là đa thức hệ số đầu $1$, bậc $k$, và
\[
 \langle\pi_j,\pi_k\rangle=0\ (j\ne k),\qquad\nu_k>0\quad(0\leq k\leq r).
\]
Với dữ liệu $f$ (hoặc $f(x_i)=y_i$), nghiệm bình phương tối thiểu trong $\mathcal P_r$ là
\begin{equation}\label{eq:ls-orthogonal-projection}
 p_r(x)=\sum_{k=0}^r\gamma_k\pi_k(x),\qquad
 \gamma_k=\frac{\langle f,\pi_k\rangle}{\nu_k},\qquad
 \|f-p_r\|^2=\|f\|^2-\sum_{k=0}^r\frac{\langle f,\pi_k\rangle^2}{\nu_k}.
\end{equation}
Trong trường hợp rời rạc, các chuẩn được hiểu trên các giá trị ở mốc.
\end{dinhly}
\begin{chungminh}
Trước hết, với $0\ne u\in\mathcal P_r$ trong trường hợp rời rạc, $r<m$ và bổ đề \ref{bd:roots} cho một chỉ số $i$ có $u(x_i)\ne0$. Vì vậy
\[
 \langle u,u\rangle=\sum_jw_ju(x_j)^2
 \geq w_i u(x_i)^2>0.
\]
Trong trường hợp tích phân, chọn $r+1$ điểm phân biệt trong $[a,b]$. Bổ đề \ref{bd:roots} cho
\[
 \#\{x\in[a,b]:u(x)=0\}\leq\deg u\leq r<r+1,
\]
nên có một điểm $c$ trong các điểm đã chọn với $u(c)\ne0$. Tính liên tục cho một đoạn $I\subseteq[a,b]$, $|I|>0$, với $|u|\geq|u(c)|/2$ và $\rho\geq\rho(c)/2$, nên
\[
 \langle u,u\rangle=\int_a^b\rho u^2dx
 \geq |I|\rho(c)|u(c)|^2/8>0.
\]
Do đó $\pi_0=1$ có hệ số đầu $1$, bậc $0$ và $\nu_0>0$.

Giả sử đã có $\pi_0,\ldots,\pi_k$ với hệ số đầu $1$, bậc $0,\ldots,k$ và trực giao đôi một. Với $u\in\mathcal P_k$, đặt $u_k=u$ và lần lượt cho $j=k,\ldots,0$,
\[
 c_j=[x^j]u_j,\qquad u_{j-1}=u_j-c_j\pi_j,\qquad
 [x^j]u_{j-1}=c_j-c_j=0.
\]
Suy ra $\deg u_{j-1}\leq j-1$, $u_{-1}=0$ và
\[
 u=\sum_{j=0}^k(u_j-u_{j-1})=\sum_{j=0}^kc_j\pi_j.
\]
Nếu $\sum_{j=0}^kd_j\pi_j=0$ mà có $d_j\ne0$, đặt $s=\max\{j:d_j\ne0\}$. Khi đó
\[
 0=[x^s]\sum_{j=0}^kd_j\pi_j
 =d_s[x^s]\pi_s+\sum_{j<s}d_j[x^s]\pi_j
 =d_s\ne0,
\]
mâu thuẫn. Vì vậy $\pi_0,\ldots,\pi_k$ là một cơ sở của $\mathcal P_k$.

Với $k=0<r$,
\[
 \pi_1=x-\alpha_0,\qquad
 \langle\pi_1,\pi_0\rangle
 =\langle x,1\rangle-\alpha_0\nu_0=0,\qquad
 [x]\pi_1=1.
\]
Xét $1\leq k<r$. Khi $0\leq j\leq k-2$, $\deg(x\pi_j)=j+1\leq k-1$; cơ sở vừa chứng minh cho
\[
 x\pi_j=\sum_{\ell=0}^{k-1}c_\ell\pi_\ell,\qquad
 \langle x\pi_k,\pi_j\rangle
 =\langle\pi_k,x\pi_j\rangle
 =\sum_{\ell=0}^{k-1}c_\ell\langle\pi_k,\pi_\ell\rangle=0.
\]
Do đó
\[
 \langle\pi_{k+1},\pi_j\rangle
 =\langle x\pi_k,\pi_j\rangle
 -\alpha_k\langle\pi_k,\pi_j\rangle
 -\beta_k\langle\pi_{k-1},\pi_j\rangle=0.
\]
Vì hệ số đầu của $x\pi_{k-1}$ và $\pi_k$ đều bằng $1$,
\[
 x\pi_{k-1}-\pi_k\in\mathcal P_{k-1},\qquad
 x\pi_{k-1}=\pi_k+\sum_{\ell=0}^{k-1}d_\ell\pi_\ell.
\]
Suy ra
\[
 \langle x\pi_k,\pi_{k-1}\rangle
 =\langle\pi_k,x\pi_{k-1}\rangle
 =\nu_k+\sum_{\ell=0}^{k-1}d_\ell\langle\pi_k,\pi_\ell\rangle
 =\nu_k,
\]
\[
 \langle\pi_{k+1},\pi_{k-1}\rangle
 =\nu_k-\beta_k\nu_{k-1}=0,\qquad
 \langle\pi_{k+1},\pi_k\rangle
 =\langle x\pi_k,\pi_k\rangle-\alpha_k\nu_k=0.
\]
Trong cả hai trường hợp $k=0$ và $k\geq1$,
\[
 [x^{k+1}]\pi_{k+1}
 =[x^{k+1}](x\pi_k-\alpha_k\pi_k-\beta_k\pi_{k-1})=1,
 \qquad \deg\pi_{k+1}=k+1\leq r.
\]
Tính dương ở đầu chứng minh cho $\nu_{k+1}>0$. Quy nạp dựng đủ $\pi_0,\ldots,\pi_r$, và lập luận cơ sở cho $k=r$ cho cơ sở của $\mathcal P_r$.

Với $p=\sum_{k=0}^rc_k\pi_k$, đặt $\gamma_k=\langle f,\pi_k\rangle/\nu_k$. Trực giao cho
\begin{align*}
 \|f-p\|^2
 &=\|f\|^2-2\sum_{k=0}^rc_k\langle f,\pi_k\rangle+\sum_{k=0}^r\nu_kc_k^2\\
 &=\|f\|^2-\sum_{k=0}^r\frac{\langle f,\pi_k\rangle^2}{\nu_k}
   +\sum_{k=0}^r\nu_k(c_k-\gamma_k)^2.
\end{align*}
Với $p_r=\sum_k\gamma_k\pi_k$,
\[
 \|f-p\|^2-\|f-p_r\|^2
 =\sum_{k=0}^r\nu_k(c_k-\gamma_k)^2\geq0.
\]
Vì mọi $\nu_k>0$,
\[
 \|f-p\|^2=\|f-p_r\|^2
 \ \Longleftrightarrow\ c_k=\gamma_k\text{ với mọi }k
 \ \Longleftrightarrow\ p=p_r.
\]
Thế $c_k=\gamma_k$ vào khai triển cho công thức sai lệch trong \eqref{eq:ls-orthogonal-projection}.
\end{chungminh}
\begin{vidu}
Lấy $\rho=1$ trên $[-1,1]$. Công thức Newton--Leibniz cho
\[
 \int_{-1}^1x^jdx=\frac{1-(-1)^{j+1}}{j+1}.
\]
Vì vậy
\[
 \pi_0=1,\quad\nu_0=2,\quad\alpha_0=0,\quad\pi_1=x,\quad
 \nu_1=2/3,\quad\alpha_1=0,\quad\beta_1=1/3,\quad\pi_2=x^2-1/3.
\]
Với $f(x)=x^2$ và mô hình affine,
\[
 \gamma_0=\frac{2/3}2=1/3,\qquad\gamma_1=\frac{\int_{-1}^1x^3dx}{2/3}=0,\qquad p_1=1/3,
\]
\[
 \int_{-1}^1(f-p_1)^2dx=\frac25-\frac49+\frac29=\frac8{45}.
\]
Với mô hình bậc $2$, $f=\pi_2+\pi_0/3$, nên $p_2=f$ và sai lệch bằng $0$.
\end{vidu}



% END EXTENSION SPLINE LS

% BEGIN EXTENSION CHAPTER 2
\chapter{Các phương pháp tính gần đúng đạo hàm và tích phân}
\label{chap:numerical-calculus}

\Needspace{9\baselineskip}
\section{Các công thức tính gần đúng đạo hàm}
Từ các giá trị $f(x_i)$, bài toán đầu tiên là tìm một tổ hợp tuyến tính xấp xỉ $f^{(s)}(x)$. Các hệ số được chọn để tổ hợp đó chính xác trên những đa thức bậc thấp; phần Taylor còn lại xác định sai số. Các đạo hàm trên đoạn đóng được hiểu là các đạo hàm liên tục tới biên. Các mốc và điểm cần tính phải thuộc đoạn xác định của hàm. Trong chương này, $h>0$.

\begin{bode}{}{dnum-taylor}
Cho $r\geq1$, $f\in C^r([\alpha,\beta])$. Với $x\in[\alpha,\beta]$,
\begin{equation}\label{eq:dnum-taylor-integral}
 f(x)=\sum_{j=0}^{r-1}\frac{f^{(j)}(\alpha)}{j!}(x-\alpha)^j
 +\int_\alpha^\beta\frac{(x-t)_+^{r-1}}{(r-1)!}f^{(r)}(t)\,dt.
\end{equation}
Ở đây $u_+^j=(\max\{u,0\})^j$ với $j\geq1$; $u_+^0$ bằng $1$ khi $u>0$ và bằng $0$ khi $u\leq0$. Với $0\leq s<r$, đạo hàm cấp $s$ của phần tích phân là
\[
 \int_\alpha^\beta\frac{(x-t)_+^{r-1-s}}{(r-1-s)!}f^{(r)}(t)\,dt.
\]
\end{bode}
\begin{chungminh}
Với $g\in C([\alpha,\beta])$, đặt
\[
 R_j(x)=\int_\alpha^x\frac{(x-t)^{j-1}}{(j-1)!}g(t)dt\quad(j\geq1).
\]
Công thức Newton--Leibniz và quy tắc đạo hàm tích phân cho
\[
 R_1'=g,\qquad
 R_j'(x)=\frac{(x-x)^{j-1}}{(j-1)!}g(x)
 +\int_\alpha^x\frac{(x-t)^{j-2}}{(j-2)!}g(t)dt=R_{j-1}(x)\quad(j\geq2).
\]
Do đó $R_r^{(s)}=R_{r-s}$ ($s<r$), $R_r^{(r)}=g$ và $R_r^{(s)}(\alpha)=0$ ($s<r$). Lấy $g=f^{(r)}$ và
\[
 H(x)=f(x)-\sum_{j=0}^{r-1}\frac{f^{(j)}(\alpha)}{j!}(x-\alpha)^j-R_r(x).
\]
Ta có $H^{(r)}=0$, $H^{(j)}(\alpha)=0$ ($0\leq j<r$). Nếu $H^{(j+1)}=0$, thì
\[
 H^{(j)}(x)=H^{(j)}(\alpha)+\int_\alpha^xH^{(j+1)}(t)dt=0.
\]
Quy nạp lùi từ $j=r-1$ tới $0$ cho $H=0$. Với $s<r$,
\[
 R_r^{(s)}(x)=R_{r-s}(x)
 =\int_\alpha^\beta\frac{(x-t)_+^{r-1-s}}{(r-1-s)!}f^{(r)}(t)dt,
\]
vì tích phân trên $(x,\beta]$ bằng $0$; giá trị tại một điểm $t=x$ không làm đổi tích phân.
\end{chungminh}

\begin{bode}{}{dnum-kernel}
Cho $r\geq1$ và phiếm hàm
\[
 \mathcal L(f)=c_0\int_\alpha^\beta f(x)\,dx
 +\sum_{i=1}^{N}\sum_{s=0}^{r-1}c_{i,s}f^{(s)}(z_i),\qquad z_i\in[\alpha,\beta].
\]
Nếu $\mathcal L(q)=0$ với mọi $q\in\mathcal P_{r-1}$, thì với $f\in C^r([\alpha,\beta])$,
\begin{equation}\label{eq:dnum-kernel}
 \mathcal L(f)=\int_\alpha^\beta K_r(t)f^{(r)}(t)\,dt,
\end{equation}
\[
 K_r(t)=c_0\frac{(\beta-t)^r}{r!}
 +\sum_{i=1}^{N}\sum_{s=0}^{r-1}
 c_{i,s}\frac{(z_i-t)_+^{r-1-s}}{(r-1-s)!}.
\]
Nếu $K_r$ có một dấu và $\int_\alpha^\beta|K_r(t)|dt>0$, thì tồn tại $\xi\in(\alpha,\beta)$ sao cho
\[
 \mathcal L(f)=f^{(r)}(\xi)\int_\alpha^\beta K_r(t)\,dt.
\]
Luôn có $|\mathcal L(f)|\leq\lVert f^{(r)}\rVert_\infty\int_\alpha^\beta|K_r(t)|\,dt$.
\end{bode}
\begin{chungminh}
Viết $f=T_{r-1}+R_r$ theo Bổ đề \ref{bd:dnum-taylor}. Với $g=f^{(r)}$, ta có $\mathcal L(T_{r-1})=0$ và
\[
 \int_\alpha^\beta R_r(x)dx
 =R_{r+1}(\beta)-R_{r+1}(\alpha)
 =\int_\alpha^\beta\frac{(\beta-t)^r}{r!}g(t)dt.
\]
Do đó
\begin{align*}
 \mathcal L(f)&=c_0\int_\alpha^\beta\frac{(\beta-t)^r}{r!}g(t)dt\\
 &\quad+\sum_{i=1}^N\sum_{s=0}^{r-1}c_{i,s}
 \int_\alpha^\beta\frac{(z_i-t)_+^{r-1-s}}{(r-1-s)!}g(t)dt\\
 &=\int_\alpha^\beta K_r(t)g(t)dt.
\end{align*}
Để chứng minh công thức giá trị trung gian, giả sử $K_r$ có một dấu và $\int_\alpha^\beta|K_r|>0$. Đặt $w=|K_r|$, $A=\int_\alpha^\beta w>0$ và $c=A^{-1}\int_\alpha^\beta wg$. Giả sử $g(t)\ne c$ trên $(\alpha,\beta)$. Theo tính liên tục và định lý giá trị trung gian, có $\sigma\in\{-1,1\}$ sao cho $\sigma(g(t)-c)>0$ trên khoảng này.

Các điểm $z_i$ chia $K_r$ thành hữu hạn đoạn đa thức. Từ $A>0$, có một điểm $t_0$ trong một đoạn như vậy với $w(t_0)>0$. Tính liên tục tại $t_0$ cho một đoạn $J$ có độ dài $|J|>0$, nằm trong $(\alpha,\beta)$, và các số $a,b>0$ sao cho
\[
 w(t)\geq a,\quad \sigma(g(t)-c)\geq b\quad(t\in J).
\]
Vì $w\geq0$ trên toàn đoạn,
\[
 0=\sigma\left(\int_\alpha^\beta wg-cA\right)
 =\int_\alpha^\beta w\sigma(g-c)
 \geq\int_J ab\,dt=ab|J|>0,
\]
mâu thuẫn. Vậy có $\xi\in(\alpha,\beta)$ với $g(\xi)=c$. Nếu $K_r=\varepsilon w$, $\varepsilon\in\{-1,1\}$, thì
\[
 \mathcal L(f)=\varepsilon\int_\alpha^\beta wg
 =\varepsilon A g(\xi)=g(\xi)\int_\alpha^\beta K_r.
\]
Cuối cùng,
\[
 |\mathcal L(f)|\leq\int_\alpha^\beta|K_r(t)||g(t)|dt
 \leq\|g\|_\infty\int_\alpha^\beta|K_r(t)|dt.
\]
\end{chungminh}

\Needspace{7\baselineskip}
\subsection{Công thức hai điểm}
Khi chỉ có hai giá trị hàm, hiệu của chúng là dữ liệu trực tiếp để xấp xỉ đạo hàm. Công thức Newton--Leibniz cho
\[
 \frac{f(x+h)-f(x)}h=\frac1h\int_x^{x+h}f'(t)dt.
\]
Sai số so với $f'(x)$ được xác định bằng cách khai triển $f$ tại $x$. Nếu hai giá trị nằm đối xứng quanh $x$, số hạng chứa $f''(x)$ bị loại khỏi hiệu.
\begin{dinhly}{}{dnum-two-point}
Đặt
\[
\begin{aligned}
 D_+f(x)&=\frac{f(x+h)-f(x)}h,\\
 D_-f(x)&=\frac{f(x)-f(x-h)}h,\\
 D_0f(x)&=\frac{f(x+h)-f(x-h)}{2h}.
\end{aligned}
\]
Nếu $f\in C^2([x,x+h])$, tồn tại $\xi_+\in(x,x+h)$ sao cho
\[
 D_+f(x)-f'(x)=\frac h2 f''(\xi_+).
\]
Nếu $f\in C^2([x-h,x])$, tồn tại $\xi_-\in(x-h,x)$ sao cho
\[
 D_-f(x)-f'(x)=-\frac h2 f''(\xi_-).
\]
Nếu $f\in C^3([x-h,x+h])$, tồn tại $\xi_0\in(x-h,x+h)$ sao cho
\[
 D_0f(x)-f'(x)=\frac{h^2}{6}f'''(\xi_0).
\]
\end{dinhly}
\begin{chungminh}
Từ \eqref{eq:dnum-taylor-integral}, hoặc thay biến $t\mapsto x-t$ cho phía trái,
\[
 f(x+h)=f(x)+hf'(x)+\int_0^h(h-t)f''(x+t)\,dt,
\]
\[
 f(x-h)=f(x)-hf'(x)+\int_0^h(h-t)f''(x-t)\,dt.
\]
Suy ra
\[
 \begin{aligned}
 D_+f(x)-f'(x)&=\frac1h\int_0^h(h-t)f''(x+t)\,dt,\\
 D_-f(x)-f'(x)&=-\frac1h\int_0^h(h-t)f''(x-t)\,dt.
 \end{aligned}
\]

Với $0\leq t\leq h$, $(h-t)/h\geq0$ và
\[
 \frac1h\int_0^h(h-t)dt=\frac1h[ht-t^2/2]_0^h=\frac h2>0.
\]
Bổ đề \ref{bd:dnum-kernel} áp dụng cho $t\mapsto f''(x+t)$ và $t\mapsto f''(x-t)$; đổi điểm $t$ thành $x+t$ hoặc $x-t$ cho hai công thức một phía.

Với phần trung tâm, khai triển tới bậc $2$ cho
\[
 f(x\pm h)=f(x)\pm hf'(x)+\frac{h^2}{2}f''(x)
 \pm\frac12\int_0^h(h-t)^2f'''(x\pm t)\,dt.
\]
Lấy hiệu và chia cho $2h$:
\[
 D_0f(x)-f'(x)=\frac1{4h}\int_{-h}^h(h-|t|)^2f'''(x+t)\,dt.
\]

Với $-h\leq t\leq h$, $(h-|t|)^2/(4h)\geq0$ và
\[
 \frac1{4h}\int_{-h}^h(h-|t|)^2dt
 =\frac1{2h}\int_0^h(h-t)^2dt
 =\frac1{2h}\left[-\frac{(h-t)^3}{3}\right]_0^h=\frac{h^2}{6}>0.
\]
Bổ đề \ref{bd:dnum-kernel} cho $\xi_0\in(x-h,x+h)$ và công thức trung tâm.
\end{chungminh}

\Needspace{7\baselineskip}
\subsection{Các công thức ba điểm}
Để loại số hạng sai số bậc $h$ ở phía phải, xét $h^{-1}(af(x)+bf(x+h)+cf(x+2h))$. Tính chính xác trên $1,t-x,(t-x)^2$ đòi hỏi
\[
 a+b+c=0,\qquad b+2c=1,\qquad b+4c=0.
\]
Trừ hai phương trình sau cho $2c=-1$, rồi $b=2$, $a=-3/2$. Với đạo hàm cấp $2$ ở ba mốc đối xứng, hệ tương ứng là $a+b+c=0$, $-a+c=0$, $a+c=2$; nghiệm là $(1,-2,1)$.
\begin{dinhly}{}{dnum-three-point}
Đặt
\[
 D_{+,3}f(x)=\frac{-3f(x)+4f(x+h)-f(x+2h)}{2h},
\]
\[
\begin{aligned}
 D_{-,3}f(x)&=\frac{3f(x)-4f(x-h)+f(x-2h)}{2h},\\
 D_2f(x)&=\frac{f(x-h)-2f(x)+f(x+h)}{h^2}.
\end{aligned}
\]
Với $f\in C^3$ trên đoạn chứa các mốc tương ứng,
\[
 D_{\pm,3}f(x)-f'(x)=-\frac{h^2}{3}f'''(\xi_\pm).
\]
Với $f\in C^4([x-h,x+h])$,
\[
 D_2f(x)-f''(x)=\frac{h^2}{12}f^{(4)}(\xi_2).
\]
Mỗi $\xi$ thuộc khoảng mở có hai đầu là mốc nhỏ nhất và lớn nhất trong công thức đang xét.
\end{dinhly}
\begin{chungminh}

Đặt $T_2(u)=f(x)+uf'(x)+u^2f''(x)/2$. Khi đó
\begin{align*}
 -3T_2(0)+4T_2(h)-T_2(2h)
 &=(-3+4-1)f(x)+(4h-2h)f'(x)\\
 &\quad+(4h^2-4h^2)f''(x)/2=2hf'(x).
\end{align*}
Thế phần dư tích phân của Bổ đề \ref{bd:dnum-taylor} vào ba giá trị cho
\[
 D_{+,3}f(x)-f'(x)=\int_0^{2h}K(t)f'''(x+t)\,dt,
\]
\[
 K(t)=\begin{cases}
 \dfrac{4(h-t)^2-(2h-t)^2}{4h}=\dfrac{t(3t-4h)}{4h},&0\leq t\leq h,\\[2mm]
 -\dfrac{(2h-t)^2}{4h},&h\leq t\leq2h.
 \end{cases}
\]
Trên $[0,h]$, $t\geq0$ và $3t-4h\leq-h<0$; trên $[h,2h]$, $-(2h-t)^2\leq0$. Do đó $K\leq0$ và
\[
 \int_0^{2h}K(t)dt=\frac1{4h}\left(4\frac{h^3}{3}-\frac{(2h)^3}{3}\right)=-\frac{h^2}{3}.
\]
Để xét phía trái, đặt $g(t)=f(x-t)$. Khi đó
\[
 D_{+,3}g(0)=-D_{-,3}f(x),\quad g'(0)=-f'(x),\quad g'''(t)=-f'''(x-t).
\]

Nếu điểm giá trị trung gian cho $g$ là $\eta\in(0,2h)$, thì
\[
 -D_{-,3}f(x)+f'(x)=-\frac{h^2}{3}g^{(3)}(\eta)
 =\frac{h^2}{3}f^{(3)}(x-\eta).
\]
Nhân với $-1$ và đặt $\xi_-=x-\eta\in(x-2h,x)$ cho công thức phía trái.


Với $T_3(u)=\sum_{j=0}^3f^{(j)}(x)u^j/j!$,
\[
 T_3(h)+T_3(-h)-2T_3(0)
 =\sum_{j=0}^3\frac{f^{(j)}(x)}{j!}\bigl(h^j+(-h)^j-2\cdot0^j\bigr)
 =h^2f''(x),
\]
trong đó $0^0=1$. Thế phần dư tích phân cấp $4$ cho
\[
 f(x+h)+f(x-h)-2f(x)=h^2f''(x)
 +\frac16\int_{-h}^h(h-|t|)^3f^{(4)}(x+t)\,dt.
\]
Chia cho $h^2$; kernel không âm và
\[
 \frac1{6h^2}\int_{-h}^h(h-|t|)^3dt
 =\frac1{3h^2}\int_0^h(h-t)^3dt
 =\frac1{3h^2}\left[-\frac{(h-t)^4}{4}\right]_0^h=\frac{h^2}{12}>0.
\]
Các công thức với $\xi$ theo bổ đề \ref{bd:dnum-kernel}.
\end{chungminh}

\Needspace{9\baselineskip}
\begin{vidu}
Với $f(t)=t^4$, $x=1$, các công thức trên cho
\[
 D_+f(1)=4+6h+4h^2+h^3,\quad D_0f(1)=4+4h^2,
\]
\[
 D_{+,3}f(1)=4-8h^2-6h^3,\qquad D_2f(1)=12+2h^2.
\]
Tại $h=1/10$, các kết quả lần lượt là $4.641$, $4.04$, $3.914$, $12.02$. Giá trị chính xác là $f'(1)=4$, $f''(1)=12$.
\end{vidu}

\Needspace{7\baselineskip}
\subsection{Các công thức \texorpdfstring{$k$}{k} điểm}
Cho $k\geq2$, $1\leq s<k$, vector mốc chuẩn $\boldsymbol\tau=(\tau_0,\ldots,\tau_{k-1})^T\in\R^k$ có các thành phần đôi một phân biệt. Ký hiệu $[t^j]p$ là hệ số của $t^j$ trong đa thức $p$. Đặt $x_i=x+h\tau_i$ và
\begin{equation}\label{eq:dnum-stencil}
 D_{s,h}f(x)=h^{-s}\sum_{i=0}^{k-1}w_if(x_i).
\end{equation}

\begin{dinhly}{}{dnum-moments}
Có duy nhất vector $\boldsymbol w=(w_0,\ldots,w_{k-1})^T$ để \eqref{eq:dnum-stencil} chính xác với mọi đa thức bậc không quá $k-1$. Vector đó thỏa
\begin{equation}\label{eq:dnum-moments}
 \sum_{i=0}^{k-1}w_i\tau_i^j=s!\,\delta_{j,s},\qquad j=0,\ldots,k-1.
\end{equation}
Nếu $L_i(t)=\prod_{j\ne i}(t-\tau_j)/(\tau_i-\tau_j)$, thì
\[
 w_i=L_i^{(s)}(0)=s![t^s]L_i(t).
\]
\end{dinhly}
\begin{chungminh}
Với $q_j(u)=(u-x)^j$,
\[
 q_j^{(s)}(x)=s!\delta_{j,s},\qquad
 D_{s,h}q_j(x)=h^{j-s}\sum_iw_i\tau_i^j.
\]
Vì $h>0$, tính chính xác trên $q_0,\ldots,q_{k-1}$ tương đương \eqref{eq:dnum-moments}. Hệ moment là $V^T\boldsymbol w=\boldsymbol b$, với $V_{i,j}=\tau_i^j$ và $b_j=s!\delta_{j,s}$. Định lý \ref{dl:vandermonde} cho
\[
 \det V=\prod_{i<j}(\tau_j-\tau_i)\ne0.
\]
Do đó tồn tại duy nhất $\boldsymbol w$. Biểu diễn \eqref{eq:cardinal} và hệ thức Horner \eqref{eq:shifted-expansion} tại $c=0$ cho
\[
 v^{(s)}(0)=\sum_i v(\tau_i)L_i^{(s)}(0)
 =\sum_i v(\tau_i)s![t^s]L_i(t)\quad(v\in\mathcal P_{k-1}).
\]

Lấy $v(t)=t^j$ ($0\leq j<k$) trong đẳng thức trên:
\[
 \sum_i\tau_i^jL_i^{(s)}(0)=s!\delta_{j,s}.
\]
Vậy $(L_i^{(s)}(0))_i$ là nghiệm của cùng hệ moment; tính duy nhất cho $w_i=L_i^{(s)}(0)=s![t^s]L_i(t)$.
\end{chungminh}

\begin{menhde}{}{dnum-general-error}
Giả sử các moment chính xác tới bậc $r-1$, với $r>s$. Đặt $\alpha=\min\{x,x_0,\ldots,x_{k-1}\}$, $\beta=\max\{x,x_0,\ldots,x_{k-1}\}$. Với $f\in C^r([\alpha,\beta])$,
\[
 f^{(s)}(x)-D_{s,h}f(x)=\int_\alpha^\beta K_{s,r}(t)f^{(r)}(t)\,dt,
\]
\[
 K_{s,r}(t)=\frac{(x-t)_+^{r-1-s}}{(r-1-s)!}
 -h^{-s}\sum_iw_i\frac{(x_i-t)_+^{r-1}}{(r-1)!}.
\]
Một chặn khác, không đòi hỏi kernel có một dấu, là
\begin{equation}\label{eq:dnum-truncation-bound}
 |D_{s,h}f(x)-f^{(s)}(x)|
 \leq\frac{h^{r-s}}{r!}\lVert f^{(r)}\rVert_\infty\sum_i|w_i||\tau_i|^r.
\end{equation}
Nếu $f\in C^{r+1}$ và $\mu_r=\sum_iw_i\tau_i^r$, thì
\[
 D_{s,h}f(x)-f^{(s)}(x)=\frac{\mu_r}{r!}h^{r-s}f^{(r)}(x)+R,
\]
\[
 |R|\leq\frac{h^{r+1-s}}{(r+1)!}\lVert f^{(r+1)}\rVert_\infty
 \sum_i|w_i||\tau_i|^{r+1}.
\]
\end{menhde}
\begin{chungminh}
Đặt $\mu_j=\sum_iw_i\tau_i^j$ và $\mathcal L(f)=f^{(s)}(x)-D_{s,h}f(x)$. Với $q_j(u)=(u-x)^j$, $0\leq j<r$,
\[
 \mathcal L(q_j)=s!\delta_{j,s}-h^{j-s}\mu_j
 =s!\delta_{j,s}(1-h^{j-s})=0.
\]
Các $q_j$ là một cơ sở của $\mathcal P_{r-1}$, nên Bổ đề \ref{bd:dnum-kernel} cho biểu diễn kernel đã nêu.

Đặt $d_i=h\tau_i$. Với $d_i\ne0$, đặt $\sigma_i=\operatorname{sgn}(d_i)$ và áp dụng Bổ đề \ref{bd:dnum-taylor} cho $u\mapsto f(x+\sigma_i u)$ trên $[0,|d_i|]$:
\[
 f(x+d_i)=\sum_{j=0}^{r-1}\frac{f^{(j)}(x)}{j!}d_i^j+R_i,
\]
\[
 R_i=\frac{\sigma_i^r}{(r-1)!}\int_0^{|d_i|}(|d_i|-u)^{r-1}f^{(r)}(x+\sigma_i u)du.
\]
Khi $d_i=0$, lấy $R_i=0$. Do đó
\[
 |R_i|\leq\frac{\|f^{(r)}\|_\infty}{(r-1)!}\int_0^{|d_i|}(|d_i|-u)^{r-1}du
 =\frac{h^r|\tau_i|^r}{r!}\|f^{(r)}\|_\infty.
\]
Nhân với $h^{-s}w_i$, cộng theo $i$:
\begin{align*}
 D_{s,h}f(x)&=\sum_{j=0}^{r-1}\frac{f^{(j)}(x)}{j!}h^{j-s}\mu_j+h^{-s}\sum_iw_iR_i\\
 &=f^{(s)}(x)+h^{-s}\sum_iw_iR_i.
\end{align*}
Suy ra
\[
 |D_{s,h}f(x)-f^{(s)}(x)|\leq h^{-s}\sum_i|w_i||R_i|
 \leq\frac{h^{r-s}}{r!}\|f^{(r)}\|_\infty\sum_i|w_i||\tau_i|^r.
\]
Nếu $f\in C^{r+1}$, khai triển tới số hạng bậc $r$ cho
\[
 D_{s,h}f(x)-f^{(s)}(x)
 =\frac{h^{r-s}\mu_r}{r!}f^{(r)}(x)+h^{-s}\sum_iw_i\widetilde R_i,
\]
\[
 |\widetilde R_i|\leq\frac{h^{r+1}|\tau_i|^{r+1}}{(r+1)!}\|f^{(r+1)}\|_\infty.
\]

Đặt $R=h^{-s}\sum_iw_i\widetilde R_i$. Khi đó
\[
 |R|\leq h^{-s}\sum_i|w_i||\widetilde R_i|
 \leq\frac{h^{r+1-s}}{(r+1)!}\|f^{(r+1)}\|_\infty\sum_i|w_i||\tau_i|^{r+1}.
\]
\end{chungminh}

\begin{menhde}{}{dnum-symmetry}
Nếu tập mốc chuẩn đối xứng, $\tau_{\sigma(i)}=-\tau_i$, thì $w_{\sigma(i)}=(-1)^s w_i$. Do đó $\mu_j=0$ khi $j+s$ lẻ.
\end{menhde}
\begin{chungminh}
Đặt $\widetilde w_i=(-1)^sw_{\sigma(i)}$. Với $0\leq j<k$,
\[
 \sum_i\widetilde w_i\tau_i^j=(-1)^s\sum_iw_{\sigma(i)}(-\tau_{\sigma(i)})^j
 =(-1)^{s+j}\sum_iw_i\tau_i^j
 =(-1)^{s+j}s!\delta_{j,s}=s!\delta_{j,s}.
\]
Tính duy nhất trong định lý \ref{dl:dnum-moments} cho $\widetilde w=w$. Với mọi $j\geq0$,
\[
 \mu_j=\sum_iw_{\sigma(i)}\tau_{\sigma(i)}^j=(-1)^{s+j}\mu_j.
\]
Nếu $j+s$ lẻ thì $2\mu_j=0$.
\end{chungminh}

\Needspace{9\baselineskip}
\begin{vidu}
Với $\boldsymbol\tau=(-2,-1,0,1,2)^T$, hai nghiệm hệ moment là
\[
 \boldsymbol w^{(1)}=\frac1{12}(1,-8,0,8,-1)^T,
 \qquad \boldsymbol w^{(2)}=\frac1{12}(-1,16,-30,16,-1)^T.
\]
Các moment tương ứng thỏa
\[
 \mu^{(1)}_j=\delta_{j,1}\ (0\leq j\leq4),\quad\mu^{(1)}_5=-4,
\]
\[
 \mu^{(2)}_j=2\delta_{j,2}\ (0\leq j\leq5),\quad\mu^{(2)}_6=-8.
\]
Vì vậy, khi các đạo hàm tiếp theo liên tục, các số hạng sai số đầu là $-h^4f^{(5)}(x)/30$ và $-h^4f^{(6)}(x)/90$. Chặn luôn đúng từ \eqref{eq:dnum-truncation-bound} là $h^4\lVert f^{(5)}\rVert_\infty/18$ và $h^4\lVert f^{(6)}\rVert_\infty/54$.
\end{vidu}

\Needspace{7\baselineskip}
Cho bảng $x_i=a+ih$, $0\leq i\leq N$, và cần đạo hàm tại $x_j$. Một công thức $n+1$ điểm cần một khối mốc $x_m,\ldots,x_{m+n}$ chứa $x_j$. Đặt
\[
 \ell=j-m,\qquad r=n-\ell,\qquad
 \tau_i=i-\ell\quad(0\leq i\leq n).
\]
Từ Định lý \ref{dl:dnum-moments}, đạo hàm cấp $s\leq n$ của đa thức nội suy trên khối này là
\begin{equation}\label{eq:dnum-block-derivative}
 p_m^{(s)}(x_j)=h^{-s}\sum_{i=0}^n w_i^{(s,\ell)}y_{m+i},\qquad
 \sum_{i=0}^nw_i^{(s,\ell)}(i-\ell)^q=s!\delta_{q,s}\quad(0\leq q\leq n).
\end{equation}
Chỉ số $s$ là cấp đạo hàm; $n$ là bậc đa thức nội suy và $n+1$ là số giá trị được sử dụng. Công thức sai số riêng cho đạo hàm cấp $1$ tại một mốc được suy ra trực tiếp như sau.

\begin{dinhly}{}{dnum-node-error}
Cho $z_0<\cdots<z_n$, $n\geq1$, $f\in C^{n+1}([z_0,z_n])$, và $p\in\mathcal P_n$ nội suy $f$ tại các $z_i$. Với mỗi $0\leq j\leq n$, có $\xi_j\in(z_0,z_n)$ sao cho
\[
 f'(z_j)-p'(z_j)=\frac{f^{(n+1)}(\xi_j)}{(n+1)!}\prod_{i\ne j}(z_j-z_i).
\]
Với khối mốc đều nói trên,
\begin{equation}\label{eq:dnum-node-uniform-error}
 p_m'(x_j)-f'(x_j)
 =\frac{(-1)^{r+1}\ell!r!}{(n+1)!}h^n f^{(n+1)}(\xi_j).
\end{equation}
\end{dinhly}
\begin{chungminh}
Đặt $W(t)=\prod_{i=0}^n(t-z_i)$. Từ $W(t)=(t-z_j)\prod_{i\ne j}(t-z_i)$,
\[
 W'(z_j)=\prod_{i\ne j}(z_j-z_i)\ne0.
\]
Lấy $A=(f'(z_j)-p'(z_j))/W'(z_j)$ và $g=f-p-AW$. Khi đó
\[
 g(z_i)=0\quad(0\leq i\leq n),\qquad g'(z_j)=f'(z_j)-p'(z_j)-AW'(z_j)=0.
\]
Rolle trên $[z_i,z_{i+1}]$ cho $\eta_i\in(z_i,z_{i+1})$ với $g'(\eta_i)=0$. Ta có
\[
 \eta_0<\cdots<\eta_{n-1},\qquad \eta_i\ne z_j\quad(0\leq i<n).
\]
Sắp $\{\eta_0,\ldots,\eta_{n-1},z_j\}$ thành $n+1$ điểm $t_0<\cdots<t_n$ trong $[z_0,z_n]$. Bổ đề \ref{bd:rolle} áp dụng cho $g'\in C^n([t_0,t_n])$ cho
\[
 \xi_j\in(t_0,t_n)\subset(z_0,z_n),\qquad
 0=(g')^{(n)}(\xi_j)=f^{(n+1)}(\xi_j)-A(n+1)!.
\]
Thế $A$:
\[
 f'(z_j)-p'(z_j)=AW'(z_j)
 =\frac{f^{(n+1)}(\xi_j)}{(n+1)!}\prod_{i\ne j}(z_j-z_i).
\]
Trên khối đều,
\[
 \prod_{i\ne\ell}(x_j-x_{m+i})
 =h^n\left(\prod_{i=0}^{\ell-1}(\ell-i)\right)
 \left(\prod_{i=\ell+1}^{n}(\ell-i)\right)
 =(-1)^rh^n\ell!r!.
\]
Thế vào công thức thứ nhất và đổi dấu cho \eqref{eq:dnum-node-uniform-error}.
\end{chungminh}

\begin{menhde}{}{dnum-forward-weights}
Với các mốc $x,x+h,\ldots,x+nh$, công thức đạo hàm cấp $1$ là
\begin{equation}\label{eq:dnum-forward-general}
 D_n^+f(x)=\frac1h\left[-H_nf(x)+\sum_{i=1}^n\frac{(-1)^{i+1}}i\binom ni f(x+ih)\right],
 \qquad H_n=\sum_{i=1}^n\frac1i.
\end{equation}
Nó chính xác tới bậc $n$ và
\[
 D_n^+f(x)-f'(x)=\frac{(-1)^{n+1}}{n+1}h^nf^{(n+1)}(\xi),\qquad\xi\in(x,x+nh).
\]
Công thức phía trái được lấy bằng cách đổi $h$ thành $-h$ trong \eqref{eq:dnum-forward-general}; trong mẫu số khi đó cũng phải dùng $-h$.
\end{menhde}
\begin{chungminh}
Tại các mốc chuẩn $0,1,\ldots,n$, đa thức cơ sở thứ $0$ là
\[
 L_0(t)=\prod_{j=1}^n(1-t/j),\qquad L_0'(0)=-\sum_{j=1}^n1/j=-H_n.
\]
Với $i\geq1$, tách nhân tử $t$ trong $L_i$ cho
\[
 L_i'(0)=\frac{\prod_{1\leq j\leq n,\ j\ne i}(-j)}{\prod_{0\leq j\leq n,\ j\ne i}(i-j)}
 =\frac{(-1)^{n-1}n!/i}{i!(-1)^{n-i}(n-i)!}
 =\frac{(-1)^{i+1}}i\binom ni.
\]

Định lý \ref{dl:dnum-moments} cho $D_n^+f=h^{-1}\sum_iL_i'(0)f(x+ih)$, chính xác trên $\mathcal P_n$. Trong \eqref{eq:dnum-node-uniform-error}, $\ell=0$, $r=n$ cho $\ell!r!/(n+1)!=1/(n+1)$, nên
\[
 D_n^+f(x)-f'(x)=\frac{(-1)^{n+1}}{n+1}h^nf^{(n+1)}(\xi).
\]
Với $g(u)=f(x-u)$,
\[
 -D_n^+g(0)
 =\frac1{-h}\left[-H_nf(x)+\sum_{i=1}^n\frac{(-1)^{i+1}}i\binom ni f(x-ih)\right],
 \qquad -g'(0)=f'(x).
\]
Đây là công thức phía trái.
\end{chungminh}

\Needspace{9\baselineskip}
\begin{vidu}
Một số vector trong \eqref{eq:dnum-forward-general} và chặn nhiễu dữ liệu $\eta h^{-1}\|\boldsymbol w\|_1$ là
\begin{center}\small
\begin{tabular}{c|c|c|c}
\toprule
$n$ & $(w_0,\ldots,w_n)$ & $D_n^+-f'$ & $\|\boldsymbol w\|_1$\\
\midrule
$3$ & $(-11,18,-9,2)/6$ & $h^3f^{(4)}(\xi)/4$ & $20/3$\\
$4$ & $(-25,48,-36,16,-3)/12$ & $-h^4f^{(5)}(\xi)/5$ & $32/3$\\
$5$ & $(-137,300,-300,200,-75,12)/60$ & $h^5f^{(6)}(\xi)/6$ & $256/15$\\
\bottomrule
\end{tabular}\end{center}
Nếu $|f^{(n+1)}|\leq M_{n+1}$, cột sai số cho chặn $M_{n+1}h^n/(n+1)$. Cột cuối cho thấy phần nhiễu cần được xét đồng thời với sai số cắt cụt.
\end{vidu}

\begin{menhde}{}{dnum-balanced-block}
Với $n+1$ mốc đều và chặn chung $|f^{(n+1)}|\leq M_{n+1}$, hệ số chặn trong \eqref{eq:dnum-node-uniform-error} nhỏ nhất khi $|\ell-r|\leq1$. Với $0\leq j\leq N$, $N\geq n$, một cách chọn khối gần cân bằng là
\[
 m=\min\{\max\{j-\lfloor n/2\rfloor,0\},N-n\},\qquad \ell=j-m.
\]
\end{menhde}
\begin{chungminh}
Đặt $F(\ell)=\ell!(n-\ell)!$, $0\leq\ell\leq n$. Với $0\leq\ell<n$,
\[
 \frac{F(\ell+1)}{F(\ell)}=\frac{\ell+1}{n-\ell}
 \begin{cases}<1,&2\ell+1<n,\\=1,&2\ell+1=n,\\>1,&2\ell+1>n.\end{cases}
\]
Do đó $F$ nhỏ nhất tại $\ell=\lfloor n/2\rfloor$ hoặc $\ell=\lceil n/2\rceil$; các chỉ số này thỏa $|\ell-(n-\ell)|\leq1$.

Đặt $q=\lfloor n/2\rfloor$. Công thức $m$ được viết thành
\[
 m=\begin{cases}0,&j<q,\\j-q,&q\leq j\leq N-n+q,\\N-n,&j>N-n+q.\end{cases}
\]
Trong ba trường hợp, tương ứng $\ell=j$, $\ell=q$, $\ell=j-N+n$; từ $0\leq j\leq N$ suy ra
\[
 0\leq m\leq N-n,\qquad0\leq\ell\leq n,\qquad j=m+\ell\in[m,m+n].
\]
Khối ở giữa có chỉ số cân bằng; hai trường hợp còn lại là các khối sát biên.
\end{chungminh}

\begin{menhde}{}{dnum-node-weights}
Với mốc chuẩn $0,1,\ldots,n$, đạo hàm tại mốc $\ell$ có các trọng số
\[
 w_\ell=H_\ell-H_{n-\ell},\qquad
 w_i=\frac{(-1)^{i-\ell}\ell!(n-\ell)!}{(\ell-i)i!(n-i)!}\quad(i\ne\ell),\qquad H_0=0.
\]
\end{menhde}
\begin{chungminh}
Với đa thức cơ sở $L_i$ của các mốc $0,\ldots,n$, định lý \ref{dl:dnum-moments} sau phép tịnh tiến cho $w_i=L_i'(\ell)$. Khi $i=\ell$, quy tắc đạo hàm tích cho
\[
 L_\ell'(\ell)=\sum_{q\ne\ell}\frac1{\ell-q}
 =\sum_{q=1}^\ell\frac1q-\sum_{q=1}^{n-\ell}\frac1q.
\]
Khi $i\ne\ell$, viết $L_i(t)=W(t)/((t-i)W'(i))$ với $W(t)=\prod_{q=0}^n(t-q)$, rồi đạo hàm tại nghiệm $\ell$ của $W$:
\[
 L_i'(\ell)=\frac{W'(\ell)}{(\ell-i)W'(i)}
 =\frac{(-1)^{n-\ell}\ell!(n-\ell)!}{(\ell-i)(-1)^{n-i}i!(n-i)!}.
\]
Vì $(-1)^{n-\ell}/(-1)^{n-i}=(-1)^{i-\ell}$, đẳng thức này cho công thức đã nêu.
\end{chungminh}

\Needspace{9\baselineskip}
\begin{vidu}
Cho $f(t)=t^4$, $h=1/4$, cần $f'(1/2)$ từ bốn mốc $1/4,1/2,3/4,1$. Với $\ell=1$, $r=2$, giải hệ moment hoặc lấy đạo hàm đa thức nội suy cho $\boldsymbol w=(-1/3,-1/2,1,-1/6)^T$:
\begin{center}
\begin{tabular}{c|cccc}
\toprule
$i$ & $0$ & $1$ & $2$ & $3$\\
\midrule
$x_{m+i}$ & $1/4$ & $1/2$ & $3/4$ & $1$\\
$\tau_i$ & $-1$ & $0$ & $1$ & $2$\\
$y_{m+i}$ & $1/256$ & $1/16$ & $81/256$ & $1$\\
$w_i y_{m+i}$ & $-1/768$ & $-1/32$ & $81/256$ & $-1/6$\\
\bottomrule
\end{tabular}\end{center}
Do đó $p_m'(1/2)=h^{-1}(15/128)=15/32$, còn $f'(1/2)=1/2$. Sai số $-1/32$ bằng $-h^3 f^{(4)}/12$. Nếu dùng bốn mốc về một phía với cùng $h$ và cùng chặn $M_4$, hệ số chặn là $1/4$ thay vì $1/12$.
\end{vidu}

\begin{menhde}{}{dnum-data}
Cho $\widetilde y_i=f(x_i)+\eta_i$, $|\eta_i|\leq\varepsilon$, và $\widetilde D=h^{-s}\sum_iw_i\widetilde y_i$. Khi đó
\[
 \max_{|\eta_i|\leq\varepsilon}|\widetilde D-D_{s,h}f(x)|
 =\varepsilon h^{-s}\sum_i|w_i|.
\]
Giả sử $w_i$, $\widetilde y_i$ và $H=h^s$ được biểu diễn chính xác, các phép nhân, cộng, chia thỏa mô hình \eqref{eq:rounding-model}, không tràn và không mất chặn sai số tương đối. Nếu tính $k$ tích rồi cộng tuần tự và chia cho $H$, với $(k+1)u<1$, thì
\[
 |\widehat D-\widetilde D|\leq\gamma_{k+1}h^{-s}\sum_i|w_i||\widetilde y_i|.
\]
\end{menhde}
\begin{chungminh}
Ta có
\[
 \widetilde D-D_{s,h}f=h^{-s}\sum_iw_i\eta_i,
 \quad |\widetilde D-D_{s,h}f|\leq\varepsilon h^{-s}\sum_i|w_i|.
\]
Với $\eta_i=\varepsilon\operatorname{sgn}(w_i)$,
\[
 h^{-s}\sum_iw_i\eta_i=\varepsilon h^{-s}\sum_i|w_i|,
\]
nên chặn đạt được.

Viết các phép tính thực tế là
\[
 t_i=w_i\widetilde y_i(1+\delta_i),\quad s_0=t_0,\quad
 s_j=(s_{j-1}+t_j)(1+e_j)\ (1\leq j<k),\quad
 \widehat D=H^{-1}s_{k-1}(1+z).
\]

Với tích rỗng bằng $1$, ta có $s_0=w_0\widetilde y_0(1+\delta_0)$. Nếu
\[
 s_{j-1}=\sum_{i=0}^{j-1}w_i\widetilde y_i(1+\delta_i)
 \prod_{\ell=\max\{1,i\}}^{j-1}(1+e_\ell),
\]
thì
\begin{align*}
 s_j&=(s_{j-1}+w_j\widetilde y_j(1+\delta_j))(1+e_j)\\
 &=\sum_{i=0}^jw_i\widetilde y_i(1+\delta_i)
 \prod_{\ell=\max\{1,i\}}^j(1+e_\ell).
\end{align*}
Quy nạp tới $j=k-1$ và nhân với $H^{-1}(1+z)$ cho
\[
 \widehat D=H^{-1}\sum_{i=0}^{k-1}w_i\widetilde y_i\theta_i,
\quad
 \theta_i=(1+\delta_i)(1+z)\prod_{j=\max\{1,i\}}^{k-1}(1+e_j).
\]
Mỗi tích có không quá $k+1$ thừa số, từng sai số có trị tuyệt đối không quá $u$. Bổ đề \ref{bd:gamma} cho $|\theta_i-1|\leq\gamma_{k+1}$. Do đó
\[
 |\widehat D-\widetilde D|\leq H^{-1}\sum_i|w_i||\widetilde y_i||\theta_i-1|
 \leq\gamma_{k+1}h^{-s}\sum_i|w_i||\widetilde y_i|.
\]
\end{chungminh}

\begin{menhde}{}{dnum-step}
Dưới các giả thiết của mệnh đề \ref{md:dnum-data}, nếu sai số cắt cụt bị chặn bởi $Ch^p$, $|f(x_i)|\leq F$, thì
\[
 |\widehat D-f^{(s)}(x)|\leq Ch^p+Bh^{-s},\qquad
 B=\Bigl(\sum_i|w_i|\Bigr)\bigl(\varepsilon+\gamma_{k+1}(F+\varepsilon)\bigr).
\]
Với $C,B,p,s>0$ không phụ thuộc $h$, trên $0<h\leq H_0$ chặn bên phải đạt nhỏ nhất tại
\[
 h_{\mathrm{opt}}=\min\left\{H_0,\left(\frac{sB}{pC}\right)^{1/(p+s)}\right\}.
\]
\end{menhde}
\begin{chungminh}
Từ $|\widetilde y_i|\leq|f(x_i)|+|\eta_i|\leq F+\varepsilon$ và Mệnh đề \ref{md:dnum-data},
\begin{align*}
 |\widehat D-f^{(s)}(x)|
 &\leq|D_{s,h}f(x)-f^{(s)}(x)|+|\widetilde D-D_{s,h}f(x)|+|\widehat D-\widetilde D|\\
 &\leq Ch^p+\varepsilon h^{-s}\sum_i|w_i|
 +\gamma_{k+1}h^{-s}(F+\varepsilon)\sum_i|w_i|\\
 &=Ch^p+Bh^{-s}.
\end{align*}
Với $\phi(h)=Ch^p+Bh^{-s}$ và $h_*=(sB/(pC))^{1/(p+s)}$,
\[
 \phi'(h)=h^{-s-1}(pCh^{p+s}-sB)
 \begin{cases}<0,&0<h<h_*,\\=0,&h=h_*,\\>0,&h>h_* .\end{cases}
\]
Vì vậy
\[
 \min_{0<h\leq H_0}\phi(h)=
 \begin{cases}\phi(H_0),&H_0<h_*,\\\phi(h_*),&H_0\geq h_* ,\end{cases}
\]
cho $h_{\mathrm{opt}}=\min\{H_0,h_*\}$. Nếu $h_*\leq H_0$, $pCh_*^{p+s}=sB$, nên
\[
 Bh_*^{-s}=\frac psCh_*^p,\qquad
 \phi(h_*)=\left(1+\frac ps\right)Ch_*^p.
\]
\end{chungminh}

\Needspace{9\baselineskip}
\begin{vidu}
Các chặn trên đoạn chứa mọi mốc có dạng
\begin{center}
\begin{tabular}{c|c|c|c}
\toprule
Công thức & $s$ & Chặn cắt cụt & $\sum_i|w_i|$\\
\midrule
$D_+,D_-$ & $1$ & $\frac12\lVert f''\rVert_\infty h$ & $2$\\
$D_0$ & $1$ & $\frac16\lVert f'''\rVert_\infty h^2$ & $1$\\
$D_{+,3},D_{-,3}$ & $1$ & $\frac13\lVert f'''\rVert_\infty h^2$ & $4$\\
$D_2$ & $2$ & $\frac1{12}\lVert f^{(4)}\rVert_\infty h^2$ & $4$\\
\bottomrule
\end{tabular}
\end{center}
Nếu chỉ xét nhiễu dữ liệu của công thức $D_0$, lấy $C=\lVert f'''\rVert_\infty/6$, $B=\varepsilon$ cho $h_*=(3\varepsilon/\lVert f'''\rVert_\infty)^{1/3}$ khi $\varepsilon>0$ và $\lVert f'''\rVert_\infty>0$.
\end{vidu}

\Needspace{9\baselineskip}
Giả sử chặn đã xác định là $\Phi(h)=Ch^p+Bh^{-s}$ trên $0<h\leq H_0$, với $C,B,p,s>0$. Mệnh đề \ref{md:dnum-step} cho $h_{\mathrm{opt}}$. Do đó
\[
 \text{có }h\in(0,H_0]\text{ thỏa }\Phi(h)\leq\varepsilon
 \quad\Longleftrightarrow\quad \Phi(h_{\mathrm{opt}})\leq\varepsilon.
\]
Đây là điều kiện đối với chặn đã có. Các giá trị $C$, $B$ phải chứa chặn đạo hàm, nhiễu dữ liệu và làm tròn tương ứng.

\Needspace{9\baselineskip}
\begin{vidu}
Với $f(t)=t^3$, công thức trung tâm hai điểm có sai số cắt cụt đúng bằng $h^2$. Nếu sai số mỗi giá trị không quá $\eta=10^{-6}$ và tạm xét số học chính xác, chặn tổng là $h^2+10^{-6}/h$:
\begin{center}
\begin{tabular}{c|c|c|c}
\toprule
$h$ & $h^2$ & $10^{-6}/h$ & Chặn tổng\\
\midrule
$10^{-1}$ & $10^{-2}$ & $10^{-5}$ & $0.01001$\\
$10^{-2}$ & $10^{-4}$ & $10^{-4}$ & $0.0002$\\
$10^{-3}$ & $10^{-6}$ & $10^{-3}$ & $0.001001$\\
$(5\cdot10^{-7})^{1/3}$ & $\approx6.2996\cdot10^{-5}$ & $\approx1.2599\cdot10^{-4}$ & $\approx1.8899\cdot10^{-4}$\\
\bottomrule
\end{tabular}\end{center}
Bước tối ưu là $(5\cdot10^{-7})^{1/3}$. Chặn nhỏ nhất vẫn lớn hơn $10^{-4}$; vì thế chặn này chưa chứng nhận được yêu cầu $\varepsilon=10^{-4}$ với bất kỳ bước nào.
\end{vidu}

\section{Các công thức tính gần đúng tích phân xác định}
Các công thức, cách chọn lưới và cách xử lý bảng dữ liệu theo bài giảng \cite{integral-slides}.
Cho $a<b$, $\ell=b-a$ và $I(f)=\int_a^b f(x)\,dx$. Với các mốc chuẩn $\tau_i\in[0,1]$, xét
\begin{equation}\label{eq:quad-standard}
 Q(f)=\ell\sum_{i=0}^{k-1}\lambda_i f(a+\ell\tau_i).
\end{equation}
Phép đổi biến $x=a+\ell t$ cho $I(f)=\ell\int_0^1 f(a+\ell t)\,dt$. Vì vậy, việc chọn trọng số được đưa về tích phân trên $[0,1]$.

\begin{dinhnghia}{}{quad-degree}
Bậc chính xác đại số của $Q$ là số nguyên $q\geq0$ sao cho $Q$ chính xác với mọi đa thức bậc không quá $q$ và có ít nhất một đa thức bậc $q+1$ mà $I(f)\ne Q(f)$.
\end{dinhnghia}

\begin{menhde}{}{quad-moments}
Công thức \eqref{eq:quad-standard} chính xác tới bậc $q$ khi và chỉ khi
\begin{equation}\label{eq:quad-moments}
 \sum_{i=0}^{k-1}\lambda_i\tau_i^j=\frac1{j+1},\qquad j=0,\ldots,q.
\end{equation}
Với $k$ mốc thực đôi một phân biệt, hệ moment tới bậc $k-1$ có duy nhất nghiệm, kể cả khi các mốc nằm ngoài $[0,1]$:
\[
 \lambda_i=\int_0^1 L_i(t)\,dt,\qquad
 L_i(t)=\prod_{j\ne i}\frac{t-\tau_j}{\tau_i-\tau_j}.
\]
Nếu $L_i(t)=\sum_{j=0}^{k-1}c_{i,j}t^j$, thì $\lambda_i=\sum_{j=0}^{k-1}c_{i,j}/(j+1)$.
\end{menhde}
\begin{chungminh}
Đặt $Q_0(v)=\sum_i\lambda_iv(\tau_i)$. Với $v(t)=\sum_{j=0}^qc_jt^j$,
\begin{align*}
 \int_0^1v(t)dt-Q_0(v)
 &=\sum_{j=0}^qc_j\left(\left[\frac{t^{j+1}}{j+1}\right]_0^1-\sum_i\lambda_i\tau_i^j\right)\\
 &=\sum_{j=0}^qc_j\left(\frac1{j+1}-\sum_i\lambda_i\tau_i^j\right).
\end{align*}
Biểu thức này bằng $0$ với mọi $(c_0,\ldots,c_q)$ khi và chỉ khi từng ngoặc bằng $0$, tức \eqref{eq:quad-moments}. Với $q=k-1$, hệ là $V^T\boldsymbol\lambda=\boldsymbol\mu$, $V_{ij}=\tau_i^j$, $\mu_j=1/(j+1)$. Định lý \ref{dl:vandermonde} cho
\[
 \det V^T=\det V=\prod_{i<j}(\tau_j-\tau_i)\ne0.
\]
Theo \eqref{eq:cardinal}, $t^j=\sum_i\tau_i^jL_i(t)$ với $j<k$, nên
\[
 \sum_i\tau_i^j\int_0^1L_i(t)dt=\int_0^1t^jdt=\frac1{j+1}.
\]
Vậy $(\int_0^1L_i)_i$ là nghiệm duy nhất của hệ. Với $L_i(t)=\sum_{j=0}^{k-1}c_{i,j}t^j$,
\[
 \lambda_i=\sum_{j=0}^{k-1}c_{i,j}\int_0^1t^jdt=\sum_{j=0}^{k-1}\frac{c_{i,j}}{j+1}.
\]
\end{chungminh}

\begin{menhde}{}{quad-kernel}
Giả sử \eqref{eq:quad-moments} đúng với $j=0,\ldots,r-1$, $r\geq1$. Với $f\in C^r([a,b])$,
\begin{equation}\label{eq:quad-kernel}
 I(f)-Q(f)=\ell^{r+1}\int_0^1K_r(t)f^{(r)}(a+\ell t)\,dt,
\end{equation}
\[
 K_r(t)=\frac{(1-t)^r}{r!}-\sum_i\lambda_i\frac{(\tau_i-t)_+^{r-1}}{(r-1)!}.
\]
Đặt $C_r=\int_0^1|K_r(t)|dt$. Khi đó
\[
 |I(f)-Q(f)|\leq C_r\ell^{r+1}\lVert f^{(r)}\rVert_\infty.
\]
Nếu kernel có một dấu và $C_r>0$, thì có $\xi\in(a,b)$ sao cho
\begin{equation}\label{eq:quad-kernel-mean}
 I(f)-Q(f)=\frac{\ell^{r+1}}{r!}
 \left(\frac1{r+1}-\sum_i\lambda_i\tau_i^r\right)f^{(r)}(\xi).
\end{equation}
\end{menhde}
\begin{chungminh}
Đặt $g(t)=f(a+\ell t)$ và $\mathcal L(v)=\int_0^1v(t)dt-\sum_i\lambda_iv(\tau_i)$. Với $0\leq j<r$,
\[
 \mathcal L(t^j)=\frac1{j+1}-\sum_i\lambda_i\tau_i^j=0.
\]
Tính tuyến tính cho $\mathcal L(\mathcal P_{r-1})=0$. Bổ đề \ref{bd:dnum-kernel} và quy tắc dây chuyền cho
\begin{align*}
 I(f)-Q(f)&=\ell\mathcal L(g)
 =\ell\int_0^1K_r(t)g^{(r)}(t)dt\\
 &=\ell^{r+1}\int_0^1K_r(t)f^{(r)}(a+\ell t)dt.
\end{align*}
Do đó
\[
 |I(f)-Q(f)|\leq\ell^{r+1}\|f^{(r)}\|_\infty\int_0^1|K_r(t)|dt.
\]
Từ $\int_0^1(\tau_i-t)_+^{r-1}dt=\tau_i^r/r$,
\[
 \int_0^1K_r(t)dt=\frac1{r!}\left(\frac1{r+1}-\sum_i\lambda_i\tau_i^r\right).
\]
Nếu $K_r$ có một dấu và $C_r>0$, Bổ đề \ref{bd:dnum-kernel} cho $t_*\in(0,1)$ với
\[
 I(f)-Q(f)=\ell^{r+1}f^{(r)}(a+\ell t_*)\int_0^1K_r(t)dt.
\]
Thế tích phân kernel và đặt $\xi=a+\ell t_*$ cho \eqref{eq:quad-kernel-mean}.
\end{chungminh}

\Needspace{7\baselineskip}
\subsection{Công thức điểm giữa}
Thay hàm trên một đoạn bằng giá trị tại trung điểm $c=(a+b)/2$ cho $M(f)=\ell f(c)$. Trên đoạn chuẩn, mốc là $\tau_0=1/2$ và trọng số là $\lambda_0=1$.
\begin{dinhly}{}{quad-midpoint}
Công thức
\[
 M(f)=\ell f\left(\frac{a+b}{2}\right)
\]
có bậc chính xác $1$. Với $f\in C^2([a,b])$,
\[
 I(f)-M(f)=\frac{\ell^3}{24}f''(\xi),\qquad \xi\in(a,b).
\]
Trên lưới đều $x_j=a+jh$, $h=(b-a)/N$, $N\geq1$, công thức ghép là
\begin{equation}\label{eq:quad-midpoint-composite}
 M_N(f)=h\sum_{j=0}^{N-1}f\left(x_j+\frac h2\right),
 \quad I(f)-M_N(f)=\frac{(b-a)h^2}{24}f''(\xi_N).
\end{equation}
\end{dinhly}
\begin{chungminh}
Với $\tau_0=1/2$, $\lambda_0=1$,
\[
 (\lambda_0\tau_0^j)_{j=0}^2=(1,1/2,1/4),\qquad
 \left(\int_0^1t^jdt\right)_{j=0}^2=(1,1/2,1/3).
\]
Mệnh đề \ref{md:quad-moments} cho bậc chính xác $1$. Kernel cấp $2$ là
\[
 K_2(t)=\frac{(1-t)^2}{2}-(1/2-t)_+
 =\begin{cases}t^2/2,&0\leq t\leq1/2,\\(1-t)^2/2,&1/2\leq t\leq1.\end{cases}
\]
Ta có $K_2\geq0$ và
\[
 \int_0^1K_2(t)dt
 =\left[\frac{t^3}{6}\right]_0^{1/2}
 +\left[-\frac{(1-t)^3}{6}\right]_{1/2}^{1}=\frac1{24}>0.
\]
Mệnh đề \ref{md:quad-kernel} cho sai số đơn.

Trên mỗi ô $(x_j,x_{j+1})$, đặt $\mathcal K(u)=h^2K_2((u-x_j)/h)\geq0$. Đổi biến và cộng sai số cục bộ cho
\[
 I(f)-h\sum_{j=0}^{N-1}f(x_j+h/2)=\int_a^b\mathcal K(u)f''(u)du,
\]
\[
 \int_a^b\mathcal K(u)du
 =\sum_{j=0}^{N-1}h^3\int_0^1K_2(t)dt
 =\frac{Nh^3}{24}=\frac{(b-a)h^2}{24}>0.
\]
Bổ đề \ref{bd:dnum-kernel} cho $\xi_N\in(a,b)$ và \eqref{eq:quad-midpoint-composite}.
\end{chungminh}

\Needspace{7\baselineskip}
\subsection{Công thức hình thang}
Đa thức nội suy tuyến tính ở hai đầu đoạn là
\[
 p_1(x)=\frac{b-x}{\ell}f(a)+\frac{x-a}{\ell}f(b),\qquad
 \int_a^b p_1(x)dx=\frac\ell2\bigl(f(a)+f(b)\bigr).
\]
\begin{dinhly}{}{quad-trapezoid}
Công thức
\[
 T(f)=\frac\ell2\bigl(f(a)+f(b)\bigr)
\]
có bậc chính xác $1$. Với $f\in C^2([a,b])$,
\[
 I(f)-T(f)=-\frac{\ell^3}{12}f''(\xi),\qquad\xi\in(a,b).
\]
Với $N\geq1$, $h=(b-a)/N$, $x_j=a+jh$, công thức ghép trên lưới đều là
\begin{equation}\label{eq:quad-trapezoid-composite}
 T_N(f)=h\left(\frac{f(a)+f(b)}2+\sum_{j=1}^{N-1}f(x_j)\right),
\end{equation}
\[
 I(f)-T_N(f)=-\frac{(b-a)h^2}{12}f''(\xi_N).
\]
\end{dinhly}
\begin{chungminh}
Với $\boldsymbol\tau=(0,1)^T$, $\boldsymbol\lambda=(1/2,1/2)^T$,
\[
 \left(\sum_i\lambda_i\tau_i^j\right)_{j=0}^2=(1,1/2,1/2),\qquad
 \left(\int_0^1t^jdt\right)_{j=0}^2=(1,1/2,1/3).
\]
Mệnh đề \ref{md:quad-moments} cho bậc chính xác $1$. Kernel cấp $2$ thỏa
\[
 K_2(t)=\frac{(1-t)^2}{2}-\frac{1-t}{2}=-\frac{t(1-t)}2\leq0,
\]
\[
 \int_0^1K_2(t)dt=-\frac12[t^2/2-t^3/3]_0^1=-\frac1{12}<0.
\]
Mệnh đề \ref{md:quad-kernel} cho sai số đơn. Tổng các công thức cục bộ là
\begin{align*}
 \sum_{j=0}^{N-1}\frac h2(f(x_j)+f(x_{j+1}))
 &=\frac h2\left(f(x_0)+f(x_N)+2\sum_{j=1}^{N-1}f(x_j)\right)=T_N(f).
\end{align*}
Trên $(x_j,x_{j+1})$, đặt $\mathcal K(u)=h^2K_2((u-x_j)/h)\leq0$. Đổi biến trong sai số cục bộ cho
\[
 I(f)-T_N(f)=\int_a^b\mathcal K(u)f''(u)du,
\]
\[
 \int_a^b\mathcal K(u)du=Nh^3\int_0^1K_2(t)dt
 =-\frac{Nh^3}{12}=-\frac{(b-a)h^2}{12}<0.
\]
Bổ đề \ref{bd:dnum-kernel} cho công thức với $\xi_N\in(a,b)$.
\end{chungminh}

\Needspace{7\baselineskip}
\subsection{Công thức Simpson}
Trên đoạn chuẩn, tích phân ba đa thức cơ sở tại $0,1/2,1$ lần lượt là
\[
 \int_0^1(2t-1)(t-1)dt=\frac16,\quad
 \int_0^1 4t(1-t)dt=\frac23,\quad
 \int_0^1t(2t-1)dt=\frac16.
\]
Tích phân đa thức nội suy bậc $2$ vì thế cho các trọng số $1,4,1$.
\begin{dinhly}{}{quad-simpson}
Công thức
\[
 S(f)=\frac\ell6\left(f(a)+4f\left(\frac{a+b}{2}\right)+f(b)\right)
\]
có bậc chính xác $3$. Với $f\in C^4([a,b])$,
\[
 I(f)-S(f)=-\frac{\ell^5}{2880}f^{(4)}(\xi),\qquad\xi\in(a,b).
\]
Nếu $N\geq2$ chẵn, $h=(b-a)/N$, $x_j=a+jh$, công thức ghép là
\begin{equation}\label{eq:quad-simpson-composite}
 S_N(f)=\frac h3\left[f(x_0)+f(x_N)+4\sum_{j=1}^{N/2}f(x_{2j-1})
 +2\sum_{j=1}^{N/2-1}f(x_{2j})\right],
\end{equation}
\[
 I(f)-S_N(f)=-\frac{(b-a)h^4}{180}f^{(4)}(\xi_N).
\]
\end{dinhly}
\begin{chungminh}
Với $\tau=(0,1/2,1)$, $\lambda=(1/6,2/3,1/6)$,
\begin{align*}
 \left(\sum_i\lambda_i\tau_i^j\right)_{j=0}^3
 &=\left(1,\frac23\frac12+\frac16,\frac23\frac14+\frac16,\frac23\frac18+\frac16\right)\\
 &=(1,1/2,1/3,1/4),\qquad \sum_i\lambda_i\tau_i^4=\frac5{24}\ne\frac15.
\end{align*}
Do đó bậc chính xác bằng $3$. Kernel cấp $4$ là
\begin{align*}
 K_4(t)&=\frac{(1-t)^4}{24}-\frac{(1/2-t)_+^3}{9}-\frac{(1-t)^3}{36}\\
 &=\begin{cases}
 \dfrac{t^3(t-2/3)}{24},&0\leq t\leq1/2,\\[1mm]
 \dfrac{(1-t)^3(1/3-t)}{24},&1/2\leq t\leq1.
 \end{cases}
\end{align*}
Với $0\leq t\leq1/2$, $t^3\geq0$ và $t-2/3\leq-1/6$; với $1/2\leq t\leq1$, $(1-t)^3\geq0$ và $1/3-t\leq-1/6$. Do đó $K_4\leq0$, và
\[
 \int_0^1K_4(t)dt=\frac1{24}\left(\frac15-\frac5{24}\right)=-\frac1{2880}.
\]
Mệnh đề \ref{md:quad-kernel} cho sai số đơn.


Cộng trên các khối $[x_{2j},x_{2j+2}]$:
\begin{align*}
 \sum_{j=0}^{N/2-1}\frac h3(f(x_{2j})+4f(x_{2j+1})+f(x_{2j+2}))
 &=\frac h3\left[f(x_0)+f(x_N)+4\sum_{j=0}^{N/2-1}f(x_{2j+1})\right.\\
 &\qquad\left. +2\sum_{j=1}^{N/2-1}f(x_{2j})\right]=S_N(f).
\end{align*}
Trên $(x_{2j},x_{2j+2})$, đặt $\mathcal K(u)=(2h)^4K_4((u-x_{2j})/(2h))\leq0$. Đổi biến trong các sai số cục bộ cho
\[
 I(f)-S_N(f)=\int_a^b\mathcal K(u)f^{(4)}(u)du,
\]
\[
 \int_a^b\mathcal K(u)du
 =\frac N2(2h)^5\left(-\frac1{2880}\right)
 =-\frac{Nh^5}{180}=-\frac{(b-a)h^4}{180}<0.
\]
Bổ đề \ref{bd:dnum-kernel} cho $\xi_N\in(a,b)$.

\end{chungminh}

\Needspace{9\baselineskip}
\begin{vidu}
Với $f(x)=x^4$ trên $[0,1]$, $I(f)=1/5$ và
\begin{center}
\begin{tabular}{c|c|c}
\toprule
Công thức & Giá trị & $I-Q$\\
\midrule
Điểm giữa & $1/16$ & $11/80$\\
Hình thang & $1/2$ & $-3/10$\\
Simpson & $5/24$ & $-1/120$\\
\bottomrule
\end{tabular}
\end{center}
Ở đây $f^{(4)}=24$, nên sai số Simpson bằng $-24/2880=-1/120$. Trên lưới đều $N$ chẵn, sai số Simpson ghép là $-2/(15N^4)$.
\end{vidu}

\begin{menhde}{}{quad-nonuniform}
Cho $N\geq1$, phân hoạch $a=t_0<\cdots<t_N=b$, $\ell_j=t_{j+1}-t_j$. Cộng các công thức đơn trên mỗi đoạn cho $M_{\mathcal T}$, $T_{\mathcal T}$, $S_{\mathcal T}$; trong công thức Simpson, điểm thứ ba là trung điểm thực của từng đoạn. Nếu $f\in C^2([a,b])$, thì
\[
 I-M_{\mathcal T}=\frac1{24}\left(\sum_j\ell_j^3\right)f''(\xi_M),\quad
 I-T_{\mathcal T}=-\frac1{12}\left(\sum_j\ell_j^3\right)f''(\xi_T),
\]
Nếu $f\in C^4([a,b])$, thì
\[
 I-S_{\mathcal T}=-\frac1{2880}\left(\sum_j\ell_j^5\right)f^{(4)}(\xi_S),
\]
với $\xi_M,\xi_T,\xi_S\in(a,b)$ trong từng công thức tương ứng.
\end{menhde}
\begin{chungminh}
Gọi $K_M,K_T,K_S$ là các kernel chuẩn đã tính. Các cấp và tích phân của chúng là
\[
 (r_M,c_M)=(2,1/24),\quad(r_T,c_T)=(2,-1/12),\quad(r_S,c_S)=(4,-1/2880),
\]
với $c_A=\int_0^1K_A$, $K_M\geq0$, $K_T,K_S\leq0$. Với $A\in\{M,T,S\}$, đặt trên $(t_j,t_{j+1})$
\[
 \mathcal K_A(u)=\ell_j^{r_A}K_A((u-t_j)/\ell_j).
\]
Đổi biến trong từng sai số đơn và cộng theo $j$ cho
\[
 I(f)-A_{\mathcal T}(f)=\int_a^b\mathcal K_A(u)f^{(r_A)}(u)du,
 \qquad\int_a^b\mathcal K_A(u)du=c_A\sum_j\ell_j^{r_A+1}.
\]
Các $\ell_j>0$ giữ dấu của $K_A$ và
\[
 \int_a^b|\mathcal K_A(u)|du=|c_A|\sum_j\ell_j^{r_A+1}>0.
\]
Bổ đề \ref{bd:dnum-kernel} cho từng $\xi_A\in(a,b)$ và các công thức sai số. Cuối cùng, với $H=\max_j\ell_j$,
\[
 \sum_j\ell_j^3\leq H^2\sum_j\ell_j=H^2(b-a),\qquad
 \sum_j\ell_j^5\leq H^4\sum_j\ell_j=H^4(b-a).
\]
\end{chungminh}

\Needspace{7\baselineskip}
\subsection{Các công thức Newton--Cotes}
Tích phân đa thức nội suy trên một bộ mốc cách đều cho công thức Newton--Cotes. Có hai lựa chọn mốc chuẩn:
\[
 \tau_i=\frac{i}{k-1},\quad i=0,\ldots,k-1,\quad k\geq2
 \qquad\text{(đóng)},
\]
\[
 \tau_i=\frac{i+1}{k+1},\quad i=0,\ldots,k-1,\quad k\geq1
 \qquad\text{(mở)}.
\]
Các trọng số đều được xác định bằng mệnh đề \ref{md:quad-moments}. Đặt $W(t)=\prod_{i=0}^{k-1}(t-\tau_i)$.

\begin{bode}{}{quad-alternating}
Cho $L>0$ là số nguyên hoặc số nguyên cộng $1/2$. Giả sử $F$ liên tục trên $[0,L]$, có một dấu nghiêm ngặt $\sigma(-1)^j$ trên $(j,j+1)\cap(0,L)$, với $\sigma\in\{-1,1\}$. Giả sử $|F(u+1)|<|F(u)|$ với $0<u<L-1$ không là số nguyên. Khi đó $\sigma\int_0^LF(u)du>0$.
\end{bode}
\begin{chungminh}
Đặt $J=\sigma\int_0^LF$. Với mọi cặp khoảng đầy đủ nằm trong $[0,L]$,
\[
 A_j=\sigma\int_{2j}^{2j+2}F(u)du
 =\int_0^1\bigl(|F(2j+v)|-|F(2j+1+v)|\bigr)dv>0.
\]
Dấu $>$ dùng $|F(2j+v)|>|F(2j+1+v)|$ với $0<v<1$. Nếu $L$ nguyên,
\[
 J=\begin{cases}
 \sum_{j=0}^{q-1}A_j>0,&L=2q,\ q\geq1,\\
 \sum_{j=0}^{q-1}A_j+\int_0^1|F(2q+v)|dv>0,&L=2q+1.
 \end{cases}
\]
Nếu $L=2q+1/2$, thì
\[
 J=\sum_{j=0}^{q-1}A_j+\int_0^{1/2}|F(2q+v)|dv>0.
\]
Nếu $L=2q+3/2$, phần còn lại sau $q$ cặp là
\begin{align*}
 J-\sum_{j=0}^{q-1}A_j
 &=\int_0^{1/2}\bigl(|F(2q+v)|-|F(2q+1+v)|\bigr)dv\\
 &\quad+\int_{1/2}^1|F(2q+v)|dv>0.
\end{align*}
Trong hai trường hợp có $q=0$, tổng rỗng bằng $0$ và tích phân còn lại dương theo giả thiết dấu nghiêm ngặt.
\end{chungminh}

\begin{dinhly}{}{quad-nc-degree}
Với $k$ mốc Newton--Cotes đóng hoặc mở,
\[
 \text{bậc chính xác}=
 \begin{cases}k-1,&k\text{ chẵn},\\k,&k\text{ lẻ}.\end{cases}
\]
Các trọng số đối xứng: $\lambda_{k-1-i}=\lambda_i$.
\end{dinhly}
\begin{chungminh}

Vì $\tau_{k-1-i}=1-\tau_i$, đổi chỉ số $j=k-1-q$ trong tích cho
\[
 L_{k-1-i}(1-t)=\prod_{q\ne i}\frac{\tau_q-t}{\tau_q-\tau_i}
 =\prod_{q\ne i}\frac{t-\tau_q}{\tau_i-\tau_q}=L_i(t).
\]
Do đó
\[
 \lambda_{k-1-i}=\int_0^1L_{k-1-i}(u)du
 =\int_0^1L_{k-1-i}(1-t)dt=\lambda_i.
\]
Cũng bằng phép đổi chỉ số này,
\[
 W(1-t)=\prod_{i=0}^{k-1}(\tau_{k-1-i}-t)
 =(-1)^k\prod_{i=0}^{k-1}(t-\tau_i)=(-1)^kW(t).
\]

Mệnh đề \ref{md:quad-moments} cho tính chính xác tới bậc $k-1$. Đặt $J=\int_0^1W(t)dt$. Đổi biến $t=1-u$ cho $J=\int_0^1W(1-u)du=(-1)^kJ$; nếu $k$ lẻ, $2J=0$, nên $J=0$. Với $p\in\mathcal P_k$ và hệ số đầu $c=[t^k]p$, đa thức
\[
 H(t)=p(t)-I_{\boldsymbol\tau}p(t)-cW(t)
\]
có bậc không quá $k-1$ và $k$ nghiệm $\tau_i$, nên $H=0$ theo bổ đề \ref{bd:roots}. 
Từ $p=I_{\boldsymbol\tau}p+cW$, $Q(p)=Q(I_{\boldsymbol\tau}p)$ và $I_{\boldsymbol\tau}p\in\mathcal P_{k-1}$,
\[
 I(p)-Q(p)=I(I_{\boldsymbol\tau}p)-Q(I_{\boldsymbol\tau}p)+cJ=0.
\]

Cần xác định các tích phân không triệt tiêu ở bậc kế tiếp. Xét biến $z$ với các mốc nguyên. Trường hợp đóng, đặt $n=k-1$,
\[
 U(z)=\prod_{j=0}^{n}(z-j),\qquad0\leq z\leq n,
\]
\[
 \frac{U(z+1)}{U(z)}
 =\frac{(z+1)\prod_{j=0}^{n-1}(z-j)}{(z-n)\prod_{j=0}^{n-1}(z-j)}=\frac{z+1}{z-n}
\]
ở mọi $z$ không là mốc. Trường hợp mở, đặt
\[
 U(z)=\prod_{j=1}^{k}(z-j),\qquad0\leq z\leq k+1,
\]
\[
 \frac{U(z+1)}{U(z)}
 =\frac{z\prod_{j=1}^{k-1}(z-j)}{(z-k)\prod_{j=1}^{k-1}(z-j)}=\frac z{z-k}.
\]

Đặt $L=n/2$ đối với bộ đóng và $L=(k+1)/2$ đối với bộ mở. 
Đổi chỉ số trong từng tích:
\[
 U(n-z)=\prod_{j=0}^n(n-j-z)=\prod_{q=0}^n(q-z)=(-1)^kU(z)\quad\text{(đóng)},
\]
\[
 U(k+1-z)=\prod_{j=1}^k(k+1-j-z)=\prod_{q=1}^k(q-z)=(-1)^kU(z)\quad\text{(mở)}.
\]
Nếu $k$ chẵn, hai đẳng thức này là $U(2L-z)=U(z)$. Trên $(j,j+1)$,
\[
 \operatorname{sgn}U(z)=(-1)^{n-j}=-(-1)^j\ \text{(đóng)},\qquad
 \operatorname{sgn}U(z)=(-1)^{k-j}=(-1)^j\ \text{(mở)}.
\] Với $0<z<L-1$,
\[
 \left|\frac{U(z+1)}{U(z)}\right|
 =\frac{z+1}{n-z}<1\quad\text{(đóng)},\qquad
 \left|\frac{U(z+1)}{U(z)}\right|=\frac z{k-z}<1\quad\text{(mở)}.
\]

Bổ đề \ref{bd:quad-alternating} cho $\sigma\int_0^LU(z)dz>0$, với $\sigma=-1$ cho bộ đóng và $\sigma=1$ cho bộ mở. Đối xứng cho
\[
 \int_0^{2L}U(z)dz
 =\int_0^LU(z)dz+\int_0^LU(2L-z)dz
 =2\int_0^LU(z)dz.
\]
Sau các phép đổi biến dương dưới đây,
\[
 \int_0^1W(t)dt<0\quad\text{(đóng)},\qquad
 \int_0^1W(t)dt>0\quad\text{(mở)},\qquad k\text{ chẵn}.
\]
Các phép đổi biến được xác định bởi
\[
 W(t)=n^{-k}U(nt)\quad\text{(đóng)},\qquad
 W(t)=(k+1)^{-k}U((k+1)t)\quad\text{(mở)};
\]
\[
 \int_0^1W(t)dt=n^{-k-1}\int_0^nU(z)dz\quad\text{(đóng)},
\]
\[
 \int_0^1W(t)dt=(k+1)^{-k-1}\int_0^{k+1}U(z)dz\quad\text{(mở)}.
\] Vì $Q(W)=\sum_i\lambda_iW(\tau_i)=0$,
\[
 I(W)-Q(W)=\int_0^1W(t)dt\ne0,
\]
nên công thức không chính xác ở bậc $k$.


Với $k$ lẻ, đặt $F(z)=(z-L)U(z)$. Từ $U(2L-z)=-U(z)$,
\[
 F(2L-z)=(L-z)(-U(z))=(z-L)U(z)=F(z).
\] Trên nửa trái và $0<z<L-1$,
\[
 \left|\frac{F(z+1)}{F(z)}\right|
 =\frac{L-z-1}{L-z}\left|\frac{U(z+1)}{U(z)}\right|<1.
\]
Với $0<z<L-1$, ta có $0<L-z-1<L-z$ và, tương ứng,
\[
 0<z+1<n-z\quad\text{(đóng)},\qquad0<z<k-z\quad\text{(mở)}.
\]
Trên $(j,j+1)$ thuộc nửa trái,
\[
 \operatorname{sgn}U(z)=(-1)^{n-j}\ \text{(đóng)},\qquad
 \operatorname{sgn}U(z)=(-1)^{k-j}\ \text{(mở)},\qquad
 \operatorname{sgn}F(z)=-\operatorname{sgn}U(z).
\]

Vì $n$ chẵn trong trường hợp đóng và $k$ lẻ trong trường hợp mở,
\[
 \operatorname{sgn}F(z)=-(-1)^{n-j}=-(-1)^j\quad\text{(đóng)},
\]
\[
 \operatorname{sgn}F(z)=-(-1)^{k-j}=(-1)^j\quad\text{(mở)}.
\] Mặt khác, các phép đổi biến cho
\[
 \int_0^1(t-1/2)W(t)dt=n^{-k-2}\int_0^nF(z)dz\quad\text{(đóng)},
\]
\[
 \int_0^1(t-1/2)W(t)dt=(k+1)^{-k-2}\int_0^{k+1}F(z)dz\quad\text{(mở)}.
\]
Bổ đề \ref{bd:quad-alternating} cho $\sigma\int_0^LF(z)dz>0$ và
\[
 \int_0^{2L}F(z)dz=2\int_0^LF(z)dz,
\]
với $\sigma=-1$ cho bộ đóng, $\sigma=1$ cho bộ mở. Các hệ số đổi biến trên đều dương, nên
\[
 \int_0^1(t-1/2)W(t)dt<0\quad\text{(đóng)},\qquad
 \int_0^1(t-1/2)W(t)dt>0\quad\text{(mở)}.
\]

Với $P(t)=(t-1/2)W(t)$, $\deg P=k+1$, $Q(P)=\sum_i\lambda_iP(\tau_i)=0$ và
\[
 I(P)-Q(P)=\int_0^1(t-1/2)W(t)dt\ne0.
\]
Do đó công thức không chính xác ở bậc $k+1$. Các bậc chính xác lần lượt là $k-1$ khi $k$ chẵn, $k$ khi $k$ lẻ.
\end{chungminh}

\begin{hequa}{}{quad-nc-error}
Đặt $r=k$ nếu $k$ chẵn và $r=k+1$ nếu $k$ lẻ. Với $f\in C^r([a,b])$, sai số Newton--Cotes là \eqref{eq:quad-kernel}, và
\[
 |I(f)-Q(f)|\leq\ell^{r+1}\lVert f^{(r)}\rVert_\infty
 \int_0^1\left|\frac{(1-t)^r}{r!}
 -\sum_i\lambda_i\frac{(\tau_i-t)_+^{r-1}}{(r-1)!}\right|dt.
\]
Nếu kernel này có một dấu, dạng \eqref{eq:quad-kernel-mean} áp dụng. Trong mọi trường hợp, hằng số $\int|K_r|$ được xác định bằng tích phân các đa thức trên các đoạn tách bởi $\tau_i$ và các nghiệm thực của kernel trong từng đoạn.
\end{hequa}
\begin{chungminh}
Định lý \ref{dl:quad-nc-degree} cho $\sum_i\lambda_i\tau_i^j=1/(j+1)$ ($j<r$); Mệnh đề \ref{md:quad-kernel} cho biểu diễn và chặn. Trên mỗi khoảng không chứa mốc, $K_r$ là đa thức có hệ số đầu $(-1)^r/r!\ne0$, nên chỉ có hữu hạn nghiệm theo Bổ đề \ref{bd:roots}. Ghép các mốc, các nghiệm này và $0,1$ thành $0=z_0<\cdots<z_m=1$. Tính liên tục và định lý giá trị trung gian cho $K_r$ có một dấu trên từng $(z_j,z_{j+1})$. Bởi vậy
\[
 \int_0^1|K_r(t)|dt
 =\sum_{j=0}^{m-1}\int_{z_j}^{z_{j+1}}|K_r(t)|dt
 =\sum_{j=0}^{m-1}\left|\int_{z_j}^{z_{j+1}}K_r(t)dt\right|.
\]
Mỗi tích phân ở vế phải là tích phân của đa thức trên khoảng tương ứng, xác định bằng nguyên hàm từng đơn thức.
\end{chungminh}

\Needspace{9\baselineskip}
\begin{vidu}
Một số bộ trọng số chuẩn, tính từ hệ moment \eqref{eq:quad-moments}, là
\begin{center}
\begin{tabular}{c|c|c|c}
\toprule
Loại & $k$ & $\boldsymbol\lambda$ & Bậc chính xác\\
\midrule
Đóng & $2$ & $(1,1)^T/2$ & $1$\\
Đóng & $3$ & $(1,4,1)^T/6$ & $3$\\
Đóng & $4$ & $(1,3,3,1)^T/8$ & $3$\\
Mở & $1$ & $(1)^T$ & $1$\\
Mở & $2$ & $(1,1)^T/2$ & $1$\\
Mở & $3$ & $(2,-1,2)^T/3$ & $3$\\
\bottomrule
\end{tabular}
\end{center}
Với bộ mốc mở $\tau=(1/3,2/3)$, kernel cấp $2$ là
\[
 K_2(t)=\begin{cases}t^2/2,&0\leq t\leq1/3,\\
 (t^2-t+1/3)/2,&1/3\leq t\leq2/3,\\
 (1-t)^2/2,&2/3\leq t\leq1.
 \end{cases}
\]
Ở khoảng giữa, $t^2-t+1/3=(t-1/2)^2+1/12>0$. Do đó $K_2\geq0$ và $\int_0^1K_2=1/36$, nên $I-Q=\ell^3 f''(\xi)/36$.
\end{vidu}

\Needspace{7\baselineskip}
Trong công thức đóng, đa thức bậc $n$ dùng $k=n+1$ mốc. Với $x_i=a+ih$ và $b=a+nh$, các trọng số chuẩn $\lambda_i$ cho
\[
 Q_n(f)=nh\sum_{i=0}^n\lambda_i f(x_i).
\]
Vì thế số $n$ trong bậc đa thức khác với số $N$ các đoạn nhỏ khi ghép nhiều khối.

\begin{menhde}{}{quad-nc-456}
Các công thức đóng bậc $n=4,5,6$ có trọng số và sai số sau:
\begin{center}\small
\begin{tabular}{c|c|c|c}
\toprule
$n$ & $\boldsymbol\lambda^T$ & Bậc chính xác & $I-Q_n$\\
\midrule
$4$ & $(7,32,12,32,7)/90$ & $5$ & $-8h^7f^{(6)}(\xi)/945$\\
$5$ & $(19,75,50,50,75,19)/288$ & $5$ & $-275h^7f^{(6)}(\xi)/12096$\\
$6$ & $(41,216,27,272,27,216,41)/840$ & $7$ & $-9h^9f^{(8)}(\xi)/1400$\\
\bottomrule
\end{tabular}\end{center}
Ở mỗi hàng, $f$ có đạo hàm liên tục tới cấp được ghi và $\xi\in(a,b)$.
\end{menhde}
\begin{chungminh}
Đặt $D=90,288,840$ và $w_i=D\lambda_i$ theo từng hàng. 
Đặt $S_q=\sum_{i=0}^nw_i i^q$. Các tổng theo từng hàng là
\begin{align*}
 (S_0,\ldots,S_4)&=(90,180,480,1440,4608),&&n=4,\\
 (S_0,\ldots,S_5)&=(288,720,2400,9000,36000,150000),&&n=5,\\
 (S_0,\ldots,S_3)&=(840,2520,10080,45360),&&n=6,\\
 (S_4,S_5,S_6)&=(217728,1088640,5598720),&&n=6.
\end{align*}
Mỗi vector bằng $(Dn^q/(q+1))_{q=0}^n$, nên
\[
 \sum_{i=0}^n\lambda_i(i/n)^q=\frac{S_q}{Dn^q}=\frac1{q+1}\quad(0\leq q\leq n).
\]
Vì $\tau_i=i/n$, đây là hệ moment \eqref{eq:quad-moments}; tính duy nhất xác định các trọng số. Định lý \ref{dl:quad-nc-degree} cho bậc chính xác. Đặt $r=6$ khi $n=4,5$, và $r=8$ khi $n=6$. 
Ở moment cấp $r$, các tổng nguyên là $S_6=52800$ khi $n=4$, $S_6=643800$ khi $n=5$, $S_8=156800448$ khi $n=6$. Chia cho $Dn^r$ rồi trừ $1/(r+1)$ cho lần lượt $1/2688$, $11/52500$, $1/38880$. Do đó
\[
 -\int_0^1K_r(t)dt=\frac1{r!}\left(\sum_i\lambda_i(i/n)^r-\frac1{r+1}\right)
 =\begin{cases}
 1/1935360,&n=4,\\
 11/37800000,&n=5,\\
 1/1567641600,&n=6.
 \end{cases}
\]
Cần xác định dấu của kernel để dùng công thức giá trị trung gian. Từ các moment tới bậc $r-1$,
\[
 K_r(t)=\frac1{r!}\left[t^r-r\sum_{\tau_i<t}\lambda_i(t-\tau_i)^{r-1}\right],
 \qquad K_r(1-t)=K_r(t).
\]
Thật vậy, các moment cho
\[
 \sum_i\lambda_i(\tau_i-t)^{r-1}=\int_0^1(u-t)^{r-1}du
 =\frac{(1-t)^r-t^r}{r}.
\]
Tách tổng tại $\tau_i<t$, với $r-1$ lẻ, cho công thức thứ nhất. Đổi chỉ số $i\mapsto n-i$ trong công thức đó cho
\[
 K_r(1-t)=\frac{(1-t)^r-r\sum_{\tau_i>t}\lambda_i(\tau_i-t)^{r-1}}{r!}=K_r(t).
\] Trên $t=(j+u)/n$, $0\leq u\leq1$, đặt
\[
 H_{n,j}(u)=Dn^rr!K_r((j+u)/n)
 =D(j+u)^r-rn\sum_{i=0}^jw_i(j+u-i)^{r-1}.
\]
Với $B_m^r(u)=\binom rm u^m(1-u)^{r-m}$, khai triển nhị thức cho
\begin{align*}
 \sum_{m=a}^r\frac{\binom ma}{\binom ra}B_m^r(u)
 &=u^a\sum_{m=a}^r\binom{r-a}{m-a}u^{m-a}(1-u)^{r-m}\\
 &=u^a.
\end{align*}
Nếu $H_{n,j}(u)=\sum_{a=0}^r A_a u^a$, thì hệ số của $B_m^r$ trong $-\rho H_{n,j}$ là
\[
 b_m=-\rho\sum_{a=0}^m A_a\frac{\binom ma}{\binom ra}.
\]
Thay $A_r=D$ và, với $0\leq a<r$,
\[
 A_a=D\binom ra j^{r-a}-rn\binom{r-1}a\sum_{i=0}^jw_i(j-i)^{r-1-a}
\]
cho các vector sau. Trong những lũy thừa bậc $0$, giá trị được hiểu là $1$.
\begin{center}\scriptsize
\begin{tabular}{c|c|c|l}
\toprule
$n$ & $j$ & $\rho$ & $(b_0,\ldots,b_r)$ trong $-\rho H_{n,j}=\sum_m b_mB_m^r$\\
\midrule
$4$ & $0$ & $1/2$ & $(0,0,0,0,0,14,39)$\\
$4$ & $1$ & $1/2$ & $(39,64,100,144,176,192,192)$\\
$5$ & $0$ & $1$ & $(0,0,0,0,0,95,282)$\\
$5$ & $1$ & $1$ & $(282,469,748,1116,1472,1799,2058)$\\
$5$ & $2$ & $1$ & $(2058,2317,2508,2592,2508,2317,2058)$\\
$6$ & $0$ & $1/6$ & $(0,0,0,0,0,0,0,41,188)$\\
$6$ & $1$ & $1/6$ & $(188,335,588,1012,1696,2736,4160,5912,7872)$\\
$6$ & $2$ & $1/24$ & $(1968,2458,3000,3564,4104,4563,4887,5049,5049)$\\
\bottomrule
\end{tabular}\end{center}

Trên mỗi đoạn của bảng, $0\leq u\leq1$ và
\[
 K_r((j+u)/n)
 =-\frac1{\rho Dn^rr!}\sum_{m=0}^rb_m\binom rm u^m(1-u)^{r-m}\leq0,
\]
vì từng $b_m\geq0$. Với $t\geq1/2$, $1-t\leq1/2$ và $K_r(t)=K_r(1-t)\leq0$, nên dấu được xác định trên toàn $[0,1]$. Đặt $C=-\int_0^1K_r>0$; \eqref{eq:quad-kernel-mean} cho
\[
 I-Q_n=-C(nh)^{r+1}f^{(r)}(\xi),\qquad\xi\in(a,b).
\]
Các hệ số là
\[
 Cn^{r+1}=\begin{cases}
 4^7/1935360=8/945,&n=4,\\
 11\cdot5^7/37800000=275/12096,&n=5,\\
 6^9/1567641600=9/1400,&n=6.
 \end{cases}
\]
\end{chungminh}

\Needspace{9\baselineskip}
\begin{vidu}
Với công thức đóng bậc $4$, bảng ngang biểu diễn cách tính trọng số từ các tích và cách sử dụng chúng:
\begin{center}
\begin{tabular}{c|ccccc}
\toprule
$i$ & $0$ & $1$ & $2$ & $3$ & $4$\\
\midrule
$\tau_i$ & $0$ & $1/4$ & $1/2$ & $3/4$ & $1$\\
$90\lambda_i$ & $7$ & $32$ & $12$ & $32$ & $7$\\
$y_i$ & $y_0$ & $y_1$ & $y_2$ & $y_3$ & $y_4$\\
$90\lambda_i y_i$ & $7y_0$ & $32y_1$ & $12y_2$ & $32y_3$ & $7y_4$\\
\bottomrule
\end{tabular}\end{center}
Vì $b-a=4h$, giá trị công thức là $2h(7y_0+32y_1+12y_2+32y_3+7y_4)/45$.
\end{vidu}

\Needspace{16\baselineskip}
\begin{menhde}{}{quad-composite-general}
Cho một công thức chuẩn có kernel $K_r$, và phân hoạch $a=t_0<\cdots<t_N=b$, $\ell_j=t_{j+1}-t_j$. Đặt
\[
 Q_{\mathcal T}(f)=\sum_{j=0}^{N-1}\ell_j\sum_i\lambda_i f(t_j+\ell_j\tau_i).
\]
Khi $f\in C^r([a,b])$,
\[
 |I(f)-Q_{\mathcal T}(f)|\leq C_r\lVert f^{(r)}\rVert_\infty\sum_j\ell_j^{r+1}
 \leq C_r(b-a)H^r\lVert f^{(r)}\rVert_\infty,
 \qquad H=\max_j\ell_j.
\]
Nếu $K_r$ có một dấu và $C_r>0$, $I-Q_{\mathcal T}=\bigl(\int_0^1K_r\bigr)\bigl(\sum_j\ell_j^{r+1}\bigr)f^{(r)}(\xi)$ với $\xi\in(a,b)$.
\end{menhde}
\begin{chungminh}
Trên $(t_j,t_{j+1})$, đặt $\mathcal K(u)=\ell_j^rK_r((u-t_j)/\ell_j)$. Đổi biến $u=t_j+\ell_jt$ trong \eqref{eq:quad-kernel} cho
\[
 \ell_j^{r+1}\int_0^1K_r(t)f^{(r)}(t_j+\ell_jt)dt
 =\int_{t_j}^{t_{j+1}}\mathcal K(u)f^{(r)}(u)du.
\]
Cộng theo $j$:
\[
 I(f)-Q_{\mathcal T}(f)
 =\sum_{j=0}^{N-1}\int_{t_j}^{t_{j+1}}\mathcal K(u)f^{(r)}(u)du
 =\int_a^b\mathcal K(u)f^{(r)}(u)du.
\]
Cũng với phép đổi biến trên,
\[
 \int_a^b|\mathcal K(u)|du
 =\sum_j\ell_j^{r+1}\int_0^1|K_r(t)|dt=C_r\sum_j\ell_j^{r+1}.
\]
Vì $0<\ell_j\leq H$,
\[
 \sum_j\ell_j^{r+1}\leq H^r\sum_j\ell_j=(b-a)H^r,
\]
\begin{align*}
 |I(f)-Q_{\mathcal T}(f)|
 &\leq\int_a^b|\mathcal K(u)||f^{(r)}(u)|du\\
 &\leq C_r\|f^{(r)}\|_\infty\sum_j\ell_j^{r+1}
 \leq C_r(b-a)H^r\|f^{(r)}\|_\infty.
\end{align*}
Nếu $K_r$ có một dấu và $C_r>0$, chọn $\sigma\in\{-1,1\}$ sao cho $\sigma K_r\geq0$. Khi đó $\sigma\mathcal K\geq0$ và
\[
 \int_a^b|\mathcal K|=C_r\sum_j\ell_j^{r+1}>0,\qquad
 \int_a^b\mathcal K=\left(\int_0^1K_r\right)\sum_j\ell_j^{r+1}.
\]
Bổ đề \ref{bd:dnum-kernel} với $g=f^{(r)}$ cho $\xi\in(a,b)$ và
\[
 I(f)-Q_{\mathcal T}(f)
 =f^{(r)}(\xi)\int_a^b\mathcal K
 =\left(\int_0^1K_r\right)\left(\sum_j\ell_j^{r+1}\right)f^{(r)}(\xi).
\]
\end{chungminh}

\begin{menhde}{}{quad-data}
Với công thức đơn, dữ liệu $\widetilde y_i=f(a+\ell\tau_i)+\eta_i$, $|\eta_i|\leq\varepsilon$, thỏa
\[
 \max_{|\eta_i|\leq\varepsilon}|\widetilde Q-Q|=\ell\varepsilon\sum_i|\lambda_i|.
\]
Với công thức ghép, một chặn là $(b-a)\varepsilon\sum_i|\lambda_i|$. Nếu toàn bộ trọng số không âm và tổng bằng $1$, chặn này bằng $(b-a)\varepsilon$ và đạt được bằng nhiễu cùng dấu tại mọi mốc.

Với dữ liệu và trọng số được biểu diễn chính xác, tính $k$ tích, cộng tuần tự rồi nhân với $\ell$ theo mô hình \eqref{eq:rounding-model}, ta có
\[
 |\widehat Q-\widetilde Q|\leq\gamma_{k+1}\ell\sum_i|\lambda_i||\widetilde y_i|,
 \qquad (k+1)u<1.
\]
\end{menhde}
\begin{chungminh}
Với công thức đơn,
\[
 \widetilde Q-Q=\ell\sum_i\lambda_i\eta_i,\qquad
 |\widetilde Q-Q|\leq\ell\sum_i|\lambda_i||\eta_i|
 \leq\ell\varepsilon\sum_i|\lambda_i|.
\]
Lấy $\eta_i=\varepsilon\operatorname{sgn}(\lambda_i)$, với $\operatorname{sgn}(0)=0$, thì
\[
 \widetilde Q-Q=\ell\varepsilon\sum_i|\lambda_i|.
\]
Với công thức ghép trên các khối có độ dài $\ell_j$,
\[
 |\widetilde Q_{\mathcal T}-Q_{\mathcal T}|
 \leq\sum_j\ell_j\sum_i|\lambda_i||\eta_{j,i}|
 \leq\varepsilon\left(\sum_j\ell_j\right)\sum_i|\lambda_i|
 =(b-a)\varepsilon\sum_i|\lambda_i|.
\]
Nếu $\lambda_i\geq0$, $\sum_i\lambda_i=1$, chọn nhiễu $\varepsilon$ tại mọi mốc, kể cả mốc chung, thì
\[
 \widetilde Q_{\mathcal T}-Q_{\mathcal T}
 =\varepsilon\sum_j\ell_j\sum_i\lambda_i=(b-a)\varepsilon.
\]

Đánh số $k$ số hạng bởi $i=0,\ldots,k-1$. Theo mô hình \eqref{eq:rounding-model}, đặt
\[
 t_i=\operatorname{fl}(\lambda_i\widetilde y_i)
 =\lambda_i\widetilde y_i(1+\delta_i),\qquad |\delta_i|\leq u,
\]
\[
 s_0=t_0,\qquad s_j=(s_{j-1}+t_j)(1+\epsilon_j)
 \quad(1\leq j<k),\qquad |\epsilon_j|\leq u.
\]
Khai triển quy nạp truy hồi tổng cho
\[
 s_{k-1}=\sum_{i=0}^{k-1}\lambda_i\widetilde y_i(1+\delta_i)
 \prod_{j=\max\{1,i\}}^{k-1}(1+\epsilon_j).
\]
Cụ thể, bước $s_j=(s_{j-1}+t_j)(1+\epsilon_j)$ nhân mọi số hạng cũ với $(1+\epsilon_j)$ và thêm $t_j(1+\epsilon_j)$; công thức đúng tại $j=0$. Phép nhân cuối cho
\[
 \widehat Q=\ell s_{k-1}(1+\zeta)
 =\ell\sum_i\lambda_i\widetilde y_i(1+\theta_i),\qquad |\zeta|\leq u,
\]
\[
 1+\theta_i=(1+\delta_i)(1+\zeta)
 \prod_{j=\max\{1,i\}}^{k-1}(1+\epsilon_j).
\]
Mỗi tích có không quá $2+(k-1)=k+1$ thừa số. Bổ đề \ref{bd:gamma} cho $|\theta_i|\leq\gamma_{k+1}$, nên
\[
 |\widehat Q-\widetilde Q|
 =\left|\ell\sum_i\lambda_i\widetilde y_i\theta_i\right|
 \leq\gamma_{k+1}\ell\sum_i|\lambda_i||\widetilde y_i|.
\]
\end{chungminh}

\Needspace{9\baselineskip}
\begin{vidu}
Các trọng số của điểm giữa, hình thang và Simpson đều không âm và có tổng $1$. Với công thức mở ba điểm, $\sum_i|\lambda_i|=5/3$, nên chặn nhiễu dữ liệu bằng $5\ell\varepsilon/3$. Sự khác nhau này được tính từ chuẩn $1$ của vector trọng số, không từ bậc chính xác của công thức.
\end{vidu}

\Needspace{7\baselineskip}
Khi được phép lấy thêm giá trị hàm, số đoạn là một biến của bài toán. Các chặn ghép chuyển việc chọn $h$ thành một bất đẳng thức cho số nguyên $N$.

\begin{menhde}{}{quad-panel-count}
Cho chặn $|I-Q_N|\leq A/N^r+\nu$, $A\geq0$, $r>0$, trong đó $\nu\geq0$ chứa chặn dữ liệu và làm tròn. Nếu $\varepsilon>\nu$ và số đoạn phải là bội của số nguyên $g\geq1$, số nhỏ nhất được chặn này chứng nhận là
\begin{equation}\label{eq:quad-panel-count}
 N=g\max\left\{1,\left\lceil\frac1g\left(\frac A{\varepsilon-\nu}\right)^{1/r}\right\rceil\right\}.
\end{equation}
Với $L=b-a$, các công thức đều có các giá trị sau:
\begin{center}
\begin{tabular}{c|c|c|c}
\toprule
Công thức & $r$ & $A$ & $g$\\
\midrule
Điểm giữa & $2$ & $L^3M_2/24$ & $1$\\
Hình thang & $2$ & $L^3M_2/12$ & $1$\\
Simpson & $4$ & $L^5M_4/180$ & $2$\\
Newton--Cotes đóng bậc $n$ & $r_n$ & $C_{r_n}L^{r_n+1}n^{r_n}M_{r_n}$ & $n$\\
\bottomrule
\end{tabular}\end{center}
Ở đây $M_j\geq\|f^{(j)}\|_\infty$, $r_n=n+1$ với $n$ lẻ và $r_n=n+2$ với $n$ chẵn; $C_{r_n}=\int_0^1|K_{r_n}|$. Trong hàng Newton--Cotes, $N$ là tổng số khoảng nhỏ, mỗi khối có $n$ khoảng nhỏ.
\end{menhde}
\begin{chungminh}
Đặt $\delta=\varepsilon-\nu>0$ và $N=gq$, $q\in\mathbb N$, $q\geq1$. Khi đó
\begin{align*}
 A/N^r+\nu\leq\varepsilon
 &\ \Longleftrightarrow\ A/N^r\leq\delta
 \ \Longleftrightarrow\ N\geq(A/\delta)^{1/r}\\
 &\ \Longleftrightarrow\ q\geq g^{-1}(A/\delta)^{1/r}\\
 &\ \Longleftrightarrow\ q\geq
 \max\{1,\lceil g^{-1}(A/\delta)^{1/r}\rceil\}.
\end{align*}
Bởi vậy số nguyên dương nhỏ nhất $q$ là vế phải cuối cùng, cho \eqref{eq:quad-panel-count}, kể cả $A=0$. Thế $h=L/N$ vào các chặn đã chứng minh:
\[
 \frac{LM_2h^2}{24}=\frac{L^3M_2}{24N^2},\qquad
 \frac{LM_2h^2}{12}=\frac{L^3M_2}{12N^2},\qquad
 \frac{LM_4h^4}{180}=\frac{L^5M_4}{180N^4}.
\]
Simpson đòi hỏi $N$ chẵn, tức $g=2$; hai công thức còn lại cho $g=1$. Với Newton--Cotes, $N$ là bội của $n$, nên có $N/n$ khối, mỗi khối dài $nL/N$. Mệnh đề \ref{md:quad-composite-general} cho
\begin{align*}
 |I-Q_N|
 &\leq C_{r_n}M_{r_n}\frac Nn(nL/N)^{r_n+1}\\
 &=\frac{C_{r_n}M_{r_n}L^{r_n+1}n^{r_n}}{N^{r_n}}.
\end{align*}
Cộng $\nu$ vào mỗi chặn cho các giá trị $A,r,g$ trong bảng.
\end{chungminh}

\begin{menhde}{}{quad-rounding-budget}
Giả sử công thức ghép có trọng số không âm, tổng trọng số bằng $L=b-a$, và dùng không quá $N+1$ giá trị. Cho $N\leq N_{\max}$, $(N_{\max}+2)u<1$, $|f(x)|\leq F$, nhiễu mỗi dữ liệu không quá $\eta$. Dưới các giả thiết biểu diễn và tính tổng của Mệnh đề \ref{md:quad-data}, một chặn chung cho dữ liệu và làm tròn là
\[
 \nu=L\left[\eta+\gamma_{N_{\max}+2}(F+\eta)\right].
\]
Công thức \eqref{eq:quad-panel-count} với chặn này được sử dụng nếu số $N$ thu được không vượt $N_{\max}$.
\end{menhde}
\begin{chungminh}
Gọi các trọng số đầy đủ là $a_i\geq0$, $\sum_i a_i=L$, và $K\leq N+1$ là số giá trị. Với $\widetilde y_i=f(x_i)+\eta_i$, ta có
\[
 |\widetilde y_i|\leq F+\eta,\qquad
 |\widetilde Q-Q|=\left|\sum_i a_i\eta_i\right|
 \leq\sum_i a_i|\eta_i|\leq L\eta.
\]
Khai triển phép nhân và cộng tuần tự trong Mệnh đề \ref{md:quad-data}, với tối đa $K+1$ thừa số nhiễu cho mỗi số hạng, cho
\[
 |\widehat Q-\widetilde Q|
 \leq\gamma_{K+1}\sum_i a_i|\widetilde y_i|
 \leq\gamma_{K+1}L(F+\eta).
\]
Với $0\leq p\leq q$ và $qu<1$,
\[
 \gamma_q-\gamma_p
 =\frac{qu}{1-qu}-\frac{pu}{1-pu}
 =\frac{(q-p)u}{(1-qu)(1-pu)}\geq0.
\]
Vì $K+1\leq N+2\leq N_{\max}+2$, suy ra
\begin{align*}
 |\widehat Q-Q|
 &\leq|\widehat Q-\widetilde Q|+|\widetilde Q-Q|\\
 &\leq L\bigl[\gamma_{N_{\max}+2}(F+\eta)+\eta\bigr]=\nu.
\end{align*}
Nếu $N$ chọn bởi \eqref{eq:quad-panel-count} thỏa $N\leq N_{\max}$, chặn $A/N^r+\nu\leq\varepsilon$ giữ nguyên giả thiết $(N_{\max}+2)u<1$.
\end{chungminh}

\Needspace{9\baselineskip}
\begin{vidu}
Với $f(x)=x^4$ trên $[0,1]$, $M_2=12$, $M_4=24$. Trong số học chính xác và dữ liệu chính xác, yêu cầu $\varepsilon=10^{-4}$ cho
\begin{center}
\begin{tabular}{c|c|c|c|c}
\toprule
Công thức & $A$ & Ngưỡng thực của $N$ & $N$ chọn & Chặn đạt được\\
\midrule
Điểm giữa & $1/2$ & $\sqrt{5000}$ & $71$ & $1/10082<10^{-4}$\\
Hình thang & $1$ & $100$ & $100$ & $10^{-4}$\\
Simpson & $2/15$ & $(4000/3)^{1/4}$ & $8$ & $1/30720<10^{-4}$\\
\bottomrule
\end{tabular}\end{center}
Với Simpson, $6^4=1296<4000/3<7^4$; yêu cầu $N$ chẵn đưa $N$ từ $7$ lên $8$. Chặn đạo hàm được xác định trên toàn đoạn, nên phép chọn này có chứng nhận.
\end{vidu}

\Needspace{7\baselineskip}
Một chặn $M_r$ dùng để chọn lưới cần đúng trên toàn đoạn. Với đa thức, các hệ thức Horner cho một chặn trực tiếp; với phân thức, phương trình tích cho một truy hồi chặn.
\begin{menhde}{}{quad-derivative-bounds}
Cho $P(x)=\sum_{j=0}^d a_j(x-c)^j$ và $|x-c|\leq R$ trên $[a,b]$. Khi $0\leq r\leq d$,
\[
 \|P^{(r)}\|_\infty\leq\sum_{j=r}^d\frac{j!}{(j-r)!}|a_j|R^{j-r}.
\]
Nếu $f=P/Q$, $P,Q$ đủ trơn và $|Q(x)|\geq q_*>0$ trên $[a,b]$, đặt
\[
 B_0=\|P\|_\infty/q_*,\qquad
 B_r=\frac{\|P^{(r)}\|_\infty+\sum_{j=1}^r\binom rj\|Q^{(j)}\|_\infty B_{r-j}}{q_*}.
\]
Khi đó $\|f^{(r)}\|_\infty\leq B_r$.
\end{menhde}
\begin{chungminh}
Đạo hàm từng đơn thức cho
\[
 P^{(r)}(x)=\sum_{j=r}^d\frac{j!}{(j-r)!}a_j(x-c)^{j-r}.
\]
Với $|x-c|\leq R$,
\[
 |P^{(r)}(x)|
 \leq\sum_{j=r}^d\frac{j!}{(j-r)!}|a_j||x-c|^{j-r}
 \leq\sum_{j=r}^d\frac{j!}{(j-r)!}|a_j|R^{j-r}.
\]
Lấy supremum theo $x\in[a,b]$ cho chặn thứ nhất.

Từ $Qf=P$, cấp $0$ thỏa $|f(x)|\leq\|P\|_\infty/q_*=B_0$. Giả sử $\|f^{(j)}\|_\infty\leq B_j$ với $0\leq j<r$. Quy tắc Leibniz cho
\[
 P^{(r)}=\sum_{j=0}^r\binom rjQ^{(j)}f^{(r-j)},\qquad
 Qf^{(r)}=P^{(r)}-\sum_{j=1}^r\binom rjQ^{(j)}f^{(r-j)}.
\]
Do $|Q(x)|\geq q_*>0$,
\begin{align*}
 |f^{(r)}(x)|
 &\leq\frac{|P^{(r)}(x)|+\sum_{j=1}^r\binom rj|Q^{(j)}(x)||f^{(r-j)}(x)|}{q_*}\\
 &\leq\frac{\|P^{(r)}\|_\infty+\sum_{j=1}^r\binom rj\|Q^{(j)}\|_\infty B_{r-j}}{q_*}=B_r.
\end{align*}
Lấy supremum và quy nạp theo $r$ cho $\|f^{(r)}\|_\infty\leq B_r$.
\end{chungminh}

\Needspace{7\baselineskip}
Khi bảng được cho sẵn, bước mốc là dữ liệu đầu vào. Công thức điểm giữa cần giá trị tại trung điểm từng ô; Simpson cần một số chẵn các khoảng nhỏ đều nhau. Nếu những giá trị cần thiết chưa có, phải chọn một phân hoạch hợp lệ hoặc dựng thêm một mô hình nội suy có sai số riêng.

\begin{menhde}{}{quad-table-only}
Từ một số hữu hạn giá trị hàm tại các mốc phân biệt, không thể suy ra một chặn hữu hạn cho sai số tích phân trên đoạn chứa các mốc, hoặc cho sai số một đạo hàm cấp $s\geq1$ tại một mốc, nếu không có giả thiết bổ sung về hàm.
\end{menhde}
\begin{chungminh}
Cho $a<b$, các mốc $z_0<\cdots<z_n$ trong $[a,b]$, $n\geq0$, $p$ là đa thức nội suy dữ liệu và $W(x)=\prod_{i=0}^n(x-z_i)$. Với $c\in\R$, đặt $f_c=p+cW^2$; khi đó $f_c(z_i)=p(z_i)$. Chọn $t_0\in(a,b)\setminus\{z_0,\ldots,z_n\}$. Tính liên tục của $W^2$ và $W(t_0)^2>0$ cho một đoạn $J\subset(a,b)$, độ dài dương, sao cho
\[
 W(t)^2\geq W(t_0)^2/2\quad(t\in J).
\]
Vì $W^2\geq0$,
\[
 A=\int_a^bW(t)^2dt\geq\int_JW(t)^2dt
 \geq\frac{W(t_0)^2}{2}|J|>0.
\]
Với mọi $T\in\R$, chọn $c=(T-\int_a^bp)/A$, thì
\[
 \int_a^bf_c=\int_a^bp+cA=T.
\]

Với đạo hàm cấp $s\geq1$ tại $z_j$, đặt
\[
 V_j(x)=\prod_{i\ne j}(x-z_i),\qquad
 g_c(x)=p(x)+c(x-z_j)^sV_j(x).
\]
Với mọi $i$, $g_c(z_i)=p(z_i)$. Quy tắc Leibniz tại $z_j$ cho
\begin{align*}
 g_c^{(s)}(z_j)-p^{(s)}(z_j)
 &=c\sum_{\ell=0}^s\binom s\ell
 \left.\frac{d^\ell}{dx^\ell}(x-z_j)^s\right|_{x=z_j}V_j^{(s-\ell)}(z_j)\\
 &=cs!V_j(z_j),\qquad V_j(z_j)=\prod_{i\ne j}(z_j-z_i)\ne0.
\end{align*}
Với mọi $T\in\R$, chọn $c=(T-p^{(s)}(z_j))/(s!V_j(z_j))$ thì $g_c^{(s)}(z_j)=T$. Một phép tính chỉ phụ thuộc dữ liệu có cùng kết quả $B$ trên mọi $f_c$ hoặc mọi $g_c$, trong khi
\[
 \sup_{T\in\R}|T-B|=+\infty.
\]
Do đó không có chặn hữu hạn cho hai sai số. Mọi $f_c,g_c$ đều là đa thức, nên chúng có đạo hàm mọi cấp.
\end{chungminh}

\begin{menhde}{}{quad-table-panels}
Cho $x_i=a+ih$, $0\leq i\leq N$, $N\geq1$, và bậc khối $n\geq1$. Viết $N=qn+v$, $q\geq0$, $0\leq v<n$. Với các trọng số Newton--Cotes đóng bậc $d$, ký hiệu $\lambda_i^{(d)}$, phép ghép trên toàn bảng là
\begin{equation}\label{eq:quad-table-panels}
 Q=h n\sum_{j=0}^{q-1}\sum_{i=0}^n\lambda_i^{(n)}y_{nj+i}
 +\begin{cases}0,&v=0,\\hv\sum_{i=0}^v\lambda_i^{(v)}y_{qn+i},&v>0.\end{cases}
\end{equation}
Đặt $r_d=d+1$ với $d$ lẻ và $r_d=d+2$ với $d$ chẵn. Nếu các chặn đạo hàm cần thiết được cho, thì
\[
 |I-Q|\leq q C_{r_n}(nh)^{r_n+1}M_{r_n}
 +\begin{cases}0,&v=0,\\C_{r_v}(vh)^{r_v+1}M_{r_v},&v>0.\end{cases}
\]
\end{menhde}
\begin{chungminh}
Từ $N=qn+v$, $0\leq v<n$, $x_{qn}=a+qnh$, $x_N=x_{qn}+vh$ và
\[
 \int_a^{x_N}f
 =\sum_{j=0}^{q-1}\int_{x_{jn}}^{x_{(j+1)n}}f
 +\begin{cases}0,&v=0,\\\displaystyle\int_{x_{qn}}^{x_N}f,&v>0.\end{cases}
\]
Trên khối $j$, độ dài $nh$ và các mốc chuẩn $i/n$ cho
\[
 x_{jn}+nh(i/n)=x_{jn+i},\qquad
 Q_j=nh\sum_{i=0}^n\lambda_i^{(n)}y_{jn+i}.
\]
Nếu $v>0$, khối cuối có độ dài $vh$ và
\[
 x_{qn}+vh(i/v)=x_{qn+i},\qquad
 Q_*=vh\sum_{i=0}^v\lambda_i^{(v)}y_{qn+i}.
\]
Cộng các $Q_j$ và, nếu có, $Q_*$ cho \eqref{eq:quad-table-panels}. Đặt
\[
 E_j=\int_{x_{jn}}^{x_{(j+1)n}}f-Q_j,
 \quad E_*=\begin{cases}0,&v=0,\\\displaystyle\int_{x_{qn}}^{x_N}f-Q_*,&v>0.\end{cases}
\]
Hệ quả \ref{hq:quad-nc-error} cho $|E_j|\leq C_{r_n}(nh)^{r_n+1}M_{r_n}$ và, với $v>0$, $|E_*|\leq C_{r_v}(vh)^{r_v+1}M_{r_v}$. Bởi vậy
\begin{align*}
 |I-Q|&=\left|\sum_{j=0}^{q-1}E_j+E_*\right|
 \leq\sum_{j=0}^{q-1}|E_j|+|E_*|\\
 &\leq qC_{r_n}(nh)^{r_n+1}M_{r_n}
 +\begin{cases}0,&v=0,\\C_{r_v}(vh)^{r_v+1}M_{r_v},&v>0.\end{cases}
\end{align*}
Khi $q=0$, tổng rỗng bằng $0$ và chỉ có khối cuối.
\end{chungminh}

\Needspace{9\baselineskip}
\begin{vidu}
Một bảng đều có $N=10$ khoảng nhỏ, chọn $n=4$. Khi đó $q=2$, $v=2$:
\begin{center}
\begin{tabular}{c|c|c|c}
\toprule
Khối & Chỉ số sử dụng & Công thức & Độ dài\\
\midrule
$0$ & $0,1,2,3,4$ & Newton--Cotes bậc $4$ & $4h$\\
$1$ & $4,5,6,7,8$ & Newton--Cotes bậc $4$ & $4h$\\
Cuối & $8,9,10$ & Simpson & $2h$\\
\bottomrule
\end{tabular}\end{center}
Công thức dùng hết bảng. Chặn gồm $2C_6(4h)^7M_6+C_4(2h)^5M_4$; bậc sai số của khối cuối cần được giữ trong phép cộng.
\end{vidu}

\Needspace{7\baselineskip}
Cho hai công thức cùng loại $Q(h)$ và $Q(h/2)$. Nếu đã có khai triển
\[
 I-Q(h)=ch^r+R(h),\qquad |R(h)|\leq Bh^{r+q},\quad q>0,
\]
thì
\begin{equation}\label{eq:quad-refinement-indicator}
 \widehat E=\frac{Q(h/2)-Q(h)}{2^r-1},\qquad
 \left|(I-Q(h/2))-\widehat E\right|
 \leq\frac{1+2^{-q}}{2^r-1}Bh^{r+q}.
\end{equation}
Thật vậy,
\[
 \widehat E=\frac{ch^r(1-2^{-r})+R(h)-R(h/2)}{2^r-1}
 =ch^r2^{-r}+\frac{R(h)-R(h/2)}{2^r-1},
\]
\[
 (I-Q(h/2))-\widehat E=\frac{2^rR(h/2)-R(h)}{2^r-1}.
\]
Do đó
\begin{align*}
 \left|(I-Q(h/2))-\widehat E\right|
 &\leq\frac{2^r|R(h/2)|+|R(h)|}{2^r-1}\\
 &\leq\frac{2^rB(h/2)^{r+q}+Bh^{r+q}}{2^r-1}
 =\frac{1+2^{-q}}{2^r-1}Bh^{r+q}.
\end{align*}
Nếu chưa có chặn $R$, $\widehat E$ chỉ là một chỉ số so sánh hai lưới. Mệnh đề \ref{md:quad-table-only} cho thấy chênh lệch giữa hai phép tính từ cùng tập hữu hạn dữ liệu không tự tạo ra một chặn sai số cho mọi hàm trơn khớp dữ liệu.

\Needspace{7\baselineskip}
\subsection{Tích phân kép trên hình chữ nhật}
Trên $D=[a,b]\times[c,d]$, đặt $L_x=b-a$, $L_y=d-c$. Áp dụng một công thức theo từng biến cho tổng hai chiều
\begin{equation}\label{eq:quad-tensor}
 Q_{xy}(f)=\sum_{i=0}^{N_x}\sum_{j=0}^{N_y}\alpha_i\beta_jf(x_i,y_j),
 \qquad x_i=a+ih_x,\ y_j=c+jh_y.
\end{equation}
Với hình thang, $h_x=L_x/N_x$, $h_y=L_y/N_y$, $\alpha_i=h_x a_i$, $\beta_j=h_yb_j$, trong đó trọng số ở mỗi đầu mút là $1/2$ và trọng số bên trong là $1$. Với Simpson, $N_x,N_y$ chẵn, $\alpha_i=h_xa_i/3$, $\beta_j=h_yb_j/3$, trong đó hai đầu mút có trọng số $1$, chỉ số lẻ có trọng số $4$, chỉ số chẵn bên trong có trọng số $2$.

\begin{dinhly}{}{quad-tensor-error}
Cho $f$ liên tục trên $D$. Nếu các đạo hàm riêng cấp $2$ theo từng biến liên tục và bị chặn bởi $M_x$, $M_y$, công thức hình thang hai chiều thỏa
\[
 \left|\iint_Df-Q_{xy}(f)\right|
 \leq\frac{L_xL_y}{12}(h_x^2M_x+h_y^2M_y).
\]
Nếu các đạo hàm riêng tới cấp $4$ theo từng biến liên tục và các đạo hàm cấp $4$ bị chặn bởi $M_x$, $M_y$, công thức Simpson hai chiều thỏa
\[
 \left|\iint_Df-Q_{xy}(f)\right|
 \leq\frac{L_xL_y}{180}(h_x^4M_x+h_y^4M_y).
\]
\end{dinhly}
\begin{chungminh}
Ký hiệu $I_x,I_y$ là phép tích phân theo từng biến, $Q_x,Q_y$ là công thức ghép tương ứng. Đối với hình thang,
\[
 \sum_{i=0}^{N_x}\alpha_i=h_x(1/2+(N_x-1)+1/2)=h_xN_x=L_x.
\]
Đối với Simpson,
\[
 \sum_{i=0}^{N_x}\alpha_i
 =\frac{h_x}{3}\bigl[2+4(N_x/2)+2(N_x/2-1)\bigr]=h_xN_x=L_x.
\]
Thay $x$ bằng $y$ cho $\sum_j\beta_j=L_y$. Mọi $\alpha_i,\beta_j\geq0$ theo các trọng số đã cho. Các tổng hữu hạn và tích phân thỏa
\[
 Q_xI_yf=\sum_i\alpha_i\int_c^df(x_i,y)dy
 =\int_c^d\sum_i\alpha_if(x_i,y)dy=I_yQ_xf.
\]
Khai triển hiệu:
\begin{align*}
 I_xI_yf-Q_xQ_yf
 &=I_xI_yf-Q_xI_yf+Q_xI_yf-Q_xQ_yf\\
 &=(I_x-Q_x)I_yf+Q_x(I_y-Q_y)f.
\end{align*}
Đặt $(r,C)=(2,1/12)$ cho hình thang hoặc $(r,C)=(4,1/180)$ cho Simpson. Với $g(x)=\int_c^df(x,y)dy$, quy tắc đạo hàm tích phân cho
\[
 |g^{(r)}(x)|=\left|\int_c^d\partial_x^rf(x,y)dy\right|
 \leq\int_c^d|\partial_x^rf(x,y)|dy\leq L_yM_x.
\]
Chặn một chiều cho
\[
 |(I_x-Q_x)I_yf|\leq CL_xh_x^r\|g^{(r)}\|_\infty
 \leq CL_xL_yh_x^rM_x.
\]
Với mọi $i$, $|(I_y-Q_y)f(x_i,\cdot)|\leq CL_yh_y^rM_y$, nên
\begin{align*}
 |Q_x(I_y-Q_y)f|
 &\leq\sum_i\alpha_i|(I_y-Q_y)f(x_i,\cdot)|\\
 &\leq\sum_i\alpha_iCL_yh_y^rM_y=CL_xL_yh_y^rM_y.
\end{align*}
Do đó
\[
 |I_xI_yf-Q_xQ_yf|
 \leq CL_xL_y(h_x^rM_x+h_y^rM_y),
\]
cho hai chặn trong phát biểu.
\end{chungminh}

\begin{menhde}{}{quad-tensor-count}
Viết chặn hai chiều là $A_x/N_x^r+A_y/N_y^r+\nu$, với $\nu<\varepsilon$ và $A_x,A_y>0$. Trong hai công thức trên, $A_x=CL_x^{r+1}L_yM_x$, $A_y=CL_xL_y^{r+1}M_y$, với $(r,C)=(2,1/12)$ hoặc $(4,1/180)$. Đặt $\delta=\varepsilon-\nu$. Một lưới được chứng nhận bởi
\[
 N_x=g\max\left\{1,\left\lceil\frac1g(2A_x/\delta)^{1/r}\right\rceil\right\},\qquad
 N_y=g\max\left\{1,\left\lceil\frac1g(2A_y/\delta)^{1/r}\right\rceil\right\},
\]
với $g=1$ cho hình thang, $g=2$ cho Simpson. Ở mức số thực dương, chia đều chặn sai số cho hai chiều làm nhỏ nhất tích $N_xN_y$.
\end{menhde}
\begin{chungminh}
Các số nguyên được chọn thỏa
\[
 N_x\geq(2A_x/\delta)^{1/r},\quad
 N_y\geq(2A_y/\delta)^{1/r},\qquad N_x,N_y\in g\mathbb N.
\]
Vì $A_x,A_y>0$,
\[
 \frac{A_x}{N_x^r}\leq\frac\delta2,\qquad
 \frac{A_y}{N_y^r}\leq\frac\delta2,
\]
và $A_x/N_x^r+A_y/N_y^r+\nu\leq\delta+\nu=\varepsilon$.

Với $N_x,N_y$ thực dương, đặt $u=A_x/N_x^r$, $v=A_y/N_y^r$. Từ $(u-v)^2\geq0$ và $u+v\leq\delta$,
\[
 4uv=(u+v)^2-(u-v)^2\leq(u+v)^2\leq\delta^2,
\]
\[
 N_xN_y=(A_xA_y/(uv))^{1/r}\geq(4A_xA_y/\delta^2)^{1/r}.
\]
Hai dấu bằng cùng đạt tại $u=v=\delta/2$, tương ứng
\[
 N_x=(2A_x/\delta)^{1/r},\qquad N_y=(2A_y/\delta)^{1/r}.
\]
Đây là cực tiểu của tích trên miền số thực dương.
\end{chungminh}

\Needspace{9\baselineskip}
\begin{vidu}
Cho $f(x,y)=x^4+2y^4$ trên $[0,1]^2$ và lưới $N_x=N_y=2$. Bảng giá trị và trọng số Simpson là
\begin{center}
\begin{tabular}{c|c|ccc}
\toprule
 & Trọng số $a_i$ & $y_0=0$ & $y_1=1/2$ & $y_2=1$\\
\midrule
Trọng số $b_j$ & & $1$ & $4$ & $1$\\
$x_0=0$ & $1$ & $0$ & $1/8$ & $2$\\
$x_1=1/2$ & $4$ & $1/16$ & $3/16$ & $33/16$\\
$x_2=1$ & $1$ & $1$ & $9/8$ & $3$\\
\bottomrule
\end{tabular}\end{center}
Các tổng theo hàng $\sum_jb_jf(x_i,y_j)$ là $5/2$, $23/8$, $17/2$. Nhân với $a_i$, cộng và nhân $h_xh_y/9=1/36$ cho $Q_{xy}=5/8$. Tích phân đúng là $3/5$, nên sai số $I-Q_{xy}=-1/40$. Vì $M_x=24$, $M_y=48$, chặn từ Định lý \ref{dl:quad-tensor-error} bằng $(24+48)/(180\cdot16)=1/40$.
\end{vidu}

% END EXTENSION CHAPTER 2

\appendix
\chapter{Thuật toán}
\label{app:algorithms}
Các chỉ số lặp thuộc $\mathbb Z$. Vòng từ $r$ tới $s$ có tập chỉ số $\{i\in\mathbb Z:r\leq i\leq s\}$ theo thứ tự tăng; vòng từ $r$ về $s$ có tập chỉ số $\{i\in\mathbb Z:s\leq i\leq r\}$ theo thứ tự giảm. Nếu tập chỉ số rỗng, khối lặp không được thực hiện.

\Needspace{20\baselineskip}
\section{Sinh mốc Chebyshev}
\begin{algorithm}[H]
\caption{Sinh mốc Chebyshev}
\label{alg:chebyshev-nodes}
\begin{algorithmic}[1]
\Require $a,b\in\R$, $a<b$; số nguyên $m\geq1$.
\Ensure $\boldsymbol x=(x_0,\ldots,x_{m-1})^T\in\R^m$ theo \eqref{eq:optimal-nodes}.
\State $\mu\gets(a+b)/2$; $\rho\gets(b-a)/2$; $h\gets\lfloor m/2\rfloor$
\For{$i$ từ $0$ tới $h-1$}
  \State $z\gets\cos((2i+1)\pi/(2m))$
  \State $x_i\gets\mu-\rho z$; $x_{m-1-i}\gets\mu+\rho z$
\EndFor
\Statex
\If{$m=2h+1$}
  \State $x_h\gets\mu$
\EndIf
\Statex
\State $\boldsymbol x\gets(x_0,\ldots,x_{m-1})^T$.
\State Kết thúc.
\end{algorithmic}
\end{algorithm}

\Needspace{20\baselineskip}
\section{Tính giá trị đa thức}
\begin{algorithm}[H]
\caption{Tính giá trị đa thức bằng Horner}
\label{alg:horner-eval}
\begin{algorithmic}[1]
\Require Số nguyên $n\geq0$; $\boldsymbol a=(a_0,\ldots,a_n)^T\in\R^{n+1}$; $c\in\R$. $p(x)=\sum_{j=0}^{n}a_jx^j$.
\Ensure $v=\sum_{j=0}^{n}a_jc^j$.
\State $v\gets a_n$
\For{$j$ từ $n-1$ về $0$}
  \State $v\gets a_j+cv$
\EndFor
\Statex

\State Kết thúc.
\end{algorithmic}
\end{algorithm}

\Needspace{20\baselineskip}
\section{Chia cho nhân tử tuyến tính}
\begin{algorithm}[H]
\caption{Chia cho nhân tử tuyến tính bằng Horner}
\label{alg:horner-div}
\begin{algorithmic}[1]
\Require Số nguyên $n\geq0$; $\boldsymbol a=(a_0,\ldots,a_n)^T\in\R^{n+1}$; $c\in\R$. $p(x)=\sum_{j=0}^{n}a_jx^j$.
\Ensure $(\boldsymbol q,r)$, trong đó $p(x)=(x-c)q(x)+r$; $\boldsymbol q$ là vector hệ số của $q$.
\If{$n=0$}
  \State $(\boldsymbol q,r)\gets((0)^T,a_0)$; dừng.
\EndIf
\Statex
\State $b_n\gets a_n$
\For{$j$ từ $n-1$ về $0$}
  \State $b_j\gets a_j+cb_{j+1}$
\EndFor
\Statex
\State $(\boldsymbol q,r)\gets((b_1,\ldots,b_n)^T,b_0)$.
\State Kết thúc.
\end{algorithmic}
\end{algorithm}

\Needspace{20\baselineskip}
\section{Đạo hàm bằng Horner lặp}
\begin{algorithm}[H]
\caption{Tính đạo hàm bằng Horner lặp -- sơ lược}
\label{alg:horner-derivatives}
\begin{algorithmic}[1]
\Require Số nguyên $n,K$, $n\geq0$, $0\leq K\leq n$; $\boldsymbol a=(a_0,\ldots,a_n)^T\in\R^{n+1}$; $c\in\R$. $p(x)=\sum_{j=0}^{n}a_jx^j$.
\Ensure $\boldsymbol v=(v_0,\ldots,v_K)^T$, $v_k=p^{(k)}(c)$.
\State $\boldsymbol b\gets\boldsymbol a$; $F\gets1$
\For{$s$ từ $0$ tới $K$}
  \State Áp dụng Thuật toán \ref{alg:horner-div} với $\boldsymbol b\in\R^{n-s+1}$ tại $c$; thu được $(\boldsymbol q,r)$.
  \State $v_s\gets Fr$
  \State $\boldsymbol b\gets\boldsymbol q$; $F\gets(s+1)F$
\EndFor
\Statex
\State $\boldsymbol v\gets(v_0,\ldots,v_K)^T$.
\State Kết thúc.
\end{algorithmic}
\end{algorithm}

\begin{algorithm}[H]
\caption{Tính đạo hàm bằng Horner lặp -- chi tiết}
\label{alg:horner-derivatives-detailed}
\begin{algorithmic}[1]
\Require Số nguyên $n,K$, $n\geq0$, $0\leq K\leq n$; $\boldsymbol a=(a_0,\ldots,a_n)^T\in\R^{n+1}$; $c\in\R$. $p(x)=\sum_{j=0}^{n}a_jx^j$.
\Ensure $\boldsymbol v=(v_0,\ldots,v_K)^T$, $v_k=p^{(k)}(c)$.
\State $\boldsymbol b\gets\boldsymbol a$; $F\gets1$.
\For{$s$ từ $0$ tới $K$}
  \State $d\gets n-s$, $\boldsymbol b=(b_0,\ldots,b_d)^T$.
  \If{$d=0$}
    \State $r\gets b_0$; $\boldsymbol q\gets(0)^T$.
  \Else
    \State $t_d\gets b_d$.
    \For{$j$ từ $d-1$ về $0$}
      \State $t_j\gets b_j+ct_{j+1}$.
    \EndFor
\Statex
    \State $r\gets t_0$; $\boldsymbol q\gets(t_1,\ldots,t_d)^T$.
  \EndIf
\Statex
  \State $v_s\gets Fr$.
  \State $\boldsymbol b\gets\boldsymbol q$; $F\gets(s+1)F$.
\EndFor
\Statex
\State $\boldsymbol v\gets(v_0,\ldots,v_K)^T$.
\State Kết thúc.
\end{algorithmic}
\end{algorithm}

\Needspace{20\baselineskip}
\section{Horner mở rộng}
\begin{algorithm}[H]
\caption{Tính đạo hàm bằng Horner mở rộng}
\label{alg:horner-jet}
\begin{algorithmic}[1]
\Require Số nguyên $n,K$, $n\geq0$, $0\leq K\leq n$; $\boldsymbol a=(a_0,\ldots,a_n)^T\in\R^{n+1}$; $c\in\R$. $p(x)=\sum_{j=0}^{n}a_jx^j$.
\Ensure $\boldsymbol d=(d_0,\ldots,d_K)^T$, $d_j=p^{(j)}(c)/j!$.
\State $d_0\gets a_n$; $d_j\gets0$ với $1\leq j\leq K$.
\For{$i$ từ $n-1$ về $0$}
  \For{$j$ từ $K$ về $1$}
    \State $d_j\gets cd_j+d_{j-1}$
  \EndFor
\Statex
  \State $d_0\gets cd_0+a_i$
\EndFor
\Statex
\State $\boldsymbol d\gets(d_0,\ldots,d_K)^T$.
\State Kết thúc.
\end{algorithmic}
\end{algorithm}

\Needspace{20\baselineskip}
\section{Horner ngược}
\begin{algorithm}[H]
\caption{Nhân với nhân tử tuyến tính bằng Horner ngược}
\label{alg:inverse-horner}
\begin{algorithmic}[1]
\Require Số nguyên $n\geq0$; $\boldsymbol a=(a_0,\ldots,a_n)^T\in\R^{n+1}$; $c\in\R$. $p(x)=\sum_{j=0}^{n}a_jx^j$.
\Ensure $\boldsymbol d=(d_0,\ldots,d_{n+1})^T$, $\sum_{j=0}^{n+1}d_jx^j=(x-c)p(x)$.
\State $d_{n+1}\gets a_n$
\For{$j$ từ $n$ về $1$}
  \State $d_j\gets a_{j-1}-ca_j$
\EndFor
\Statex
\State $d_0\gets-ca_0$
\State $\boldsymbol d\gets(d_0,\ldots,d_{n+1})^T$.
\State Kết thúc.
\end{algorithmic}
\end{algorithm}

\Needspace{20\baselineskip}
\section{Tạo đa thức tích}
\begin{algorithm}[H]
\caption{Tạo đa thức tích bằng Horner ngược -- sơ lược}
\label{alg:horner-product}
\begin{algorithmic}[1]
\Require Số nguyên $m\geq0$; $\boldsymbol z\in\R^m$.
\Ensure $\boldsymbol a=(a_0,\ldots,a_m)^T$, $\sum_{j=0}^{m}a_jx^j=\prod_{i=0}^{m-1}(x-z_i)$.
\State $\boldsymbol a\gets(1)^T$
\For{$s$ từ $0$ tới $m-1$}
  \State Áp dụng Thuật toán \ref{alg:inverse-horner} với $\boldsymbol a\in\R^{s+1}$ tại $z_s$; thay $\boldsymbol a$ bằng vector kết quả.
\EndFor
\Statex

\State Kết thúc.
\end{algorithmic}
\end{algorithm}

\begin{algorithm}[H]
\caption{Tạo đa thức tích bằng Horner ngược -- chi tiết}
\label{alg:horner-product-detailed}
\begin{algorithmic}[1]
\Require Số nguyên $m\geq0$; $\boldsymbol z\in\R^m$.
\Ensure $\boldsymbol a=(a_0,\ldots,a_m)^T$, $\sum_{j=0}^{m}a_jx^j=\prod_{i=0}^{m-1}(x-z_i)$.
\State $\boldsymbol a\gets(1)^T$.
\For{$s$ từ $0$ tới $m-1$}
  \State $\boldsymbol a=(a_0,\ldots,a_s)^T\in\R^{s+1}$.
  \State $d_{s+1}\gets a_s$.
  \For{$j$ từ $s$ về $1$}
    \State $d_j\gets a_{j-1}-z_sa_j$.
  \EndFor
\Statex
  \State $d_0\gets-z_sa_0$.
  \State $\boldsymbol a\gets(d_0,\ldots,d_{s+1})^T$.
\EndFor
\Statex

\State Kết thúc.
\end{algorithmic}
\end{algorithm}

\Needspace{20\baselineskip}
\section{Nhân hai đa thức}
\begin{algorithm}[H]
\caption{Nhân hai đa thức bằng Horner -- sơ lược}
\label{alg:polynomial-horner-mul}
\begin{algorithmic}[1]
\Require Số nguyên $n,m\geq0$; $\boldsymbol a=(a_0,\ldots,a_n)^T\in\R^{n+1}$; $\boldsymbol\beta=(\beta_0,\ldots,\beta_m)^T\in\R^{m+1}$. $p(x)=\sum_{j=0}^{n}a_jx^j$, $q(x)=\sum_{k=0}^{m}\beta_kx^k$.
\Ensure $\boldsymbol r=(r_0,\ldots,r_{n+m})^T$, $\sum_{j=0}^{n+m}r_jx^j=p(x)q(x)$.
\State $\boldsymbol r\gets\beta_m\boldsymbol a$; $\ell\gets n$
\For{$k$ từ $m-1$ về $0$}
  \State Áp dụng Thuật toán \ref{alg:inverse-horner} với $\boldsymbol r\in\R^{\ell+1}$ tại $0$; thay $\boldsymbol r$ bằng vector kết quả, $\ell\gets\ell+1$.
  \For{$j$ từ $0$ tới $n$}
    \State $r_j\gets r_j+\beta_ka_j$
  \EndFor
\EndFor
\Statex

\State Kết thúc.
\end{algorithmic}
\end{algorithm}

\begin{algorithm}[H]
\caption{Nhân hai đa thức bằng Horner -- chi tiết}
\label{alg:polynomial-horner-mul-detailed}
\begin{algorithmic}[1]
\Require Số nguyên $n,m\geq0$; $\boldsymbol a=(a_0,\ldots,a_n)^T\in\R^{n+1}$; $\boldsymbol\beta=(\beta_0,\ldots,\beta_m)^T\in\R^{m+1}$. $p(x)=\sum_{j=0}^{n}a_jx^j$, $q(x)=\sum_{k=0}^{m}\beta_kx^k$.
\Ensure $\boldsymbol r=(r_0,\ldots,r_{n+m})^T$, $\sum_{j=0}^{n+m}r_jx^j=p(x)q(x)$.
\State $\boldsymbol r\gets\beta_m\boldsymbol a$; $\ell\gets n$.
\For{$k$ từ $m-1$ về $0$}
  \State $\boldsymbol r=(r_0,\ldots,r_\ell)^T\in\R^{\ell+1}$; $d_0\gets0$.
  \For{$j$ từ $1$ tới $\ell+1$}
    \State $d_j\gets r_{j-1}$.
  \EndFor
\Statex
  \For{$j$ từ $0$ tới $n$}
    \State $d_j\gets d_j+\beta_ka_j$.
  \EndFor
\Statex
  \State $\ell\gets\ell+1$; $\boldsymbol r\gets(d_0,\ldots,d_\ell)^T$.
\EndFor
\Statex

\State Kết thúc.
\end{algorithmic}
\end{algorithm}

\Needspace{20\baselineskip}
\section{Chia Horner tổng quát}
\begin{algorithm}[H]
\caption{Chia đa thức bằng Horner tổng quát}
\label{alg:generalized-horner-div}
\begin{algorithmic}[1]
\Require Số nguyên $n\geq d\geq1$; $\boldsymbol a=(a_0,\ldots,a_n)^T\in\R^{n+1}$; $\boldsymbol g=(g_0,\ldots,g_d)^T\in\R^{d+1}$, $g_d\ne0$. $p(x)=\sum_{j=0}^{n}a_jx^j$, $g(x)=\sum_{j=0}^{d}g_jx^j$.
\Ensure $(\boldsymbol q,\boldsymbol r)$, $p(x)=g(x)q(x)+r(x)$, $\deg r<d$, theo \eqref{eq:generalized-division}.
\State $\boldsymbol b\gets\boldsymbol a$
\For{$i$ từ $n$ về $d$}
  \State $b_i\gets b_i/g_d$
  \For{$s$ từ $1$ tới $d$}
    \State $b_{i-s}\gets b_{i-s}-g_{d-s}b_i$
  \EndFor
\EndFor
\Statex
\State $(\boldsymbol q,\boldsymbol r)\gets((b_d,\ldots,b_n)^T,(b_0,\ldots,b_{d-1})^T)$.
\State Kết thúc.
\end{algorithmic}
\end{algorithm}

\Needspace{20\baselineskip}
\section{Hệ số nội suy Lagrange bằng Horner}
\begin{algorithm}[H]
\caption{Tạo hệ số nội suy Lagrange bằng Horner -- sơ lược}
\label{alg:lagrange-coefficients}
\begin{algorithmic}[1]
\Require Số nguyên $n\geq0$; $\boldsymbol x=(x_0,\ldots,x_n)^T\in\R^{n+1}$ đôi một phân biệt; $\boldsymbol y=(y_0,\ldots,y_n)^T\in\R^{n+1}$.
\Ensure $\boldsymbol a=(a_0,\ldots,a_n)^T$, $\sum_{j=0}^{n}a_jx_i^j=y_i$ với $0\leq i\leq n$.
\State Áp dụng Thuật toán \ref{alg:horner-product} với $\boldsymbol x$ gồm $n+1$ mốc; thu được $\boldsymbol w\in\R^{n+2}$.
\For{$i$ từ $0$ tới $n$}
  \State Áp dụng Thuật toán \ref{alg:horner-div} với $\boldsymbol w$ tại $x_i$; thu được $(\boldsymbol q_i,r_i)$.
  \State Áp dụng Thuật toán \ref{alg:horner-eval} với $\boldsymbol q_i$ tại $x_i$; thu được $\lambda_i$.
\EndFor
\Statex
\State $\displaystyle\boldsymbol a\gets\sum_{i=0}^{n}\frac{y_i}{\lambda_i}\boldsymbol q_i$.

\State Kết thúc.
\end{algorithmic}
\end{algorithm}

\begin{algorithm}[H]
\caption{Tạo hệ số nội suy Lagrange bằng Horner -- chi tiết}
\label{alg:lagrange-coefficients-detailed}
\small
\begin{algorithmic}[1]
\Require Số nguyên $n\geq0$; $\boldsymbol x=(x_0,\ldots,x_n)^T\in\R^{n+1}$ đôi một phân biệt; $\boldsymbol y=(y_0,\ldots,y_n)^T\in\R^{n+1}$.
\Ensure $\boldsymbol a=(a_0,\ldots,a_n)^T$, $\sum_{j=0}^{n}a_jx_i^j=y_i$ với $0\leq i\leq n$.
\State $\boldsymbol w\gets(1)^T$.
\For{$s$ từ $0$ tới $n$}
  \State $\boldsymbol w=(w_0,\ldots,w_s)^T$; $d_{s+1}\gets w_s$.
  \For{$j$ từ $s$ về $1$}
    \State $d_j\gets w_{j-1}-x_sw_j$.
  \EndFor
\Statex
  \State $d_0\gets-x_sw_0$; $\boldsymbol w\gets(d_0,\ldots,d_{s+1})^T$.
\EndFor
\Statex
\State $a_j\gets0$ với $0\leq j\leq n$.
\For{$i$ từ $0$ tới $n$}
  \State $b_{n+1}\gets w_{n+1}$.
  \For{$j$ từ $n$ về $0$}
    \State $b_j\gets w_j+x_ib_{j+1}$.
  \EndFor
\Statex
  \State $r_i\gets b_0$; $\boldsymbol q_i\gets(b_1,\ldots,b_{n+1})^T$; $\lambda_i\gets b_{n+1}$.
  \For{$j$ từ $n-1$ về $0$}
    \State $\lambda_i\gets b_{j+1}+x_i\lambda_i$.
  \EndFor
\Statex
  \For{$j$ từ $0$ tới $n$}
    \State $a_j\gets a_j+(y_i/\lambda_i)b_{j+1}$.
  \EndFor
\EndFor
\Statex
\State $\boldsymbol a\gets(a_0,\ldots,a_n)^T$.
\State Kết thúc.
\end{algorithmic}
\end{algorithm}

\Needspace{20\baselineskip}
\section{Tính giá trị đa thức nội suy tại \texorpdfstring{$c$}{c} bằng công thức Lagrange}
\begin{algorithm}[H]
\caption{Tính giá trị đa thức nội suy tại $c$ bằng công thức Lagrange -- sơ lược}
\label{alg:lagrange-evaluate-short}
\begin{algorithmic}[1]
\Require Số nguyên $n\geq0$; $\boldsymbol x=(x_0,\ldots,x_n)^T\in\R^{n+1}$ đôi một phân biệt; $\boldsymbol y=(y_0,\ldots,y_n)^T\in\R^{n+1}$; $c\in\R$.
\Ensure $v=p(c)$, trong đó $p\in\mathcal P_n$ là đa thức nội suy dữ liệu $(x_i,y_i)$, $0\leq i\leq n$.
\If{$c=x_i$ với một $i\in\{0,\ldots,n\}$}
  \State $v\gets y_i$; dừng.
\EndIf
\Statex
\State $\displaystyle\ell_i\gets\prod_{\substack{0\leq j\leq n\\j\ne i}}\frac{c-x_j}{x_i-x_j}$ với $0\leq i\leq n$.
\State $\displaystyle v\gets\sum_{i=0}^{n}y_i\ell_i$.

\State Kết thúc.
\end{algorithmic}
\end{algorithm}

\begin{algorithm}[H]
\caption{Tính giá trị đa thức nội suy tại $c$ bằng công thức Lagrange -- chi tiết}
\label{alg:lagrange-evaluate}
\begin{algorithmic}[1]
\Require Số nguyên $n\geq0$; $\boldsymbol x=(x_0,\ldots,x_n)^T\in\R^{n+1}$ đôi một phân biệt; $\boldsymbol y=(y_0,\ldots,y_n)^T\in\R^{n+1}$; $c\in\R$.
\Ensure $v=p(c)$, trong đó $p\in\mathcal P_n$ là đa thức nội suy dữ liệu $(x_i,y_i)$, $0\leq i\leq n$.
\For{$i$ từ $0$ tới $n$}
  \If{$c=x_i$}
    \State $v\gets y_i$; dừng.
  \EndIf
\EndFor
\Statex
\State $v\gets0$.
\For{$i$ từ $0$ tới $n$}
  \State $\ell\gets1$.
  \For{$j$ từ $0$ tới $n$}
    \If{$j\ne i$}
      \State $\ell\gets\ell(c-x_j)/(x_i-x_j)$.
    \EndIf
  \EndFor
\Statex
  \State $v\gets v+y_i\ell$.
\EndFor
\Statex

\State Kết thúc.
\end{algorithmic}
\end{algorithm}

\Needspace{20\baselineskip}
\section{Tính giá trị đa thức nội suy tại \texorpdfstring{$c$}{c} theo dạng \texorpdfstring{$\omega_{n+1}$}{omega}}
\begin{algorithm}[H]
\caption{Tính giá trị đa thức nội suy tại $c$ theo dạng $\omega_{n+1}$ -- sơ lược}
\label{alg:lagrange-omega-evaluate-short}
\begin{algorithmic}[1]
\Require Số nguyên $n\geq0$; $\boldsymbol x=(x_0,\ldots,x_n)^T\in\R^{n+1}$ đôi một phân biệt; $\boldsymbol y=(y_0,\ldots,y_n)^T\in\R^{n+1}$; $c\in\R$. $\boldsymbol\lambda=(\lambda_0,\ldots,\lambda_n)^T$, $\lambda_i=\prod_{\substack{0\leq j\leq n\\j\ne i}}(x_i-x_j)$.
\Ensure $v=p(c)$, trong đó $p\in\mathcal P_n$ là đa thức nội suy dữ liệu $(x_i,y_i)$, $0\leq i\leq n$.
\If{$c=x_i$ với một $i\in\{0,\ldots,n\}$}
  \State $v\gets y_i$; dừng.
\EndIf
\Statex
\State $\displaystyle\omega\gets\prod_{i=0}^{n}(c-x_i)$.
\State $\displaystyle s\gets\sum_{i=0}^{n}\frac{y_i}{\lambda_i(c-x_i)}$.
\State $v\gets\omega s$.
\State Kết thúc.
\end{algorithmic}
\end{algorithm}

\begin{algorithm}[H]
\caption{Tính giá trị đa thức nội suy tại $c$ theo dạng $\omega_{n+1}$ -- chi tiết}
\label{alg:lagrange-omega-evaluate}
\begin{algorithmic}[1]
\Require Số nguyên $n\geq0$; $\boldsymbol x=(x_0,\ldots,x_n)^T\in\R^{n+1}$ đôi một phân biệt; $\boldsymbol y=(y_0,\ldots,y_n)^T\in\R^{n+1}$; $c\in\R$. $\boldsymbol\lambda=(\lambda_0,\ldots,\lambda_n)^T$, $\lambda_i=\prod_{\substack{0\leq j\leq n\\j\ne i}}(x_i-x_j)$.
\Ensure $v=p(c)$, trong đó $p\in\mathcal P_n$ là đa thức nội suy dữ liệu $(x_i,y_i)$, $0\leq i\leq n$.
\State $\omega\gets1$.
\For{$i$ từ $0$ tới $n$}
  \If{$c=x_i$}
    \State $v\gets y_i$; dừng.
  \EndIf
\Statex
  \State $\omega\gets\omega(c-x_i)$.
\EndFor
\Statex
\State $s\gets0$.
\For{$i$ từ $0$ tới $n$}
  \State $s\gets s+y_i/(\lambda_i(c-x_i))$.
\EndFor
\Statex
\State $v\gets\omega s$.
\State Kết thúc.
\end{algorithmic}
\end{algorithm}

\Needspace{20\baselineskip}
\section{Tính giá trị đa thức nội suy tại \texorpdfstring{$c$}{c} theo dạng tỷ số}
\begin{algorithm}[H]
\caption{Tính giá trị đa thức nội suy tại $c$ theo dạng tỷ số -- sơ lược}
\label{alg:barycentric-evaluate-short}
\begin{algorithmic}[1]
\Require Số nguyên $n\geq0$; $\boldsymbol x=(x_0,\ldots,x_n)^T\in\R^{n+1}$ đôi một phân biệt; $\boldsymbol y=(y_0,\ldots,y_n)^T\in\R^{n+1}$; $c\in\R$. $\boldsymbol\lambda=(\lambda_0,\ldots,\lambda_n)^T$, $\lambda_i=\prod_{\substack{0\leq j\leq n\\j\ne i}}(x_i-x_j)$.
\Ensure $v=p(c)$, trong đó $p\in\mathcal P_n$ là đa thức nội suy dữ liệu $(x_i,y_i)$, $0\leq i\leq n$.
\If{$c=x_i$ với một $i\in\{0,\ldots,n\}$}
  \State $v\gets y_i$; dừng.
\EndIf
\Statex
\State $u_i\gets1/(\lambda_i(c-x_i))$ với $0\leq i\leq n$.
\State $\displaystyle s\gets\sum_{i=0}^{n}y_iu_i$, $\displaystyle t\gets\sum_{i=0}^{n}u_i$.
\State $v\gets s/t$.
\State Kết thúc.
\end{algorithmic}
\end{algorithm}

\begin{algorithm}[H]
\caption{Tính giá trị đa thức nội suy tại $c$ theo dạng tỷ số -- chi tiết}
\label{alg:barycentric-evaluate}
\begin{algorithmic}[1]
\Require Số nguyên $n\geq0$; $\boldsymbol x=(x_0,\ldots,x_n)^T\in\R^{n+1}$ đôi một phân biệt; $\boldsymbol y=(y_0,\ldots,y_n)^T\in\R^{n+1}$; $c\in\R$. $\boldsymbol\lambda=(\lambda_0,\ldots,\lambda_n)^T$, $\lambda_i=\prod_{\substack{0\leq j\leq n\\j\ne i}}(x_i-x_j)$.
\Ensure $v=p(c)$, trong đó $p\in\mathcal P_n$ là đa thức nội suy dữ liệu $(x_i,y_i)$, $0\leq i\leq n$.
\For{$i$ từ $0$ tới $n$}
  \If{$c=x_i$}
    \State $v\gets y_i$; dừng.
  \EndIf
\EndFor
\Statex
\State $s\gets0$; $t\gets0$.
\For{$i$ từ $0$ tới $n$}
  \State $u\gets1/(\lambda_i(c-x_i))$.
  \State $s\gets s+uy_i$; $t\gets t+u$.
\EndFor
\Statex
\State $v\gets s/t$.
\State Kết thúc.
\end{algorithmic}
\end{algorithm}

% BEGIN EXTENSION NEWTON ALGORITHMS
\Needspace{20\baselineskip}
\section{Bảng tỷ sai phân và hệ số Newton}
\begin{algorithm}[H]
\caption{Dựng bảng tỷ sai phân}
\label{alg:newton-table}
\begin{algorithmic}[1]
\Require $n\geq0$; $\boldsymbol x=(x_0,\ldots,x_n)^T$ có các thành phần đôi một phân biệt; $\boldsymbol y=(y_0,\ldots,y_n)^T$.
\Ensure Các số $D_i^{(k)}=f[x_i,\ldots,x_{i+k}]$, $0\leq k\leq n$, $0\leq i\leq n-k$; $\boldsymbol d=(d_0,\ldots,d_n)^T$, $d_k=D_0^{(k)}$.
\For{$i$ từ $0$ tới $n$}
 \State $D_i^{(0)}\gets y_i$.
\EndFor
\Statex
\State $d_0\gets y_0$.
\For{$k$ từ $1$ tới $n$}
 \For{$i$ từ $0$ tới $n-k$}
  \State $D_i^{(k)}\gets(D_{i+1}^{(k-1)}-D_i^{(k-1)})/(x_{i+k}-x_i)$.
 \EndFor
 \State $d_k\gets D_0^{(k)}$.
\EndFor
\Statex
\State Kết thúc.
\end{algorithmic}
\end{algorithm}

\Needspace{20\baselineskip}
\section{Giá trị và đạo hàm trong cơ sở Newton}
\begin{algorithm}[H]
\caption{Horner trong cơ sở Newton}
\label{alg:newton-eval}
\begin{algorithmic}[1]
\Require $n\geq0$; $\boldsymbol x=(x_0,\ldots,x_{n-1})^T$ (vector rỗng nếu $n=0$); $\boldsymbol d=(d_0,\ldots,d_n)^T$; $c\in\R$. $p(x)=\sum_{k=0}^{n}d_k\prod_{j=0}^{k-1}(x-x_j)$.
\Ensure $v=p(c)$ và $u=p'(c)$.
\State $v\gets d_n$; $u\gets0$.
\For{$i$ từ $n-1$ về $0$}
 \State $u\gets v+(c-x_i)u$.
 \State $v\gets d_i+(c-x_i)v$.
\EndFor
\Statex
\State Kết thúc.
\end{algorithmic}
\end{algorithm}

\Needspace{20\baselineskip}
\section{Bổ sung một mốc nội suy Newton}
\begin{algorithm}[H]
\caption{Bổ sung mốc Newton}
\label{alg:newton-append}
\begin{algorithmic}[1]
\Require $\boldsymbol x=(x_0,\ldots,x_n)^T$ phân biệt; hệ số Newton $\boldsymbol d=(d_0,\ldots,d_n)^T$; $z\notin\{x_0,\ldots,x_n\}$; giá trị mới $w$.
\Ensure $\widetilde{\boldsymbol d}=(d_0,\ldots,d_n,d_{n+1})^T$ cho tập mốc $(x_0,\ldots,x_n,z)$ với dữ liệu mới $w$.
\State $v\gets d_n$; $W\gets1$.
\For{$i$ từ $n-1$ về $0$}
 \State $v\gets d_i+(z-x_i)v$.
\EndFor
\Statex
\For{$i$ từ $0$ tới $n$}
 \State $W\gets W(z-x_i)$.
\EndFor
\Statex
\State $d_{n+1}\gets(w-v)/W$; $\widetilde{\boldsymbol d}\gets(d_0,\ldots,d_n,d_{n+1})^T$.
\State Kết thúc.
\end{algorithmic}
\end{algorithm}

\Needspace{20\baselineskip}
\section{Chuyển hệ số Newton sang hệ số đơn thức}
\begin{algorithm}[H]
\caption{Chuyển cơ sở Newton bằng Horner ngược}
\label{alg:newton-coefficients}
\begin{algorithmic}[1]
\Require $n\geq0$; $\boldsymbol x=(x_0,\ldots,x_{n-1})^T$; $\boldsymbol d=(d_0,\ldots,d_n)^T$ như trong \eqref{eq:newton-form}.
\Ensure $\boldsymbol a=(a_0,\ldots,a_n)^T$ sao cho $p_n(x)=\sum_{j=0}^{n}a_jx^j$.
\State $a_0\gets d_n$; $m\gets0$.
\For{$i$ từ $n-1$ về $0$}
 \State $a_{m+1}\gets a_m$.
 \For{$j$ từ $m$ về $1$}
  \State $a_j\gets a_{j-1}-x_i a_j$.
 \EndFor
 \Statex
 \State $a_0\gets d_i-x_i a_0$; $m\gets m+1$.
\EndFor
\Statex
\State $\boldsymbol a\gets(a_0,\ldots,a_n)^T$.
\State Kết thúc.
\end{algorithmic}
\end{algorithm}

\Needspace{20\baselineskip}
\section{Bảng nhân Newton}
\begin{algorithm}[H]
\caption{Dựng bảng nhân bằng Horner ngược}
\label{alg:newton-product-table}
\begin{algorithmic}[1]
\Require $n\geq0$; các mốc $x_0,\ldots,x_{n-1}$ (không cần mốc khi $n=0$).
\Ensure $B=(B_{j,k})_{0\leq j,k\leq n}$, cột $k$ là vector hệ số của $W_k(x)=\prod_{i=0}^{k-1}(x-x_i)$.
\State $B_{0,0}\gets1$.
\For{$k$ từ $0$ tới $n$}
 \For{$j$ từ $k+1$ tới $n$}
  \State $B_{j,k}\gets0$.
 \EndFor
\EndFor
\Statex
\For{$k$ từ $0$ tới $n-1$}
 \State $B_{0,k+1}\gets-x_kB_{0,k}$; $B_{k+1,k+1}\gets B_{k,k}$.
 \For{$j$ từ $1$ tới $k$}
  \State $B_{j,k+1}\gets B_{j-1,k}-x_kB_{j,k}$.
 \EndFor
\EndFor
\Statex
\State Kết thúc.
\end{algorithmic}
\end{algorithm}

\Needspace{20\baselineskip}
\section{Nhân bảng nhân với vector tỷ sai phân}
\begin{algorithm}[H]
\caption{Hệ số Newton từ bảng nhân: sơ lược}
\label{alg:newton-table-coefficients-short}
\begin{algorithmic}[1]
\Require $n\geq0$; mốc $x_0,\ldots,x_{n-1}$; tỷ sai phân $\boldsymbol d=(d_0,\ldots,d_n)^T$.
\Ensure $\boldsymbol a=(a_0,\ldots,a_n)^T$, hệ số đơn thức của đa thức Newton.
\State Dựng $B$ bằng thuật toán \ref{alg:newton-product-table}.
\State $\boldsymbol a\gets B\boldsymbol d$.
\State Kết thúc.
\end{algorithmic}
\end{algorithm}
\begin{algorithm}[H]
\caption{Hệ số Newton từ bảng nhân: chi tiết}
\label{alg:newton-table-coefficients}
\begin{algorithmic}[1]
\Require $n\geq0$; mốc $x_0,\ldots,x_{n-1}$; tỷ sai phân $\boldsymbol d=(d_0,\ldots,d_n)^T$.
\Ensure $\boldsymbol a=(a_0,\ldots,a_n)^T$, hệ số đơn thức của đa thức Newton.
\State $B_{0,0}\gets1$.
\For{$k$ từ $0$ tới $n-1$}
 \State $B_{0,k+1}\gets-x_kB_{0,k}$; $B_{k+1,k+1}\gets B_{k,k}$.
 \For{$j$ từ $1$ tới $k$}
  \State $B_{j,k+1}\gets B_{j-1,k}-x_kB_{j,k}$.
 \EndFor
\EndFor
\Statex
\For{$j$ từ $0$ tới $n$}
 \State $a_j\gets0$.
 \For{$k$ từ $j$ tới $n$}
  \State $a_j\gets a_j+B_{j,k}d_k$.
 \EndFor
\EndFor
\Statex
\State $\boldsymbol a\gets(a_0,\ldots,a_n)^T$.
\State Kết thúc.
\end{algorithmic}
\end{algorithm}

\Needspace{20\baselineskip}
\section{Bảng sai phân tiến}
\begin{algorithm}[H]
\caption{Dựng bảng sai phân tiến}
\label{alg:forward-difference-table}
\begin{algorithmic}[1]
\Require Hai số nguyên $\ell\leq r$; các giá trị $y_\ell,\ldots,y_r$.
\Ensure $F_i^{(k)}=\Delta^ky_i$, $0\leq k\leq r-\ell$, $\ell\leq i\leq r-k$.
\For{$i$ từ $\ell$ tới $r$}
 \State $F_i^{(0)}\gets y_i$.
\EndFor
\Statex
\For{$k$ từ $1$ tới $r-\ell$}
 \For{$i$ từ $\ell$ tới $r-k$}
  \State $F_i^{(k)}\gets F_{i+1}^{(k-1)}-F_i^{(k-1)}$.
 \EndFor
\EndFor
\Statex
\State Kết thúc.
\end{algorithmic}
\end{algorithm}

\Needspace{20\baselineskip}
\section{Newton tiến và lùi từ một mốc chọn trước}
\begin{algorithm}[H]
\caption{Giá trị và đạo hàm từ $x_k$: sơ lược}
\label{alg:newton-from-k-short}
\begin{algorithmic}[1]
\Require $x_i=x_0+ih$, $h\ne0$; bảng $F_i^{(j)}=\Delta^jy_i$; $\sigma\in\{1,-1\}$; $k,r$ có $0\leq k\leq k+r\leq n$ nếu $\sigma=1$, $0\leq k-r\leq k\leq n$ nếu $\sigma=-1$; $c\in\R$.
\Ensure $v=P_{k,r}^{\sigma}(c)$ và $u=(P_{k,r}^{\sigma})'(c)$ theo \eqref{eq:newton-forward-k}, \eqref{eq:newton-backward-k}.
\State $t\gets(c-x_k)/h$; $z_j\gets\sigma j$ ($0\leq j<r$).
\State $d_j\gets F_{k-(1-\sigma)j/2}^{(j)}/j!$ ($0\leq j\leq r$).
\State Tính $(v,w)$ bằng thuật toán \ref{alg:newton-eval} với mốc $z_j$, hệ số $d_j$ và điểm $t$.
\State $u\gets w/h$.
\State Kết thúc.
\end{algorithmic}
\end{algorithm}
\begin{algorithm}[H]
\caption{Giá trị và đạo hàm từ $x_k$: chi tiết}
\label{alg:newton-from-k}
\begin{algorithmic}[1]
\Require $x_i=x_0+ih$, $h\ne0$; bảng $F_i^{(j)}=\Delta^jy_i$; $\sigma=1$ dùng $k,\ldots,k+r$, $\sigma=-1$ dùng $k-r,\ldots,k$; mọi chỉ số thuộc $0,\ldots,n$; $c\in\R$.
\Ensure $v=P_{k,r}^{\sigma}(c)$ và $u=(P_{k,r}^{\sigma})'(c)$.
\State $b\gets1$.
\For{$j$ từ $0$ tới $r$}
 \State $i\gets k-(1-\sigma)j/2$; $d_j\gets F_i^{(j)}/b$.
 \State $b\gets(j+1)b$.
\EndFor
\Statex
\State $t\gets(c-x_k)/h$; $v\gets d_r$; $u\gets0$.
\For{$j$ từ $r-1$ về $0$}
 \State $u\gets v+(t-\sigma j)u$.
 \State $v\gets d_j+(t-\sigma j)v$.
\EndFor
\Statex
\State $u\gets u/h$.
\State Kết thúc.
\end{algorithmic}
\end{algorithm}

\Needspace{20\baselineskip}
\section{Giá trị nội suy theo Gauss I và Gauss II}
\begin{algorithm}[H]
\caption{Tính giá trị theo công thức Gauss}
\label{alg:central-gauss}
\begin{algorithmic}[1]
\Require $m\geq1$, $h\ne0$, $x_0,c\in\R$; bảng $F_i^{(k)}=\Delta^ky_i$ trên chỉ số $-m,\ldots,m$; $\sigma=1$ cho Gauss I, $\sigma=-1$ cho Gauss II.
\Ensure $v=p_{2m}(c)$, đa thức nội suy dữ liệu $(x_0+ih,y_i)$, $-m\leq i\leq m$.
\State $t\gets(c-x_0)/h$; $A\gets t$; $v\gets F_0^{(0)}$; $b\gets1$.
\For{$r$ từ $0$ tới $m-1$}
 \State $i\gets-r-(1-\sigma)/2$.
 \State $v\gets v+F_i^{(2r+1)}A/b$.
 \State $b\gets(2r+2)b$.
 \State $v\gets v+F_{-r-1}^{(2r+2)}(t-\sigma(r+1))A/b$.
 \State $A\gets A(t^2-(r+1)^2)$; $b\gets(2r+3)b$.
\EndFor
\Statex
\State Kết thúc.
\end{algorithmic}
\end{algorithm}

\Needspace{20\baselineskip}
\section{Giá trị nội suy theo Stirling}
\begin{algorithm}[H]
\caption{Tính giá trị theo công thức Stirling}
\label{alg:central-stirling}
\begin{algorithmic}[1]
\Require $m\geq1$, $h\ne0$, $x_0,c\in\R$; bảng $F_i^{(k)}=\Delta^ky_i$ trên chỉ số $-m,\ldots,m$.
\Ensure $v=p_{2m}(c)$, đa thức nội suy dữ liệu $(x_0+ih,y_i)$, $-m\leq i\leq m$.
\State $t\gets(c-x_0)/h$; $A\gets t$; $v\gets F_0^{(0)}$; $b\gets1$.
\For{$r$ từ $0$ tới $m-1$}
 \State $v\gets v+(F_{-r-1}^{(2r+1)}+F_{-r}^{(2r+1)})A/(2b)$.
 \State $b\gets(2r+2)b$.
 \State $v\gets v+F_{-r-1}^{(2r+2)}tA/b$.
 \State $A\gets A(t^2-(r+1)^2)$; $b\gets(2r+3)b$.
\EndFor
\Statex
\State Kết thúc.
\end{algorithmic}
\end{algorithm}

\Needspace{20\baselineskip}
\section{Giá trị nội suy theo Bessel}
\begin{algorithm}[H]
\caption{Tính giá trị theo công thức Bessel}
\label{alg:central-bessel}
\begin{algorithmic}[1]
\Require $m\geq0$, $h\ne0$, $x_0,c\in\R$; bảng $F_i^{(k)}=\Delta^ky_i$ trên chỉ số $-m,\ldots,m+1$.
\Ensure $v=p_{2m+1}(c)$, đa thức nội suy dữ liệu $(x_0+ih,y_i)$, $-m\leq i\leq m+1$.
\State $t\gets(c-x_0)/h$; $B\gets1$; $v\gets0$; $b\gets1$.
\For{$r$ từ $0$ tới $m$}
 \State $v\gets v+(F_{-r}^{(2r)}+F_{1-r}^{(2r)})B/(2b)$.
 \State $v\gets v+F_{-r}^{(2r+1)}(t-1/2)B/((2r+1)b)$.
 \State $B\gets B(t+r)(t-r-1)$; $b\gets(2r+1)(2r+2)b$.
\EndFor
\Statex
\State Kết thúc.
\end{algorithmic}
\end{algorithm}

\Needspace{20\baselineskip}
\section{Dựng đa thức nội suy trung tâm bằng bảng nhân}
\begin{algorithm}[H]
\caption{Đa thức Gauss, Stirling hoặc Bessel: sơ lược}
\label{alg:central-polynomial-short}
\begin{algorithmic}[1]
\Require $m\geq0$; $x_j=x_0+jh$, $h\ne0$; $F_j^{(k)}=\Delta^ky_j$ trên các chỉ số $-m,\ldots,m$ với $\sigma\in\{\mathrm I,\mathrm{II},\mathrm S\}$, hoặc $-m,\ldots,m+1$ với $\sigma=\mathrm B$.
\Ensure $\boldsymbol a=(a_0,\ldots,a_N)^T$ sao cho $p(x_0+ht)=\sum_{j=0}^{N}a_jt^j$; $N=2m$ nếu $\sigma\ne\mathrm B$, $N=2m+1$ nếu $\sigma=\mathrm B$.
\State $N\gets2m$ nếu $\sigma\ne\mathrm B$; $N\gets2m+1$ nếu $\sigma=\mathrm B$; $r\gets1$ nếu $\sigma\in\{\mathrm I,\mathrm{II}\}$; $r\gets2$ nếu $\sigma\in\{\mathrm S,\mathrm B\}$.
\For{$\nu$ từ $1$ tới $r$}
 \State $(\delta,\eta)\gets(0,1)$ nếu $\sigma=\mathrm I$ hoặc $\nu=1$ và $\sigma\in\{\mathrm S,\mathrm B\}$; $(\delta,\eta)\gets(0,-1)$ nếu $\sigma=\mathrm{II}$ hoặc $\nu=2,\sigma=\mathrm S$; $(\delta,\eta)\gets(1,-1)$ nếu $\nu=2,\sigma=\mathrm B$.
 \State $z_0\gets\delta$; $z_k\gets\delta+\eta(-1)^{k+1}\lceil k/2\rceil$, $1\leq k\leq N$.
 \State $\ell_k\gets\min\{z_0,\ldots,z_k\}$; $d_k\gets F_{\ell_k}^{(k)}/k!$, $0\leq k\leq N$.
 \State Dựng $B^{(\nu)}$ từ $z_0,\ldots,z_{N-1}$ bằng thuật toán \ref{alg:newton-product-table}.
 \State $\boldsymbol a^{(\nu)}\gets B^{(\nu)}\boldsymbol d$.
\EndFor
\Statex
\State $\boldsymbol a\gets r^{-1}\sum_{\nu=1}^{r}\boldsymbol a^{(\nu)}$.
\State Kết thúc.
\end{algorithmic}
\end{algorithm}

\begin{algorithm}[H]
\caption{Đa thức Gauss, Stirling hoặc Bessel: chi tiết}
\label{alg:central-polynomial}
\begin{algorithmic}[1]
\Require $m\geq0$, $h\ne0$; bảng $F_j^{(k)}=\Delta^ky_j$ và lựa chọn $\sigma$ như trong thuật toán \ref{alg:central-polynomial-short}.
\Ensure $\boldsymbol a=(a_0,\ldots,a_N)^T$, hệ số của $p(x_0+ht)$, với $N$ như trong thuật toán \ref{alg:central-polynomial-short}.
\State $N\gets2m$; $r\gets1$.
\If{$\sigma\in\{\mathrm S,\mathrm B\}$}
 \State $r\gets2$.
\EndIf
\If{$\sigma=\mathrm B$}
 \State $N\gets2m+1$.
\EndIf
\Statex
\For{$j$ từ $0$ tới $N$}
 \State $a_j\gets0$.
\EndFor
\Statex
\For{$\nu$ từ $1$ tới $r$}
 \State $\delta\gets0$; $\eta\gets1$.
 \If{$\sigma=\mathrm{II}$ hoặc $\nu=2$}
  \State $\eta\gets-1$.
 \EndIf
 \If{$\sigma=\mathrm B$ và $\nu=2$}
  \State $\delta\gets1$.
 \EndIf
 \State $z_0\gets\delta$; $\ell_0\gets\delta$; $d_0\gets F_{\delta}^{(0)}$; $f\gets1$; $B_{0,0}\gets1$.
 \For{$k$ từ $1$ tới $N$}
  \State $z_k\gets\delta+\eta(-1)^{k+1}\lceil k/2\rceil$; $\ell_k\gets\min\{\ell_{k-1},z_k\}$; $f\gets kf$; $d_k\gets F_{\ell_k}^{(k)}/f$.
  \State $B_{0,k}\gets-z_{k-1}B_{0,k-1}$; $B_{k,k}\gets B_{k-1,k-1}$.
  \For{$j$ từ $1$ tới $k-1$}
   \State $B_{j,k}\gets B_{j-1,k-1}-z_{k-1}B_{j,k-1}$.
  \EndFor
 \EndFor
 \Statex
 \For{$j$ từ $0$ tới $N$}
  \For{$k$ từ $j$ tới $N$}
   \State $a_j\gets a_j+B_{j,k}d_k/r$.
  \EndFor
 \EndFor
\EndFor
\Statex
\State $\boldsymbol a\gets(a_0,\ldots,a_N)^T$.
\State Kết thúc.
\end{algorithmic}
\end{algorithm}

\section{Nội suy hàm ngược tại giá trị đích}
\begin{algorithm}[H]
\caption{Nội suy ngược: sơ lược}
\label{alg:inverse-interpolation-short}
\begin{algorithmic}[1]
\Require $n\geq0$; các $y_i$ đôi một phân biệt; $\boldsymbol x=(x_0,\ldots,x_n)^T$, $\boldsymbol y=(y_0,\ldots,y_n)^T$; $y_\ast\in\R$.
\Ensure $\widehat x=q_n(y_\ast)$, với $q_n(y_i)=x_i$, $\deg q_n\leq n$.
\State Dựng $\boldsymbol d$ bằng thuật toán \ref{alg:newton-table} cho mốc $\boldsymbol y$ và dữ liệu $\boldsymbol x$.
\State Tính $\widehat x$ bằng thuật toán \ref{alg:newton-eval} với mốc $\boldsymbol y$, hệ số $\boldsymbol d$, điểm $y_\ast$.
\State Kết thúc.
\end{algorithmic}
\end{algorithm}
\begin{algorithm}[H]
\caption{Nội suy ngược: chi tiết}
\label{alg:inverse-interpolation}
\begin{algorithmic}[1]
\Require $n\geq0$; các $y_i$ đôi một phân biệt; $\boldsymbol x=(x_0,\ldots,x_n)^T$, $\boldsymbol y=(y_0,\ldots,y_n)^T$; $y_\ast\in\R$.
\Ensure $\widehat x=q_n(y_\ast)$, với $q_n(y_i)=x_i$, $\deg q_n\leq n$.
\For{$i$ từ $0$ tới $n$}
 \State $d_i\gets x_i$.
\EndFor
\Statex
\For{$k$ từ $1$ tới $n$}
 \For{$i$ từ $n$ về $k$}
  \State $d_i\gets(d_i-d_{i-1})/(y_i-y_{i-k})$.
 \EndFor
\EndFor
\Statex
\State $\widehat x\gets d_n$.
\For{$i$ từ $n-1$ về $0$}
 \State $\widehat x\gets d_i+(y_\ast-y_i)\widehat x$.
\EndFor
\Statex
\State Kết thúc.
\end{algorithmic}
\end{algorithm}

\Needspace{20\baselineskip}
\section{Lặp nội suy ngược với sai số được kiểm soát}
\begin{algorithm}[H]
\caption{Lặp nội suy ngược: sơ lược}
\label{alg:inverse-iteration-short}
\begin{algorithmic}[1]
\Require $x_0\in\R$, $h\ne0$, $y_0,y_\ast\in\R$, $d\ne0$; hệ số $\boldsymbol a=(a_0,\ldots,a_r)^T$ của $H(t)$ trong phân rã $p(x_0+ht)=y_0+dt+H(t)$; đoạn đóng $I$ với $\varphi(t)=(y_\ast-y_0-H(t))/d$ thỏa $\varphi(I)\subset I$, $\max_I|H'|\leq q|d|$, $0\leq q<1$; $t_0\in I$, $\varepsilon>0$, $K\geq1$.
\Ensure $x=x_0+ht_k$; $E\geq|x-x_\ast|$ cho nghiệm của đa thức trong $x_0+hI$; trạng thái $s\in\{0,1\}$, $s=1$ khi $E\leq\varepsilon$.
\State $t\gets t_0$; $s\gets0$.
\For{$k$ từ $1$ tới $K$}
 \State Tính $z=H(t)$ bằng thuật toán \ref{alg:horner-eval}.
 \State $w\gets(y_\ast-y_0-z)/d$; $E\gets |h|q|w-t|/(1-q)$; $t\gets w$.
 \If{$E\leq\varepsilon$}
  \State $s\gets1$; $x\gets x_0+ht$; dừng.
 \EndIf
\EndFor
\Statex
\State $x\gets x_0+ht$.
\State Kết thúc.
\end{algorithmic}
\end{algorithm}
\begin{algorithm}[H]
\caption{Lặp nội suy ngược: chi tiết}
\label{alg:inverse-iteration}
\begin{algorithmic}[1]
\Require $x_0\in\R$, $h\ne0$, $y_0,y_\ast\in\R$, $d\ne0$; $H(t)=\sum_{j=0}^{r}a_jt^j$ trong phân rã $p(x_0+ht)=y_0+dt+H(t)$; đoạn đóng $I$ với $\varphi(t)=(y_\ast-y_0-H(t))/d$ thỏa $\varphi(I)\subset I$, $\max_I|H'|\leq q|d|$, $0\leq q<1$; $t_0\in I$, $\varepsilon>0$, $K\geq1$.
\Ensure $x=x_0+ht_k$; $E\geq|x-x_\ast|$ cho nghiệm của đa thức trong $x_0+hI$; $s=1$ khi $E\leq\varepsilon$, $s=0$ nếu chưa đạt chặn này.
\State $t\gets t_0$; $s\gets0$.
\For{$k$ từ $1$ tới $K$}
 \State $z\gets a_r$.
 \For{$j$ từ $r-1$ về $0$}
  \State $z\gets a_j+tz$.
 \EndFor
 \Statex
 \State $w\gets(y_\ast-y_0-z)/d$; $E\gets|h|q|w-t|/(1-q)$.
 \State $t\gets w$.
 \If{$E\leq\varepsilon$}
  \State $s\gets1$; $x\gets x_0+ht$; dừng.
 \EndIf
\EndFor
\Statex
\State $x\gets x_0+ht$.
\State Kết thúc.
\end{algorithmic}
\end{algorithm}

% INSERT ALGORITHM NEWTON-A: after alg:newton-table.
\Needspace{20\baselineskip}
\section{Hệ số Newton tiến và lùi từ cùng bảng}
\begin{algorithm}[H]
\caption{Lấy hai vector hệ số Newton trên các mốc bất kỳ}
\label{alg:newton-two-directions}
\begin{algorithmic}[1]
\Require $n\geq0$; $\boldsymbol x=(x_0,\ldots,x_n)^T$ có mốc đôi một phân biệt; $\boldsymbol y=(y_0,\ldots,y_n)^T$.
\Ensure Các vector $\boldsymbol d^+,\boldsymbol d^-$ của \eqref{eq:newton-arbitrary-forward}, \eqref{eq:newton-arbitrary-backward}.
\State Áp dụng Thuật toán \ref{alg:newton-table} với $\boldsymbol x,\boldsymbol y$; thu được các $D_i^{(k)}$.
\State $d_k^+\gets D_0^{(k)}$; $d_k^-\gets D_{n-k}^{(k)}$, $0\leq k\leq n$.
\State $\boldsymbol d^+\gets(d_0^+,\ldots,d_n^+)^T$; $\boldsymbol d^-\gets(d_0^-,\ldots,d_n^-)^T$.
\State Kết thúc.
\end{algorithmic}
\end{algorithm}

% INSERT ALGORITHM NEWTON-B: after divided-difference and finite-difference table modules.
\Needspace{20\baselineskip}
\section{Cập nhật một đường biên của bảng}
\begin{algorithm}[H]
\caption{Thêm một mốc vào bảng tỷ sai phân}
\label{alg:newton-dd-update}
\begin{algorithmic}[1]
\Require Các mốc $x_0,\ldots,x_n$ và bảng $D_i^{(k)}$ đã có; $x_{n+1}$ khác mọi mốc cũ, $y_{n+1}\in\R$.
\Ensure Bảng mới chứa $D_{n+1-k}^{(k)}$, $0\leq k\leq n+1$, và giữ các phần tử cũ, theo \eqref{eq:newton-dd-update}.
\State $E\gets y_{n+1}$; $D_{n+1}^{(0)}\gets E$.
\For{$k$ từ $1$ tới $n+1$}
  \State $E\gets(E-D_{n+1-k}^{(k-1)})/(x_{n+1}-x_{n+1-k})$.
  \State $D_{n+1-k}^{(k)}\gets E$.
\EndFor
\Statex
\State $d_{n+1}^+\gets E$; $d_k^-\gets D_{n+1-k}^{(k)}$, $0\leq k\leq n+1$.
\State Kết thúc.
\end{algorithmic}
\end{algorithm}

\begin{algorithm}[H]
\caption{Thêm một mốc vào bảng sai phân cách đều}
\label{alg:newton-difference-update}
\begin{algorithmic}[1]
\Require Bảng $\Delta^ky_i$ trên $x_i=x_0+ih$, $0\leq i\leq n$, $h\ne0$; giá trị mới $y_{n+1}$ tại $x_0+(n+1)h$.
\Ensure Đường biên mới $\Delta^ky_{n+1-k}=\nabla^ky_{n+1}$ theo \eqref{eq:newton-difference-update}; giữ mọi sai phân cũ.
\State $F\gets y_{n+1}$; $\Delta^0y_{n+1}\gets F$.
\For{$k$ từ $1$ tới $n+1$}
  \State $F\gets F-\Delta^{k-1}y_{n+1-k}$.
  \State $\Delta^ky_{n+1-k}\gets F$; $\nabla^ky_{n+1}\gets F$.
\EndFor
\Statex
\State Kết thúc.
\end{algorithmic}
\end{algorithm}

% INSERT ALGORITHM NEWTON-C: after Newton value module.
\Needspace{20\baselineskip}
\section{Bổ sung một giá trị khi biết chặn bậc}
\begin{algorithm}[H]
\caption{Giá trị còn thiếu với giả thiết bậc không quá $n-1$}
\label{alg:newton-missing-value}
\begin{algorithmic}[1]
\Require $n\geq1$, $x_0,\ldots,x_n$ đôi một phân biệt; chỉ số $j$ thiếu giá trị; các $y_i$ với $i\ne j$; giả thiết dữ liệu thuộc một đa thức của $\mathcal P_{n-1}$.
\Ensure $y_j$ duy nhất theo Mệnh đề \ref{md:newton-missing-data}.
\State $(i_0,\ldots,i_{n-1})\gets(0,\ldots,j-1,j+1,\ldots,n)$.
\State $\boldsymbol z\gets(x_{i_0},\ldots,x_{i_{n-1}})^T$; $\boldsymbol v\gets(y_{i_0},\ldots,y_{i_{n-1}})^T$.
\State Áp dụng Thuật toán \ref{alg:newton-table} với $n-1,\boldsymbol z,\boldsymbol v$; thu được $\boldsymbol d$.
\State $\boldsymbol\xi\gets(z_0,\ldots,z_{n-2})^T$ (vector rỗng khi $n=1$).
\State Áp dụng Thuật toán \ref{alg:newton-eval} với $\boldsymbol\xi,\boldsymbol d,c=x_j$; thu được $y_j=P(x_j)$.
\State Kết thúc.
\end{algorithmic}
\end{algorithm}

% INSERT ALGORITHM NEWTON-D: after alg:newton-eval.
\Needspace{20\baselineskip}
\section{Đạo hàm mọi cấp trong cơ sở Newton}
\begin{algorithm}[H]
\caption{Horner với vector đạo hàm đã chuẩn hóa}
\label{alg:newton-derivative-jet}
\begin{algorithmic}[1]
\Require $n,s\geq0$; $\boldsymbol x=(x_0,\ldots,x_{n-1})^T$; $\boldsymbol d=(d_0,\ldots,d_n)^T$ của \eqref{eq:newton-form}; $c\in\R$.
\Ensure $D_r=p_n^{(r)}(c)$, $0\leq r\leq s$, theo \eqref{eq:newton-derivative-jet}.
\State $(J_0,\ldots,J_s)\gets(d_n,0,\ldots,0)$.
\For{$i$ từ $n-1$ về $0$}
  \For{$r$ từ $s$ về $1$}
    \State $J_r\gets(c-x_i)J_r+J_{r-1}$.
  \EndFor
  \Statex
  \State $J_0\gets d_i+(c-x_i)J_0$.
\EndFor
\Statex
\State $D_r\gets r!J_r$, $0\leq r\leq s$.
\State Kết thúc.
\end{algorithmic}
\end{algorithm}

% INSERT ALGORITHM NEWTON-E: before central formula modules.
\Needspace{20\baselineskip}
\section{Trích xuất một cửa sổ mốc gần điểm cần tính}
\begin{algorithm}[H]
\caption{Lấy $k$ mốc gần nhất trên lưới đều}
\label{alg:newton-nearest-window}
\begin{algorithmic}[1]
\Require $M\geq1$, $1\leq k\leq M$, $x_i=a+ih$, $0\leq i<M$, $h>0$; $\boldsymbol y=(y_0,\ldots,y_{M-1})^T$; $c\in\R$.
\Ensure Cửa sổ $(z_j,v_j)=(x_{s+j},y_{s+j})$, $0\leq j<k$, theo Mệnh đề \ref{md:newton-nearest-window}; khi bằng khoảng cách, ưu tiên chỉ số nhỏ hơn.
\State $q\gets(c-a)/h$; $s\gets\min\{M-k,\max\{0,\lceil q-k/2\rceil\}\}$.
\For{$j$ từ $0$ tới $k-1$}
  \State $z_j\gets a+(s+j)h$; $v_j\gets y_{s+j}$.
\EndFor
\Statex
\State $\boldsymbol z\gets(z_0,\ldots,z_{k-1})^T$; $\boldsymbol v\gets(v_0,\ldots,v_{k-1})^T$.
\State Kết thúc.
\end{algorithmic}
\end{algorithm}

% INSERT ALGORITHM INVERSE-A: before inverse interpolation.
\Needspace{20\baselineskip}
\section{Phát hiện khoảng đổi dấu của bảng}
\begin{algorithm}[H]
\caption{Nghiệm tại mốc, khoảng đổi dấu và đoạn hằng -- sơ lược}
\label{alg:inv-sign-scan}
\begin{algorithmic}[1]
\Require $n\geq0$, $x_0<\cdots<x_n$, $\boldsymbol y=(y_0,\ldots,y_n)^T$, giá trị đích $y_\ast$.
\Ensure $Z,C,F$ theo Mệnh đề \ref{md:inv-sampled-branches}; $C$ ghi mọi khoảng đổi dấu liền nhau của bảng.
\State $Z\gets\{i:0\leq i\leq n,\ y_i=y_\ast\}$.
\State $C\gets\{i:0\leq i<n,\ (y_i-y_\ast)(y_{i+1}-y_\ast)<0\}$.
\State $F\gets\{i:0\leq i<n,\ y_i=y_{i+1}=y_\ast\}$.
\State Kết thúc.
\end{algorithmic}
\end{algorithm}

\begin{algorithm}[H]
\caption{Nghiệm tại mốc, khoảng đổi dấu và đoạn hằng -- chi tiết}
\label{alg:inv-sign-scan-detailed}
\begin{algorithmic}[1]
\Require $n\geq0$, $x_0<\cdots<x_n$, $\boldsymbol y=(y_0,\ldots,y_n)^T$, giá trị đích $y_\ast$.
\Ensure $Z,C,F$ theo Mệnh đề \ref{md:inv-sampled-branches}.
\State $Z\gets\varnothing$; $C\gets\varnothing$; $F\gets\varnothing$.
\For{$i$ từ $0$ tới $n-1$}
  \If{$y_i=y_\ast$}
    \State $Z\gets Z\cup\{i\}$.
  \EndIf
  \Statex
  \If{$(y_i-y_\ast)(y_{i+1}-y_\ast)<0$}
    \State $C\gets C\cup\{i\}$.
  \EndIf
  \Statex
  \If{$y_i=y_{i+1}=y_\ast$}
    \State $F\gets F\cup\{i\}$.
  \EndIf
  \Statex
\EndFor
\Statex
\If{$y_n=y_\ast$}
  \State $Z\gets Z\cup\{n\}$.
\EndIf
\Statex
\State Kết thúc.
\end{algorithmic}
\end{algorithm}

\Needspace{20\baselineskip}
\section{Các khối đơn điệu cực đại của bảng}
\begin{algorithm}[H]
\caption{Tách bảng theo dấu của sai phân cấp một}
\label{alg:inv-monotone-blocks}
\begin{algorithmic}[1]
\Require $n\geq1$; $\boldsymbol y=(y_0,\ldots,y_n)^T$.
\Ensure $\mathcal B$ chứa các bộ $(a,b,\delta)$ mô tả khối cực đại $[a,b]$ với dấu sai phân $\delta\in\{-1,0,1\}$; theo Mệnh đề \ref{md:inv-sampled-monotonicity}.
\State $\mathcal B\gets\varnothing$; $a\gets0$; $\delta\gets\operatorname{sgn}(y_1-y_0)$.
\For{$i$ từ $1$ tới $n-1$}
  \State $\sigma\gets\operatorname{sgn}(y_{i+1}-y_i)$.
  \If{$\sigma\ne\delta$}
    \State $\mathcal B\gets\mathcal B\cup\{(a,i,\delta)\}$; $a\gets i$; $\delta\gets\sigma$.
  \EndIf
  \Statex
\EndFor
\Statex
\State $\mathcal B\gets\mathcal B\cup\{(a,n,\delta)\}$.
\State Kết thúc.
\end{algorithmic}
\end{algorithm}

% INSERT ALGORITHM INVERSE-B: after polynomial coefficient conversion and inverse scan.
\Needspace{20\baselineskip}
\section{Dựng phần dư cho hai công thức lặp nội suy ngược}
\begin{algorithm}[H]
\caption{Phần dư Newton tiến hoặc lùi và chứng nhận co -- sơ lược}
\label{alg:inv-remainder-certificate}
\begin{algorithmic}[1]
\Require Bảng sai phân cách đều, $h>0$, $k$ với $A=y_{k+1}-y_k\ne0$; $\alpha=(y_\ast-y_k)/A\in[0,1]$; dấu $\sigma\in\{+1,-1\}$; bậc $d\geq1$ với mọi mốc cần dùng thuộc bảng.
\Ensure Hệ số $\boldsymbol a$ của $H_\sigma$, $q$, trạng thái $v\in\{0,1\}$; $v=1$ chứng nhận \eqref{eq:inv-coefficient-certificate}.
\State $\xi_j\gets\sigma j$, $0\leq j<d$; $\eta_0\gets0$; $\eta_1\gets0$.
\State $\eta_r\gets\begin{cases}\Delta^ry_k/r!,&\sigma=+1,\\\Delta^ry_{k+1-r}/r!,&\sigma=-1,\end{cases}$ $2\leq r\leq d$.
\State Áp dụng Thuật toán \ref{alg:newton-coefficients} với $\boldsymbol\xi,\boldsymbol\eta$; thu được $\boldsymbol a=(a_0,\ldots,a_d)^T$.
\State $B\gets\sum_{r=0}^d|a_r|$; $q\gets|A|^{-1}\sum_{r=1}^dr|a_r|$; $\delta\gets\min(\alpha,1-\alpha)$.
\State $v\gets\mathbf1_{\{B\leq|A|\delta,\ q<1\}}$.
\State Kết thúc.
\end{algorithmic}
\end{algorithm}

\begin{algorithm}[H]
\caption{Phần dư Newton tiến hoặc lùi và chứng nhận co -- chi tiết}
\label{alg:inv-remainder-certificate-detailed}
\small
\begin{algorithmic}[1]
\Require Bảng sai phân cách đều, $h>0$, $k$ với $A=y_{k+1}-y_k\ne0$; $\alpha=(y_\ast-y_k)/A\in[0,1]$; dấu $\sigma\in\{+1,-1\}$; bậc $d\geq1$ với mọi mốc cần dùng thuộc bảng.
\Ensure Hệ số $\boldsymbol a$ của $H_\sigma$, $q$, trạng thái $v\in\{0,1\}$; $v=1$ chứng nhận \eqref{eq:inv-coefficient-certificate}.
\State $\eta_0\gets0$; $\eta_1\gets0$; $F\gets1$.
\For{$j$ từ $0$ tới $d-1$}
  \State $\xi_j\gets\sigma j$.
\EndFor
\Statex
\For{$r$ từ $2$ tới $d$}
  \State $F\gets rF$.
  \If{$\sigma=+1$}
    \State $\eta_r\gets\Delta^ry_k/F$.
  \Else
    \State $\eta_r\gets\Delta^ry_{k+1-r}/F$.
  \EndIf
  \Statex
\EndFor
\Statex
\State $a_0\gets\eta_d$; $m\gets0$.
\For{$i$ từ $d-1$ về $0$}
  \State $a_{m+1}\gets a_m$.
  \For{$j$ từ $m$ về $1$}
    \State $a_j\gets a_{j-1}-\xi_i a_j$.
  \EndFor
  \Statex
  \State $a_0\gets\eta_i-\xi_i a_0$; $m\gets m+1$.
\EndFor
\Statex
\State $B\gets0$; $Q\gets0$.
\For{$r$ từ $0$ tới $d$}
  \State $B\gets B+|a_r|$; $Q\gets Q+r|a_r|$.
\EndFor
\Statex
\State $\boldsymbol a\gets(a_0,\ldots,a_d)^T$; $q\gets Q/|A|$; $\delta\gets\min(\alpha,1-\alpha)$.
\State $v\gets\mathbf1_{\{B\leq|A|\delta,\ q<1\}}$.
\State Kết thúc.
\end{algorithmic}
\end{algorithm}

\Needspace{20\baselineskip}
\section{Hai hướng lặp trên một khoảng đổi dấu}
\begin{algorithm}[H]
\caption{Nội suy ngược theo hướng tiến hoặc lùi -- sơ lược}
\label{alg:inv-two-directions}
\begin{algorithmic}[1]
\Require Bảng sai phân trên $x_i=x_0+ih$, $h>0$; $k$ với $A=y_{k+1}-y_k\ne0$, $\alpha=(y_\ast-y_k)/A\in[0,1]$; $\sigma\in\{+1,-1\}$, bậc $d\geq1$ hợp lệ; $\varepsilon>0$, $K\geq1$.
\Ensure $v=0$ nếu chứng nhận hệ số chưa đạt; nếu $v=1$, $(x,E,s)$ với $|x-x_*^\sigma|\leq E$ và $s=1$ khi $E\leq\varepsilon$.
\State Áp dụng Thuật toán \ref{alg:inv-remainder-certificate} với $k,\alpha,\sigma,d$; thu được $\boldsymbol a,q,v$.
\If{$v=1$}
  \State $(x_c,y_c,z_0,I)\gets\begin{cases}(x_k,y_k,\alpha,[0,1]),&\sigma=+1,\\(x_{k+1},y_{k+1},\alpha-1,[-1,0]),&\sigma=-1.\end{cases}$
  \State Áp dụng Thuật toán \ref{alg:inverse-iteration-short} với $x_c,h,y_c,y_\ast,A,\boldsymbol a,I,q,z_0,\varepsilon,K$; thu được $(x,E,s)$.
\EndIf
\Statex
\State Kết thúc.
\end{algorithmic}
\end{algorithm}

\begin{algorithm}[H]
\caption{Nội suy ngược theo hướng tiến hoặc lùi -- chi tiết}
\label{alg:inv-two-directions-detailed}
\small
\begin{algorithmic}[1]
\Require Bảng sai phân trên $x_i=x_0+ih$, $h>0$; $k$ với $A=y_{k+1}-y_k\ne0$, $\alpha=(y_\ast-y_k)/A\in[0,1]$; $\sigma\in\{+1,-1\}$, bậc $d\geq1$ hợp lệ; $\varepsilon>0$, $K\geq1$.
\Ensure $v=0$ nếu chứng nhận hệ số chưa đạt; nếu $v=1$, $(x,E,s)$ với $|x-x_*^\sigma|\leq E$ và $s=1$ khi $E\leq\varepsilon$.
\State $\eta_0\gets0$; $\eta_1\gets0$; $F\gets1$; $s\gets0$.
\For{$r$ từ $2$ tới $d$}
  \State $F\gets rF$; $\eta_r\gets\begin{cases}\Delta^ry_k/F,&\sigma=+1,\\\Delta^ry_{k+1-r}/F,&\sigma=-1.\end{cases}$
\EndFor
\Statex
\State $a_0\gets\eta_d$; $m\gets0$.
\For{$i$ từ $d-1$ về $0$}
  \State $\xi\gets\sigma i$; $a_{m+1}\gets a_m$.
  \For{$j$ từ $m$ về $1$}
    \State $a_j\gets a_{j-1}-\xi a_j$.
  \EndFor
  \Statex
  \State $a_0\gets\eta_i-\xi a_0$; $m\gets m+1$.
\EndFor
\Statex
\State $B\gets0$; $Q\gets0$.
\For{$r$ từ $0$ tới $d$}
  \State $B\gets B+|a_r|$; $Q\gets Q+r|a_r|$.
\EndFor
\Statex
\State $q\gets Q/|A|$; $\delta\gets\min(\alpha,1-\alpha)$; $v\gets\mathbf1_{\{B\leq|A|\delta,\ q<1\}}$.
\If{$v=1$}
  \State $(x_c,y_c,t)\gets\begin{cases}(x_k,y_k,\alpha),&\sigma=+1,\\(x_{k+1},y_{k+1},\alpha-1),&\sigma=-1.\end{cases}$
  \For{$\ell$ từ $1$ tới $K$}
    \State $z\gets a_d$.
    \For{$j$ từ $d-1$ về $0$}
      \State $z\gets a_j+tz$.
    \EndFor
    \Statex
    \State $w\gets(y_\ast-y_c-z)/A$; $E\gets hq|w-t|/(1-q)$; $t\gets w$.
    \If{$E\leq\varepsilon$}
      \State $s\gets1$; Dừng lặp.
    \EndIf
    \Statex
  \EndFor
  \Statex
  \State $x\gets x_c+ht$.
\EndIf
\Statex
\State Kết thúc.
\end{algorithmic}
\end{algorithm}

% END EXTENSION NEWTON ALGORITHMS

% BEGIN EXTENSION SPLINE LS ALGORITHMS
\Needspace{20\baselineskip}
\section{Giải hệ tam đường chéo đối xứng}
\begin{algorithm}[H]
\caption{Giải hệ tam đường chéo bằng khử và thế ngược}
\label{alg:spl-thomas}
\begin{algorithmic}[1]
\Require $q\geq1$; $d_i$ ($0\leq i<q$), $e_i$ ($0\leq i<q-1$), $\mathbf r\in\R^q$; ma trận $T_{ii}=d_i$, $T_{i,i+1}=T_{i+1,i}=e_i$ đối xứng xác định dương.
\Ensure $\mathbf z=(z_0,\ldots,z_{q-1})^T$, $T\mathbf z=\mathbf r$.
\State $D_0\gets d_0$; $\rho_0\gets r_0$.
\For{$i$ từ $1$ tới $q-1$}
  \State $D_i\gets d_i-e_{i-1}^2/D_{i-1}$.
  \State $\rho_i\gets r_i-e_{i-1}\rho_{i-1}/D_{i-1}$.
\EndFor
\Statex
\State $z_{q-1}\gets\rho_{q-1}/D_{q-1}$.
\For{$i$ từ $q-2$ về $0$}
  \State $z_i\gets(\rho_i-e_iz_{i+1})/D_i$.
\EndFor
\Statex
\State Kết thúc.
\end{algorithmic}
\end{algorithm}

\Needspace{20\baselineskip}
\section{Giá trị nội suy tuyến tính từng đoạn}
\begin{algorithm}[H]
\caption{Nội suy tuyến tính tại một điểm}
\label{alg:spl-linear-evaluate}
\begin{algorithmic}[1]
\Require $N\geq1$; $x_0<\cdots<x_N$; $y_0,\ldots,y_N$; $c\in[x_0,x_N]$.
\Ensure $v=L(c)$ theo \eqref{eq:spl-linear}.
\State $k\gets0$.
\For{$i$ từ $0$ tới $N-1$}
  \If{$x_i\leq c$}
    \State $k\gets i$.
  \EndIf
\EndFor
\Statex
\State $t\gets(c-x_k)/(x_{k+1}-x_k)$.
\State $v\gets y_k+t(y_{k+1}-y_k)$.
\State Kết thúc.
\end{algorithmic}
\end{algorithm}

\Needspace{20\baselineskip}
\section{Moment của spline tự nhiên}
\begin{algorithm}[H]
\caption{Dựng moment của spline tự nhiên -- sơ lược}
\label{alg:spl-natural}
\begin{algorithmic}[1]
\Require $N\geq1$; $x_0<\cdots<x_N$; $\mathbf y=(y_0,\ldots,y_N)^T$.
\Ensure $\mathbf M=(M_0,\ldots,M_N)^T$ thỏa \eqref{eq:spl-interior-system} và $M_0=M_N=0$.
\State $h_i\gets x_{i+1}-x_i$, $\delta_i\gets(y_{i+1}-y_i)/h_i$ với $0\leq i<N$; $M_0\gets0$; $M_N\gets0$.
\If{$N\geq2$}
  \State $d_k\gets2(h_k+h_{k+1})$, $r_k\gets6(\delta_{k+1}-\delta_k)$ với $0\leq k<N-1$.
  \State $e_k\gets h_{k+1}$ với $0\leq k<N-2$.
  \State Áp dụng Thuật toán \ref{alg:spl-thomas} cho $q=N-1$, $\boldsymbol d,\boldsymbol e,\mathbf r$; nhận $\mathbf z$.
  \State $M_i\gets z_{i-1}$ với $1\leq i<N$.
\EndIf
\Statex
\State Kết thúc.
\end{algorithmic}
\end{algorithm}

\begin{algorithm}[H]
\caption{Dựng moment của spline tự nhiên -- chi tiết}
\label{alg:spl-natural-detailed}
\begin{algorithmic}[1]
\Require $N\geq1$; $x_0<\cdots<x_N$; $\mathbf y=(y_0,\ldots,y_N)^T$.
\Ensure $\mathbf M=(M_0,\ldots,M_N)^T$ thỏa \eqref{eq:spl-interior-system} và $M_0=M_N=0$.
\For{$i$ từ $0$ tới $N-1$}
  \State $h_i\gets x_{i+1}-x_i$; $\delta_i\gets(y_{i+1}-y_i)/h_i$.
\EndFor
\Statex
\State $M_0\gets0$; $M_N\gets0$.
\If{$N\geq2$}
  \State $D_1\gets2(h_0+h_1)$; $\rho_1\gets6(\delta_1-\delta_0)$.
  \For{$i$ từ $2$ tới $N-1$}
    \State $D_i\gets2(h_{i-1}+h_i)-h_{i-1}^2/D_{i-1}$.
    \State $\rho_i\gets6(\delta_i-\delta_{i-1})-h_{i-1}\rho_{i-1}/D_{i-1}$.
  \EndFor
  \Statex
  \State $M_{N-1}\gets\rho_{N-1}/D_{N-1}$.
  \For{$i$ từ $N-2$ về $1$}
    \State $M_i\gets(\rho_i-h_iM_{i+1})/D_i$.
  \EndFor
\EndIf
\Statex
\State Kết thúc.
\end{algorithmic}
\end{algorithm}

\Needspace{20\baselineskip}
\section{Moment của spline ấn định đạo hàm biên}
\begin{algorithm}[H]
\caption{Dựng moment với đạo hàm biên -- sơ lược}
\label{alg:spl-clamped}
\begin{algorithmic}[1]
\Require $N\geq1$; $x_0<\cdots<x_N$; $y_0,\ldots,y_N$; $\alpha,\beta\in\R$.
\Ensure $\mathbf M$ thỏa \eqref{eq:spl-interior-system}--\eqref{eq:spl-clamped-system}.
\State $h_i\gets x_{i+1}-x_i$, $\delta_i\gets(y_{i+1}-y_i)/h_i$ với $0\leq i<N$.
\State $d_0\gets2h_0$; $d_N\gets2h_{N-1}$; $r_0\gets6(\delta_0-\alpha)$; $r_N\gets6(\beta-\delta_{N-1})$.
\State $d_i\gets2(h_{i-1}+h_i)$, $r_i\gets6(\delta_i-\delta_{i-1})$ với $1\leq i<N$; $e_i\gets h_i$ với $0\leq i<N$.
\State Áp dụng Thuật toán \ref{alg:spl-thomas} cho $q=N+1$, $\boldsymbol d,\boldsymbol e,\mathbf r$; nhận $\mathbf M$.
\State Kết thúc.
\end{algorithmic}
\end{algorithm}

\begin{algorithm}[H]
\caption{Dựng moment với đạo hàm biên -- chi tiết}
\label{alg:spl-clamped-detailed}
\begin{algorithmic}[1]
\Require $N\geq1$; $x_0<\cdots<x_N$; $y_0,\ldots,y_N$; $\alpha,\beta\in\R$.
\Ensure $\mathbf M$ thỏa \eqref{eq:spl-interior-system}--\eqref{eq:spl-clamped-system}.
\For{$i$ từ $0$ tới $N-1$}
  \State $h_i\gets x_{i+1}-x_i$; $\delta_i\gets(y_{i+1}-y_i)/h_i$.
\EndFor
\Statex
\State $D_0\gets2h_0$; $\rho_0\gets6(\delta_0-\alpha)$.
\For{$i$ từ $1$ tới $N-1$}
  \State $D_i\gets2(h_{i-1}+h_i)-h_{i-1}^2/D_{i-1}$.
  \State $\rho_i\gets6(\delta_i-\delta_{i-1})-h_{i-1}\rho_{i-1}/D_{i-1}$.
\EndFor
\Statex
\State $D_N\gets2h_{N-1}-h_{N-1}^2/D_{N-1}$.
\State $\rho_N\gets6(\beta-\delta_{N-1})-h_{N-1}\rho_{N-1}/D_{N-1}$.
\State $M_N\gets\rho_N/D_N$.
\For{$i$ từ $N-1$ về $0$}
  \State $M_i\gets(\rho_i-h_iM_{i+1})/D_i$.
\EndFor
\Statex
\State Kết thúc.
\end{algorithmic}
\end{algorithm}

\Needspace{20\baselineskip}
\section{Hệ số từng đoạn của spline bậc ba}
\begin{algorithm}[H]
\caption{Dựng hệ số Horner từng đoạn}
\label{alg:spl-coefficients}
\begin{algorithmic}[1]
\Require $N\geq1$; $x_0<\cdots<x_N$; $y_0,\ldots,y_N$; moment $M_0,\ldots,M_N$ của spline đã xác định.
\Ensure $a_i,b_i,c_i,d_i$ trong \eqref{eq:spl-horner}, $0\leq i<N$.
\For{$i$ từ $0$ tới $N-1$}
  \State $h_i\gets x_{i+1}-x_i$; $\delta_i\gets(y_{i+1}-y_i)/h_i$.
  \State $a_i\gets y_i$; $b_i\gets\delta_i-h_i(2M_i+M_{i+1})/6$.
  \State $c_i\gets M_i/2$; $d_i\gets(M_{i+1}-M_i)/(6h_i)$.
\EndFor
\Statex
\State Kết thúc.
\end{algorithmic}
\end{algorithm}

\Needspace{20\baselineskip}
\section{Giá trị spline bậc ba}
\begin{algorithm}[H]
\caption{Chọn đoạn và tính giá trị bằng Horner}
\label{alg:spl-evaluate}
\begin{algorithmic}[1]
\Require $N\geq1$; $x_0<\cdots<x_N$; các hệ số $a_i,b_i,c_i,d_i$ trong \eqref{eq:spl-horner}; $\xi\in[x_0,x_N]$.
\Ensure $v=S(\xi)$.
\State $k\gets0$.
\For{$i$ từ $0$ tới $N-1$}
  \If{$x_i\leq\xi$}
    \State $k\gets i$.
  \EndIf
\EndFor
\Statex
\State $t\gets\xi-x_k$; $v\gets d_k$.
\State $v\gets c_k+tv$; $v\gets b_k+tv$; $v\gets a_k+tv$.
\State Kết thúc.
\end{algorithmic}
\end{algorithm}

\Needspace{20\baselineskip}
\section{Phân tích Cholesky}
\begin{algorithm}[H]
\caption{Dựng nhân tử tam giác dưới Cholesky}
\label{alg:ls-cholesky-factor}
\begin{algorithmic}[1]
\Require $q\geq1$; $G\in\R^{q\times q}$ đối xứng xác định dương.
\Ensure $L$ tam giác dưới, $L_{kk}>0$, $G=LL^T$.
\State $L_{ij}\gets0$ với $0\leq i<j<q$.
\For{$k$ từ $0$ tới $q-1$}
  \State $L_{kk}\gets\sqrt{G_{kk}-\sum_{j=0}^{k-1}L_{kj}^2}$.
  \For{$i$ từ $k+1$ tới $q-1$}
    \State $L_{ik}\gets(G_{ik}-\sum_{j=0}^{k-1}L_{ij}L_{kj})/L_{kk}$.
  \EndFor
\EndFor
\Statex
\State Kết thúc.
\end{algorithmic}
\end{algorithm}

\Needspace{20\baselineskip}
\section{Giải hệ đối xứng xác định dương}
\begin{algorithm}[H]
\caption{Giải hệ bằng Cholesky -- sơ lược}
\label{alg:ls-cholesky-solve}
\begin{algorithmic}[1]
\Require $q\geq1$; $G\in\R^{q\times q}$ đối xứng xác định dương; $\mathbf b\in\R^q$.
\Ensure $\mathbf a\in\R^q$, $G\mathbf a=\mathbf b$.
\State Áp dụng Thuật toán \ref{alg:ls-cholesky-factor} cho $G$; nhận $L$.
\For{$i$ từ $0$ tới $q-1$}
  \State $z_i\gets(b_i-\sum_{j=0}^{i-1}L_{ij}z_j)/L_{ii}$.
\EndFor
\Statex
\For{$i$ từ $q-1$ về $0$}
  \State $a_i\gets(z_i-\sum_{j=i+1}^{q-1}L_{ji}a_j)/L_{ii}$.
\EndFor
\Statex
\State Kết thúc.
\end{algorithmic}
\end{algorithm}

\begin{algorithm}[H]
\caption{Giải hệ bằng Cholesky -- chi tiết}
\label{alg:ls-cholesky-solve-detailed}
\begin{algorithmic}[1]
\Require $q\geq1$; $G\in\R^{q\times q}$ đối xứng xác định dương; $\mathbf b\in\R^q$.
\Ensure $\mathbf a\in\R^q$, $G\mathbf a=\mathbf b$.
\State $L_{ij}\gets0$ với $0\leq i<j<q$.
\For{$k$ từ $0$ tới $q-1$}
  \State $L_{kk}\gets\sqrt{G_{kk}-\sum_{j=0}^{k-1}L_{kj}^2}$.
  \For{$i$ từ $k+1$ tới $q-1$}
    \State $L_{ik}\gets(G_{ik}-\sum_{j=0}^{k-1}L_{ij}L_{kj})/L_{kk}$.
  \EndFor
\EndFor
\Statex
\For{$i$ từ $0$ tới $q-1$}
  \State $z_i\gets(b_i-\sum_{j=0}^{i-1}L_{ij}z_j)/L_{ii}$.
\EndFor
\Statex
\For{$i$ từ $q-1$ về $0$}
  \State $a_i\gets(z_i-\sum_{j=i+1}^{q-1}L_{ji}a_j)/L_{ii}$.
\EndFor
\Statex
\State Kết thúc.
\end{algorithmic}
\end{algorithm}

\Needspace{20\baselineskip}
\section{Bình phương tối thiểu rời rạc có trọng số}
\begin{algorithm}[H]
\caption{Giải phương trình chuẩn -- sơ lược}
\label{alg:ls-normal}
\begin{algorithmic}[1]
\Require $m,q\geq1$; $A\in\R^{m\times q}$, $\operatorname{rank}A=q$; $\mathbf y\in\R^m$; $w_i>0$ với $0\leq i<m$.
\Ensure $\mathbf a$ cực tiểu $\sum_iw_i((A\mathbf a)_i-y_i)^2$.
\State $G_{jk}\gets\sum_{i=0}^{m-1}w_iA_{ij}A_{ik}$ với $0\leq j,k<q$.
\State $b_j\gets\sum_{i=0}^{m-1}w_iA_{ij}y_i$ với $0\leq j<q$.
\State Áp dụng Thuật toán \ref{alg:ls-cholesky-solve} cho $G,\mathbf b$; nhận $\mathbf a$.
\State Kết thúc.
\end{algorithmic}
\end{algorithm}

\begin{algorithm}[H]
\caption{Giải phương trình chuẩn -- chi tiết}
\label{alg:ls-normal-detailed}
\begin{algorithmic}[1]
\Require $m,q\geq1$; $A\in\R^{m\times q}$, $\operatorname{rank}A=q$; $\mathbf y\in\R^m$; $w_i>0$ với $0\leq i<m$.
\Ensure $\mathbf a$ cực tiểu $\sum_iw_i((A\mathbf a)_i-y_i)^2$.
\For{$j$ từ $0$ tới $q-1$}
  \State $b_j\gets\sum_{i=0}^{m-1}w_iA_{ij}y_i$.
  \For{$k$ từ $0$ tới $j$}
    \State $G_{jk}\gets\sum_{i=0}^{m-1}w_iA_{ij}A_{ik}$; $G_{kj}\gets G_{jk}$.
  \EndFor
\EndFor
\Statex
\State $L_{ij}\gets0$ với $0\leq i<j<q$.
\For{$k$ từ $0$ tới $q-1$}
  \State $L_{kk}\gets\sqrt{G_{kk}-\sum_{j=0}^{k-1}L_{kj}^2}$.
  \For{$i$ từ $k+1$ tới $q-1$}
    \State $L_{ik}\gets(G_{ik}-\sum_{j=0}^{k-1}L_{ij}L_{kj})/L_{kk}$.
  \EndFor
\EndFor
\Statex
\For{$i$ từ $0$ tới $q-1$}
  \State $z_i\gets(b_i-\sum_{j=0}^{i-1}L_{ij}z_j)/L_{ii}$.
\EndFor
\Statex
\For{$i$ từ $q-1$ về $0$}
  \State $a_i\gets(z_i-\sum_{j=i+1}^{q-1}L_{ji}a_j)/L_{ii}$.
\EndFor
\Statex
\State Kết thúc.
\end{algorithmic}
\end{algorithm}

\Needspace{20\baselineskip}
\section{Hiệu chuẩn bằng mô hình affine}
\begin{algorithm}[H]
\caption{Hồi quy affine bằng tọa độ quy tâm}
\label{alg:ls-affine}
\begin{algorithmic}[1]
\Require $m\geq2$; $x_i,y_i\in\R$, $w_i>0$ với $0\leq i<m$; ít nhất hai $x_i$ khác nhau.
\Ensure $a,b$ cực tiểu $\sum_iw_i(a+bx_i-y_i)^2$.
\State $s\gets\sum_{i=0}^{m-1}w_i$; $\bar x\gets\sum_iw_ix_i/s$; $\bar y\gets\sum_iw_iy_i/s$.
\State $U\gets\sum_{i=0}^{m-1}w_i(x_i-\bar x)^2$.
\State $V\gets\sum_{i=0}^{m-1}w_i(x_i-\bar x)(y_i-\bar y)$.
\State $b\gets V/U$; $a\gets\bar y-b\bar x$.
\State Kết thúc.
\end{algorithmic}
\end{algorithm}

\Needspace{20\baselineskip}
\section{Bình phương tối thiểu theo đa thức trực giao}
\begin{algorithm}[H]
\caption{Dựng cơ sở trực giao rời rạc và các hệ số chiếu}
\label{alg:ls-orthogonal}
\begin{algorithmic}[1]
\Require $m\geq1$, $0\leq r<m$; các $x_i$ đôi một phân biệt; $y_i\in\R$, $w_i>0$ với $0\leq i<m$.
\Ensure $\gamma_0,\ldots,\gamma_r$, $\alpha_k,\beta_k$ ($0\leq k<r$) xác định nghiệm \eqref{eq:ls-orthogonal-projection} qua truy hồi \eqref{eq:ls-orthogonal-recurrence}.
\State $P_{-1,i}\gets0$; $P_{0,i}\gets1$ với $0\leq i<m$.
\For{$k$ từ $0$ tới $r$}
  \State $\nu_k\gets\sum_{i=0}^{m-1}w_iP_{k,i}^2$; $\gamma_k\gets\sum_iw_iy_iP_{k,i}/\nu_k$.
  \If{$k<r$}
    \State $\alpha_k\gets\sum_{i=0}^{m-1}w_ix_iP_{k,i}^2/\nu_k$; $\beta_k\gets\begin{cases}0,&k=0,\\\nu_k/\nu_{k-1},&k>0.\end{cases}$
    \For{$i$ từ $0$ tới $m-1$}
      \State $P_{k+1,i}\gets(x_i-\alpha_k)P_{k,i}-\beta_kP_{k-1,i}$.
    \EndFor
  \EndIf
\EndFor
\Statex
\State Kết thúc.
\end{algorithmic}
\end{algorithm}

% BEGIN SPL-SUPPLEMENT-ALGORITHMS
\Needspace{20\baselineskip}
\section{Đa thức cục bộ từ các đạo hàm đầu trái}
\begin{algorithm}[H]
\caption{Dựng một đoạn và truyền các đạo hàm}
\label{alg:spl-local-jet}
\begin{algorithmic}[1]
\Require $1\leq p\leq4$; $h>0$; $s^{(0)}=y$, $s^{(1)},\ldots,s^{(p-1)}$; giá trị đầu phải $z$.
\Ensure $a_0,\ldots,a_p$ của $P(t)=\sum_{j=0}^pa_jt^j$ với $P^{(j)}(0)=s^{(j)}$ ($j<p$), $P(h)=z$; $r^{(k)}=P^{(k)}(h)$ ($1\leq k<p$).
\State $a_j\gets s^{(j)}/j!$ ($0\leq j<p$); $v\gets a_{p-1}$.
\For{$j$ từ $p-2$ về $0$}
 \State $v\gets a_j+hv$.
\EndFor
\Statex
\State $a_p\gets(z-v)/h^p$.
\State $r^{(k)}\gets\sum_{j=k}^p j!a_jh^{j-k}/(j-k)!$ ($1\leq k<p$).
\State Kết thúc.
\end{algorithmic}
\end{algorithm}

\Needspace{20\baselineskip}
\section{Hàm ghép trơn bậc không quá bốn}
\begin{algorithm}[H]
\caption{Dựng spline từ đạo hàm đầu trái: sơ lược}
\label{alg:spl-initial-jet-short}
\begin{algorithmic}[1]
\Require $N\geq1$, $1\leq p\leq4$; $x_0<\cdots<x_N$; $y_0,\ldots,y_N$; $\alpha_1,\ldots,\alpha_{p-1}$.
\Ensure Các hệ số $a_{i,j}$ của spline trong định lý \ref{dl:spl-initial-jet}: $S_i(x)=\sum_{j=0}^pa_{i,j}(x-x_i)^j$.
\State $s_0^{(k)}\gets\alpha_k$ ($1\leq k<p$).
\For{$i$ từ $0$ tới $N-1$}
 \State $h\gets x_{i+1}-x_i$; $s_i^{(0)}\gets y_i$.
 \State Dựng $a_{i,0},\ldots,a_{i,p}$ và $s_{i+1}^{(k)}$ bằng thuật toán \ref{alg:spl-local-jet} với $h,s_i^{(k)},y_{i+1}$.
\EndFor
\Statex
\State Kết thúc.
\end{algorithmic}
\end{algorithm}
\begin{algorithm}[H]
\caption{Dựng spline từ đạo hàm đầu trái: chi tiết}
\label{alg:spl-initial-jet}
\begin{algorithmic}[1]
\Require $N\geq1$, $1\leq p\leq4$; $x_0<\cdots<x_N$; $y_0,\ldots,y_N$; $\alpha_1,\ldots,\alpha_{p-1}$.
\Ensure $S_i(x)=\sum_{j=0}^pa_{i,j}(x-x_i)^j$ nội suy dữ liệu, thuộc $C^{p-1}$ và $S^{(k)}(x_0)=\alpha_k$ ($1\leq k<p$).
\State $s_0^{(k)}\gets\alpha_k$ ($1\leq k<p$).
\For{$i$ từ $0$ tới $N-1$}
 \State $h\gets x_{i+1}-x_i$; $a_{i,0}\gets y_i$.
 \For{$j$ từ $1$ tới $p-1$}
  \State $a_{i,j}\gets s_i^{(j)}/j!$.
 \EndFor
 \Statex
 \State $v\gets a_{i,p-1}$.
 \For{$j$ từ $p-2$ về $0$}
  \State $v\gets a_{i,j}+hv$.
 \EndFor
 \Statex
 \State $a_{i,p}\gets(y_{i+1}-v)/h^p$.
 \For{$k$ từ $1$ tới $p-1$}
  \State $s_{i+1}^{(k)}\gets0$.
  \For{$j$ từ $k$ tới $p$}
   \State $s_{i+1}^{(k)}\gets s_{i+1}^{(k)}+j!a_{i,j}h^{j-k}/(j-k)!$.
  \EndFor
 \EndFor
\EndFor
\Statex
\State Kết thúc.
\end{algorithmic}
\end{algorithm}

\Needspace{20\baselineskip}
\section{Giá trị và đạo hàm của spline từ hệ số cục bộ}
\begin{algorithm}[H]
\caption{Chọn đoạn và dùng Horner đồng thời}
\label{alg:spl-general-evaluate}
\begin{algorithmic}[1]
\Require $N\geq1$, $1\leq p\leq4$; $x_0<\cdots<x_N$; các hệ số cục bộ $S_i(x)=\sum_{j=0}^pa_{i,j}(x-x_i)^j$ của spline; $\xi\in[x_0,x_N]$.
\Ensure $v=S(\xi)$; $u$ là đạo hàm của đoạn được chọn tại $\xi$; nếu $p\geq2$, $u=S'(\xi)$; nếu $p=1$ tại mốc trong đoạn, lấy đạo hàm từ phía phải; tại đầu mút, lấy đạo hàm một phía.
\State $i\gets0$.
\While{$i<N-1$ và $\xi\geq x_{i+1}$}
 \State $i\gets i+1$.
\EndWhile
\Statex
\State $t\gets\xi-x_i$; $v\gets a_{i,p}$; $u\gets0$.
\For{$j$ từ $p-1$ về $0$}
 \State $u\gets v+tu$; $v\gets a_{i,j}+tv$.
\EndFor
\Statex
\State Kết thúc.
\end{algorithmic}
\end{algorithm}

\Needspace{20\baselineskip}
\section{Spline bậc ba với hai moment biên}
\begin{algorithm}[H]
\caption{Moment spline với $S''(x_0)=A$, $S''(x_N)=B$: sơ lược}
\label{alg:spl-prescribed-moments-short}
\begin{algorithmic}[1]
\Require $N\geq1$; $x_0<\cdots<x_N$; $y_0,\ldots,y_N$; $A,B\in\R$.
\Ensure $M_0,\ldots,M_N$ của spline trong định lý \ref{dl:spl-prescribed-moments}.
\State $M_0\gets A$; $M_N\gets B$; $h_i\gets x_{i+1}-x_i$, $\delta_i\gets(y_{i+1}-y_i)/h_i$ ($0\leq i<N$).
\If{$N\geq2$}
 \State $q\gets N-1$; $d_j\gets2(h_j+h_{j+1})$, $r_j\gets6(\delta_{j+1}-\delta_j)$ ($0\leq j<q$).
 \State $e_j\gets h_{j+1}$ ($0\leq j<q-1$).
 \State $r_0\gets r_0-h_0A$; $r_{q-1}\gets r_{q-1}-h_{N-1}B$.
 \State Tính $\boldsymbol z$ bằng thuật toán \ref{alg:spl-thomas} cho $d_j,e_j,\boldsymbol r$.
 \State $M_i\gets z_{i-1}$ ($1\leq i<N$).
\EndIf
\Statex
\State Kết thúc.
\end{algorithmic}
\end{algorithm}
\begin{algorithm}[H]
\caption{Moment spline với $S''(x_0)=A$, $S''(x_N)=B$: chi tiết}
\label{alg:spl-prescribed-moments}
\begin{algorithmic}[1]
\Require $N\geq1$; $x_0<\cdots<x_N$; $y_0,\ldots,y_N$; $A,B\in\R$.
\Ensure $M_0,\ldots,M_N$ của spline trong định lý \ref{dl:spl-prescribed-moments}.
\State $M_0\gets A$; $M_N\gets B$.
\For{$i$ từ $0$ tới $N-1$}
 \State $h_i\gets x_{i+1}-x_i$; $\delta_i\gets(y_{i+1}-y_i)/h_i$.
\EndFor
\Statex
\If{$N\geq2$}
 \State $q\gets N-1$.
 \For{$j$ từ $0$ tới $q-1$}
  \State $d_j\gets2(h_j+h_{j+1})$; $r_j\gets6(\delta_{j+1}-\delta_j)$.
 \EndFor
 \Statex
 \For{$j$ từ $0$ tới $q-2$}
  \State $e_j\gets h_{j+1}$.
 \EndFor
 \Statex
 \State $r_0\gets r_0-h_0A$; $r_{q-1}\gets r_{q-1}-h_{N-1}B$.
 \State $D_0\gets d_0$; $\rho_0\gets r_0$.
 \For{$j$ từ $1$ tới $q-1$}
  \State $D_j\gets d_j-e_{j-1}^2/D_{j-1}$; $\rho_j\gets r_j-e_{j-1}\rho_{j-1}/D_{j-1}$.
 \EndFor
 \Statex
 \State $M_{N-1}\gets\rho_{q-1}/D_{q-1}$.
 \For{$j$ từ $q-2$ về $0$}
  \State $M_{j+1}\gets(\rho_j-e_jM_{j+2})/D_j$.
 \EndFor
\EndIf
\Statex
\State Kết thúc.
\end{algorithmic}
\end{algorithm}
% END SPL-SUPPLEMENT-ALGORITHMS

% BEGIN LS-SUPPLEMENT-ALGORITHMS
\Needspace{20\baselineskip}
\section{Sai lệch trên bảng dữ liệu}
\begin{algorithm}[H]
\caption{Tổng bình phương và sai số trung bình phương}
\label{alg:ls-error-measures}
\begin{algorithmic}[1]
\Require $m\geq1$; $y_i,\widehat y_i\in\R$, $w_i>0$ ($0\leq i<m$).
\Ensure $r_i=y_i-\widehat y_i$, $Q=\sum_iw_ir_i^2$, $E_2=Q/\sum_iw_i$, $E=\sqrt{E_2}$.
\State $Q\gets0$; $s\gets0$.
\For{$i$ từ $0$ tới $m-1$}
 \State $r_i\gets y_i-\widehat y_i$; $Q\gets Q+w_ir_i^2$; $s\gets s+w_i$.
\EndFor
\Statex
\State $E_2\gets Q/s$; $E\gets\sqrt{E_2}$.
\State Kết thúc.
\end{algorithmic}
\end{algorithm}

\Needspace{20\baselineskip}
\section{Hàm thực nghiệm và sai số của nó}
\begin{algorithm}[H]
\caption{Bình phương tối thiểu và sai số: sơ lược}
\label{alg:ls-fit-errors-short}
\begin{algorithmic}[1]
\Require $m,q\geq1$; $A\in\R^{m\times q}$ có hạng cột $q$; $\boldsymbol y\in\R^m$; $w_i>0$ ($0\leq i<m$).
\Ensure $\boldsymbol a$ cực tiểu $\sum_iw_i((A\boldsymbol a)_i-y_i)^2$; $\widehat y_i=(A\boldsymbol a)_i$; $Q,E_2,E$ trong \eqref{eq:ls-mean-square-error}.
\State Tính $\boldsymbol a$ bằng thuật toán \ref{alg:ls-normal} cho $A,\boldsymbol y,w_i$.
\State $\widehat y_i\gets\sum_{j=0}^{q-1}A_{ij}a_j$ ($0\leq i<m$).
\State Tính $Q,E_2,E$ bằng thuật toán \ref{alg:ls-error-measures} cho $y_i,\widehat y_i,w_i$.
\State Kết thúc.
\end{algorithmic}
\end{algorithm}
\begin{algorithm}[H]
\caption{Bình phương tối thiểu và sai số: chi tiết}
\label{alg:ls-fit-errors}
\begin{algorithmic}[1]
\Require $m,q\geq1$; $A\in\R^{m\times q}$ có hạng cột $q$; $\boldsymbol y\in\R^m$; $w_i>0$ ($0\leq i<m$).
\Ensure $\boldsymbol a$ cực tiểu $\sum_iw_i((A\boldsymbol a)_i-y_i)^2$; $\widehat y_i=(A\boldsymbol a)_i$; $Q,E_2,E$ trong \eqref{eq:ls-mean-square-error}.
\For{$j$ từ $0$ tới $q-1$}
 \State $b_j\gets\sum_{i=0}^{m-1}w_iA_{ij}y_i$.
 \For{$k$ từ $0$ tới $j$}
  \State $G_{jk}\gets\sum_{i=0}^{m-1}w_iA_{ij}A_{ik}$; $G_{kj}\gets G_{jk}$.
 \EndFor
\EndFor
\Statex
\For{$j$ từ $0$ tới $q-1$}
 \State $L_{jj}\gets\sqrt{G_{jj}-\sum_{k=0}^{j-1}L_{jk}^2}$.
 \For{$i$ từ $j+1$ tới $q-1$}
  \State $L_{ij}\gets(G_{ij}-\sum_{k=0}^{j-1}L_{ik}L_{jk})/L_{jj}$.
 \EndFor
\EndFor
\Statex
\For{$j$ từ $0$ tới $q-1$}
 \State $u_j\gets(b_j-\sum_{k=0}^{j-1}L_{jk}u_k)/L_{jj}$.
\EndFor
\Statex
\For{$j$ từ $q-1$ về $0$}
 \State $a_j\gets(u_j-\sum_{k=j+1}^{q-1}L_{kj}a_k)/L_{jj}$.
\EndFor
\Statex
\State $Q\gets0$; $s\gets0$.
\For{$i$ từ $0$ tới $m-1$}
 \State $\widehat y_i\gets\sum_{j=0}^{q-1}A_{ij}a_j$; $r_i\gets y_i-\widehat y_i$.
 \State $Q\gets Q+w_ir_i^2$; $s\gets s+w_i$.
\EndFor
\Statex
\State $E_2\gets Q/s$; $E\gets\sqrt{E_2}$.
\State Kết thúc.
\end{algorithmic}
\end{algorithm}

\Needspace{20\baselineskip}
\section{Mô hình mũ sau phép biến đổi logarit}
\begin{algorithm}[H]
\caption{Khớp $a\exp(\sum_{j=1}^rb_j\varphi_j(x))$ trên thang logarit}
\label{alg:ls-exponential-transform}
\begin{algorithmic}[1]
\Require $r\geq0$, $m\geq r+1$; $s\in\{-1,1\}$, $sy_i>0$, $w_i>0$; các giá trị $\varphi_j(x_i)$; thiết kế $A_{i0}=1$, $A_{ij}=\varphi_j(x_i)$ có hạng cột đầy đủ.
\Ensure $a,b_j$ cực tiểu \eqref{eq:ls-exponential-objective}; $Q_{\log},E_{2,\log},E_{\log}$ trên thang logarit; $Q_y,E_{2,y},E_y$ trên thang dữ liệu gốc.
\State $z_i\gets\ln(sy_i)$ ($0\leq i<m$).
\State Tính $\boldsymbol\theta=(c,b_1,\ldots,b_r)^T$, $Q_{\log},E_{2,\log},E_{\log}$ bằng thuật toán \ref{alg:ls-fit-errors-short} cho $A,\boldsymbol z,w_i$.
\State $a\gets s e^c$.
\For{$i$ từ $0$ tới $m-1$}
 \State $\widehat y_i\gets s\exp(\sum_{j=0}^rA_{ij}\theta_j)$.
\EndFor
\Statex
\State Tính $Q_y,E_{2,y},E_y$ bằng thuật toán \ref{alg:ls-error-measures} cho $y_i,\widehat y_i,w_i$.
\State Kết thúc.
\end{algorithmic}
\end{algorithm}

\Needspace{20\baselineskip}
\section{Mô hình lũy thừa sau phép biến đổi logarit}
\begin{algorithm}[H]
\caption{Khớp $ax^b$ trên thang logarit}
\label{alg:ls-power-transform}
\begin{algorithmic}[1]
\Require $m\geq2$; $x_i>0$ có ít nhất hai giá trị khác nhau; $s\in\{-1,1\}$, $sy_i>0$; $w_i>0$ ($0\leq i<m$).
\Ensure $a,b$ cực tiểu tiêu chuẩn trong hệ quả \ref{hq:ls-power-transform}; $Q_{\log},E_{2,\log},E_{\log}$ và $Q_y,E_{2,y},E_y$.
\State $u_i\gets\ln x_i$; $z_i\gets\ln(sy_i)$ ($0\leq i<m$).
\State Tính $c,b$ bằng thuật toán \ref{alg:ls-affine} cho $u_i,z_i,w_i$.
\State $a\gets s e^c$.
\For{$i$ từ $0$ tới $m-1$}
 \State $\widehat z_i\gets c+bu_i$; $\widehat y_i\gets s e^{\widehat z_i}$.
\EndFor
\Statex
\State Tính $Q_{\log},E_{2,\log},E_{\log}$ bằng thuật toán \ref{alg:ls-error-measures} cho $z_i,\widehat z_i,w_i$.
\State Tính $Q_y,E_{2,y},E_y$ bằng thuật toán \ref{alg:ls-error-measures} cho $y_i,\widehat y_i,w_i$.
\State Kết thúc.
\end{algorithmic}
\end{algorithm}

\Needspace{20\baselineskip}
\section{Mô hình logarit sau phép biến đổi mũ}
\begin{algorithm}[H]
\caption{Khớp $\ln(\sum_jb_j\varphi_j(x))$ và kiểm tra miền xác định}
\label{alg:ls-logarithmic-transform}
\begin{algorithmic}[1]
\Require $m,q\geq1$; $A_{ij}=\varphi_{j+1}(x_i)$ có hạng cột $q$; $y_i\in\R$, $w_i>0$ ($0\leq i<m$).
\Ensure $\boldsymbol b$ cực tiểu \eqref{eq:ls-logarithmic-transform}; $Q_z,E_{2,z},E_z$ trên thang biến đổi; $\sigma\in\{0,1\}$, $\sigma=1$ khi $(A\boldsymbol b)_i>0$ với mọi $i$; nếu $\sigma=1$, có $g_i=\ln((A\boldsymbol b)_i)$ và $Q_y,E_{2,y},E_y$.
\State $z_i\gets e^{y_i}$ ($0\leq i<m$).
\State Tính $\boldsymbol b$, $\widehat z_i$, $Q_z,E_{2,z},E_z$ bằng thuật toán \ref{alg:ls-fit-errors-short} cho $A,\boldsymbol z,w_i$.
\State $\sigma\gets0$.
\If{$\min_i\widehat z_i>0$}
 \State $\sigma\gets1$; $g_i\gets\ln\widehat z_i$ ($0\leq i<m$).
 \State Tính $Q_y,E_{2,y},E_y$ bằng thuật toán \ref{alg:ls-error-measures} cho $y_i,g_i,w_i$.
\EndIf
\Statex
\State Kết thúc.
\end{algorithmic}
\end{algorithm}
% END LS-SUPPLEMENT-ALGORITHMS

% END EXTENSION SPLINE LS ALGORITHMS

% BEGIN EXTENSION CHAPTER 2 ALGORITHMS
\Needspace{20\baselineskip}
\section{Trọng số đạo hàm từ các mốc chuẩn}
\begin{algorithm}[H]
\caption{Dựng trọng số đạo hàm bằng Horner -- sơ lược}
\label{alg:dnum-weights}
\begin{algorithmic}[1]
\Require Số nguyên $k\geq2$, $1\leq s<k$; $\boldsymbol\tau=(\tau_0,\ldots,\tau_{k-1})^T\in\R^k$ có các thành phần đôi một phân biệt.
\Ensure $\boldsymbol w\in\R^k$ thỏa $\sum_{i=0}^{k-1}w_i\tau_i^j=s!\delta_{j,s}$ với $0\leq j<k$, theo \eqref{eq:dnum-moments}.
\State Áp dụng Thuật toán \ref{alg:horner-product} với $\boldsymbol\tau$; thu được $\boldsymbol a\in\R^{k+1}$, $\sum_{j=0}^ka_jt^j=\prod_{i=0}^{k-1}(t-\tau_i)$.
\For{$i$ từ $0$ tới $k-1$}
  \State Áp dụng Thuật toán \ref{alg:horner-div} với $\boldsymbol a$ tại $\tau_i$; thu được $\boldsymbol q_i=(q_{i,0},\ldots,q_{i,k-1})^T$ và phần dư $0$.
  \State Áp dụng Thuật toán \ref{alg:horner-eval} với $\boldsymbol q_i$ tại $\tau_i$; thu được $\lambda_i$.
  \State $w_i\gets s!q_{i,s}/\lambda_i$.
\EndFor
\Statex
\State $\boldsymbol w\gets(w_0,\ldots,w_{k-1})^T$.
\State Kết thúc.
\end{algorithmic}
\end{algorithm}

\begin{algorithm}[H]
\caption{Dựng trọng số đạo hàm bằng Horner -- chi tiết}
\label{alg:dnum-weights-detailed}
\small
\begin{algorithmic}[1]
\Require Số nguyên $k\geq2$, $1\leq s<k$; $\boldsymbol\tau=(\tau_0,\ldots,\tau_{k-1})^T\in\R^k$ có các thành phần đôi một phân biệt.
\Ensure $\boldsymbol w\in\R^k$ thỏa $\sum_{i=0}^{k-1}w_i\tau_i^j=s!\delta_{j,s}$ với $0\leq j<k$.
\State $\boldsymbol a\gets(1)^T$.
\For{$r$ từ $0$ tới $k-1$}
  \State $d_{r+1}\gets a_r$.
  \For{$j$ từ $r$ về $1$}
    \State $d_j\gets a_{j-1}-\tau_ra_j$.
  \EndFor
  \Statex
  \State $d_0\gets-\tau_ra_0$; $\boldsymbol a\gets(d_0,\ldots,d_{r+1})^T$.
\EndFor
\Statex
\State $F\gets\prod_{j=1}^{s}j$.
\For{$i$ từ $0$ tới $k-1$}
  \State $b_k\gets a_k$.
  \For{$j$ từ $k-1$ về $0$}
    \State $b_j\gets a_j+\tau_ib_{j+1}$.
  \EndFor
  \Statex
  \State $\lambda_i\gets b_k$.
  \For{$j$ từ $k-1$ về $1$}
    \State $\lambda_i\gets b_j+\tau_i\lambda_i$.
  \EndFor
  \Statex
  \State $w_i\gets F b_{s+1}/\lambda_i$.
\EndFor
\Statex
\State $\boldsymbol w\gets(w_0,\ldots,w_{k-1})^T$.
\State Kết thúc.
\end{algorithmic}
\end{algorithm}

\Needspace{20\baselineskip}
\section{Đạo hàm hai điểm và ba điểm}
\begin{algorithm}[H]
\caption{Tính các công thức đạo hàm hai điểm và ba điểm}
\label{alg:dnum-small-stencils}
\begin{algorithmic}[1]
\Require $h>0$; $\boldsymbol y=(y_{-2},y_{-1},y_0,y_1,y_2)^T$, trong đó $y_j=f(x+jh)$.
\Ensure $(D_+,D_-,D_0,D_{+,3},D_{-,3},D_2)$ là các công thức của Định lý \ref{dl:dnum-two-point} và \ref{dl:dnum-three-point} tại $x$.
\State $D_+\gets(y_1-y_0)/h$.
\State $D_-\gets(y_0-y_{-1})/h$.
\State $D_0\gets(y_1-y_{-1})/(2h)$.
\State $D_{+,3}\gets(-3y_0+4y_1-y_2)/(2h)$.
\State $D_{-,3}\gets(3y_0-4y_{-1}+y_{-2})/(2h)$.
\State $D_2\gets(y_{-1}-2y_0+y_1)/h^2$.
\State Kết thúc.
\end{algorithmic}
\end{algorithm}

\Needspace{20\baselineskip}
\section{Đạo hàm từ một vector trọng số}
\begin{algorithm}[H]
\caption{Tính công thức đạo hàm $k$ điểm}
\label{alg:dnum-apply}
\begin{algorithmic}[1]
\Require Số nguyên $k\geq2$, $1\leq s<k$; $h>0$; $\boldsymbol w=(w_0,\ldots,w_{k-1})^T$, $\boldsymbol y=(y_0,\ldots,y_{k-1})^T\in\R^k$, trong đó $y_i=f(x+h\tau_i)$.
\Ensure $D=h^{-s}\sum_{i=0}^{k-1}w_i y_i$, theo \eqref{eq:dnum-stencil}.
\State $v\gets0$.
\For{$i$ từ $0$ tới $k-1$}
  \State $v\gets v+w_i y_i$.
\EndFor
\Statex
\State $D\gets v/h^s$.
\State Kết thúc.
\end{algorithmic}
\end{algorithm}

\Needspace{20\baselineskip}
\section{Trọng số tích phân từ các mốc chuẩn}
\begin{algorithm}[H]
\caption{Dựng trọng số tích phân bằng Horner -- sơ lược}
\label{alg:quad-weights}
\begin{algorithmic}[1]
\Require Số nguyên $k\geq1$; $\boldsymbol\tau=(\tau_0,\ldots,\tau_{k-1})^T\in\R^k$ có các thành phần đôi một phân biệt.
\Ensure $\boldsymbol\lambda\in\R^k$ thỏa $\sum_{i=0}^{k-1}\lambda_i\tau_i^j=1/(j+1)$ với $0\leq j<k$, theo \eqref{eq:quad-moments}.
\State Áp dụng Thuật toán \ref{alg:horner-product} với $\boldsymbol\tau$; thu được $\boldsymbol a\in\R^{k+1}$, $\sum_{j=0}^ka_jt^j=\prod_{i=0}^{k-1}(t-\tau_i)$.
\For{$i$ từ $0$ tới $k-1$}
  \State Áp dụng Thuật toán \ref{alg:horner-div} với $\boldsymbol a$ tại $\tau_i$; thu được $\boldsymbol q_i=(q_{i,0},\ldots,q_{i,k-1})^T$ và phần dư $0$.
  \State Áp dụng Thuật toán \ref{alg:horner-eval} với $\boldsymbol q_i$ tại $\tau_i$; thu được $v_i$.
  \State $\lambda_i\gets v_i^{-1}\sum_{j=0}^{k-1}q_{i,j}/(j+1)$.
\EndFor
\Statex
\State $\boldsymbol\lambda\gets(\lambda_0,\ldots,\lambda_{k-1})^T$.
\State Kết thúc.
\end{algorithmic}
\end{algorithm}

\begin{algorithm}[H]
\caption{Dựng trọng số tích phân bằng Horner -- chi tiết}
\label{alg:quad-weights-detailed}
\small
\begin{algorithmic}[1]
\Require Số nguyên $k\geq1$; $\boldsymbol\tau=(\tau_0,\ldots,\tau_{k-1})^T\in\R^k$ có các thành phần đôi một phân biệt.
\Ensure $\boldsymbol\lambda\in\R^k$ thỏa $\sum_{i=0}^{k-1}\lambda_i\tau_i^j=1/(j+1)$ với $0\leq j<k$.
\State $\boldsymbol a\gets(1)^T$.
\For{$r$ từ $0$ tới $k-1$}
  \State $d_{r+1}\gets a_r$.
  \For{$j$ từ $r$ về $1$}
    \State $d_j\gets a_{j-1}-\tau_ra_j$.
  \EndFor
  \Statex
  \State $d_0\gets-\tau_ra_0$; $\boldsymbol a\gets(d_0,\ldots,d_{r+1})^T$.
\EndFor
\Statex
\For{$i$ từ $0$ tới $k-1$}
  \State $b_k\gets a_k$.
  \For{$j$ từ $k-1$ về $0$}
    \State $b_j\gets a_j+\tau_ib_{j+1}$.
  \EndFor
  \Statex
  \State $v_i\gets b_k$.
  \For{$j$ từ $k-1$ về $1$}
    \State $v_i\gets b_j+\tau_i v_i$.
  \EndFor
  \Statex
  \State $u_i\gets0$.
  \For{$j$ từ $0$ tới $k-1$}
    \State $u_i\gets u_i+b_{j+1}/(j+1)$.
  \EndFor
  \Statex
  \State $\lambda_i\gets u_i/v_i$.
\EndFor
\Statex
\State $\boldsymbol\lambda\gets(\lambda_0,\ldots,\lambda_{k-1})^T$.
\State Kết thúc.
\end{algorithmic}
\end{algorithm}

\Needspace{20\baselineskip}
\section{Điểm giữa, hình thang và Simpson trên một đoạn}
\begin{algorithm}[H]
\caption{Tính ba công thức tích phân trên một đoạn}
\label{alg:quad-single}
\begin{algorithmic}[1]
\Require $a<b$; $\boldsymbol y=(y_0,y_1,y_2)^T$ với $y_0=f(a)$, $y_1=f((a+b)/2)$, $y_2=f(b)$.
\Ensure $(M,T,S)$ là các công thức điểm giữa, hình thang và Simpson trên $[a,b]$.
\State $\ell\gets b-a$.
\State $M\gets\ell y_1$.
\State $T\gets\ell(y_0+y_2)/2$.
\State $S\gets\ell(y_0+4y_1+y_2)/6$.
\State Kết thúc.
\end{algorithmic}
\end{algorithm}

\Needspace{20\baselineskip}
\section{Điểm giữa và hình thang ghép trên lưới đều}
\begin{algorithm}[H]
\caption{Tính công thức điểm giữa và hình thang ghép}
\label{alg:quad-midpoint-trapezoid-composite}
\begin{algorithmic}[1]
\Require Số nguyên $N\geq1$; $a<b$; $h=(b-a)/N$; $\boldsymbol y=(y_0,\ldots,y_N)^T$, $\boldsymbol m=(m_0,\ldots,m_{N-1})^T$, với $y_i=f(a+ih)$, $m_i=f(a+(i+1/2)h)$.
\Ensure $(M_N,T_N)$ thỏa \eqref{eq:quad-midpoint-composite} và \eqref{eq:quad-trapezoid-composite}.
\State $u\gets0$; $v\gets(y_0+y_N)/2$.
\For{$i$ từ $0$ tới $N-1$}
  \State $u\gets u+m_i$.
  \If{$i\geq1$}
    \State $v\gets v+y_i$.
  \EndIf
\EndFor
\Statex
\State $(M_N,T_N)\gets(hu,hv)$.
\State Kết thúc.
\end{algorithmic}
\end{algorithm}

\Needspace{20\baselineskip}
\section{Simpson ghép trên lưới đều}
\begin{algorithm}[H]
\caption{Tính công thức Simpson ghép}
\label{alg:quad-simpson-composite}
\begin{algorithmic}[1]
\Require Số nguyên chẵn $N\geq2$; $a<b$; $h=(b-a)/N$; $\boldsymbol y=(y_0,\ldots,y_N)^T$, với $y_i=f(a+ih)$.
\Ensure $S_N$ thỏa \eqref{eq:quad-simpson-composite}.
\State $v\gets y_0+y_N$.
\For{$i$ từ $1$ tới $N-1$}
  \If{$i$ lẻ}
    \State $v\gets v+4y_i$.
  \Else
    \State $v\gets v+2y_i$.
  \EndIf
\EndFor
\Statex
\State $S_N\gets hv/3$.
\State Kết thúc.
\end{algorithmic}
\end{algorithm}

\Needspace{20\baselineskip}
\section{Ghép một công thức tích phân trên phân hoạch bất kỳ}
\begin{algorithm}[H]
\caption{Tính công thức tích phân ghép từ vector trọng số}
\label{alg:quad-general-composite}
\begin{algorithmic}[1]
\Require Số nguyên $N,k\geq1$; $\boldsymbol t=(t_0,\ldots,t_N)^T$, $t_0<\cdots<t_N$; $\boldsymbol\tau\in[0,1]^k$, $\boldsymbol\lambda\in\R^k$; các vector $\boldsymbol y^{(j)}\in\R^k$ với $y_i^{(j)}=f(t_j+(t_{j+1}-t_j)\tau_i)$, $0\leq j<N$.
\Ensure $Q_{\mathcal T}=\sum_{j=0}^{N-1}(t_{j+1}-t_j)\sum_{i=0}^{k-1}\lambda_i y_i^{(j)}$, theo Mệnh đề \ref{md:quad-composite-general}.
\State $v\gets0$.
\For{$j$ từ $0$ tới $N-1$}
  \State $\ell_j\gets t_{j+1}-t_j$.
  \State $u_j\gets\sum_{i=0}^{k-1}\lambda_i y_i^{(j)}$.
  \State $v\gets v+\ell_j u_j$.
\EndFor
\Statex
\State $Q_{\mathcal T}\gets v$.
\State Kết thúc.
\end{algorithmic}
\end{algorithm}

\Needspace{20\baselineskip}
\section{Đạo hàm cấp một tại mọi mốc của bảng đều}
\begin{algorithm}[H]
\caption{Tính đạo hàm cấp một từ các khối $n+1$ mốc}
\label{alg:dnum-table-first}
\begin{algorithmic}[1]
\Require Số nguyên $N\geq n\geq1$; $h>0$; $\boldsymbol y=(y_0,\ldots,y_N)^T$, $y_i=f(a+ih)$.
\Ensure $\boldsymbol d=(d_0,\ldots,d_N)^T$, $d_j=p_m'(a+jh)$ theo Mệnh đề \ref{md:dnum-balanced-block}, với $p_m$ nội suy các cặp $(a+(m+i)h,y_{m+i})$, $0\leq i\leq n$.
\For{$j$ từ $0$ tới $N$}
  \State $m\gets\min\{\max\{j-\lfloor n/2\rfloor,0\},N-n\}$.
  \State $\ell\gets j-m$; $r\gets n-\ell$.
  \State $w_\ell\gets\sum_{q=1}^{\ell}1/q-\sum_{q=1}^r1/q$.
  \For{$i$ từ $0$ tới $n$}
    \If{$i\ne\ell$}
      \State $w_i\gets(-1)^{i-\ell}\ell!r!/((\ell-i)i!(n-i)!)$.
    \EndIf
  \EndFor
  \Statex
  \State $d_j\gets h^{-1}\sum_{i=0}^nw_i y_{m+i}$.
\EndFor
\Statex
\State $\boldsymbol d\gets(d_0,\ldots,d_N)^T$.
\State Kết thúc.
\end{algorithmic}
\end{algorithm}

\Needspace{20\baselineskip}
\section{Chọn bước đạo hàm và kiểm tra độ chính xác}
\begin{algorithm}[H]
\caption{Chọn bước theo chặn cắt cụt và nhiễu}
\label{alg:dnum-step-certificate}
\begin{algorithmic}[1]
\Require $C,B,p,s,H_0,\varepsilon>0$; chặn sai số $Ch^p+Bh^{-s}$ hợp lệ với $0<h\leq H_0$.
\Ensure $h$ tối thiểu hóa chặn; $E=Ch^p+Bh^{-s}$; $\chi=1$ nếu $E\leq\varepsilon$, $\chi=0$ nếu $E>\varepsilon$.
\State $h\gets(sB/(pC))^{1/(p+s)}$.
\If{$h>H_0$}
  \State $h\gets H_0$.
\EndIf
\Statex
\State $E\gets Ch^p+Bh^{-s}$; $\chi\gets0$.
\If{$E\leq\varepsilon$}
  \State $\chi\gets1$.
\EndIf
\Statex
\State Kết thúc.
\end{algorithmic}
\end{algorithm}

\Needspace{20\baselineskip}
\section{Chọn số đoạn từ chặn sai số tích phân}
\begin{algorithm}[H]
\caption{Chọn số đoạn nguyên thỏa độ chính xác và điều kiện bội}
\label{alg:quad-count}
\begin{algorithmic}[1]
\Require $A\geq0$, $r>0$, $0\leq\nu<\varepsilon$; số nguyên $g\geq1$; chặn $|I-Q_N|\leq A/N^r+\nu$.
\Ensure $N$ là bội nguyên dương nhỏ nhất của $g$ thỏa $A/N^r+\nu\leq\varepsilon$; $E=A/N^r+\nu$.
\State $v\gets(A/(\varepsilon-\nu))^{1/r}$.
\State $N\gets g\lceil v/g\rceil$.
\If{$N<g$}
  \State $N\gets g$.
\EndIf
\Statex
\State $E\gets A/N^r+\nu$.
\State Kết thúc.
\end{algorithmic}
\end{algorithm}

\Needspace{20\baselineskip}
\section{Trọng số Newton--Cotes theo bậc đa thức}
\begin{algorithm}[H]
\caption{Dựng trọng số Newton--Cotes đóng bậc $n$ -- sơ lược}
\label{alg:quad-nc-degree-weights}
\begin{algorithmic}[1]
\Require Số nguyên $n\geq1$.
\Ensure $\boldsymbol\lambda=(\lambda_0,\ldots,\lambda_n)^T$, $\sum_i\lambda_i(i/n)^j=1/(j+1)$ với $0\leq j\leq n$.
\State $\boldsymbol\tau\gets(0,1/n,\ldots,(n-1)/n,1)^T$.
\State Áp dụng Thuật toán \ref{alg:quad-weights} với $k=n+1$, vector $\boldsymbol\tau$; thu được $\boldsymbol\lambda$.
\State Kết thúc.
\end{algorithmic}
\end{algorithm}

\begin{algorithm}[H]
\caption{Dựng trọng số Newton--Cotes đóng bậc $n$ -- chi tiết}
\label{alg:quad-nc-degree-weights-detailed}
\small
\begin{algorithmic}[1]
\Require Số nguyên $n\geq1$.
\Ensure $\boldsymbol\lambda=(\lambda_0,\ldots,\lambda_n)^T$, $\sum_i\lambda_i(i/n)^j=1/(j+1)$ với $0\leq j\leq n$.
\State $k\gets n+1$; $\tau_i\gets i/n$ với $0\leq i<k$; $\boldsymbol a\gets(1)^T$.
\For{$r$ từ $0$ tới $k-1$}
  \State $d_{r+1}\gets a_r$.
  \For{$j$ từ $r$ về $1$}
    \State $d_j\gets a_{j-1}-\tau_ra_j$.
  \EndFor
  \Statex
  \State $d_0\gets-\tau_ra_0$; $\boldsymbol a\gets(d_0,\ldots,d_{r+1})^T$.
\EndFor
\Statex
\For{$i$ từ $0$ tới $k-1$}
  \State $b_k\gets a_k$.
  \For{$j$ từ $k-1$ về $0$}
    \State $b_j\gets a_j+\tau_ib_{j+1}$.
  \EndFor
  \Statex
  \State $v_i\gets b_k$.
  \For{$j$ từ $k-1$ về $1$}
    \State $v_i\gets b_j+\tau_iv_i$.
  \EndFor
  \Statex
  \State $u_i\gets\sum_{j=0}^{k-1}b_{j+1}/(j+1)$; $\lambda_i\gets u_i/v_i$.
\EndFor
\Statex
\State $\boldsymbol\lambda\gets(\lambda_0,\ldots,\lambda_n)^T$.
\State Kết thúc.
\end{algorithmic}
\end{algorithm}

\Needspace{20\baselineskip}
\section{Ghép Newton--Cotes khi bảng còn một khối ngắn}
\begin{algorithm}[H]
\caption{Tính tích phân từ toàn bộ một bảng đều}
\label{alg:quad-table-panels}
\begin{algorithmic}[1]
\Require Số nguyên $N,n\geq1$; $h>0$; $\boldsymbol y=(y_0,\ldots,y_N)^T$; các vector trọng số đóng $\boldsymbol\lambda^{(d)}\in\R^{d+1}$, $1\leq d\leq n$.
\Ensure $Q$ thỏa \eqref{eq:quad-table-panels}, dùng toàn bộ giá trị $y_i=f(a+ih)$.
\State $q\gets\lfloor N/n\rfloor$; $v\gets N-qn$.
\State $Q\gets hn\sum_{j=0}^{q-1}\sum_{i=0}^n\lambda_i^{(n)}y_{nj+i}$.
\If{$v>0$}
  \State $Q\gets Q+hv\sum_{i=0}^v\lambda_i^{(v)}y_{qn+i}$.
\EndIf
\Statex
\State Kết thúc.
\end{algorithmic}
\end{algorithm}

\Needspace{20\baselineskip}
\section{Tích phân kép bằng một tổng tensor}
\begin{algorithm}[H]
\caption{Tính tích phân kép từ các trọng số theo từng chiều}
\label{alg:quad-tensor}
\begin{algorithmic}[1]
\Require Số nguyên $N_x,N_y\geq1$; $\boldsymbol\alpha\in\R^{N_x+1}$, $\boldsymbol\beta\in\R^{N_y+1}$; các vector $\boldsymbol f_i=(f_{i,0},\ldots,f_{i,N_y})^T$, $0\leq i\leq N_x$, với $f_{i,j}=f(x_i,y_j)$.
\Ensure $Q_{xy}=\sum_i\sum_j\alpha_i\beta_jf_{i,j}$, theo \eqref{eq:quad-tensor}.
\State $v\gets0$.
\For{$i$ từ $0$ tới $N_x$}
  \State $u_i\gets\sum_{j=0}^{N_y}\beta_jf_{i,j}$.
  \State $v\gets v+\alpha_i u_i$.
\EndFor
\Statex
\State $Q_{xy}\gets v$.
\State Kết thúc.
\end{algorithmic}
\end{algorithm}

\Needspace{20\baselineskip}
\section{Chọn lưới tích phân kép theo độ chính xác}
\begin{algorithm}[H]
\caption{Chọn số đoạn theo hai chiều -- sơ lược}
\label{alg:quad-tensor-count}
\begin{algorithmic}[1]
\Require $A_x,A_y\geq0$, $r>0$, $0\leq\nu<\varepsilon$; số nguyên $g\geq1$; chặn $|I-Q_{xy}|\leq A_x/N_x^r+A_y/N_y^r+\nu$.
\Ensure $N_x,N_y$ là bội nguyên dương của $g$; $E=A_x/N_x^r+A_y/N_y^r+\nu\leq\varepsilon$.
\State $\delta\gets\varepsilon-\nu$.
\State Áp dụng Thuật toán \ref{alg:quad-count} với $A=A_x$, cấp $r$, nhiễu $0$, độ chính xác $\delta/2$, bội $g$; thu được $N_x$.
\State Áp dụng Thuật toán \ref{alg:quad-count} với $A=A_y$, cấp $r$, nhiễu $0$, độ chính xác $\delta/2$, bội $g$; thu được $N_y$.
\State $E\gets A_x/N_x^r+A_y/N_y^r+\nu$.
\State Kết thúc.
\end{algorithmic}
\end{algorithm}

\begin{algorithm}[H]
\caption{Chọn số đoạn theo hai chiều -- chi tiết}
\label{alg:quad-tensor-count-detailed}
\begin{algorithmic}[1]
\Require $A_x,A_y\geq0$, $r>0$, $0\leq\nu<\varepsilon$; số nguyên $g\geq1$; chặn $|I-Q_{xy}|\leq A_x/N_x^r+A_y/N_y^r+\nu$.
\Ensure $N_x,N_y$ là bội nguyên dương của $g$; $E=A_x/N_x^r+A_y/N_y^r+\nu\leq\varepsilon$.
\State $\delta\gets\varepsilon-\nu$.
\State $N_x\gets g\lceil(2A_x/\delta)^{1/r}/g\rceil$; $N_y\gets g\lceil(2A_y/\delta)^{1/r}/g\rceil$.
\If{$N_x<g$}
  \State $N_x\gets g$.
\EndIf
\Statex
\If{$N_y<g$}
  \State $N_y\gets g$.
\EndIf
\Statex
\State $E\gets A_x/N_x^r+A_y/N_y^r+\nu$.
\State Kết thúc.
\end{algorithmic}
\end{algorithm}

\Needspace{20\baselineskip}
\section{So sánh hai lưới tích phân có chặn phần dư}
\begin{algorithm}[H]
\caption{Tính chỉ số sai số và giá trị ngoại suy từ hai lưới}
\label{alg:quad-refinement-indicator}
\begin{algorithmic}[1]
\Require $h,r,q>0$, $B\geq0$; hai giá trị $Q_h,Q_{h/2}$ thỏa $I-Q_t=ct^r+R(t)$, $|R(t)|\leq Bt^{r+q}$ với $t=h,h/2$ và cùng $c$.
\Ensure $e=(Q_{h/2}-Q_h)/(2^r-1)$, $Q_*=Q_{h/2}+e$; $|I-Q_*|\leq\Delta$.
\State $e\gets(Q_{h/2}-Q_h)/(2^r-1)$.
\State $Q_*\gets Q_{h/2}+e$.
\State $\Delta\gets(1+2^{-q})Bh^{r+q}/(2^r-1)$.
\State Kết thúc.
\end{algorithmic}
\end{algorithm}

% END EXTENSION CHAPTER 2 ALGORITHMS

\backmatter
\begin{thebibliography}{99}
\raggedright
\bibitem{outline} Đề cương học phần MI3042, \textit{Phương pháp số}, phiên bản 26/07/2025, tài liệu trong thư mục \texttt{doc}.
\bibitem{slides} Hà Thị Ngọc Yến, \textit{Xấp xỉ hàm số bằng đa thức nội suy}, bài giảng năm 2026, tài liệu trong thư mục \texttt{doc}.
\bibitem{newton-slides} Hà Thị Ngọc Yến, \textit{Đa thức nội suy Newton}, bài giảng năm 2026, tài liệu trong thư mục \texttt{doc}.
\bibitem{newton-slides-2022} Hà Thị Ngọc Yến, \textit{Đa thức nội suy Newton}, bài giảng tháng 9/2022, tệp \texttt{15\_noisuynewton\_.1m.pdf} trong thư mục \texttt{doc}.
\bibitem{central-slides} Hà Thị Ngọc Yến, \textit{Nội suy trung tâm}, bài giảng tháng 9/2022, tệp \texttt{16\_noisuytrungtam\_.2m.pdf} trong thư mục \texttt{doc}.
\bibitem{inverse-slides} Hà Thị Ngọc Yến, \textit{Nội suy ngược}, bài giảng năm 2019, tệp \texttt{18\_noisuynguoc\_.1m.pdf} trong thư mục \texttt{doc}.
\bibitem{integral-slides} Hà Thị Ngọc Yến, \textit{Tích phân}, tệp \texttt{17\_tichphan\_.1m.pdf} trong thư mục \texttt{doc}.
\bibitem{ls-slides} Hà Thị Ngọc Yến, \textit{Bình phương tối thiểu}, tệp \texttt{19\_binhphuongtoithieu\_.2m.pdf} trong thư mục \texttt{doc}.
\bibitem{spline-slides} Hà Thị Ngọc Yến, \textit{Spline}, tệp \texttt{20\_spline\_.5m.pdf} trong thư mục \texttt{doc}.
\bibitem{central-nptel} A. K. Lal, \textit{Gauss's and Stirling's Formulas}, bài giảng NPTEL, IIT Kanpur. \url{https://archive.nptel.ac.in/content/storage2/courses/122104018/node114.html}.
\bibitem{exercises} Hà Thị Ngọc Yến, \textit{Bài tập Phương pháp số MI3042}, năm 2026.
\bibitem{suli} Endre Süli, \textit{Numerical Computation I}, Oxford University Computing Laboratory, 1998. \url{https://people.maths.ox.ac.uk/suli/nc.pdf}.
\bibitem{berrut} Jean-Paul Berrut và Lloyd N. Trefethen, \textit{Barycentric Lagrange Interpolation}, SIAM Review 46(3), 501--517, 2004. DOI: \href{https://doi.org/10.1137/S0036144502417715}{10.1137/S0036144502417715}.
\bibitem{oracle-754} Oracle, \textit{Numerical Computation Guide}, Chương 2: IEEE Arithmetic. \url{https://docs.oracle.com/cd/E19059-01/stud.10/819-0499/ncg_math.html}.
\bibitem{bindel} David Bindel, \textit{CS 5220: Floating Point Arithmetic}, Cornell University, 2024. \url{https://www.cs.cornell.edu/courses/cs5220/2024fa/slides/20-fp.html}.
\end{thebibliography}
\end{document}
